%% PERSIAN EDITION of main.tex. One file, as the English book is, and the
%% same preamble but for the three things Unicode right-to-left typesetting
%% forces: the fonts, the margins, and xepersian at the foot.
%%
%% XeLaTeX only -- pdflatex cannot build this file. The English book's
%% pdfTeX font stack (fontenc, inputenc, mathpazo, helvet) does not survive
%% the move, so the text fonts are set through fontspec after xepersian
%% below. Only the math survives untouched: eulervm over TeX Gyre Pagella
%% is the same Palatino-with-Euler pairing mathpazo+eulervm gives the
%% English edition, in the OpenType clone XeLaTeX can load.
%%
%% Build:  latexmk -xelatex main-fa.tex
%%
%% What is translated and what is not: prose, headings, box titles, theorem
%% names and code comments are Persian; notation, macro names, labels,
%% \ref targets, interface and operation names, code keywords and the
%% bibliography stay in Latin script, being notation rather than prose.
%% See GLOSSARY-fa.md for the terminology decisions.
%%
%% Book, one-sided: this is read on a screen, so there is no verso to
%% balance and the running head is the same on every page.
\documentclass[10pt,oneside,openany]{book}

%% Math. The English edition sets Euler over Palatino with mathpazo plus
%% eulervm, and eulervm cannot come along: it is built on virtual fonts,
%% and XeTeX's native-font code cannot map a glyph through one. The run
%% dies at shipout -- "Internal error: bad native font flag in
%% `map_char_to_glyph'", reproducibly, on the first page whose math it
%% reaches, with no error a reader of the source could have predicted.
%% The design survives the move all the same, because Euler exists as an
%% OpenType font: unicode-math loads Euler-Math.otf below, and the
%% mathematics is the mathematics of the English page.
\usepackage{unicode-math}
%% 5.5in x 8.5in trade-paperback trim. The measure is set by the alphabet,
%% not by the margin: 108mm of Palatino at 10pt is 2.31 lowercase alphabets,
%% inside the 2-3 band and close to its centre. On A4 at 12pt it was 2.71.
%% ...and the 17mm inner margin moves from the left to the RIGHT, which is
%% the spine side of a right-to-left book. geometry sets margins
%% physically and bidi does not swap them, so this has to be said here.
\usepackage[paperwidth=5.5in,paperheight=8.5in,
  width=108mm,top=19mm,bottom=21mm,right=17mm,
  marginparwidth=0pt,headsep=7mm]{geometry}
%% amssymb goes with eulervm -- unicode-math supplies its symbols from the
%% math font, and this book uses exactly three of them: \mathbb, \square
%% and \varnothing. Only \square needs saying, below; it is the one
%% unicode-math gives another name.
\usepackage{amsmath,amsthm}
\usepackage{xcolor}

%% ---------------------------------------------------------------- palette
%% One accent, as a book should have. The four box types no longer tell
%% themselves apart by hue -- they tell themselves apart by *shape*:
%%   functionality  filled body, filled title band, rule    (heaviest)
%%   procedure      no body fill, tinted title band, rule   (lighter)
%%   definition     thin frame, no fill                     (outline)
%%   theorem-like   solid fill, no frame                    (solid)
%% which is the same distinction the accent used to carry, drawn in ink
%% instead of in colour.
\definecolor{accent}{HTML}{C8791A}      % chapter titles, running heads, rules, links
\definecolor{accentsoft}{HTML}{A8630F}  % operation names, secondary links
\definecolor{boxrule}{HTML}{E0A85C}     % box rules and frames
\definecolor{boxtint}{HTML}{F6E2C4}     % functionality title band
\definecolor{boxback}{HTML}{FDF6EC}     % functionality body
\definecolor{boxtitle}{HTML}{7A4A0C}    % box title lettering
\definecolor{proctint}{HTML}{EFEAE2}    % procedure title band (warm neutral)
\definecolor{procback}{HTML}{F8F6F3}    % procedure body (paler neutral still)
\definecolor{procrule}{HTML}{B9A98F}    % procedure rule
\definecolor{thmtint}{HTML}{F7DFB6}     % theorem-like fill, the attached's solid box
\definecolor{commenttint}{HTML}{8A8378} % code comments

\usepackage[most]{tcolorbox}
\usepackage{enumitem}
\setlist{leftmargin=1.6em, topsep=0.5ex, itemsep=0.25ex, parsep=0pt}
\usepackage{array}
\usepackage{booktabs}
\usepackage{microtype}
\usepackage{graphicx}
%% Diagrams were drawn for a 152mm measure and the book gives them 108mm.
%% max width scales only the ones that actually overflow, and leaves the
%% rest at their natural size -- so nothing is shrunk that need not be.
\usepackage[export]{adjustbox}
\usepackage{tikz}
\usetikzlibrary{positioning, fit, calc, decorations.pathreplacing}
%% Fit every diagram to the measure, except where one is sized by hand.
%% \fitdiagramfalse turns the hook off for a single picture; without the
%% switch, a hand-sized diagram gets clamped by the hook first and the two
%% fight -- which is exactly how the parallel-composition chain came to be
%% scaled to nothing and vanish from the page.
\newif\iffitdiagram \fitdiagramtrue
\BeforeBeginEnvironment{tikzpicture}{\iffitdiagram\begin{adjustbox}{max width=\linewidth}\fi}
\AfterEndEnvironment{tikzpicture}{\iffitdiagram\end{adjustbox}\fi}
\usepackage{algpseudocode}
\usepackage{multicol}
\usepackage{titlesec}
\usepackage{etoolbox}
\usepackage{tocloft}
\renewcommand{\cfttoctitlefont}{\Large\bfseries}
\setlength{\cftbeforetoctitleskip}{0pt}
\setlength{\cftaftertoctitleskip}{1.2ex}
%% Chapters are the spine of the contents: set bold, no leader, so the eye
%% runs down the chapter list and the sections stay subordinate.
\renewcommand{\cftchapfont}{\bfseries}
\renewcommand{\cftchappagefont}{\bfseries}
\renewcommand{\cftchapleader}{\cftdotfill{\cftnodots}}
\setlength{\cftbeforechapskip}{1.0ex plus 0.2ex}
\setlength{\cftchapindent}{0pt}
\setlength{\cftchapnumwidth}{2.1em}
\renewcommand{\cftsecfont}{\normalfont}
\renewcommand{\cftsecpagefont}{\normalfont}
\renewcommand{\cftsecleader}{\normalfont\cftdotfill{\cftdotsep}}
\setlength{\cftbeforesecskip}{0.15ex}
\setlength{\cftsecindent}{2.1em}
\setlength{\cftsecnumwidth}{2.6em}
%% the paper runs past page 99, so the page column must fit three digits
%% without the dot leader touching them
\makeatletter
\renewcommand{\@pnumwidth}{2.1em}
\makeatother
%% unicode: without it the Persian pdftitle and every Persian bookmark
%% come out as mojibake in the viewer's tab and outline pane.
\usepackage[colorlinks=true,linkcolor=accent,citecolor=accent,urlcolor=accentsoft,
  unicode=true,pdftitle={UC برای بازی‌بازان},pdfdisplaydoctitle=true]{hyperref}
%% pdftitle/pdfdisplaydoctitle: without these a browser-embedded viewer shows the bare
%% filename in its tab/toolbar instead of the paper's own title.

%% ------------------------------------------------- Persian typesetting
%% xepersian loads bidi, and bidi must be loaded after every package it
%% patches. It ships patches for tcolorbox, titlesec, multicol, tocloft,
%% hyperref, array, float and graphicx -- eight packages this book uses
%% and all eight loaded above -- which is why this line sits here, at the
%% foot of the preamble, and not with the other \usepackage lines.
\usepackage{xepersian}

%% Persian text is Vazirmatn (SIL OFL; fonts/, with OFL.txt beside it),
%% pinned by path rather than by name. Two reasons, both learned the hard
%% way: fontconfig does not index a repository, so a name lookup finds
%% nothing here even when the font is installed system-wide; and a build
%% that depends on what a machine happens to have installed is not a
%% reproducible build. The static faces are used rather than the variable
%% font, XeTeX having no way to instance a weight axis.
%%
%% PersianTeXNoNumbers is the one font feature that is a decision rather
%% than plumbing: it leaves digits in Latin form where xepersian's default
%% mapping would set them as Persian numerals. This book is read beside
%% its own notation, and a definition numbered 4.7 in the contents, cited
%% as 4.7 in a sentence and appearing as 4.7 inside a formula has to be
%% the same 4.7 in all three places. Scale matches Vazirmatn's x-height to
%% Pagella's, which is what stops a Latin term inside a Persian sentence
%% from reading a size too small.
\settextfont[Path=fonts/,Extension=.ttf,
  UprightFont=Vazirmatn-Regular,BoldFont=Vazirmatn-Bold,
  Ligatures=PersianTeXNoNumbers,Scale=0.95]{Vazirmatn}
\setpersiansansfont[Path=fonts/,Extension=.ttf,
  UprightFont=Vazirmatn-Medium,BoldFont=Vazirmatn-Bold,
  Ligatures=PersianTeXNoNumbers,Scale=0.95]{Vazirmatn}
%% Latin runs inside Persian prose -- every piece of notation, every
%% interface name, every citation -- stay Palatino, so the two scripts sit
%% together on the page as they do in the English edition.
\setlatintextfont[BoldFont=texgyrepagella-bold.otf,
  ItalicFont=texgyrepagella-italic.otf,
  BoldItalicFont=texgyrepagella-bolditalic.otf]{texgyrepagella-regular.otf}
\setlatinsansfont[BoldFont=texgyreheros-bold.otf,
  ItalicFont=texgyreheros-italic.otf,Scale=0.90]{texgyreheros-regular.otf}
\setlatinmonofont[BoldFont=texgyrecursor-bold.otf,Scale=0.90]{texgyrecursor-regular.otf}

%% Pseudocode is set left-to-right, as pseudocode is in every Persian
%% technical text: keywords, identifiers and operators are all Latin, and
%% the only Persian in a listing is the trailing comment, which XeTeX
%% reverses by itself inside an LTR block. Reversing the listings instead
%% would put every \State's number on the right and read every assignment
%% backwards, which is not a thing anyone does to code.
\BeforeBeginEnvironment{algorithmic}{\begin{LTR}}
\AfterEndEnvironment{algorithmic}{\end{LTR}}
%% Diagrams need no such wrapper, and one was tried and removed. A
%% tikzpicture is a box laid out in its own absolute coordinates, so an
%% RTL page does not mirror it; the only direction that matters inside is
%% each node's own text, which XeTeX resolves per run. Wrapping the
%% picture in LTR instead fights the fitting hook above for which of the
%% two closes first, and loses.

%% Theorem-like names. xepersian already supplies \chaptername,
%% \contentsname, \bibname and \proofname; these are the ones amsthm asks
%% the document for.
%% Euler, as OpenType. Loaded here rather than beside unicode-math so it
%% is set after xepersian has had its say about families.
\setmathfont{Euler-Math.otf}
%% \square is \unset throughout this book -- the marker for a table entry
%% never written -- and it came from amssymb. unicode-math has the glyph
%% under its own name.
\AtBeginDocument{\ifdefined\mdlgwhtsquare\let\square\mdlgwhtsquare\fi}

\renewcommand{\proofname}{اثبات}

%% ------------------------------------------- running head, and its links
%% The head carries the two affordances the old bottom folio carried, now
%% attached to the two things already in it. The chapter line on the left
%% links to the top of whichever chapter/section is active on the page (the
%% last one that started on it, or the nearest before it); the page number
%% on the right links one heading further back. Both read off one \mark per
%% page carrying *two* hyperref anchors -- current and previous -- the way
%% LaTeX's own \leftmark/\rightmark pair carries two values in one mark.
%% Read with \botmark, not \rightmark: \rightmark is LaTeX's *first*-mark-
%% on-the-page convention, the wrong one when two headings share a page.
%%
%% The mark is written from \refstepcounter (filtered to chapter/section),
%% not from \titleformat's before/after-code. That is not a style choice:
%% titlesec fully expands (\edef's) whatever a heading's <label>/<before-
%% code>/[after-code] contains at essentially first-use time, so a plain
%% macro reference placed there is frozen forever at whatever it held then
%% -- confirmed by direct test, and confirmed immune to the standard fixes
%% (\protect at the call site still froze; making the macro itself
%% \protected broke hyperref's own legitimate expansion of it instead).
%% \@currentHref and \mark survive there only because hyperref already sets
%% \@currentHref *before* titlesec's machinery runs -- verified by checking
%% its value at the moment \refstepcounter fires -- so hooking
%% \refstepcounter directly sidesteps titlesec's cached zone.
%%
%% The head's *text* cannot come from the mark: at \refstepcounter time the
%% title has not been read yet. It comes from \chaptermark, which the class
%% calls just after, into a global. That is sound here only because a
%% chapter always opens a fresh page and the opening page carries no head,
%% so the global is never stale where it is read.
\makeatletter
\newcommand{\nav@prevanchor}{toc}
\newcommand{\nav@headtext}{}
\apptocmd{\refstepcounter}{%
  \ifstrequal{#1}{chapter}{\nav@record}{%
  \ifstrequal{#1}{section}{\nav@record}{}}%
}{}{\ClassError{main}{Could not patch \string\refstepcounter\space for
  navigation}{}}
\newcommand{\nav@record}{%
  \protected@edef\nav@this{\@currentHref}%
  \mark{{\nav@this}{\nav@prevanchor}}%
  \global\let\nav@prevanchor\nav@this}
\newcommand{\nav@current}{\expandafter\@firstoftwo\botmark\@empty\@empty}
\newcommand{\nav@previous}{\expandafter\@secondoftwo\botmark\@empty\@empty}
\newcommand{\nav@setheadtext}[1]{%
  \gdef\nav@headtext{%
    \ifnum\c@secnumdepth>\m@ne
      \if@mainmatter \@chapapp\ \lr{\thechapter.}\ \fi
    \fi #1}}
\renewcommand{\chaptermark}[1]{\nav@setheadtext{#1}}
%% \chapterhead overrides the head for a chapter whose title carries
%% mathematics: maths in a 7.6pt head asks cmss and cmtt for scriptscript
%% at 4.56pt, a size neither has a design for. The contents keep the real
%% title; only the head is spelled out. Place it just after \chapter.
\newcommand{\chapterhead}[1]{\nav@setheadtext{#1}}
%% ...and one for the unnumbered ones. \chapterhead prepends "Chapter N."
%% from the counter, which for a \chapter* would name whichever chapter
%% happened to come last. These set the head text and nothing else.
%% Without them a starred chapter's continuation pages carry a bare page
%% number in the corner and no head at all, which is what the contents
%% and the introduction were doing.
\newcommand{\chapterheadstar}[1]{\gdef\nav@headtext{#1}}
\renewcommand{\ps@headings}{%
  \let\@mkboth\@gobbletwo
  \let\@oddfoot\@empty \let\@evenfoot\@empty
  \def\@oddhead{%
    \hyperlink{\nav@current}{\textcolor{accent}{%
      \fontsize{7.6}{9}\selectfont\persianfont\nav@headtext}}%
    \hfil
    \hyperlink{\nav@previous}{\textcolor{accent}{%
      \fontsize{9.5}{11}\selectfont\persianfont\bfseries\thepage}}}%
  \let\@evenhead\@oddhead}
%% A chapter's own opening page, and the contents, carry the number alone.
\renewcommand{\ps@plain}{%
  \let\@mkboth\@gobbletwo
  \let\@oddhead\@empty \let\@evenhead\@empty
  \def\@oddfoot{\reset@font\hfil
    \hyperlink{\nav@previous}{\textcolor{accent}{%
      \fontsize{9.5}{11}\selectfont\persianfont\bfseries\thepage}}\hfil}%
  \let\@evenfoot\@oddfoot}
\makeatother
%% ---------------------------------------------------------- float packing
%% The code boxes float, and LaTeX's stock parameters were written for a
%% page with far less display material than this one. On A4 they were
%% harmless; at 180mm of text height they strand a box that is merely large
%% -- \topfraction 0.7 forbids it from any page carrying text, and
%% \floatpagefraction 0.5 then gives it a page of its own with nothing
%% else on it. Raising the ceilings lets a box share a page with whatever
%% text will fit beside it, which is what stops the holes.
%% Float separation. The code boxes carry their own 2.4ex of air above and
%% below (panel's before/after skip), and the float machinery then adds
%% \textfloatsep's stock 20pt on top of it -- so every floating box sat in
%% about 32pt of white, two and a half lines, where the box already looked
%% separated at twelve. These are the float side only; the box keeps its own
%% skip, which is what the 31 nofloat boxes rely on.
\setlength{\textfloatsep}{6pt plus 2pt minus 2pt}
\setlength{\intextsep}{6pt plus 2pt minus 2pt}
\setlength{\floatsep}{8pt plus 2pt minus 2pt}

\renewcommand{\topfraction}{0.95}
\renewcommand{\bottomfraction}{0.95}
\renewcommand{\textfraction}{0.06}
\renewcommand{\floatpagefraction}{0.92}   % a float page must be nearly full,
                                          % or the box goes back on a text page
\setcounter{topnumber}{3}
\setcounter{bottomnumber}{2}
\setcounter{totalnumber}{5}

\pagestyle{headings}
\mark{{toc}{toc}}

%% -------------------------------------------------------- screen comfort
\linespread{1.06}
%% a narrower measure leaves TeX less room to break dense inline math; give it a
%% final-pass allowance rather than rewording (that belongs to the prose pass)
\tolerance=1200
\emergencystretch=2em
\setlength{\parskip}{0.15ex}
%% Widows and orphans. Not 10000: an infinite penalty makes TeX pay for the
%% fix in interword spacing instead. These are high enough that a stray line
%% has to be the only alternative. Safe here because article one-sided is
%% \raggedbottom, so a short page costs white space at the foot, not a
%% stretched column.
\clubpenalty=9996
\widowpenalty=9999
\displaywidowpenalty=9999

%% ------------------------------------------------------------- headings
%% Chapter openers carry the accent; section headings do not. One coloured
%% thing per opening is the whole restraint here -- a page of amber
%% subheads would be the decoration this design is trying to avoid.
\titleformat{\chapter}[display]
  {\normalfont}
  {\fontsize{9}{11}\selectfont\persianfont\bfseries\chaptertitlename\ \thechapter}
  {0.6ex}
  %% \bfseries, not \itshape: Persian has no italic tradition and Vazirmatn
  %% no italic face, so an italic Persian title is either a fake slant or a
  %% silent substitution. Bold is what a Persian book emphasises with.
  {\Huge\bfseries\color{accent}}
\titlespacing*{\chapter}{0pt}{6ex}{7ex}
\titleformat{name=\chapter,numberless}[display]
  {\normalfont}{}{0pt}{\Huge\bfseries\color{accent}}
\titlespacing*{name=\chapter,numberless}{0pt}{6ex}{7ex}

\titleformat{\section}{\normalfont\large\bfseries}{\thesection}{0.6em}{}
\titleformat{\subsection}{\normalfont\normalsize\bfseries}{\thesubsection}{0.5em}{}
\titlespacing*{\section}{0pt}{3.0ex plus .6ex minus .2ex}{1.2ex plus .2ex}
\titlespacing*{\subsection}{0pt}{2.2ex plus .5ex minus .2ex}{0.9ex plus .2ex}

%% ------------------------------------------------------- figure captions
%% The five diagrams sit inline in \begin{center}, not in floats: each is
%% read at the point its prose reaches it and must not drift. They still
%% need a counter, because \label on a bare \begin{center} latches onto
%% whichever counter was stepped last -- which is how ``Figure 10'' came to
%% be printed from a section number, linking to a definition thirteen pages
%% from the diagram. \figcaption steps a counter of its own, so \ref
%% resolves to the figure and the hyperlink lands on it.
%% Numbered straight through, not per section: there are six of them in
%% 147 pages, and a reader hunting "Figure 4" should not have to know which
%% section it fell in.
%% The caption sits inside the diagram's \begin{center}, where \centering is
%% in force -- which centres every line of a multi-line caption and leaves the
%% last one stranded. The \parbox restores ordinary justified setting inside a
%% measure narrower than the text block, which is what a caption wants.
\newcounter{diagram}
\newcommand{\figcaption}[2]{%
  \par\smallskip
  \refstepcounter{diagram}\label{#1}%
  \parbox{0.92\textwidth}{\footnotesize
    \textcolor{accent}{شکل \thediagram:}\ #2\par}}

%% ------------------------------------------------ AI-authorship disclosure
%% Every paper written with an AI collaborator carries this block: which
%% models drove the writing, which humans supplied the prompts, and which
%% humans reviewed the result. Copy it, with fresh values, into the
%% preamble of any such paper.
\newcommand{\AIModelsUsed}{\lr{Claude Opus 5 (1M context)}}
\newcommand{\PromptedBy}{\lr{Pooya Farshim}}
\newcommand{\ReviewedBy}{هیچ‌کس}
\newcommand{\PaperDate}{\today}
\newcommand{\AIDisclosure}{%
  \begingroup
  \small\centering
  \begin{tabular}{@{}rl@{}}
    \textbf{هوش مصنوعی:} & \AIModelsUsed \\
    \textbf{پرامپت‌ها:} & \PromptedBy \\
    \textbf{بازبینان:} & \ReviewedBy \\
    \textbf{تاریخ:} & \PaperDate \\
  \end{tabular}
  \par
  \endgroup
}

%% ------------------------------------------ translator's disclosure
%% This edition owes a line the English one does not. The translation is
%% machine-made and has been read by no Persian-speaking cryptographer, so
%% its terminology is a proposal and not settled usage; where a term is
%% load-bearing the English edition is the reference. Kept as a macro
%% beside \AIDisclosure so the two are never separated.
\newcommand{\TranslationDisclosure}{%
  \begingroup
  \footnotesize\centering
  \begin{minipage}{0.86\textwidth}
  \setlength{\parindent}{0pt}
  \textbf{یادداشت مترجم.} این متن ترجمهٔ ماشینی
  \lr{\emph{UC for Gamers}} است و هیچ رمزنگارِ فارسی‌زبانی آن را نخوانده
  است. واژه‌گزینی‌ها پیشنهاد است، نه کاربردِ جاافتاده؛ فهرست آن‌ها در
  \lr{\texttt{GLOSSARY-fa.md}} آمده است. نمادگذاری، نام‌ها، برچسب‌ها،
  کلیدواژه‌های کد و کتاب‌نامه به خطّ لاتین مانده‌اند، چون نمادند و نه
  نوشتار. مرجع، همان نسخهٔ انگلیسی
  \lr{\texttt{main.tex}} است: هر جا چیزی وزن حمل می‌کند، آنجا را ببینید.
  \end{minipage}
  \par
  \endgroup
}

%% ------------------------------------------------------------- notation
\newcommand{\Ps}{\mathbf{P}}
\newcommand{\PID}{\mathit{pid}}
\renewcommand{\emptyset}{\varnothing}
\newcommand{\admits}{\mathbf{N}}
\newcommand{\Apid}{\mathbf{A}}
\newcommand{\Zpid}{\mathbf{Z}}
\newcommand{\Stdpid}{\mathbf{Std}}
\newcommand{\pars}{\mathit{par}}
\newcommand{\uses}{\mathbf{U}}
\newcommand{\serves}{\op{serves}}
\newcommand{\WF}{\op{wf}}
\newcommand{\Cs}{\mathbf{C}}
\newcommand{\Corr}{\mathcal{C}}
\newcommand{\Blk}{\mathcal{B}}
\newcommand{\F}{\mathcal{F}}
\newcommand{\G}{\mathcal{G}}
\newcommand{\Zenv}{\mathcal{Z}}
\newcommand{\Adv}{\mathcal{A}}
\newcommand{\Dum}{\mathcal{D}}
\newcommand{\Sim}{\mathcal{S}}
\newcommand{\bd}[1]{\mathbf{#1}}
\newcommand{\advtg}{\op{Adv}}
%% property advantage, indexed by finalization:
%%   \padv{Fin}{\Gm,\sigma,\Zenv,\mathcal{X}}  --- drop the last for the dummy
\newcommand{\padv}[2]{\advtg^{#1}_{#2}}
\newcommand{\Aset}{\mathbb{A}}
\newcommand{\SimSet}{\mathbb{S}}
\newcommand{\ZenvSet}{\mathbb{Z}}
%% \textnormal keeps interface names upright inside italic theorem statements:
%% Palatino has no small-caps italic, and LaTeX would substitute upright anyway.
\newcommand{\op}[1]{\textnormal{\textsc{#1}}}
%% clickable operation names: \opdef marks a definition, \opl links to it
\newcommand{\opdef}[1]{\hypertarget{op:#1}{\op{#1}}}
\newcommand{\opl}[1]{\hyperlink{op:#1}{\textcolor{accentsoft}{\op{#1}}}}
\newcommand{\fopdef}[1]{\hypertarget{op:#1}{\op{Full}\op{#1}}}
\newcommand{\fopl}[1]{\hyperlink{op:#1}{\textcolor{accentsoft}{\op{Full}\op{#1}}}}
\newcommand{\IDs}{\op{IDs}}
\newcommand{\Ids}{\mathcal{I}}
\newcommand{\id}{id}
%% value atoms -- ok, rej, done, true, false -- share \op's small caps but
%% not its color, so a bare name and a bare value no longer look alike:
%% \atomtint reuses the comment color, one quiet tone for everything that
%% is not the operative code.
\newcommand{\atom}[1]{\textcolor{commenttint}{\op{#1}}}
\newcommand{\rej}{\atom{rej}}
\newcommand{\ok}{\atom{ok}}
\newcommand{\true}{\atom{true}}
\newcommand{\false}{\atom{false}}
\newcommand{\parent}{\mathord{\uparrow}}
\newcommand{\Gm}{\mathcal{G}m}
\newcommand{\Xset}{\mathbb{X}}
\newcommand{\Win}{\op{Win}}
\newcommand{\Bad}{\op{Bad}}
\newcommand{\done}{\atom{done}}
\newcommand{\Frand}{\mathcal{F}_{\mathsf{Rand}}}
\newcommand{\Fstore}{\mathcal{F}_{\mathsf{Store}}}
\newcommand{\Fio}{\mathcal{F}_{\mathsf{IO}}}
\newcommand{\Fdiffuse}{\mathcal{F}_{\mathsf{Diffuse}}}
\newcommand{\Fsig}{\mathcal{F}_{\mathsf{Sig}}}
\newcommand{\Gclock}{\mathcal{G}_{\mathsf{Clock}}}
\newcommand{\Gpki}{\mathcal{G}_{\mathsf{PKI}}}
\newcommand{\Fnet}{\mathcal{F}_{\mathsf{Net}}}
\newcommand{\Fac}{\mathcal{F}_{\mathsf{AC}}}
\newcommand{\Facz}{\mathcal{F}^{\circ}_{\mathsf{AC}}}
\newcommand{\Fpki}{\mathcal{F}_{\mathsf{PKI}}}
\newcommand{\DS}{\mathsf{DS}}
\newcommand{\psig}{\pi_{\mathsf{Sig}}}
\newcommand{\vk}{\mathit{vk}}
\newcommand{\msg}{\mathit{msg}}
\newcommand{\unset}{\square}
\newcommand{\Keys}{\mathcal{K}}
\newcommand{\Msgs}{\mathcal{M}}
\newcommand{\Sigs}{\Sigma}
\newcommand{\none}{\bot}
\newcommand{\San}{\opl{San}}
\newcommand{\Clean}{\op{Clean}}
\newcommand{\V}[1]{\mathtt{#1}}
\newcommand{\bits}{\{0,1\}^*}
\newcommand{\Op}{\op{Op}}
\newcommand{\ncl}[1]{\langle #1 \rangle}

%% ------------------------------------------------- algorithmic settings
\algrenewcommand\algorithmicrequire{\textbf{require}}
\algnewcommand\Req[1]{\State \textbf{require}\ #1}
\algrenewcommand{\algorithmiccomment}[1]{\hfill{\color{commenttint}\slshape /\!\!/ #1}}
\algrenewcommand\algorithmicindent{1.1em}
%% invisible block that just indents its body
\algdef{SE}[INDENTBLOCK]{Indent}{EndIndent}{}{}
\algtext*{Indent}
\algtext*{EndIndent}
%% ``on <event>'' block
\algdef{SE}[ONBLOCK]{On}{EndOn}[1]{\textbf{on}\ #1}{}
%% block ends are carried by indentation alone
\algtext*{EndOn}
\algtext*{EndIf}
%% carry algorithmic line numbers from one block to the next
\newcounter{contline}
\newcommand{\algcont}{\setcounter{ALG@line}{\value{contline}}}
\newcommand{\algsave}{\setcounter{contline}{\value{ALG@line}}}
%% column widths for aligned code
\newlength{\algretwd}\settowidth{\algretwd}{\textbf{return}}
\newcommand{\algand}{\makebox[\algretwd][r]{$\wedge$\,}}
\newlength{\clausewd}
\newlength{\reswd}

%% ----------------------------------------------------------- box panels
%% Four displays, one hue, told apart by shape. Geometry is shared so they
%% sit on the page as one family; weight is what separates them.
\tcbset{panel/.style={%
  enhanced,
  boxrule=0pt, leftrule=0pt,
  sharp corners, arc=0pt, outer arc=0pt,
  left=4pt, right=4pt, top=5pt, bottom=5pt,
  before skip=1.8ex, after skip=1.8ex}}
%% Kept for anything that genuinely cannot be brought back to the measure.
%% Nothing uses it at present: statements align with the prose because they
%% hold prose, and the code boxes buy their width back in font size instead.
\tcbset{widepanel/.style={panel,
  grow to left by=9mm, grow to right by=9mm}}

%% the optional argument takes extra tcolorbox keys for one box -- e.g.
%% \begin{interface}[nofloat]{...} pins a box in place instead of floating it
%% functionality: the heaviest of the four -- filled body, filled title
%% band, and now a bounding box all round, at the same 0.9pt as defbox
%% below, so the listing reads as one enclosed unit rather than a title
%% band floating over open text.
\newtcolorbox{interface}[2][]{%
  panel, unbreakable, float=htbp, every float=\centering,
  fontupper=\footnotesize,
  colback=boxback, colframe=boxrule, colbacktitle=boxtint, coltitle=boxtitle,
  boxrule=0.9pt, leftrule=0.9pt, titlerule=0pt, fonttitle=\small\bfseries,
  title={#2}, #1}

%% procedure: same geometry, one step lighter -- no body fill, a neutral
%% title band, the same bounding box as interface. Predicates, wrappers
%% and everything that is not a functionality.
\newtcolorbox{procedure}[1]{%
  panel, unbreakable, float=htbp, every float=\centering,
  fontupper=\footnotesize,
  colback=procback, colframe=procrule, colbacktitle=proctint, coltitle=black!62,
  boxrule=0.9pt, leftrule=0.9pt, titlerule=0pt, fonttitle=\small\bfseries,
  title={#1}}

\newcommand{\params}[1]{{\footnotesize\normalfont\color{black!72} #1}}
\newcommand{\opsig}[1]{\textcolor{accent}{\underline{\smash{#1}}}}
\newcommand{\heading}[1]{\medskip\par\noindent\textcolor{accent}{\textbf{#1.}}\ }

\newtheoremstyle{accentdef}{1.8ex plus .4ex}{1.8ex plus .4ex}%
  {\normalfont}{0pt}{\bfseries}{.}{0.5em}{}
%% Body font upright where the English edition sets italic: the four box
%% styles already tell a claim from a definition by fill and frame, which
%% is the distinction the italic was carrying.
\newtheoremstyle{accentplain}{1.8ex plus .4ex}{1.8ex plus .4ex}%
  {\normalfont}{0pt}{\bfseries}{.}{0.5em}{}

\theoremstyle{accentdef}
\newtheorem{definition}{تعریف}[chapter]
\newtheorem{remark}[definition]{نکته}
\theoremstyle{accentplain}
\newtheorem{proposition}[definition]{گزاره}
\newtheorem{theorem}[definition]{قضیه}
\newtheorem{corollary}[definition]{نتیجه}
\newtheorem{lemma}[definition]{لم}

%% Composite numbers stay left-to-right as a unit. "4.17" is two digit runs
%% with a neutral dot between them, so an RTL paragraph reorders it to
%% "17.4" -- which is what Definition 4.17 printed in its own box, in the
%% contents, and at every \ref pointing to it. \lr binds each number into
%% one run. Built from the raw counters rather than from each other, so
%% that \thesubsection does not nest one \lr inside another.
\renewcommand{\thesection}{\lr{\thechapter.\arabic{section}}}
\renewcommand{\thesubsection}{\lr{\thechapter.\arabic{section}.\arabic{subsection}}}
\renewcommand{\thedefinition}{\lr{\thechapter.\arabic{definition}}}

%% Definitions are framed and unfilled; results are filled and unframed.
%% That is the distinction the attached book draws, and it is the right one
%% -- a definition is a boundary, a theorem is a claim. Proofs stay outside
%% the box in both cases.
\tcbset{defbox/.style={panel, breakable,
  colback=white, colframe=boxrule, boxrule=0.9pt, leftrule=0.9pt}}
\tcbset{thmbox/.style={panel, breakable,
  colback=thmtint, colframe=thmtint, boxrule=0pt, leftrule=0pt}}
\tcolorboxenvironment{definition}{defbox}
\tcolorboxenvironment{remark}{defbox}
\tcolorboxenvironment{theorem}{thmbox}
\tcolorboxenvironment{proposition}{thmbox}
\tcolorboxenvironment{corollary}{thmbox}
\tcolorboxenvironment{lemma}{thmbox}

%% ------------------------------------------------------------- abstract
%% book has no abstract environment. In a book the abstract is a front-
%% matter block rather than a floating summary, so it is set narrower than
%% the measure and unheaded -- it faces the title, and needs no label.
%% Vertically centred on its own page, and carrying the folio at the foot
%% like every other page that opens something -- \pagestyle{headings} would
%% otherwise put a bare page number in the top corner here, the running head
%% having no chapter to name yet.
\newenvironment{abstract}
  {\thispagestyle{plain}\null\vfill
   \begin{center}%
   \begin{minipage}{0.90\textwidth}\small\setlength{\parindent}{0pt}}
  {\par\end{minipage}\end{center}\vfill\null}

%% ------------------------------------------- the composition square
%% Drawn once, used twice: on the cover and as Figure 1. The cover is
%% the reason it is a macro -- duplicating sixty lines of TikZ would
%% guarantee the two drift apart.
\newcommand{\compositionsquare}{%
\begin{tikzpicture}[
  sysnode/.style={draw=accent, rounded corners=1pt, inner sep=2pt, font=\scriptsize, minimum height=4.5mm},
  wire/.style={draw=black!70},
  shell/.style={draw=accent, rounded corners=2pt, inner sep=2.5mm},
  absorbed/.style={draw=black!55, dashed, rounded corners=2pt, inner sep=2.5mm},
  role/.style={font=\small},
  rel/.style={font=\normalsize},
]

%% ===== Row 1: real world (sigma = pi, X = Adv) =====
\node[sysnode] (A-rest) {$\rho\setminus\varphi$};
\node[sysnode, right=3mm of A-rest] (A-pi) {$\pi$};
\draw[wire] (A-rest) -- (A-pi);
\node[shell, fit=(A-rest)(A-pi)] (A-shell) {};
\node[role, above=6mm of A-shell] (A-Z) {$\Zenv$};
\draw[wire] (A-Z) -- (A-shell.north);
\node[role, below=6mm of A-shell] (A-Adv) {$\Adv$};
\draw[wire] (A-shell.south) -- (A-Adv);

\node[rel, right=11mm of A-shell] (eq1) {$=$};

\node[sysnode, right=11mm of eq1] (Ap-rest) {$\rho\setminus\varphi$};
\node[role, above=2mm of Ap-rest] (Ap-Z) {$\Zenv$};
\draw[wire] (Ap-Z) -- (Ap-rest);
\node[absorbed, fit=(Ap-Z)(Ap-rest)] (Ap-env) {};
\node[sysnode, below=6mm of Ap-env] (Ap-pi) {$\pi$};
\draw[wire] (Ap-env.south) -- (Ap-pi);
\node[role, below=6mm of Ap-pi] (Ap-Adv) {$\Adv$};
\draw[wire] (Ap-pi) -- (Ap-Adv);

%% ===== Row 2: ideal world (sigma = varphi, X = Sim); rho[varphi/varphi] = rho =====
\node[sysnode, below=40mm of A-rest] (B-rest) {$\rho\setminus\varphi$};
\node[sysnode, right=3mm of B-rest] (B-phi) {$\varphi$};
\draw[wire] (B-rest) -- (B-phi);
\node[shell, fit=(B-rest)(B-phi)] (B-shell) {};
\node[role, above=6mm of B-shell] (B-Z) {$\Zenv$};
\draw[wire] (B-Z) -- (B-shell.north);
\node[role, below=6mm of B-shell] (B-Sim) {$\Sim$};
\draw[wire] (B-shell.south) -- (B-Sim);

\node[rel] at (eq1 |- B-shell) (eq2) {$=$};

\node[sysnode] at (Ap-rest |- B-rest) (Bp-rest) {$\rho\setminus\varphi$};
\node[role, above=2mm of Bp-rest] (Bp-Z) {$\Zenv$};
\draw[wire] (Bp-Z) -- (Bp-rest);
\node[absorbed, fit=(Bp-Z)(Bp-rest)] (Bp-env) {};
\node[sysnode, below=6mm of Bp-env] (Bp-phi) {$\varphi$};
\draw[wire] (Bp-env.south) -- (Bp-phi);
\node[role, below=6mm of Bp-phi] (Bp-Sim) {$\Sim$};
\draw[wire] (Bp-phi) -- (Bp-Sim);

%% ===== vertical relations, routed clear of the cells' own wiring =====
\draw[<->]
  (A-shell.west) .. controls +(-16mm,-6mm) and +(-16mm,6mm) .. (B-shell.west)
  node[midway, left=2mm, rel, align=center] {$\approx_{\varepsilon}$\\[1pt]{\tiny\color{accentsoft}(conclusion)}};

\draw[<->]
  (Ap-env.east) .. controls +(18mm,-10mm) and +(18mm,10mm) .. (Bp-env.east)
  node[midway, right=2mm, rel, align=center] {$\approx_{\varepsilon}$\\[1pt]{\tiny\color{accentsoft}(hypothesis)}};

\end{tikzpicture}
}
\title{\textbf{UC برای بازی‌بازان}}
\author{}
\date{}

%% No open notes remain in the text. The three note macros above are kept
%% only so that new ones can be dropped in without touching the preamble.

\begin{document}

%% ---------------------------------------------------------------- cover
%% The composition square is the book's argument in one picture -- three
%% of the four composition results are corollaries of the absorption it
%% draws -- so it carries the cover rather than an ornament. \thispagestyle
%% is empty: a cover wants no folio and no running head.
\begin{titlepage}
\thispagestyle{empty}
\centering
%% One block, centred on the page. Equal stretch above and below, fixed
%% gaps within: the three elements are one composition, so the air belongs
%% outside them rather than between them.
\vspace*{\stretch{1}}

{\fontsize{34}{40}\selectfont\bfseries\color{accent} UC برای بازی‌بازان\par}

\vspace{1.6ex}

{\fontsize{13}{16}\selectfont\color{accentsoft}\lr{\itshape UC for Gamers}\par}

\vspace{9ex}

\adjustbox{max width=0.86\textwidth, max totalheight=0.40\textheight}{%
  \compositionsquare}

\vspace{9ex}

\AIDisclosure

\vspace*{\stretch{1}}
\end{titlepage}


\begin{abstract}
ترکیب‌پذیری جهانی را، هر جا که نوشته می‌شود، به آمیزه‌ای از نوشتار و کد
می‌نویسند؛ و همین است که پژوهشِ UC را سخت‌ابطال می‌کند و آزمودن یا
وارسیِ صوری‌اش را سخت‌تر. ما، به سبکِ کدمحورِ امنیتِ ویژگی‌محور، زبانی
دقیق برای برنامه‌نویسیِ کارکردهای UC می‌نشانیم: هویت و کنترلِ دسترسی
داده‌اند، پوشانه‌ها آن‌ها را اعمال می‌کنند، فساد دفتری است که هر واسط
می‌خواندش، و یک اجرا هندلری است که تنها یک ژتون را دست‌به‌دست می‌کند.
آن‌گاه تقلید را روی رده‌های صریحِ ماشین و با کران‌های عینی تعریف می‌کنیم
--- بی‌پارامترِ امنیتی، بی‌مجانبی --- و جعبه‌ابزارِ ترکیب از پیِ آن
می‌آید: جانشانی، ترکیبِ موازی، و قضیهٔ تمامیت برای حریفِ بدل --- هر سه
نتیجهٔ یک لمِ جذب، لمی که پروتکل یا حریف را به درونِ محیط می‌بَرد و
بازگروه‌بندیِ دقیقِ یک اجرای واحد است، نه مقایسهٔ دو اجرا --- و ترایایی
مفت، از خودِ متر. ویژگی‌های بازی‌محور نیز، چنان که خواهیم دید، خودشان
کارکردند و در امتدادِ تقلید منتقل می‌شوند. یک کارکردِ امضا، یک ساعتِ
جهانی و یک شبکهٔ با تأخیرِ کراندار سبک را در کار نشان می‌دهند. امیدمان
این است که این کار به به‌کارگیریِ روش‌های صوری در UC کمک کند و سرانجام
به اثباتی صوری از UC برسد.
\end{abstract}

%% Contents before the introduction: a book hands the reader the map, then
%% the argument for reading it. (In the article it ran the other way round,
%% because an abstract and a contents page were the whole of the front
%% matter and the introduction was the first real thing.)
\clearpage
\chapterheadstar{فهرست مطالب}
\hypertarget{toc}{}%
{\setlength{\parskip}{0pt}\tableofcontents}

%% ==================================================================
%% Introduction. Unnumbered on purpose: the body's sections carry the
%% definition numbering (\newtheorem{definition}{Definition}[chapter]),
%% so numbering this one would renumber every statement in the paper and
%% every reference to one in the audit record. \phantomsection before
%% \addcontentsline is what makes the contents entry link here rather
%% than to the previous anchor.
%% ==================================================================
\phantomsection
\addcontentsline{toc}{chapter}{مقدمه}
\chapter*{مقدمه}
\chapterheadstar{مقدمه}

ترکیب‌پذیری جهانی را، تقریباً هر جا که نوشته می‌شود، به آمیزه‌ای از
نوشتار و کد می‌نویسند. واسط‌های یک کارکرد کدند؛ اینکه چه کسی می‌تواند
آن‌ها را فرا بخواند، یک جمله است. مدلِ فساد یک بند است؛ و اینکه واسطِ یک
طرفِ فاسد در واقع چه برمی‌گرداند، به خواننده واگذار می‌شود. هر یک از این
شکاف‌ها کوچک است، و داورِ انسانی بی‌آنکه بفهمد پُرش می‌کند. ماشین نمی‌تواند،
و همین است که اثباتِ UC را سخت‌ابطال می‌کند و وارسیِ صوری‌اش را سخت‌تر:
بسیاری از آنچه باید وارسی شود، هرگز نوشته نشده بود.

این مقاله آن بخش‌ها را به‌صورتِ کد می‌نویسد. هویت و کنترلِ دسترسی
داده‌هایی می‌شوند که کارکرد در میدان‌های نام‌دار حمل می‌کند. پوشانه‌ها
اعمالشان می‌کنند: نگهبان تصمیم می‌گیرد چه فراخوان‌هایی به هسته می‌رسند،
خموش‌کننده تصمیم می‌گیرد هسته چه ادعایی می‌تواند بکند، و واسطه‌گر
فراخوانِ یک طرفِ فاسد را به حریف می‌سپارد. فساد دفتری است که هر واسط
می‌خواندش، دو بار، در نقاطی که خودِ کد نامشان می‌برد. یک اجرا هندلری است
که تنها یک ژتون را دست‌به‌دست می‌کند. هیچ‌چیزِ اینجا مدلی تازه نیست ---
همان مدلِ کانتی~\cite{Canetti20} است، با چند قراردادش که صریح شده‌اند ---
اما بیانِ آن‌ها به کد و نه به نوشتار، همان چیزی است که می‌گذارد باقیِ
مقاله اثبات شود و نه استدلال.

\heading{یک لم} آن‌گاه تقلید روی رده‌های صریحِ ماشین تعریف می‌شود، با
کران‌های عینی: بی‌پارامترِ امنیتی، بی‌مجانبی، کرانی که یک عدد است. آنچه
از پی می‌آید جعبه‌ابزارِ ترکیب است، و شکلِ آن ادعای اصلیِ ساختاریِ این
مقاله است. جانشانی، ترکیبِ موازی و تمامیتِ حریفِ بدل، هر سه، نتیجهٔ یک لمِ
\emph{جذب}‌اند، لمی که پروتکل یا حریف را به درونِ محیط می‌بَرد. جذب مقایسهٔ
دو اجرا نیست، بازگروه‌بندیِ دقیقِ یک اجراست: همان اجرا، با کمانک‌گذاریِ
دیگر، پس مزیتِ دو سو یک عدد است و نه دو عدد با کرانی مشترک. ترایایی
استثناست و از گونهٔ دیگری است، چون بر نامساویِ مثلثی روی متر تکیه دارد و
نه بر هیچ بازگروه‌بندی؛ به همین سبب همراهِ متر بیان می‌شود، پیش از
ماشین‌آلاتی که آن سه دیگر لازم دارند.

\heading{بازی‌ها کارکردند} یک ویژگی به سبکِ بازی‌محورِ بلر و
راگوی~\cite{BR06} --- جعل‌ناپذیری، زنده‌بودن، اصالت --- چنان که پیدا
می‌شود، خودش هم به‌صورتِ کارکرد بیان‌شدنی است، با مبارزطلبِ بازی به‌عنوانِ
ماشینی معمولی و حکمش به‌عنوانِ واسطی معمولی. ویژگی‌هایی که این‌گونه نوشته
شوند در امتدادِ تقلید منتقل می‌شوند: اگر پروتکلی شیئی ایده‌آل را تقلید کند،
ویژگی‌های آن شیء به آن فرود می‌آیند، زیرِ فرض‌هایی که مقاله بیانشان می‌کند
و در هر کاربرد وارسی‌شان می‌کند و مفروضشان نمی‌گیرد.

\heading{نخست چه بخوانیم} دستگاه از آنچه شمارِ قضیه‌ها نشان می‌دهد
بزرگ‌تر است، و همهٔ آن باربر نیست. پنج تعریف بیشترِ بارِ مقاله را حمل
می‌کنند: یک \emph{کارکرد} و پنج میدانش؛ یک \emph{اجرا} و اینکه اشغالِ یک
هویت چه معنایی دارد؛ یک \emph{دسته}، که چند ماشین را چون یک اشغال‌گر
می‌راند و جذب از آن ساخته شده است؛ خودِ \emph{جذب}؛ و \emph{تناظرِ
بازگروه‌بندی} که دو اجرا را ماشین‌به‌ماشین جفت می‌کند. خواننده‌ای که این
پنج را داشته باشد تیرکِ مقاله را دارد، و نتایجِ ترکیب به‌ترتیب از آن‌ها
می‌آیند. کارکردهای عینی --- یک طرحِ امضا، یک ساعتِ جهانی و یک شبکهٔ با
تأخیرِ کراندار به‌پیرویِ کاتز، مائورر، تاکمان و زیکاس~\cite{KMTZ13} و
بادرچر، مائورر، چودی و زیکاس~\cite{BMTZ17}، و یک کانالِ اصالت‌دار که روی
آن‌ها محقق می‌شود --- در پایان می‌آیند و هستند تا سبک را در کار نشان
دهند، نه چون چیزی پیش از آن‌ها به آن‌ها نیاز داشته باشد.

\heading{این به چه کار می‌آید} پژوهشگرِ UC باید قراردادهای اینجا را همان
قراردادهایی بیابد که خودش به کار می‌برد، اما نوشته‌شده. مخاطبِ دوم آن کسی
است که این دقتِ افزوده برای اوست. هر کس که ابزاری روی UC می‌سازد --- یک
نوع‌سنج برای کدِ کارکرد، یک رمزگذاری در دستیارِ اثبات، یک داربستِ آزمون ---
مشخصه‌ای لازم دارد که پرسش‌هایی را حل کند که ارائهٔ غیررسمی می‌تواند به
خوانندهٔ خود واگذارد، و حلِ همان‌ها بیشترِ کاری است که این مقاله می‌کند:
نگهبان چه فراخوانی را می‌پذیرد، هستهٔ خموش‌شده چه ادعایی می‌تواند بکند،
دفترِ فساد کِی خوانده می‌شود. اینکه این اندازه برای حملِ اثباتی
ماشین‌وارسی‌شده از UC بس باشد یا نه، اینجا برقرار نشده است، و ما ادعایش
نمی‌کنیم. همین است که چرا این چارچوب این‌گونه نوشته شده است.

%% ==================================================================
\chapter{کارکردها و سیستم‌ها}\label{sec:functionalities}

ما در چارچوبِ ترکیب‌پذیریِ جهانیِ کانتی~\cite{Canetti20} کار می‌کنیم، که
چند قراردادش را در زیر صریح می‌کنیم.

شناسهٔ فرایند نامِ یک نمونهٔ در حالِ اجرای پروتکل است و سه بخش دارد:
\[
  \PID \;=\; (F,\, s,\, i) ,
\]
یک \emph{نامِ کارکرد} $F$، یک \emph{شناسهٔ نشست} $s$، و یک \emph{شمارهٔ
نمونه} $i$. نام می‌گوید چه کدی اجرا می‌شود، نشست می‌گوید نمونه به کدام
اجرای پروتکل تعلق دارد، و شماره چند نمونه از یک نام را در یک نشست از هم
جدا می‌کند. برای سه بخشِ $\id.\PID$ می‌نویسیم $\id.F$، $\id.s$ و $\id.i$،
و به همین شکل $\F.F$، $\F.s$، $\F.i$ برای بخش‌های $\F.\PID$؛ و $\id.\PID
= \F.\PID$ به‌تنهایی یعنی هر سه یکی‌اند.

نام‌ها عضوِ $\bits \cup \{A,Z,C\}$ و شناسه‌های طرف عضوِ $\bits \cup
\{A,Z\}$اند، که در آن $A$، $Z$ و $C$ برای حریف، محیط و دفترِ فساد
رزرو شده‌اند. چون رزروند، هیچ‌یک از این سه عضوِ $\bits$ نیست، پس میدانی که
اعلام شده $\bits$ را در بر دارد هیچ‌یک از آن‌ها را در بر ندارد؛ و $C$ نامِ
هیچ طرفی نیست، چون دفتر ماشینی است که طرف‌ها فرایش می‌خوانند، نه ماشینی
که طرفی آن را براند. هر سه به‌صورتِ یک نمونهٔ واحد در یک نشستِ واحد اجرا
می‌شوند، که ثابتاً نشستِ $0$ و نمونهٔ $0$ است، پس شناسه‌های فرایندشان
$(A,0,0)$، $(Z,0,0)$ و $(C,0,0)$ است. ما در تمامِ کتاب این‌ها را به همان
نامِ خالیِ $A$، $Z$ و $C$ کوتاه می‌کنیم، چون آن دو مؤلفهٔ دیگر خبری حمل
نمی‌کنند.

آزمون‌ها تقریباً همیشه فقط روی نام‌اند. هر جا می‌نویسیم $\id.F = A$
مقصودمان مؤلفهٔ نام است، هرگز کلِ سه‌گانه؛ و همین برای $Z$، برای $C$، و
برای هر کارکردی که در آزمونِ عضویت نامش می‌آید.

دو تا از میدان‌های زیر نام‌هایی از این دست را جمع می‌کنند، و چیزهای
متفاوتی جمع می‌کنند. مجموعهٔ طرف‌ها \emph{شناسه‌های طرف} را در بر دارد،
عضوهای $\bits \cup \{A,Z\}$؛ و مجموعهٔ فراخواننده‌ها کلِ \emph{شناسه‌های
فرایند} را، چون اینکه چه کسی می‌تواند فرا بخواند پرسشی است در بابِ نمونه‌ها
و نه فقط در بابِ کد. پس مجموعهٔ فراخواننده‌ها را با این کوتاه‌نویسی‌ها
می‌نویسیم:
\[
  \begin{gathered}
  \Apid := \{\PID : \PID.F = A\}, \qquad
  \Zpid := \{\PID : \PID.F = Z\}, \\[3pt]
  \Stdpid := \{\PID : \PID.F \in \bits\}
  \end{gathered}
\]
برای شناسه‌های فرایندِ حریف، محیط، و هر چیزِ استاندارد. هر جا نمادگذاریِ
قدیمی‌تر ``$A$ و $Z$'' را می‌پذیرفت، اکنون $\Apid \cup \Zpid$ را می‌پذیرد،
که از آن تنها $(A,0,0)$ و $(Z,0,0)$ هرگز رخ می‌دهند و اجرا از هر یک یکی
فراهم می‌کند؛ و هر جا ``$\bits$'' را می‌پذیرفت، اکنون $\Stdpid$ را
می‌پذیرد، و آنجا هر نشست و هر نمونه‌ای واقعاً رخ می‌دهد. مجموعه‌های طرف
همان $\{A,Z\}$ِ خالی را نگه می‌دارند، چون آن‌ها شناسهٔ طرف‌اند و نه شناسهٔ
فرایند، پس اعلامی که پیش‌ازاین $\Ps := \admits$ خوانده می‌شد باید دو
شقّه شود.

کارکرد هم به همین شکل نام می‌گیرد، با مؤلفهٔ نامش که به‌صورتِ زیرنویس
نوشته می‌شود: $\Fsig$ کارکردی است با $\Fsig.F = \op{Sig}$، که نشست و
نمونه‌اش ناگفته مانده‌اند، پس این نمادگذاری به‌جای شناسهٔ فرایندِ
$(\op{Sig},s,i)$ می‌ایستد با $s$ و $i$ی نامعین. هر جا نمونهٔ خاصی مقصود
باشد آن‌ها را جداگانه می‌دهیم؛ هر جا فقط کد مهم باشد، هیچ نوشته
نمی‌شوند.

با ناگفته گذاشتنشان چیزِ کمی از دست می‌رود، چون کدِ خودِ کارکرد به‌سختی
هرگز آن‌ها را می‌خواند. محلی یا جهانی، هسته برای طرفی که به آن خدمت
می‌دهد و حالتی که نگه می‌دارد نوشته می‌شود؛ اینکه به کدام نشست تعلق دارد،
و کدام نمونه از نامِ خودش است، کارِ اجراست و نه کارِ او. در این مقاله $s$
تنها در سه جا رخ می‌دهد --- سطرِ~\ref{ln:session} از $\opl{Guard}$،
پوشانه‌هایی که آن را فرا می‌خوانند، و آزمونِ ایستای تعریفِ~\ref{def:wfsys}
--- و $i$ هرگز خوانده نمی‌شود. شمارهٔ نمونه جای خود را از راهی دیگر
به‌دست می‌آورد: همان است که می‌گذارد دو نمونه از یک نام در یک نشست، آن
شناسه‌های فرایندِ متمایزی را حمل کنند که تعریفِ~\ref{def:sys} می‌خواهد.
فصلِ~\ref{sec:reloc} این مشاهده را به یک حکم بدل می‌کند: تغییرنامِ نشست‌ها
و نمونه‌ها یک یکریختیِ اجراهاست، و تقلیدی که در یک شناسهٔ فرایند اثبات
شود در هر شناسهٔ دیگری برقرار است، زیرِ یک شرطِ بهداشتی در بابِ اینکه چه
کسی را می‌توان از میانِ نشست‌ها فرا خواند، شرطی که هر کارکردِ اینجا
برآورده‌اش می‌کند.

\begingroup\tcbset{statementbox/.append style={breakable=false}}
\begin{definition}[هویت‌ها و کارکردها]\label{def:func}
یک \emph{هویت} زوجِ $\id = (\PID,P)$ از یک شناسهٔ فرایند و یک شناسهٔ طرف
است، و یک \emph{مجموعهٔ هویت} هر مجموعهٔ $\Ids$ از این‌هاست --- شناسهٔ
فرایند سه‌گانه است، و دیگر مجموعه‌ای از زوج‌های نام نیست. یک
\emph{کارکردِ} $\F$ این پنج میدان را حمل می‌کند
\begin{center}
\footnotesize
\begin{tabular}{@{}ll l@{}}
\toprule
$\F.\PID$   & یک شناسهٔ فرایند                       & \\
$\F.\Ps$    & مجموعهٔ طرف‌های \emph{خدمت‌گیرنده}       & واسط‌هایی که وجود دارند \\
$\F.\admits$  & شناسه‌های فرایندِ فراخواننده‌ای که می‌پذیرد & چه کسی می‌تواند آن واسط‌ها را فرا بخواند \\
$\F.\uses$  & کارکردهایی که \emph{به کار می‌برد}      & چه چیزی باید باشد تا او بتواند بدود \\
$\F.\pars$  & هر پارامترِ دیگر                       & کد جز موارد بالا چه می‌خواند \\
\bottomrule
\end{tabular}
\end{center}
و این مجموعهٔ هویت را \emph{اشغال می‌کند}
\[
  \IDs(\F) \;:=\; \bigl\{\, (\F.\PID,\, P) \;:\; P \in \F.\Ps \,\bigr\} .
\]
\end{definition}
\endgroup

مجموعه‌های هویت دو چیز می‌گویند: اینکه گردایه‌ای از کارکردها چه هویت‌هایی
را اشغال می‌کند، و اینکه یک ماشین چه هویت‌هایی را می‌تواند ادعا کند. این
پنج میدان در سطرِ پارامترِ هر جعبهٔ زیر تکرار می‌شوند، بی نامِ صاحبشان،
چون سرِ جعبه آن را می‌دهد. تنها دو تای اول $\IDs(\F)$ را تعیین می‌کنند:
شناسهٔ فرایند بینِ همهٔ هویت‌هایش مشترک است و طرف‌های خدمت‌گیرنده مؤلفهٔ
دوم را می‌دهند.

میدانِ چهارم وابستگی را ثبت می‌کند. هر درایهٔ $\F.\uses$ زوجی است
$(\F',R)$: $\F'$ نامِ کارکردی است که باید حاضر باشد تا $\F$ بدود، و $R$
محمولی است که می‌گوید $\F$ از هر چه آنجا یافت شود چه می‌خواهد. نام‌بردن از
یک زیرروال به‌تنهایی به‌ندرت بس است --- زیرروالی که به هیچ‌یک از طرف‌های
فراخواننده‌اش خدمت نمی‌داد، یا هیچ‌یک از فراخواننده‌هایش را نمی‌پذیرفت،
بی‌مصرف بود --- پس می‌نویسیم
\[
  \serves(\F,\G) \;:\iff\; \G.\Ps \supseteq \F.\Ps \ \wedge\ \F.\PID \in \G.\admits
\]
برای آن خواسته‌ای که بیش از همه پیش می‌آید: $\G$ به هر طرفی که $\F$ خدمت
می‌دهد خدمت می‌دهد، و $\F$ را در میانِ فراخواننده‌هایش می‌پذیرد. بی عطفِ
دوم، $\opl{Guard}$ هر فراخوان از $\F$ به $\G$ را رد می‌کرد.

آنچه $\uses$ ثبت می‌کند همان است که \emph{کدِ} $\F$ نامش را می‌برد.
فراخوان‌هایی که پوشانه‌اش به‌جای او می‌گذارد --- به $\Corr$ در
سطرهای~\ref{ln:stat} و~\ref{ln:stattwo} از بخشِ~\ref{sec:fullinterface}، و
به $\Adv$ از طریقِ $\opl{Mediate}$ --- بینِ هر کارکردِ استاندارد مشترکند و
کنار گذاشته می‌شوند، و این هیچ هزینه‌ای ندارد: هر دو هدف به $\bits \cup
\{A,Z\}$ خدمت می‌دهند و $\Stdpid$ را می‌پذیرند، پس $\serves(\F,\Corr)$ و
$\serves(\F,\Adv)$ برای هر $\F$ِ استاندارد، هر میدانی هم که خودش داشته
باشد، برقرارند و درایه‌ای برای هیچ‌یک هرگز نمی‌توانست شکست بخورد.

\begin{definition}[سیستم‌ها]\label{def:sys}
یک کارکرد \emph{استاندارد} است اگر $\F.F \notin \{A,Z,C\}$ و هر یک از
هسته‌هایش چنان که در بخشِ~\ref{sec:fullinterface} آمده پوشانده شده باشد ---
جز $\op{Leak}$، که چنان که در بخشِ~\ref{sec:leakage} آمده پوشانده می‌شود، و
$\op{Initialize}$، که یک رویه است و هیچ پوشانده نمی‌شود. یک \emph{سیستم}
(یا \emph{جهانِ}) $\pi$ مجموعه‌ای متناهی از کارکردهای استاندارد است که
شناسه‌های فرایندِ متمایز حمل می‌کنند: اگر $\F,\F' \in \pi$ و $\F \neq \F'$
آن‌گاه $\F.\PID \neq \F'.\PID$. سیستم هر چه اعضایش اشغال کنند اشغال
می‌کند،
\[
  \IDs(\pi) \;:=\; \bigcup_{\F \in \pi} \IDs(\F) \;=\; \bigl\{\, (\F.\PID,\, P) \;:\; \F \in \pi,\ P \in \F.\Ps \,\bigr\} .
\]
\end{definition}

سیستم همان است که طراحِ پروتکل می‌نویسدش --- ما \emph{پروتکل} و
\emph{سیستم} را به‌جای هم به کار می‌بریم، و \emph{ایده‌آل} و \emph{واقعی}
نامِ سبک‌های نوشتنِ یکی از آن‌هایند (در زیر)، نه دو ردهٔ جدا. سه ماشین
هیچ بخشی از آن نیستند: محیطِ $\Zenv$، حریفِ $\Adv$، و دفترِ فسادِ $\Corr$
از بخشِ~\ref{sec:corruption}. هیچ‌یک استاندارد نیست، پس بنا بر
تعریفِ~\ref{def:sys} هیچ سیستمی هیچ‌یک از آن‌ها را در بر ندارد و
$\IDs(\Zenv)$، $\IDs(\Adv)$ و $\IDs(\Corr)$ بیرونِ هر سیستمی می‌مانند؛
به‌جایش اجرای بخشِ~\ref{sec:execution} هر سه را فراهم می‌کند. برای $\Zenv$
و $\Adv$ همین است که می‌گذارد هر یک هویتِ خودش را ادعا کند و در همان حال
روی سیستمِ زیرِ آزمون خموش باشد. برای $\Corr$ همین است که دفتر را از
عملیاتِ مجموعه‌ایِ بخشِ~\ref{sec:repl} و فصلِ~\ref{sec:blockers} برکنار
نگه می‌دارد: هیچ مسدودگری برایش ساخته نمی‌شود، هیچ جانشانی‌ای جایش را
نمی‌برد، و هر مقایسه روی یک و همان دفتر می‌دود. رزروکردنِ آن سه نام آن‌ها
را از هر $\F.\admits$ که به‌صورتِ $\Stdpid$ اعلام شده هم بیرون نگه
می‌دارد، پس هیچ کارکردی دفتر را در میانِ فراخواننده‌هایش نمی‌پذیرد --- و
این درست است، چون دفتر تنها فراخوانده است.

تمایزِ شناسه‌های فرایند کارِ واقعی می‌کند. سیستم با آنچه در هر شناسه‌ای که
به کار می‌برد می‌نشاند تعیین می‌شود، پس عضویت را تنها شناسه حل می‌کند. و
هر هویت آن شناسه را در مؤلفهٔ اولش ثبت می‌کند، پس اعضای متمایز
$\IDs(\F) \cap \IDs(\F') = \emptyset$ دارند: اجتماعِ بالا مجزاست، و هیچ
هویتی از یک سیستم دو بار اشغال نمی‌شود.

\begin{definition}[سیستم‌های خوش‌ساخت]\label{def:wfsys}
بنویسیم $\pi^{+} := \pi \cup \{\Zenv,\Adv,\Corr\}$ برای یک سیستم به همراهِ
آن سه ماشینی که هر اجرا فراهم می‌کند. سیستمِ $\pi$ \emph{خوش‌ساخت} است، و
می‌نویسیم $\WF(\pi)$، اگر
\[
  \begin{gathered}
  \forall\, \F \in \pi^{+} \ \ \forall\, (\F',R) \in \F.\uses \ \ \ \exists\, \G \in \pi^{+} \;:\; \\[3pt]
  \begin{aligned}
  \G.F = \F'.F \ \ &\wedge\ \ \bigl(\G.s = \F.s \ \vee\ \op{global}(\G)\bigr) \\
  &\wedge\ \ R(\F,\G) .
  \end{aligned}
  \end{gathered}
\]
\end{definition}

سه چیز در بابِ این. شرط‌های روی $\G$ درست همان‌هایی هستند که
$\opl{Guard}$ هنگامِ گذاشتنِ فراخوان تحمیل می‌کند --- نامِ پذیرفته‌شده،
نشستِ جور، مگر آنکه فراخوانده مشترک باشد، و میدان‌های برازنده --- پس
$\WF(\pi)$ درست همین را می‌گوید که فراخوان‌هایی که اعضای $\pi$ برای
گذاشتنشان ساخته شده‌اند می‌توانند موفق شوند. جورشدن با \emph{نام} است، نه
با خودِ کارکرد، و همین است که می‌گذارد یک وابستگی از جانشانیِ زیرسیستم جان
سالم ببرد: درایه‌ای که نامِ $\varphi_q$ را می‌برد با هر چه در آن نام
بنشیند برآورده می‌شود، ایده‌آل یا واقعی، تا وقتی باقیِ شرط‌ها را برآورد. و
نیمهٔ خواسته میدان‌های فراخواننده و فراخوانده را به هم گره می‌زند؛ $\WF$
بی آن تقریباً هیچ نمی‌گفت.

هر دو سور روی $\pi^{+}$ دامنه دارند، و هر بزرگ‌شدن جای خود را به‌دست
می‌آورد: در سمتِ راست، درایه‌ای که نامِ یکی از آن سه ماشینِ فراهم‌شده را
می‌برد اصلاً برآورده‌شدنی می‌شود، چون هیچ‌یک عضوِ هیچ سیستمی نیست؛ و در
سمتِ چپ، درایه‌های خودِ آن‌ها وارسی می‌شود --- بی‌درنگ همین‌جا، چون
$\serves(\Zenv,\Adv)$ و $\serves(\Adv,\Zenv)$ برقرارند و $\Corr.\uses$
تهی است. آنچه $\WF$ به‌تنهایی نمی‌کند، جان‌بردن از ساخت‌هایی است که در پی
می‌آیند: افزودنِ عضو می‌تواند تعهد بیاورد، و جانشانیِ زیرسیستم هر خواسته
را در اشغال‌گرِ تازه از نو می‌سنجد. گزارهٔ~\ref{prop:wfblk}
(فصلِ~\ref{sec:blockers}) اولی را حل می‌کند و قضیهٔ~\ref{thm:wfrepl}
دومی را.

\heading{شرط‌ها، یک‌جا} خوش‌ساختی یکی از چند شرطِ ایستایی است که این مقاله
بر سیستم‌ها، هسته‌ها و ماشین‌ها می‌گذارد، و هر یک همان جا که نخست لازم
می‌شود تعریف می‌شود --- که خواننده را بی هیچ راهی می‌گذارد تا ببیند کدام
نتیجه کدام شرط را می‌خواند. این جدول یک بار همه را جمع می‌کند. هیچ‌یک از
آن‌ها در طولِ اجرا خوانده نمی‌شود: $\opl{Exec}$ میدان می‌خواند، نه شرط،
چنان که گزارهٔ~\ref{prop:inert} به‌طورِ خاص برای $\uses$ ثبت می‌کند.

\begin{center}
\small
%% ragged, not justified: a 0.20\textwidth column cannot be justified
%% without visible stretching, and it was 12 of the 17 underfull boxes
\begin{tabular}{@{}>{\raggedright\arraybackslash}p{0.20\textwidth} >{\raggedright\arraybackslash}p{0.30\textwidth} >{\raggedright\arraybackslash}p{0.40\textwidth}@{}}
\toprule
شرط & چه چیزی را مقید می‌کند & چه کسی می‌خواندش \\
\midrule
$\WF(\pi)$، تعریفِ~\ref{def:wfsys}
  & یک سیستم: هر درایهٔ $\uses$ شاهدی دارد --- فراخوان‌هایی که اعضایش برای \emph{گذاشتنشان} ساخته شده‌اند می‌توانند موفق شوند
  & قضیهٔ~\ref{thm:wfrepl}، گزارهٔ~\ref{prop:wfblk}، تعریفِ~\ref{def:prop} \\[0.6ex]
جانشانیِ خوش‌ساخت، تعریفِ~\ref{def:wf}
  & سیستمِ ورودی در برابرِ آنچه می‌ماند: هیچ شناسهٔ فرایندِ مشترکی
  & گزارهٔ~\ref{prop:repl}، قضیهٔ~\ref{thm:wfrepl}، گزارهٔ~\ref{prop:absorb}، قضیهٔ~\ref{thm:comp}، قضیهٔ~\ref{thm:parcomp}، گزارهٔ~\ref{prop:proptransfer} \\[0.6ex]
سازگار، تعریفِ~\ref{def:compat}
  & دو سیستم: شناسهٔ فرایندِ مشترک یکسان خدمت می‌دهد و یکسان می‌پذیرد
  & قضیهٔ~\ref{thm:blk}، ادعای دوم \\[0.6ex]
نشست‌یکنوا، تعریفِ~\ref{def:sessunif}
  & یک سیستم یا محیط: اینکه چه کسی می‌تواند به عضوی \emph{برسد} به نشستِ او وابسته نیست
  & لمِ~\ref{lem:reloc}، گزارهٔ~\ref{prop:relocemul}، نتیجهٔ~\ref{cor:onepid}، نتیجهٔ~\ref{cor:moveone} \\[0.6ex]
تنگ در $I$، نکتهٔ~\ref{rem:tight}
  & یک سیستم: اعضای بیرونِ $I$ فراخواننده‌ها را با شناسهٔ فرایند سنجاق می‌کنند
  & تعریفِ~\ref{def:prop} \\[0.6ex]
هیچ فراخوانِ پاسخگو نمی‌گذارد، تعریفِ~\ref{def:respcall}
  & هسته‌های یک سیستم
  & قضیهٔ~\ref{thm:dummy}، لمِ~\ref{lem:absorb} (نیمهٔ واگذارنده)، گزارهٔ~\ref{prop:propslotabs}، نتیجهٔ~\ref{cor:proptransfer2} \\[0.6ex]
بی‌اعتنا به رد، تعریفِ~\ref{def:refobl}
  & ماشینی در جایگاهِ حریف
  & همان چهار \\[0.6ex]
هویت‌اتمی، تعریفِ~\ref{def:atomic}
  & همهٔ کد؛ قراردادی جاری و نه آزمونی به‌ازای هر سیستم
  & لمِ~\ref{lem:reloc}، گزارهٔ~\ref{prop:rename} \\[0.6ex]
نشت‌خوش‌ساخت، بخشِ~\ref{sec:leakage}
  & یک تحقق
  & هیچ‌کس: ثبت شده، به کار نرفته \\
\bottomrule
\end{tabular}
\end{center}

دو چیز که این جدول آشکار می‌کند. دو سطرِ اول یک نام دارند و دو شرطِ
متفاوتند: $\WF$ در بابِ یک سیستم به‌تنهایی است، و \emph{جانشانیِ}
خوش‌ساخت در بابِ اینکه سیستمِ ورودی کجا فرود می‌آید. و $\WF$ و
نشست‌یکنوایی دوگانِ یکدیگرند --- فراخوان‌هایی که عضوی برای گذاشتنشان ساخته
شده، در برابرِ اینکه چه کسی می‌تواند به او برسد --- که وسوسه می‌کند
آن‌ها را در یک شرط تا کنیم. جدا نگه داشته شده‌اند چون هیچ نتیجه‌ای هر دو
را لازم ندارد، و چون یک تعریف در هر حال نمی‌توانست هر دو را حمل کند: از
بازیِ تعریفِ~\ref{def:prop} خواسته می‌شود که خوش‌ساخت باشد و عامداً
نشست‌یکنوا \emph{نیست}، چنان که نکتهٔ~\ref{rem:gmsessunif} توضیح می‌دهد.

می‌گوییم $\F$ \emph{محلی} است اگر $\F.\admits = \{\PID : \PID.F =
F^{\parent}\}$ برای نامِ والدِ (فراخوانندهٔ) معینِ $F^{\parent}$. (این را
به خودِ کارکرد وامی‌گذاریم که بگوید آیا $\Apid$ و $\Zpid$ می‌توانند فرایش
بخوانند یا نه.) مفهوم‌های میانی ممکن‌اند: برپاسازی‌ای که از دو پروتکلِ
نام‌دار در دسترس است $\{\PID : \PID.F \in \{F_1,F_2\}\}$ را می‌پذیرد؛ و
یک $\admits$ می‌تواند کلِ شناسه‌های فرایند را سنجاق کند و نه نام‌ها را، و
نمونه‌های اشغال‌شدهٔ فراخواننده‌هایش را رو نویسد ---
نکتهٔ~\ref{rem:tight} می‌گوید چرا اعضای خصوصیِ یک تحقق باید چنین کنند. در
سرِ دیگرِ طیف می‌گوییم $\F$ \emph{جهانی} است، به معنای برپاسازیِ مشترکِ
UCِ تعمیم‌یافته~\cite{CDPW07}، اگر
\[
  \op{global}(\F) \;:\iff\; \Stdpid \subseteq \F.\admits \ \ \vee\ \ \F.F \in \{A,Z,C\} ,
\]
که یک مفهوم است و نه دو: عطفِ اول حالتِ معمولی است، هر کارکردِ استانداردی
می‌تواند فرایش بخواند، و عطفِ دوم آن سه ماشینی را می‌پوشاند که اجرا فراهم
می‌کند --- و هر چه در جایگاهشان بایستد --- که به دلیلی دیگر جهانی‌اند و
هر جا $\op{global}$ آزموده شود همان‌گونه با آن‌ها رفتار می‌شود.

این تمایز همان چیزی است که مؤلفهٔ نشست برای آن است. کارکردِ محلی بخشی از
یک اجرا از یک پروتکل است، پس فراخوان تنها از نشستِ خودش می‌تواند به آن
برسد: سطرِ~\ref{ln:session} از $\opl{Guard}$ می‌پرسد $\id'.s = \id.s$.
شمارهٔ نمونه آزموده نمی‌شود، پس یک نشست می‌تواند هر چند نمونه از یک نام که
بخواهد در بر داشته باشد و آن‌ها می‌توانند یکدیگر را فرا بخوانند. کارکردِ
جهانی بینِ اجراها مشترک است --- یک نشست از آن، در دسترسِ هر نشستی --- پس
آن آزمون برای او بخشیده می‌شود، و همهٔ نشست‌های $\F$ می‌توانند یک نمونهٔ
مشترک را فرا بخوانند. کارکردِ ایده‌آلِ جهانی را با $\op{G}$ در آغاز نام
می‌دهیم و واقعیِ جهانی را با $\gamma$ در آغاز، چنان که در
$\F_{\op{G}\op{Sig}}$؛ و $\G$ِ خوش‌نویس همان می‌ماند که بوده، کارکردِ
دومی در کنارِ $\F$. آن سه ماشینی که اجرا فراهم می‌کند جهانی‌اند: $\Adv$ و
$\Corr$ همین حالا با بندِ اول، و $\Zenv$ با بندِ دوم، که آنجاست چون
$\Zenv.\admits = \Apid \cup \Zpid$ عامداً تنگ است و با این حال محیط باید
از هر نشستی در دسترس باشد.


\heading{مشخصهٔ واقعی در برابرِ ایده‌آل} کارکردهای ایده‌آل و پروتکل‌های
واقعی از چند جهت متفاوتند. اول، کارکردهای ایده‌آل وضعِ فسادِ طرف‌های
موجود در مجموعهٔ طرف‌هایشان را می‌دانند. دوم، می‌توانند اجرا را به
شبیه‌ساز بسپارند و می‌توانند واسط‌های ویژهٔ شبیه‌ساز عرضه کنند (مثلاً
$\Fdiffuse$ را ببینید). سرانجام، می‌توانند حالت را بینِ طرف‌ها مشترک کنند
--- و همین اشتراک نکتهٔ اصلی است: نگهبانِ فصلِ~\ref{sec:wrappers} هر
فراخوان را به یک طرف گره می‌زند، پس طرف‌ها تنها از راهِ کارکردهایی با هم
سخن می‌گویند که حالت را بینشان مشترک می‌کنند، و پروتکلِ واقعی با همتاهایش
از راهِ زیرروال‌های ایده‌آل سخن می‌گوید و از هیچ راهِ دیگری. اینجا و در
زیر، $\Fdiffuse$ و $\Fio$ نامِ کارکردهایی است که نمی‌دهیمشان: یک کانالِ
پخش و یک بازپخشِ ورودی--خروجی. هیچ‌چیزِ این مقاله به کدِ آن‌ها بند نیست.
سرچشمهٔ تصادفِ $\Frand$ و انبارِ $\Fstore$ داده شده‌اند، در
فصل‌های~\ref{sec:frand} و~\ref{sec:fstore}؛ آن‌ها کوچک‌ترین مثال‌های این
مقاله‌اند و همان جایی هستند که خوانندهٔ ناآشنا با این سبک باید از آنجا
شروع کند.

واسط‌های هستهٔ پروتکل‌های \emph{واقعیِ} خوش‌مشخصه هیچ‌یک از این ویژگی‌ها
را ندارند، جز آنچه در بلوک‌های ساختمانیِ ایده‌آل‌شده و پوشانه‌های کنترلِ
دسترسی مستتر است؛ و تحقق‌های واقعی مشخص نمی‌کنند که نشت هنگامِ فساد چگونه
حساب می‌شود، چون آن محتوای نوارهای کار
است.\footnote{رهیافتی منظم‌تر ذخیره‌سازی را به‌صورتِ کارکردی ایده‌آل مدل
می‌کرد که پاک‌کردنِ امن عرضه می‌کند؛ ما از آن لایهٔ افزوده می‌پرهیزیم تا
نمادگذاری سبک بماند.}

\heading{مقدارِ آغازین} $\op{Initialize}$ یک رویه است و نه یک واسط: یک بار
می‌دود، وقتی کارکرد نخستین بار فعال می‌شود، و متغیرهایی را که باقیِ کد به
کار می‌برد برپا می‌کند. هیچ فراخوانی نمی‌گذارد --- اجرای
بخشِ~\ref{sec:execution} هر ماشین را پیش از نصبِ توزیع‌گرش مقدارِ آغازین
می‌دهد، پس چیزی نبود که فراخوانی را خدمت کند --- هیچ فراخواننده‌ای
نمی‌تواند به آن برسد، هیچ شناسهٔ فراخوانندهٔ ادعایی نمی‌گیرد، و کارکردی که
هیچ اعلام نکند تهیِ آن را می‌گیرد. پروتکل‌های واقعی هیچ مقداردهِ جهانی از
این دست حمل نمی‌کنند، تنها مقداردهِ ویژهٔ یک طرف در یک فرایند، که
$(\PID,P).\op{Initialize}()$ نوشته می‌شود و در آن نمادگذاری را بدرفتاری
می‌کنیم: آن‌ها خصوصی‌اند، نه واسط‌هایی که هر کس بتواند خطابشان کند.

%% ==================================================================
\chapter{نگهبان‌ها، خموش‌کننده‌ها و واسطه‌گرها}\label{sec:wrappers}

این فصل سه تکه ماشین‌آلات را تعیین می‌کند. محمولِ $\opl{Guard}$
فراخوان‌هایی را که به یک ماشین می‌رسند مقید می‌کند؛ پوشانهٔ
$\opl{Silence}$ فراخوان‌هایی را که او می‌گذارد؛ و پوشانهٔ $\opl{Mediate}$
فراخوانی را که طرفِ فاسد نمی‌تواند خودش پاسخش دهد به حریف می‌سپارد. واسطِ
کاملِ بخشِ~\ref{sec:fullinterface} ترکیبِ همین سه است.

بنویسیم $\Op$ برای واسطِ هستهٔ عمومیِ $\F$. ما تنها هسته‌ها را مشخص
می‌کنیم؛ پوشانهٔ زیر هر یک را به واسطی کامل بدل می‌کند، و فراخواننده به
$\fopl{Op}$ می‌رسد، هرگز به $\Op$. آزمون‌های کنترلِ دسترسی در محمولِ
$\opl{Guard}$ جمع شده‌اند، که کارکرد نیست بلکه الگوریتم است: هویتِ هدفِ
$\id$، هویتِ فراخوانندهٔ ادعاییِ $\id'$ و کارکردِ $\F$ را می‌گیرد و یک
بولی برمی‌گرداند. دو قراردادِ ارائه. جعبه‌ها دو وزن دارند: پُر نشانهٔ
کارکرد است، که هویت اشغال می‌کند، و تُهی نشانهٔ محمول یا پوشانه یا رویه،
که هیچ اشغال نمی‌کند. و فراخوان چنین نوشته می‌شود: $\id.\Op(in)$
\textbf{as} $\id'$ که در آن $\id'$ هویتی است که در گذاشتنِ آن ادعا شده، و
چنین دریافت می‌شود: $\id.\Op(in)$ \textbf{from} $\id'$ که در آن $\id'$
هویتی است که فراخوان زیرِ نامِ آن ادعا شده بود. اِعمال در پوشانه نشسته
است، همان جا که سطرِ~\ref{ln:guard} آن را \textbf{require} می‌کند، پس
فراخوانی که او ردش کند هرگز به هسته نمی‌رسد.
سطرِ~\ref{ln:interface} آزمون می‌کند که واسط وجود دارد؛
سطرِ~\ref{ln:caller}، که نامِ فراخواننده پذیرفته است و برای طرفِ درست
عمل می‌کند؛ سطرِ~\ref{ln:session}، که فراخواننده و فراخوانده در یک نشستند،
مگر آنکه فراخوانده بینِ نشست‌ها مشترک باشد.


\begin{procedure}{محمولِ $\opdef{Guard}$}
\opsig{$\opl{Guard}_{\F}\bigl(\id,\id'\bigr)$}:
\begin{algorithmic}[1]
\State \textbf{return}\  $\id.\PID = \F.\PID \ \wedge\  \id.P \in \F.\Ps$
       \label{ln:interface}
  \Comment{واسط‌هایی که وجود دارند}
\State \algand\  $\id'.\PID \in \F.\admits \ \wedge\  \id'.P = \id.P$
       \label{ln:caller}
  \Comment{فراخواننده‌های پذیرفته}
\State \algand\  $\bigl(\id'.s = \id.s \ \vee\ \op{global}(\F)\bigr)$
       \label{ln:session}
  \Comment{یک نشست، مگر آنکه مشترک باشد}
\end{algorithmic}
\end{procedure}

$\opl{Guard}$ فراخوان‌هایی را که به ماشین می‌رسند مقید می‌کند؛ پوشانهٔ
$\opl{Silence}$ در زیر، فراخوان‌هایی را که او \emph{می‌گذارد}. این پوشانه
یک مجموعهٔ هویتِ $\Ids$ و یک برنامه می‌گیرد، برنامه را می‌دواند، و هر
فراخوانی که صادر کند را می‌رباید. فراخوانی که هویتی در $\Ids$ ادعا کند با
$\rej$ رد می‌شود، بی‌آنکه هرگز گذاشته شود؛ هر فراخوانِ دیگری گذاشته
می‌شود و پاسخش پس داده می‌شود. پس $\Ids$ درست همان جایی است که برنامهٔ
پوشانده‌شده در آن خموش است، و می‌تواند به‌عنوانِ هر چیزِ بیرونِ آن سخن
بگوید.

دامنه، همان برنامهٔ پوشانده‌شده است و نه فراتر از آن. فراخوانی که او
می‌گذارد ژتون را به ماشینی دیگر می‌سپارد، و آنچه \emph{آن} ماشین می‌گوید
با پوشانهٔ خودش حکم می‌شود، نه با این یکی. اگر جز این بود، کارکرد
زیرروال‌هایش را ترایا خموش می‌کرد و هیچ‌یک نمی‌توانست به‌عنوانِ خودش سخن
بگوید. هر کاربردِ $\opl{Silence}$ در زیر همین شکل را دارد: واسطی کامل که
هستهٔ زیرِ خدمتش را می‌پوشاند.

خواندنِ این به‌عنوانِ هندلر، به معنای اثرهای جبری~\cite{PP09,BP15}، به‌کار
می‌آید. فراخوانی که برنامه صادر می‌کند عملی از اثرِ محیطی است، و $k$ همان
ادامه‌ای است که چشم‌به‌راهِ پاسخِ آن است. پوشانهٔ $\opl{Silence}$ تصمیم
می‌گیرد با چه چیزی از سر گرفته شود: $\rej$ برای ادعای خموش‌شده، و در غیرِ
این صورت پاسخِ فراخوانی که واقعاً گذاشته شد. بندِ \textbf{return} نتیجهٔ
خودِ برنامه را دست‌نخورده رد می‌کند، پس خموش‌سازی آنچه برنامه می‌تواند
بگوید را تغییر می‌دهد و هیچ‌چیزِ دیگر را. این قیاس تنها شکلِ پوشانه را
تعیین می‌کند.

\begin{procedure}{پوشانهٔ $\opdef{Silence}$ \\
\params{مجموعهٔ هویتِ $\Ids$، \quad برنامهٔ $\id.\Op(in)$ از $\id'$}}
\opsig{$\opl{Silence}\bigl(\Ids,\ \id.\Op(in) \text{ from } \id' \bigr)$}:
\settowidth{\clausewd}{$\id''.\fopl{Op}(in'')\ \text{as}\ \id''' \,;\ k$}
\begin{algorithmic}[1]
\State \textbf{handle}\ \ $\id.\Op(in)$ from $\id'$\ \ \textbf{with}
%% \clausewd aligns the \mapsto column, which is the alignment that carries
%% meaning. \reswd used to pad every result out to the width of the longest,
%% which at this measure pushed the trailing condition off the line. The
%% conditions right-align instead, so they line up with each other at the
%% right edge and the results take their natural width.
%% The longest of these needs 284.5pt and the listing gives 280.0pt, so the
%% two fixed gaps -- the sub-indent and the space before the condition --
%% are \enspace rather than \quad. Four points each, eight in all, and the
%% clause alignment is untouched.
\State \enspace \makebox[\clausewd][l]{$\textbf{return}\ out$}
       $\ \mapsto\ $ $out$
\State \enspace \makebox[\clausewd][l]{$\id''.\fopl{Op}(in'')\ \text{as}\ \id''' \,;\ k$}
       $\ \mapsto\ $ $k(\rej)$\enspace\hfill اگر $\id''' \in \Ids$
\State \enspace \makebox[\clausewd][l]{}
       $\ \mapsto\ $ $k\bigl(\id''.\fopl{Op}(in'')\ \text{as}\ \id'''\bigr)$\enspace\hfill وگرنه
\end{algorithmic}
\end{procedure}

با داشتنِ $\opl{Silence}$ می‌توانیم بگوییم بر فراخوانی که طرفِ فاسد
نمی‌تواند خودش پاسخش دهد چه می‌گذرد. هر واسطِ کاملِ زیر چنین فراخوانی را
دو بار به حریف می‌سپارد، در راهِ ورود و در راهِ خروج، و هر دو بار در همان
یک شکل: پوشانهٔ $\opl{Mediate}$. این پوشانه یک مجموعهٔ فسادِ $\Cs$، عملِ
$\id.\Op$ که در حالِ خدمت است، مقدارِ $x$ که باید سپرده شود، و
فراخوانندهٔ ادعاییِ $\id'$ِ عملی که به‌جایش می‌ایستد را می‌گیرد. اینکه
پاسخ از کجا می‌آید در درونِ خودش تعیین شده است: هر عملی که باشد، پاسخ‌های
طرفِ فاسد از $(A,\id.P)$ می‌آیند. یک چیز در بابِ آن آرگومانی معمولی
نیست: هم \textbf{require} و هم \textbf{return}ِ $\opl{Mediate}$ از عملی
برمی‌گردند که در کدِ آن ایستاده‌اند، و نه فقط از خودِ آن سطر. وقتی آزمونِ
$\opl{Mediate}$ شکست بخورد او هیچ نمی‌کند و آن عمل ادامه می‌دهد.

این قلاب یک فراخوانِ معمولی روی واسطِ کاملِ $\Adv$ است، که هویتِ $\id$ِ
واسطِ در حالِ خدمت را ادعا می‌کند. پس با $\opl{Guard}_{\Adv}$ روبه‌رو
می‌شود، و از آن می‌گذرد چون طرفِ ادعاییِ $\id.P$ همان طرفِ هدف است؛ و
حریفی که پاسخ می‌دهد در بندِ خموش‌کنندهٔ خودِ $\Adv$ است، که تنها
$(A,\id.P)$ را می‌پذیرد. پس حریف در حالِ پاسخ‌دادن برای $\id.P$ می‌تواند
به‌عنوانِ خودش در $\id.P$ سخن بگوید و به‌عنوانِ هیچ‌چیزِ دیگر، و
$\opl{Mediate}$ به خموش‌سازیِ خودش نیازی ندارد.

\begin{procedure}{پوشانهٔ فسادِ $\opdef{Mediate}$ \\
\params{مجموعهٔ فسادِ $\Cs$، \quad عملِ $\id.\Op$، \quad مقدارِ $x$، \quad فراخوانندهٔ ادعاییِ $\id'$}}
\opsig{$\opl{Mediate}\bigl(\Cs,\ \id.\Op,\ x,\ \id'\bigr)$}:
\begin{algorithmic}[1]
\If{$(\id'.F \neq A \ \wedge\ \id.P \in \Cs)$}
  \State \textbf{return} $(A,\id.P).\fopl{A}\bigl(\id.\Op,\ x\bigr) \text{ as } \id$
\EndIf
\end{algorithmic}
\end{procedure}

%% ==================================================================
\chapter{مدلِ اجرا}\label{ch:execmodel}

%% Six short chapters merged here. Each was two to four pages, and each
%% bought a half-empty page on the way out; together they are one
%% argument -- the register that records corruption, what a functionality
%% reveals when corrupted, the two machines the execution supplies, the
%% interface they all present, and the execution that dispatches to it.
%% Corruption leads because leakage is about what corruption exposes.

%% ==================================================================

%% ==================================================================
\section{فساد}\label{sec:corruption}

فساد واسطی نیست که پروتکل‌های واقعی محققش کنند (پروتکلِ واقعی وقتی
«دستور» بگیرد که فاسد شود چه می‌کند؟)، اما برای تقلید باید وجود داشته
باشد. می‌توانیم فرض کنیم پروتکل‌های واقعی همیشه فساد را رد می‌کنند، اما
آن‌گاه هرگز به‌عنوانِ فاسد ثبت نمی‌شوند.

ما فساد را به‌صورتِ کارکردی جهانی به نامِ $\Corr$ مدل می‌کنیم. چند ماشین
واسطِ کاملشان یک‌جا داده شده و از هسته به‌دست نیامده --- محیط، حریف، و
عملِ نشتِ بخشِ~\ref{sec:leakage} --- اما $\Corr$ تنها ماشینی است که نه
$\opl{Silence}$ را به کار می‌برد و نه $\opl{Mediate}$، و گزارشِ وضع دلیلش
است: آن پوشانه در هر فراخوان وضعِ فساد را می‌خواند، که همان چیزی است که
این واسط پیاده‌اش می‌کند، پس پوشاندنش دوری بود. پس هر واسطِ زیر تنها
آزمون‌هایی را که لازم دارد به‌دست حمل می‌کند، و $\fopl{Status}$ به‌طورِ
خاص \emph{همیشه} مجموعهٔ فاسدها را برمی‌گرداند، بی‌آنکه نخست بپرسد آیا آن
طرف فاسد است.


\begin{interface}{دفترِ فسادِ $\Corr$ \\
\params{$\PID := C$, \quad $\Ps := \bits \cup \{A,Z\}$, \quad
$\admits := \Stdpid \cup \Apid \cup \Zpid$, \quad $\uses := \emptyset$, \quad $\pars := \none$}}

\opsig{$\op{Initialize}()$}:
\begin{algorithmic}[1]
\setcounter{contline}{0}\algcont
\State $\Cs \gets \emptyset$
\algsave
\end{algorithmic}
\vspace{1.4ex}
\opsig{$\id.\fopdef{Corrupt}()$} \quad از $\id'$
\begin{algorithmic}[1]
\algcont
  \Req{$\opl{Guard}_{\Corr}\bigl(\id,\id'\bigr)$}
  \Req{$\id'.F \in \{A,Z\}$} \label{ln:corronlya}
    \Comment{تنها بیرون، هرگز زیرروال}
  \Req{$\id.P \in \bits$}
    \Comment{نام‌های رزرو درستکار می‌مانند}
  \State $\Cs \gets \Cs \cup \{\id.P\}$ \label{ln:corrwrite}
    \State \textbf{return} $\ok$
\algsave
\end{algorithmic}
\vspace{1.4ex}
\opsig{$\id.\fopdef{Status}()$} \quad از $\id'$
\begin{algorithmic}[1]
\algcont
  \Req{$\opl{Guard}_{\Corr}\bigl(\id,\id'\bigr)$}
  \State \textbf{return} $\Cs$
\algsave
\end{algorithmic}
\vspace{1.4ex}

\opsig{$\id.\fopl{Leak}()$} \quad از $\id'$
\begin{algorithmic}[1]
\algcont
  \Req{$\opl{Guard}_{\Corr}\bigl(\id,\id'\bigr)$}
  \Req{$\id'.F = A$}
  \State \textbf{return} $\none$
    \Comment{ثابت: دفتر هیچ رازی نگه نمی‌دارد}
\algsave
\end{algorithmic}
\end{interface}

نشتِ خودِ $\Corr$ ثابت است، پس آزمونِ فسادی که بخشِ~\ref{sec:leakage} بر
هر عملِ نشتِ دیگری تحمیل می‌کند اینجا چیزی برای محافظت نداشت؛ با این حال
نگهبان نگه داشته می‌شود، تا حریف دفتر را در همان طرفی بخواند که برایش عمل
می‌کند و در هیچ طرفِ دیگر.

\heading{خواندنِ دفتر} $\Cs$ حالتِ خودِ $\Corr$ است، پس هیچ ماشینِ دیگری
رونوشتی از آن ندارد. هر جا $\Cs$ در کدِ کارکردی جز خودِ $\Corr$ ظاهر
می‌شود، فهمِ مطلب این است که نخست واکشی شده است، با
\[
  \Cs \;\gets\; (C,\id.P).\fopl{Status}() \ \text{ as } \id ,
\]
و ما آن سطر را هر جا که تنها چیزی را تکرار می‌کرد که کدِ پیرامون خودش
روشن می‌کند حذف می‌کنیم. در درونِ $\Corr$ چیزی برای واکشی نیست: $\Cs$
مستقیم خوانده و نوشته می‌شود، و $\op{Corrupt}$ در جا به‌روزش می‌کند.

\begin{remark}[چرا $\fopl{Status}$ واسطه‌گری نمی‌شود]\label{rem:statusmed}
اگر مثلِ عمل‌های دیگر پوشانده می‌شد، $\fopl{Status}$ درست همان جا که مهم
است بی‌مصرف می‌بود. فرض کنید $\Zenv$ به‌عنوانِ $\id'$ دست دراز کند تا
وضعِ فسادِ $\id'.P$ را بپرسد. اگر $\id'.P$ فاسد باشد، دروازه فراخوان را
رد می‌کند و $\opl{Mediate}$ آن را به حریف می‌سپارد --- پس پاسخِ «آیا
$\id'.P$ فاسد است؟» هر چیزی است که حریف بگوید، و می‌تواند بگوید نه. و چون
هر فراخواننده‌ای دفتر را می‌خواند تا تصمیم بگیرد با طرفِ فاسد چگونه رفتار
کند، پاسخی که حریف مهارش کند ماشین‌آلاتِ واسطه‌گری را به میلِ خود خاموش
می‌کرد.

پس $\fopl{Status}$ هیچ واسطه‌گری و هیچ خموش‌سازی نمی‌کند: پس از نگهبان،
$\Cs$ را برمی‌گرداند، حالتِ خودِ دفتر را. این همان یک عمل است که پاسخش
باید از دفتر بیاید، و چون تنها‌خواندنی است چیزی به حریف نمی‌دهد که پیش‌تر
نتوانسته باشد استنتاج کند. هستهٔ $\op{Status}$ِ پیش‌نویس‌های قدیمی‌تر هم با
آن می‌رود، چون پوشانه‌ای نمانده که از آن برکنار بماند.
\end{remark}

فساد کارِ بیرون است. سطرِ~\ref{ln:corronlya} $\id'.F \in \{A,Z\}$ را
می‌پذیرد، پس حریف و محیط هر یک می‌توانند طرفی را در $\Cs$ بگذارند و هیچ
کارکردِ استانداردی اصلاً نمی‌تواند چیزی آنجا بگذارد. نگهبانِ بالای آن،
واسطِ فاسدکننده را به همان طرفی که فاسدش می‌کند گره می‌زند و می‌پرسد
$\id'.P = \id.P$: برای فاسدکردنِ $P$، حریف باید در $(A,P)$ بدود، و محیط
باید $(Z,P)$ را ادعا کند، که سطرِ~\ref{ln:zparty} می‌گذارد ریشه‌اش این
کار را برای هر طرفی بکند. هر دو می‌توانند بخوانند هم: $\Apid \cup \Zpid
\subseteq \Corr.\admits$، و $\fopl{Status}$ کلِ $\Cs$ را برمی‌گرداند، پس
هر یک می‌تواند در هر نقطه‌ای بپرسد و همیشه کلِ مجموعه را می‌آموزد. و تنها
طرف‌های اصیل می‌توانند فاسد شوند: $\fopl{Corrupt}$ می‌پرسد $\id.P \in
\bits$، پس نام‌های رزرو بنا بر ساخت درستکار می‌مانند.

هم‌اندازگی درست درمی‌آید، و آنچه تأمینش می‌کند \emph{خواندن} است و نه
هیچ تقسیمِ اینکه چه کسی می‌تواند بنویسد. ماشینِ نشسته در $(A,P)$ در یک
اجرای آزمایشِ تقلید $\Adv$ است و در اجرای دیگر $\Sim$، و هیچ‌چیز آن دو را
به فاسدکردنِ همان طرف‌ها وادار نمی‌کند. اما $\Zenv$ در هر دو سو همان یک
ماشین است و در هر دو $\Cs$ را می‌خواند، پس اشغال‌گرِ جایگاه اگر مجموعهٔ
دیگری را فاسد کند، محیطی که تنها بپرسد او را می‌گیرد. پس هم‌اندازگی
نتیجه است و نه فرض: هیچ محیطِ مجازی آن دو مجموعهٔ فساد را با مزیتی بیش از
آنچه تعریفِ~\ref{def:uc} به شبیه‌ساز می‌دهد از هم تشخیص نمی‌دهد. و اینکه
محیط می‌تواند برای خودش فاسد کند این را تنها تیزتر می‌کند، چون فسادهایی که
او می‌کند در هر دو سو از خودِ اوست؛ آنچه برای گرفتن می‌ماند ابتکارِ خودِ
اشغال‌گر است.

فساد تنها با نامِ طرف اندیس‌گذاری می‌شود، پس $P$ نمی‌تواند در یک شناسهٔ
فرایند فاسد باشد و در دیگری درستکار.\footnote{می‌شد به‌جایش یک
\emph{هویت} را فاسد کرد --- $\Cs$ هویت‌ها را جمع می‌کرد و هر آزمونِ $\id.P
\in \Cs$ به $\id \in \Cs$ بدل می‌شد --- به بهای «پخش‌شدنِ» فسادها به
زیرپروتکل‌ها. قراردادِ طرف‌گستر پاکیزه‌تر است، چون یک خانوادهٔ کامل را
یک‌جا آلوده می‌کند.}
\section{نشت}\label{sec:leakage}

\begin{interface}{کارکردِ $\F$، عملِ نشت \\
\params{$\PID$, \quad $\Ps$, \quad $\admits$, \quad $\uses$, \quad $\pars$}}
\opsig{$\id.\fopdef{Leak}()$} \quad از $\id'$
\begin{algorithmic}[1]
\Req{$\id.\PID = \F.\PID \ \wedge\ \id.P \in \F.\Ps$} \label{ln:leakexists}
  \Comment{واسط وجود دارد}
\Req{$\id'.F = A \ \wedge\ \id'.P = \id.P$} \label{ln:leakadv}
  \Comment{$\Adv$، برای این طرف}
\State $\Cs \gets (C,\id.P).\fopl{Status}() \text{ as } \id$
\Req{$\id.P \in \Cs$} \label{ln:leakcorr}
  \Comment{تنها طرفِ فاسد نشت می‌کند}
\State \textbf{return} $\opl{Silence}\bigl(\{\id'' : \id'' \ne \id \},\ \id.\op{Leak}() \text{ from } \id'\bigr)$
\end{algorithmic}
\end{interface}

این سه خواسته همان‌هایی هستند که پوشانهٔ عمومیِ
بخشِ~\ref{sec:fullinterface} تحمیل می‌کند، و از نو بیان شده‌اند چون نشت
به شیوهٔ معمول از هسته به‌دست نمی‌آید. سطرِ~\ref{ln:leakexists} نیمهٔ
وجودِ $\opl{Guard}$ است؛ نیمهٔ فراخواننده عامداً ضعیف‌تر است، چون $\Adv$
باید بتواند از کارکردی نشت بگیرد که او را به‌عنوانِ فراخواننده نمی‌پذیرد،
پس سطرِ~\ref{ln:leakadv} تنها می‌خواهد که حریف برای طرفی عمل کند که دارد
می‌خواندش. سطرِ~\ref{ln:leakcorr} دروازهٔ فساد است. بی آن، حالتِ طرفِ
درستکار به میل خواندنی می‌شد، و همین است که سطرِ~\ref{ln:corrgate} از
واسطِ کامل هر جای دیگری ردش می‌کند. خموش‌کننده همان خموش‌کنندهٔ
سطرِ~\ref{ln:call} است، پس هستهٔ نشتی که فراخوان می‌گذارد، مثلِ هر هستهٔ
دیگری هویتِ خودش را ادعا می‌کند.

برای کارکردِ ایده‌آل، $\opl{Leak}$ بخشی از مشخصه است؛ پیام‌هایی که تا
اینجا دریافت شده و محتوای فهرست‌های درونی نمونه‌های معمولند. برای
پروتکل‌های واقعی می‌توان آن را چنان نشاند که هیچ نشت نکند، که یعنی مدلی
از بی‌نشتیِ کامل؛ و چاره همان کارکردهای ایده‌آلی هستند که می‌گویند کدام
بخش‌های حالت نشت می‌کنند.

پروتکلی \emph{نشت‌خوش‌ساخت} است اگر بدنهٔ هر ماشین بینِ فراخوانی‌ها پاک
باشد، هر تصادفِ نمونه‌برداری‌شده از راهِ $\Frand$ کشیده شود و هر حالتی که
از فراخوانی‌ای به فراخوانیِ دیگر می‌پاید در $\Fstore$ نگه داشته شود و نه
روی نوارِ ماشین --- که هر یک آن را از راهِ $\opl{Leak}$ِ خودش رو می‌کند
--- و ورودی‌ها و خروجی‌ها تنها آنجا نشت کنند که پروتکل آن‌ها را از راهِ
کارکردِ ورودی--خروجیِ $\Fio$ می‌گذراند، که پیش‌فرض نیست. آن‌گاه نوارهای
کار هرگز نشت نمی‌کنند، چون پاکیِ بینِ فراخوانی‌ها چیزی رویشان نمی‌گذارد.
هیچ نتیجهٔ زیر هیچ‌یک از این‌ها را مفروض نمی‌گیرد؛ این شرطی است که یک تحقق
باید برآورده کند تا $\opl{Leak}$ِ آن معنایی داشته باشد، و ما ثبتش می‌کنیم
و به کارش نمی‌بریم.

$\Zenv$ نمی‌تواند این واسط را فرا بخواند، پس برای تمایزناپذیری لازم نیست
محقق شود. با این حال باید شبیه‌سازی‌شدنی باشد، چون حریف می‌تواند فرایش
بخواند.

%% ==================================================================
\section{محیط}\label{sec:env}

ما کارکردِ ویژه‌ای به نامِ $\Zenv$ تعریف می‌کنیم، با یک واسطِ یگانه به‌ازای
هر $\id$، یعنی $\id.\fopl{Z}$، که هستهٔ $\id.\op{Z}$ را نگهبانی و خموش
می‌کند. میدانِ $\Zenv.\admits = \Apid \cup \Zpid$ سزاوارِ تأکید است: حریف
و خودِ محیط را می‌پذیرد و هیچ‌چیزِ دیگر را، پس هیچ کارکردِ استانداردی
نمی‌تواند $\Zenv$ را فرا بخواند. هر چیزِ فراخ‌تری می‌گذاشت پروتکلی محیط را
به‌عنوانِ یکی از زیرروال‌های خودش فرا بخواند --- یک «پروتکلِ باز»، که
زیرروال‌هایش نامعین می‌مانند. دو چیزِ متمایز این را ممنوع می‌کنند. تنگیِ
$\Zenv.\admits$ خودِ فراخوان را رد می‌کند: شناسهٔ فرایندِ فراخوانندهٔ
استاندارد در $\Stdpid$ است، که هیچ کجا با $\Apid \cup \Zpid$ برخورد
نمی‌کند، پس سطرِ~\ref{ln:caller} از $\opl{Guard}$، هر چه فراخواننده اعلام
کند، ردش می‌کند. تعریفِ~\ref{def:wfsys} نیمهٔ دیگرِ بازبودن را رد می‌کند،
یعنی عضوی که نامِ وابستگی‌ای را می‌برد که هیچ عضوی فراهمش نمی‌کند. اینجا
اولی است که کار می‌کند: نام‌بردن از $\Zenv$ $\WF$ را نمی‌لغزاند، چون
$\Zenv \in \pi^{+}$ چنین درایه‌ای را برآورده‌شدنی می‌کند، پس اگر آن میدان
فراخ‌تر بود، سیستمی می‌توانست خوش‌ساخت باشد و با این حال باز.

آن $\Zpid$ در $\Zenv.\admits$ به دلیلی دیگر آنجاست: می‌گذارد $\Zenv$
خودش را فرا بخواند. واسط‌های $\op{Z}$ِ محیط در طرف‌های مختلف جدایند، و
پذیرفتنِ $\Zpid$ همان است که می‌گذارد آن‌ها به یکدیگر برسند.

اینکه کدامشان می‌تواند به کدام برسد را مجموعهٔ خموش‌شده حل می‌کند، و آنجا
کاری بیش از دور نگه داشتنِ $\Zenv$ از سیستم‌ها می‌کند. هر ماشینِ دیگری به
جایی که می‌دود گره خورده است: کارکردِ استاندارد تنها هویتِ خودش را
می‌تواند ادعا کند، و $\Adv$ روی $\{\id'' \neq \id\}$ خموش است. هیچ‌چیز
$\Zenv$ را به‌تنهایی گره نمی‌زد، چون $\opl{Silence}$ تنها $\id'$ِ ادعایی
را آزمون می‌کند و $\opl{Guard}$ تنها هدف را، پس $(Z,P)$ و $(Z,P')$
جایگزین‌پذیر می‌شدند و مؤلفهٔ طرفِ یک هویتِ $\Zenv$ هیچ خبری حمل نمی‌کرد.
سطرِ~\ref{ln:zparty} این گره را می‌دهد، و تا آنجا که کاربردها می‌گذارند
شل: واسطی در $\id$ می‌تواند $\id'$ را ادعا کند هرگاه $\id'.P = \id.P$، یا
$\id'.P = Z$، یا $\id.P = Z$. به زبانِ ساده، \emph{برای طرفِ خودت عمل کن،
یا برای ریشه؛ و ریشه می‌تواند برای هر کسی عمل کند.}

هر دو راهِ فرار لازمند. بی $\id.P = Z$ ریشه در تنگنا می‌ماند، و ریشه همان
جایی است که $\opl{Exec}$ از آن آغاز می‌کند و تنها واسطی است که می‌تواند
به دیگران برسد. بی $\id'.P = Z$ برگ‌ها در تنگنا می‌مانند: تنها راهِ
بازگشتِ حریف به محیط $(A,P) \to (Z,P)$ است، و $\opl{Mediate}$ $\Adv$ را
در $(A,\id.P)$ بیدار می‌کند، برای طرفی که پاسخش را می‌دهد؛ پس محیطی که
نیاز داشته باشد آنچه حریف در $P$ گفته را با آنچه در $P'$ گفته بسنجد، هرگز
نمی‌توانست آن دو را کنارِ هم بیاورد --- این واسط‌ها جدایند، و نمی‌توانند
از راهِ حالتِ مشترک هم آن را یک‌کاسه کنند --- و ردهٔ تعریفِ~\ref{def:uc}
واقعاً کوچک می‌شد. با هر دو راهِ فرار، همبندی ستاره‌ای دوسویه است که در
$(Z,Z)$ ریشه دارد، و ما چیزی نمی‌شناسیم که محیطی نامقید بتواند بکند و این
ستاره نتواند. سطرِ~\ref{ln:zadv} در همین روح مقید می‌کند: جلوی ادعای
$(A,P)$ به‌دستِ $\Zenv$ را می‌گیرد، و از این راه جلوی گذشتنش از دروازهٔ
فسادِ سطرِ~\ref{ln:corrgate} و راندنِ واسطِ طرفِ \emph{درستکار} را ---
توانی که حریفی که او نقشش را بازی می‌کرد هم ندارد، چنان که
نکتهٔ~\ref{rem:advpars} در بابِ شبیه‌سازی که روی کمتری خموش شده ثبت
می‌کند. آنچه آن سطر دیگر لازم نیست بکند، دور نگه داشتنِ $\Zenv$ از
$\fopl{Corrupt}$ است، که سطرِ~\ref{ln:corronlya} او را زیرِ نامِ خودش به
آن می‌پذیرد.

$\Apid \cup \Zpid$ با هم همان است که کاربردها لازم دارند و نه بیشتر.
پذیرفتنِ $\Apid$ ضروری است چون محیط هیچ راهِ دیگری به اجرا ندارد: وقتی
$\Zenv$ ژتون را واگذارد، دوباره تنها زمانی می‌دود که چیزی فرایش بخواند، و
تنها حریف می‌تواند؛ و همان کانال است که $\Zenv$ از آن هر چه حریف بخواهد
بگویدش می‌شنود. پذیرفتنِ $\Zpid$ ضروری است چون محیط یک واسط به‌ازای هر
طرف است و نه یک ماشین، و بی آن ریشه به هیچ‌یک از آن‌ها نمی‌رسید. جز این
دو هیچ ضروری نیست: کارکردی که ناچار بود $\Zenv$ را فرا بخواند با آن مثلِ
زیرروالی رفتار می‌کرد که خودش تعریفش نکرده، و $\Corr$ فسادها را به
هیچ‌کس گزارش نمی‌دهد --- حریف می‌داند چه فاسد کرده، چون خودش کرده --- و
همین است که چرا $\Zenv$ باید مجموعهٔ فساد را \emph{بپرسد}. پس محیط تنها
از آن دو ماشینی در دسترس است که بیرونِ هر سیستمی‌اند.


\begin{interface}{محیطِ $\Zenv$ \\
\params{$\PID := Z$, \quad $\Ps :=  \bits \cup \{A,Z\}$, \quad $\admits := \Apid \cup \Zpid$,
\quad $\uses := \{(\Adv,\serves)\}$, \quad $\pars := \Ids$}}
\opsig{$\id.\fopdef{Z}(in)$} \quad از $\id'$
\begin{algorithmic}[1]
  \Req{$\opl{Guard}_{\Zenv}\bigl(\id,\id'\bigr)$}
  \State $\overline{\Ids} \gets \Ids \cup \{\id'' : \id''.F = A \}$ \label{ln:zadv}
    \Comment{$\Adv$ دور از دسترس}
  \State $\phantom{\overline{\Ids} \gets {}} \cup\ \{\id'' : \id.P \neq Z \ \wedge\ \id''.P \notin \{\id.P, Z\}\}$
    \label{ln:zparty}
    \Comment{طرفِ خود، یا ریشه}
  \State $out \gets \opl{Silence}\bigl(\overline{\Ids},\ \id.\op{Z}(in) \text{ from } \id' \bigr)$
  \State \textbf{return} $out$
\end{algorithmic}
\end{interface}


%% ==================================================================
\section{حریف}\label{sec:adv}

حریفِ $\Adv$ همان شکل را دارد: یک واسطِ یگانه به‌ازای هر $\id$، یعنی
$\id.\fopl{A}$، که هستهٔ $\id.\op{A}$ را نگهبانی و خموش می‌کند. آنچه او را
از محیط جدا می‌کند این است که چه کسی می‌تواند به او برسد. آنجا که
$\Zenv.\admits$ تنها $\Apid \cup \Zpid$ است، میدانِ $\Adv.\admits =
\Stdpid \cup \Apid \cup \Zpid$ هر شناسهٔ فرایندی را در بر دارد که هرگز
فراخوانی می‌گذارد. تنها شناسهٔ $\Corr$ نیست، و $\Corr$ هیچ فراخوانی
نمی‌گذارد. این فراخی یک خواسته است و نه یک راحتی: قلاب‌های واسطه‌گریِ
فصلِ~\ref{sec:wrappers}، و فراخوان‌هایی که کارکردهای ایده‌آل از راهِ
کوتاه‌نویسیِ $\Adv(\cdot)$ می‌گذارند، هر دو فراخوان‌هایی هستند که ماشینی
در حالِ عمل برای طرفی، روی حریف می‌گذارد.

این واسط‌ها همان چیزی هستند که دَرِ پشتی در این نمادگذاری چنین می‌نماید.
حریف ژتونی از خودش ندارد و نمی‌تواند فراخوانی را آغاز کند، پس هر حمله‌ای
که ترتیب می‌دهد از آنجا شروع می‌شود که ماشینِ دیگری مهار را به او
می‌سپارد: واسطِ یک طرف از راهِ $\opl{Mediate}$، یا مستقیماً محیط، که
اصلاً همان راهی است که محیط به حریف دستور می‌دهد. آنچه او در پاسخ
می‌تواند بگوید تنگ است: مجموعهٔ خموش‌شده $\{\id'' : \id'' \ne \id\}$ است،
پس واسطی از $\Adv$ تنها می‌تواند همان هویتی را ادعا کند که در آن می‌دود.

در درونِ واسطی با هویتِ $\id$، ما $\Adv(in)$ را برای فراخوانِ
$(A,\id.P).\fopl{A}(in)$ به‌عنوانِ $\id$ می‌نویسیم، و $\Zenv(in)$ را به
همین شکل. ادعا باید $\id$ باشد و نه هیچ‌چیزِ دیگر، چون
سطرِ~\ref{ln:call} هسته را روی هر چیزی جز هویتِ خودش خموش می‌کند؛ پس حریف
فراخواننده را به‌عنوانِ شناسهٔ ادعایی می‌آموزد و میدانِ جداگانه‌ای برایش
لازم ندارد. هستهٔ او به‌صورتِ $\id.\op{A}(in)$ از $\id'$ وارد می‌شود، همان
شکلی که $\opl{Silence}$ اعلام می‌کند و همان شکلی که هر هستهٔ دیگری در آن
وارد می‌شود. مؤلفهٔ طرفِ هدف $\id.P$ است، پس فراخوان این را هم ثبت می‌کند
که به‌نیابتِ چه کسی گذاشته شده است.

\begin{interface}{حریفِ $\Adv$ \\
\params{$\PID := A$, \quad $\Ps := \bits \cup \{A,Z\}$, \quad
$\admits := \Stdpid \cup \Apid \cup \Zpid$, \quad $\uses := \{(\Zenv,\serves)\}$, \quad
$\pars := \none$}}
\opsig{$\id.\fopdef{A}(in)$} \quad از $\id'$
\begin{algorithmic}[1]
  \Req{$\opl{Guard}_{\Adv}\bigl(\id,\id'\bigr)$}
  \State $out \gets \opl{Silence}\bigl(\{\id'' : \id'' \ne \id \} ,\ \id.\op{A}(in) \text{ from } \id' \bigr)$
    \label{ln:advsil}
  \State \textbf{return} $out$
\end{algorithmic}
\end{interface}

%% ==================================================================
\section{واسطِ کامل}\label{sec:fullinterface}

هر چه پوشانه لازم دارد اکنون سرِ جایش است: $\opl{Guard}$ و
$\opl{Mediate}$ و $\opl{Silence}$ از فصلِ~\ref{sec:wrappers}، و دفترِ
$\Corr$ از بخشِ~\ref{sec:corruption} که $\Cs$ را می‌دهد. واسطِ کاملِ یک
کارکردِ استاندارد ترکیبِ همین‌هاست.

\begin{interface}{کارکردِ $\F$ \\
\params{$\PID$, \quad $\Ps$, \quad $\admits$, \quad $\uses$, \quad $\pars$}}
\opsig{$\id.\fopdef{Op}(in)$} \quad از $\id'$
\begin{algorithmic}[1]
\Req{$\opl{Guard}_{\F}\bigl(\id,\id'\bigr)$} \label{ln:guard}
\State $\Cs \gets (C,\id.P).\fopl{Status}() \text{ as } \id$ \label{ln:stat}
\Req{$\neg \bigl(\id'.F = A \ \wedge\ \id.P \not\in \Cs\bigr)$}
  \label{ln:corrgate}
  \Comment{$\Adv$ تنها در طرفِ فاسد}
\State $\opl{Mediate}\bigl(\Cs,\ \id.\Op,\ in,\ \id'\bigr)$ \label{ln:medin}
\State $out \gets \opl{Silence}\bigl(\{\id'' : \id'' \ne \id \} ,\ \id.\Op(in) \text{ from } \id' \bigr)$ \label{ln:call}
\State $\Cs \gets (C,\id.P).\fopl{Status}() \text{ as } \id$ \label{ln:stattwo}
\State $\opl{Mediate}\bigl(\Cs,\ \id.\Op,\ out,\ \id'\bigr)$ \label{ln:medout}
\State \textbf{return} $out$ \label{ln:ret}
\end{algorithmic}
\end{interface}

سطرهای~\ref{ln:stat}--\ref{ln:medin} و~\ref{ln:stattwo}--\ref{ln:medout}
فساد را می‌گردانند، و تنها جایی هستند که واسطِ کامل با دفتر مشورت می‌کند.
دو بار خوانده می‌شود چون $\Cs$ با پیشرفتِ اجرا تغییر می‌کند و هر خواندن
تنها یک عکسِ فوری است: آزمون‌های پیش از سطرِ~\ref{ln:stattwo} اولی را به
کار می‌برند و آزمونِ سطرِ~\ref{ln:medout} دومی را. عضویت یکنواست، چون
$\opl{Corrupt}$ هرگز جز افزودن نمی‌کند، پس طرفی که در خواندنِ اول درستکار
است می‌تواند در خواندنِ دوم فاسد باشد، اما هرگز برعکس.

سطرِ~\ref{ln:corrgate} یک ترکیب را رد می‌کند: حریف در طرفِ درستکار، که
آنجا کاری ندارد. هر چیزِ دیگری می‌گذرد، و $\opl{Mediate}$ باقی را حل
می‌کند. پس طرفِ درستکار برای همه جز $\Adv$ خطاب‌پذیر است و هستهٔ خودش را
می‌دواند؛ و طرفِ فاسد برای همه خطاب‌پذیر است، اما تنها $\Adv$ به هسته‌اش
می‌رسد و هر فراخوانندهٔ دیگری را $\Adv$ به‌جای آن طرف پاسخ می‌دهد. محیط
حالتِ ویژهٔ خودش را لازم ندارد: چه هویتِ خودش را ادعا کند و چه هویتِ
بیرونیِ پذیرفته‌ای را، طرفِ فاسد در هر حال از راهِ $\Adv$ پاسخش می‌دهد.

می‌ماند طرفِ فاسدی که هر کسی جز $\Adv$ خطابش کند، که فراخوانش را نمی‌شود
همین‌طور دواند، پس $\Adv$ به‌جای آن طرف پاسخ می‌دهد.
سطرهای~\ref{ln:medin} و~\ref{ln:medout} همان یک پوشانه در دو لحظه‌اند.
اولی طرفی را می‌گیرد که در لحظهٔ ورود از پیش فاسد است و $in$ را واگذار
می‌کند، پس هسته هرگز نمی‌دود. دومی طرفی را می‌گیرد که در حینِ دویدنِ
هسته‌اش در سطرِ~\ref{ln:call} فاسد شد و $out$ را واگذار می‌کند، که دیگر
مالِ آن طرف نیست که برگرداندش. یکنوایی باعث می‌شود آن دو، حالت‌ها را
دقیقاً یک بار بپوشانند: طرفی که در اولی فاسد است پیش‌تر برگشته، و هیچ طرفی
در اولی فاسد و در دومی درستکار نیست. هر دو قلاب $\Adv$ را در $(A,\id.P)$
از راهِ واسطِ کاملش خطاب می‌کنند، پس حریفی که پاسخ می‌دهد مثلِ هر
فراخواندهٔ دیگری نگهبانی و خموش شده است. سطرِ~\ref{ln:ret} مقدارِ خودِ
هسته را برمی‌گرداند، که در دو حالت رخ می‌دهد: برای طرفی که سرتاسر درستکار
است، و برای طرفِ فاسدی که خودِ $\Adv$ خطابش می‌کند، که دروازه می‌پذیردش و
هیچ‌یک از دو قلاب نمی‌گیردش، چون هر دو $\id'.F \neq A$ را آزمون می‌کنند.

\begin{remark}[فسادِ میانِ فراخوان پیش‌دستانه نیست]\label{rem:midcall}
هیچ‌چیز هسته‌ای را که پیش‌تر در حالِ دویدن است قطع نمی‌کند. فرض کنید
هستهٔ سطرِ~\ref{ln:call} حریف را فرا بخواند --- از راهِ $\Adv(\cdot)$، یا
از راهِ زیرروالی که واسطه‌گری می‌کند --- و حریف پیش از برگرداندنِ مقداری،
$\id.P$ را فاسد کند. دفتر این را بی‌درنگ ثبت می‌کند، چون $\Cs$ حالتِ خودِ
اوست، اما هسته نه خبر می‌شود و نه متوقف: با هر چه حریف برگردانده از سر
می‌گیرد و تا پایان می‌دود. تنها در لحظهٔ بازگشتِ آن است که
سطرِ~\ref{ln:stattwo} $\Cs$ را از نو می‌خواند، $\id.P$ را در آن می‌یابد، و
$out$ را در سطرِ~\ref{ln:medout} به $\opl{Mediate}$ می‌سپارد تا جایگزین
شود.

پس حریف در طولِ چنین فراخوانی دو بار ژتون را در دست دارد، به این ترتیب:
یک بار در درونِ هسته، برای فراخوانی که هسته گذاشت، و یک بار پس از بازگشتِ
هسته، برای واسطه‌گریِ خروجیِ آن. هیچ لحظه‌ای در میانه نمی‌گذارد آن فساد
اثر کند، چون طرفی که در میانِ فراخوان فاسد می‌شود نمی‌تواند آنچه پیش‌تر
کرده را نکرده کند. فراخوان از نو هم دروازه‌بانی نمی‌شود:
سطرِ~\ref{ln:corrgate} در خواندنِ اول تصمیم گرفته شد، و فراخوانی که در
درستکاریِ $\id.P$ پذیرفته شده پذیرفته می‌ماند. پس آنچه پوشانه تضمین
می‌کند از سقط‌کردن ضعیف‌تر است، و بس است: هیچ‌چیزی که طرفِ تازه‌فاسد حساب
کرده باشد بی‌واسطه‌گری به فراخوانندهٔ او نمی‌رسد. فراخوان‌هایی که
\emph{پس از} فساد می‌رسند، به‌جایش در راهِ ورود گرفته می‌شوند، با
سطرِ~\ref{ln:medin}.
\end{remark}

%% ==================================================================
\section{اجراها}\label{sec:execution}

در یک اجرا، دقیقاً یک ماشین ژتونِ فعال‌سازی را در دست دارد. فراخوانِ $out
\gets \id.\Op(in)$ کوتاه‌نویسیِ این است که فراخواننده با یک ادامه معلق
شود و ژتون به واسطِ $\Op$ از $\F$ در هویتِ $\id$ سپرده شود. هر فراخوان
میانِ ماشین‌ها یک واسطِ کامل را خطاب می‌کند، و $\fopl{Op}$ همان است که
$\opl{Exec}$ توزیعش می‌کند؛ هسته تنها به‌دستِ پوشانه‌ای که خدمتش می‌کند
فراخوانده می‌شود، در درونِ ماشینی که حملش می‌کند، و هرگز از آن‌سوی یک
سپردنِ ژتون. فراخوانی که به هویتی خطاب شود که هیچ ماشینی اشغالش نکرده با
$\rej$ پاسخ می‌گیرد. یک نمونه می‌تواند چند ادامهٔ معلق داشته باشد و
می‌تواند در حالِ تعلیق دوباره وارد شود، به‌ویژه به‌دستِ $\Adv$.

\begin{definition}[ماشین‌ها و اشغال]\label{def:occ}
یک \emph{ماشین} هر چیزی است که اجرا می‌دواندش: اعضای سیستم، و هر چه اجرا
در کنارِ آن‌ها فراهم می‌کند. ماشین مجموعهٔ متناهی‌ای از هویت‌ها را
\emph{اشغال می‌کند} --- برای کارکرد، همان $\IDs(\F)$ از
تعریفِ~\ref{def:func} --- و $\opl{Exec}$ درست همان فراخوان‌هایی را به او
توزیع می‌کند که به هویت‌های اشغال‌شده‌اش خطاب شده‌اند. اشغال‌گرانِ یک اجرا
مجموعه‌های هویتِ دوبه‌دو مجزا حمل می‌کنند، و $\opl{Exec}$ برای هیچ
خانوادهٔ دیگری تعریف نشده است. \emph{اشغال‌گرِ جایگاهِ حریف} ماشینی است که
هویت‌های شناسهٔ فرایندِ $A$ را اشغال می‌کند و هیچ هویتِ دیگری را؛ حریفِ
بخشِ~\ref{sec:adv} یکی از آن‌هاست، و هر شبیه‌سازِ فصلِ~\ref{sec:uc} نیز.
یک \emph{فعال‌سازی} بازهٔ ماکسیمالی است که در طولِ آن یک ماشین ژتون را در
دست دارد.
\end{definition}

هر فراخوان یک شناسهٔ ادعاییِ $\id'$ حمل می‌کند. کارکردهای استاندارد تنها
هویتِ همان واسطی را می‌توانند ادعا کنند که فراخوان را می‌گذارد؛ $\Zenv$ و
$\Adv$ می‌توانند هویت‌های دیگری را ادعا کنند، مشروط به آزمون‌های
فصلِ~\ref{sec:wrappers}. آن آزمون‌ها همان نگهبان‌اند، که $\id'$ را در
برابرِ هویتِ $\id$ِ واسط و در برابرِ پارامترهای کارکرد می‌سنجد، و ما
آن‌ها را با گزاره‌های \textbf{require} مشخص می‌کنیم. نگهبانی که شکست
بخورد مقدارِ ممتازِ $\rej$ را به فرستنده برمی‌گرداند، و $\rej$ یک رد است
و نه یک مقدار: هیچ ماشینی نمی‌تواند آن را به‌عنوانِ خروجیِ خودش
برگرداند. این قاعده اِعمال می‌شود و مفروض گرفته نمی‌شود --- هسته‌ای که
کدش $\rej$ برمی‌گرداند، و پاسخِ واسطه‌گری‌شدهٔ $\rej$ از جایگاهِ حریف،
به‌جایش به‌صورتِ $\none$ تحویل داده می‌شوند --- پس فراخواننده‌ای که $\rej$
می‌گیرد می‌آموزد که نگهبانی فراخوانش را رد کرده، و هیچ‌چیزِ دیگری
نمی‌تواند یکی بسازد. به‌ویژه $\rej \neq \none$.

تصادف به‌ازای هر ماشین است: هر ماشین نوارِ سکهٔ خودش را حمل می‌کند،
نامتناهی، یکنواخت، و مستقل از نوارِ هر ماشینِ دیگر، که تنها وقتی کدش
نمونه‌برداری می‌کند خوانده می‌شود. احتمال‌ها روی یک اجرا، روی این نوارها
به‌طورِ توأم‌اند. دسته‌کردنِ ماشین‌ها در یک اشغال‌گر --- چنان که
فصلِ~\ref{sec:uc} و بخشِ~\ref{app:dummy} می‌کنند --- نوارها را یکی
نمی‌کند: سکه‌های یک رونوشتِ میهمان همان‌اند که پیش از جذب بودند.

جز سیستمِ $\pi$، اجرا سه ماشین حمل می‌کند: محیطِ $\Zenv$، حریفِ $\Adv$، و
دفترِ فسادِ $\Corr$، که حالتش $\Cs$ را هر پوشانه‌ای می‌خواند. بنا بر
فصلِ~\ref{sec:functionalities} هیچ‌یک عضوِ $\pi$ نیست، و همین
سطرِ~\ref{ln:execinit} را یکه‌معنا می‌کند --- دفتر یک بار مقدارِ آغازین
می‌گیرد، به‌عنوانِ آرگومان، و نه بارِ دوم به‌عنوانِ عضو --- و $\pi$ را
همان شیئی نگه می‌دارد که طراح نوشت. رویه آن سه را مقدارِ آغازین می‌دهد، و
سپس محیط را در $(Z,Z)$ فعال می‌کند، چنان که گویی خودش را فرا خوانده است.
آن نخستین فعال‌سازی مثلِ هر فعال‌سازیِ دیگری از واسطِ کاملِ $\Zenv$
می‌گذرد، پس محیط از همان آغاز خموش است، و $\opl{Guard}_{\Zenv}$
به‌سادگی می‌گذرد، چون هویتِ ادعایی و هدف یکی‌اند. هر چه محیط برگرداند
همان برآمدِ کار است.

\begin{procedure}{اجرای $\opdef{Exec}$ \\
\params{سیستمِ $\pi$، \quad محیطِ $\Zenv$، \quad حریفِ $\Adv$، \quad دفترِ فسادِ $\Corr$}}
\opsig{$\opl{Exec}(\pi,\Zenv,\Adv,\Corr)$}:
\settowidth{\clausewd}{$\id.\fopl{Op}(in)\ \text{from}\ \id' \,;\ k$}
\settowidth{\reswd}{$k\bigl(\id.\fopl{Op}(in)\ \text{from}\ \id'\bigr)$}
\begin{algorithmic}[1]
\State $\Corr.\op{Initialize}()$; \ $\F.\op{Initialize}()$ for every $\F \in \pi$ \label{ln:execinit}
\State \textbf{handle}\ \ $(Z,Z).\fopl{Z}()$ from $(Z,Z)$\ \ \textbf{with} \label{ln:execstart}
  \Comment{$\Zenv$ نخست فعال می‌شود و ژتون را دارد}
\State \quad \makebox[\clausewd][l]{$\textbf{return}\ out$}
       $\ \mapsto\ $ \makebox[\reswd][l]{$out$}
\State \quad \makebox[\clausewd][l]{$\id.\fopl{Op}(in)\ \text{from}\ \id' \,;\ k$}
       $\ \mapsto\ $ \makebox[\reswd][l]{$k\bigl(\id.\fopl{Op}(in)\ \text{from}\ \id'\bigr)$} \label{ln:execcall}
\end{algorithmic}
\end{procedure}

اگر آن را به‌عنوانِ هندلر بخوانیم~\cite{PP09,BP15}، اجرا دقیقاً یک عمل
دارد: گذاشتنِ فراخوان. سطرِ~\ref{ln:execstart} محیط را فعال می‌کند، و این
تنها فعال‌سازی‌ای است که ماشینی آن را نگذاشته است.
سطرِ~\ref{ln:execcall} از آن پس هر فراخوانی را می‌گرداند، و ادامهٔ $k$
همان جایی است که ژتون به آن بازمی‌گردد: معلق‌کردنِ فراخواننده، سپردنِ
ژتون به $\id$، و از سر گرفتنِ فراخواننده با پاسخ، درست همان معنایی است که
فراخواندنِ $k$ روی آن پاسخ دارد. دقیقاً یک ماشین در سرتاسر ژتون را در دست
دارد. بندِ \textbf{return} وقتی رسیده می‌شود که برنامهٔ محیط بایستد، و
خروجیِ آن خروجیِ اجراست.

%% ==================================================================

%% ==================================================================
\chapter{تقلید و ترکیبِ UC}\label{sec:uc}

\section{جانشانیِ زیرسیستم}\label{sec:repl}

پروتکل را معمولاً در برابرِ زیرروالی می‌نویسند که بعداً پروتکلِ دیگری
فراهمش می‌کند. جانشین‌کردنِ یکی به‌جای دیگری، عملی مجموعه‌ای روی
سیستم‌هاست.

\begin{definition}[جانشانیِ زیرسیستم]\label{def:repl}
بگذارید $\rho$ سیستمی باشد و $\varphi \subseteq \rho$ زیرسیستمی از آن.
برای سیستمِ $\pi$، \emph{جانشانیِ $\varphi$ با $\pi$ در $\rho$} این است:
\[
  \rho[\pi/\varphi] \;:=\; (\rho \setminus \varphi) \cup \pi .
\]
\end{definition}

تا اینجا هیچ‌چیز جلوی $\pi$ِ ورودی را نمی‌گیرد که روی شناسهٔ فرایندی فرود
بیاید که پروتکلِ پیرامون از پیش به کارش می‌برد، و این فراخوانی را با
خطابی دوپهلو رها می‌کرد. ما این را ممنوع می‌کنیم و به تعمیرش نمی‌رویم.

\begin{definition}[جانشانیِ خوش‌ساخت]\label{def:wf}
جانشانیِ $\varphi$ با $\pi$ در $\rho$ \emph{خوش‌ساخت} است اگر هیچ شناسهٔ
فرایندی هم در $\rho \setminus \varphi$ و هم در $\pi$ رخ ندهد، یعنی اگر
\[
  \{\, \F.\PID : \F \in \rho \setminus \varphi \,\} \;\cap\; \{\, \F.\PID : \F \in \pi \,\} \;=\; \emptyset .
\]
\end{definition}

\begin{proposition}[جانشانی]\label{prop:repl}
بگذارید $\rho$ و $\pi$ سیستم باشند، $\varphi \subseteq \rho$، و جانشانیِ
$\varphi$ با $\pi$ در $\rho$ خوش‌ساخت باشد. آن‌گاه $\rho[\pi/\varphi]$ یک
سیستم است. و اگر علاوه بر آن $\IDs(\pi) = \IDs(\varphi)$، آن‌گاه
$\IDs(\rho[\pi/\varphi]) = \IDs(\rho)$.
\end{proposition}

\begin{proof}
شناسه‌های فرایندِ $\rho[\pi/\varphi]$ همان‌های $\rho \setminus \varphi$ به
همراهِ همان‌های $\pi$ هستند. هر سیستم درونِ خودش شناسه‌های متمایز حمل
می‌کند و خوش‌ساختی آن دو گردایه را مجزا می‌کند، پس هیچ شناسه‌ای تکرار
نمی‌شود؛ اعضا استاندارد و به‌شمارِ متناهی‌اند، که هر دو از $\rho$ و $\pi$
به ارث رسیده، پس $\rho[\pi/\varphi]$ یک سیستم است. آن ادعای اول تنها بر
خوش‌ساختیِ جانشانی تکیه دارد و $\IDs(\pi) = \IDs(\varphi)$ هیچ کجای آن
وارد نمی‌شود. هویت‌ها فرضِ جداگانه‌ای لازم ندارند: هویت شناسهٔ فرایندش را
در مؤلفهٔ اول حمل می‌کند، پس اعضایی با شناسه‌های متمایز مجموعه‌های هویتِ
مجزا دارند. و برای خودِ هویت‌ها،
\[
  \begin{aligned}
  \IDs(\rho[\pi/\varphi]) &\;=\; \IDs(\rho \setminus \varphi) \cup \IDs(\pi) \\
  &\;=\; \IDs(\rho \setminus \varphi) \cup \IDs(\varphi) \;=\; \IDs(\rho),
  \end{aligned}
\]
که در گامِ میانی $\IDs(\pi) = \IDs(\varphi)$ و در گامِ آخر $\varphi
\subseteq \rho$ را به کار می‌برد.
\end{proof}

دو نیمهٔ این گزاره در زیر به دو شکلِ متفاوت به کار می‌روند. اینکه
$\rho[\pi/\varphi]$ یک سیستم است می‌گذارد فصلِ~\ref{sec:uc} اصلاً از آن
سخن بگوید، و تنها بر خوش‌ساختی تکیه دارد. اینکه هویت‌هایش همان‌های $\rho$
است $\IDs(\pi) = \IDs(\varphi)$ را لازم دارد، و هیچ نتیجه‌ای در
فصلِ~\ref{sec:uc} آن را مفروض نمی‌گیرد: مجازبودن در آنجا برای زوجی از
سیستم‌ها و از راهِ اجتماعِ مجموعه‌های هویتشان بیان می‌شود، و
گزارهٔ~\ref{prop:absorb} و قضیهٔ~\ref{thm:comp} رو می‌گویند که هیچ نسبتی
میانِ آن دو به کار نرفته است. هر جا آن تساوی خواسته باشد، مسدودسازی
فراهمش می‌کند.

تعریفِ~\ref{def:wf} شناسه‌های فرایند را حل می‌کند و
گزارهٔ~\ref{prop:repl} هویت‌ها را. هیچ‌یک به وابستگی نمی‌پردازد، و
وابستگی همان جایی است که جانشانی می‌تواند بد پیش برود: نام‌هایی که $\rho
\setminus \varphi$ به آن‌ها تکیه دارد پس از آنکه $\varphi$ جایش را به
$\pi$ بدهد هنوز اشغال‌شده‌اند، اما هر خواسته در اشغال‌گرِ تازه از نو
سنجیده می‌شود. دو شرط این شکاف را می‌بندند.

به‌زبانِ غیررسمی، قضیهٔ~\ref{thm:wfrepl}ِ زیر می‌گوید جانشانیِ خوش‌ساخت
$\WF$ را نگه می‌دارد، هرگاه درایه‌های $\pi$ِ ورودی در جهانِ تازه
برآورده‌شدنی باشند، و هر جا درایه‌ای از $\rho \setminus \varphi$ شاهدش را
در درونِ $\varphi$ می‌یافت، عضوی از $\pi$ به‌جای $\varphi$ بایستد. فرضِ
(ii) همان معنای این است که $\pi$ جانشینِ آمادهٔ $\varphi$ باشد، و تقلید
آن را فراهم نمی‌کند.

\begin{theorem}[جانشانی خوش‌ساختی را نگه می‌دارد]\label{thm:wfrepl}
بگذارید $\rho$ و $\pi$ سیستم باشند، $\varphi \subseteq \rho$، جانشانیِ
$\varphi$ با $\pi$ در $\rho$ به معنای تعریفِ~\ref{def:wf} خوش‌ساخت باشد، و
$\WF(\rho)$. فرض کنید
\begin{enumerate}[label=(\roman*),leftmargin=2.2em]
  \item هر درایهٔ هر $\F \in \pi$ شاهدی در $\bigl(\rho[\pi/\varphi]\bigr)^{+}$ دارد،
        به معنای تعریفِ~\ref{def:wfsys}؛ و
  \item هرگاه $\F \in \rho \setminus \varphi$ درایه‌ای $(\F',R) \in \F.\uses$ داشته باشد که
        \emph{برخی} شاهدش به معنای تعریفِ~\ref{def:wfsys} در $\varphi$ باشد، آن‌گاه برخی
        $\G \in \pi$ در $\G.F = \F'.F$، $\ \G.s = \F.s \vee \op{global}(\G)$ و $R(\F,\G)$
        صادق باشند.
\end{enumerate}
آن‌گاه $\WF\bigl(\rho[\pi/\varphi]\bigr)$.
\end{theorem}

\begin{proof}
بنویسیم $\rho' := \rho[\pi/\varphi] = (\rho \setminus \varphi) \cup \pi$،
که بنا بر ادعای اولِ گزارهٔ~\ref{prop:repl} یک سیستم است --- ادعای دوم
اینجا در دست نیست، چون $\IDs(\pi) = \IDs(\varphi)$ فرضِ این قضیه نیست.
$\F \in \rho'^{+}$ و درایهٔ $(\F',R) \in \F.\uses$ را بگیرید؛ باید شاهدی
در $\rho'^{+}$ بنشانیم. بر حسبِ اینکه $\F$ از کجا می‌آید، سه حالت هست.

اگر $\F \in \pi$، فرضِ (i) درست همان ادعاست.

اگر $\F$ یکی از $\Zenv$، $\Adv$، $\Corr$ باشد، درایه‌هایش یک‌بار برای
همیشه و مستقل از سیستم حل شده‌اند: $\Corr.\uses$ تهی است، و درایه‌های
$\Zenv$ و $\Adv$ نامِ یکدیگر را می‌برند، با $\Adv, \Zenv \in \rho'^{+}$ و
$\serves$ که بینشان برقرار است، چنان که فصلِ~\ref{sec:functionalities}
ثبت می‌کند. هر سه جهانی‌اند، پس عطفِ نشست بخشیده می‌شود.

اگر $\F \in \rho \setminus \varphi$، آن‌گاه $\F \in \rho^{+}$، پس
$\WF(\rho)$ شاهدی $\G \in \rho^{+}$ برای این درایه فراهم می‌کند. یا $\G
\notin \varphi$، که در آن صورت $\G$ در $(\rho \setminus \varphi) \cup
\{\Zenv,\Adv,\Corr\} \subseteq \rho'^{+}$ است و دوباره خدمت می‌کند: همان
ماشین است با همان میدان‌ها، پس $\G.F$، $\G.s$ و هر چه $R$ می‌خواند تغییر
نکرده، و $\F$ هم تغییر نکرده است. یا $\G \in \varphi$، و فرضِ (ii) عضوی از
$\pi$ فراهم می‌کند که همان سه شرط را به‌جای او برمی‌آورد. در هر دو حالت
درایه شاهدی در $\rho'^{+}$ دارد.
\end{proof}

فرضِ (ii) همان معنای این است که $\pi$ جانشینِ آمادهٔ $\varphi$ باشد، و
تقلید آن را فراهم نمی‌کند. تعریفِ~\ref{def:uc} رفتار را در مرز مقایسه
می‌کند؛ در بابِ میدان‌های اعلام‌شدهٔ ماشین‌های پشتِ آن مرز هیچ نمی‌گوید، و
$\pi$ای که $\varphi$ را بی‌نقص تقلید کند می‌تواند باز هم در پذیرفتنِ
فراخواننده‌ای که $\varphi$ می‌پذیرفت، یا خدمت‌دادن به طرفی که $\varphi$
خدمت می‌داد، شکست بخورد. این آزمون ارزان است --- روی درایه‌هایی از $\rho
\setminus \varphi$ دامنه دارد که به درونِ $\varphi$ اشاره می‌کردند و هیچ
چیزِ دیگر --- اما آزمون است، و جایش کنارِ خوش‌ساختیِ جانشانی است و نه در
درونِ قضیهٔ ترکیب.

%% ==================================================================
%% The chapter body needs a heading of its own now that a section
%% precedes it: replacement has to be defined before the composition
%% theorem that substitutes with it.
\section{تقلید و جذب}\label{sec:emabs}


هر چه اجرا لازم دارد اکنون تعیین شده است، پس می‌توانیم بگوییم یک سیستم کِی
دستِ‌کم به‌اندازهٔ دیگری امن است. آن دو زیرِ یک و همان محیط نشانده می‌شوند،
پس نخست باید تعیین شود که او چه چیزی را می‌تواند ادعا کند؛ و این نقشِ
$\Zenv.\pars$ در زیر است.

مقدارِ $\opl{Exec}(\pi,\Zenv,\Adv,\Corr)$ هر چیزی است که محیط برمی‌گرداند.
ما محیط را چنان می‌گیریم که در ریشه یک بیت برمی‌گرداند --- حدسش در بابِ
اینکه در کدام سیستم نشانده شده --- و اجرا را چنان می‌خوانیم که دقیقاً
وقتی $1$ برمی‌گرداند که محیط برگرداند. هیچ‌چیزِ دیگری در چارچوب آن مقدار
را مقید نمی‌کند، پس این شرطی است بر محیط‌هایی که روی آن‌ها سور می‌زنیم، و
نه نتیجه‌ای. آن‌گاه شکافِ میانِ دو اجرا یک عدد است، و نه یک شباهتِ
غیررسمی.

\begin{definition}[مزیتِ UC]\label{def:adv}
بگذارید $\pi$ و $\varphi$ سیستم باشند. برای محیطِ $\Zenv$، حریفِ $\Adv$ و
شبیه‌سازِ $\Sim$،
\[
  \begin{gathered}
  \advtg^{\op{uc}}_{\pi,\varphi}(\Zenv,\Adv,\Sim)
  \;:=\;
  \Bigl|\, \Pr\bigl[\opl{Exec}(\pi,\Zenv,\Adv,\Corr) = 1\bigr] \\[3pt]
  \;-\; \Pr\bigl[\opl{Exec}(\varphi,\Zenv,\Sim,\Corr) = 1\bigr] \,\Bigr| ,
  \end{gathered}
\]
که احتمال‌ها روی سکه‌های هر ماشینِ آن دو اجرا هستند، و $\opl{Exec} = 1$
کوتاه‌نویسیِ پیشامدی است که اجرا با برگرداندنِ $1$ از سوی محیط می‌ایستد.
\end{definition}

آن دو اجرا در سیستم و در ماشینی که هویت‌های حریف را اشغال می‌کند تفاوت
دارند: $\Adv$ در یک سو، $\Sim$ در سوی دیگر. هر دو با همان یک جفت تعریف
نوع‌بندی می‌شوند.

\begin{definition}[حریف‌ها]\label{def:advcls}
یک \emph{حریف} ماشینی است که $\IDs(\Adv)$ را اشغال می‌کند و هیچ هویتِ
دیگری را، و واسط‌های حریف را حمل می‌کند --- پوشانهٔ بخشِ~\ref{sec:adv}،
نگهبان و خموش‌کننده، و هر واسطِ دیگری که به ماشینِ آن بخش داده شده
(فصلِ~\ref{app:responsive} یکی می‌افزاید) --- به دورِ هستهٔ دلخواهی در هر
یک از آن‌ها. ماشینِ بخشِ~\ref{sec:adv} یکی از آن‌هاست؛ بدلِ
بخشِ~\ref{app:dummy} دیگری.
\end{definition}

\begin{definition}[شبیه‌سازها]\label{def:sim}
یک \emph{شبیه‌ساز} حریفی است که دستِ‌کم $\Stdpid \cup \Zpid$ را می‌پذیرد
--- شناسه‌های فرایندِ هر کارکردِ استاندارد به همراهِ شناسهٔ محیط.
\end{definition}

پس هر شبیه‌سازی حریف است، و همین است که می‌گذارد یکی از آن‌ها در ردهٔ
حریف بایستد: گزارهٔ~\ref{prop:trans} و قضیهٔ~\ref{thm:parcomp}
شبیه‌سازها را به‌عنوانِ حریف می‌خوانند، و شمول‌های بودجه‌ای
بخشِ~\ref{sec:concrete} آن دو رده را مستقیم می‌سنجند. اشغال و
فراخواننده‌های پذیرفته همان معنای ایستادن در جایگاهِ حریف‌اند، پس هر چه
بخشِ~\ref{sec:adv} در بابِ ماشینِ نشسته در $(A,P)$ حل می‌کند بی‌تغییر
برای شبیه‌ساز هم برقرار است: هیچ فراخوانی از خودش آغاز نمی‌کند، و
خموش‌کننده‌اش تنها می‌گذارد هویتی را ادعا کند که در آن می‌دود.
گزارهٔ~\ref{prop:regroup} تنها اشغالِ هر چه به آن داده شود را می‌خواهد؛
قضیهٔ~\ref{thm:dummy} پوشانه و فراخواننده‌های پذیرفته را هم به کار
می‌برد، و نکتهٔ~\ref{rem:advpars} توضیح می‌دهد چرا هیچ‌یک محدودیت نیست.

متر آن سه ماشین را به‌عنوانِ آرگومان می‌گیرد، پس تقلید را می‌شود با
سور‌زدن روی رده‌هایی از آن‌ها و کران‌گذاشتن بر نتیجه بیان کرد. اینکه
اصلاً چه محیط‌هایی به کار می‌آیند را همان زوجی که مقایسه می‌شود حل
می‌کند.

\begin{remark}[تقلید در بابِ جایگاهِ حریف چه چیزی را تعیین می‌کند]\label{rem:advpars}
تعریفِ~\ref{def:admis}ِ زیر $\Zenv.\pars$ را مقید می‌کند و در بابِ
ماشین‌های جایگاهِ حریف هیچ نمی‌گوید؛ تعریف‌های~\ref{def:advcls}
و~\ref{def:sim} هر چه لازم است می‌گویند، و آن کمتر از پنج میدان است.
میدان‌های $\PID$ و $\Ps$ با اشغال مفت می‌آیند: $\IDs(\F) = \{(\F.\PID,P)
: P \in \F.\Ps\}$، پس اشغالِ $\IDs(\Adv)$ هر دو را سنجاق می‌کند، و
$\Sim.\PID = (A,0,0)$ در سطرِ~\ref{ln:interface} و $\op{global}(\Sim)$ در
سطرِ~\ref{ln:session} را می‌دهد؛ و سنجاق‌کردنِ $\Ps$ همان است که خواسته
است، چون اشغال‌گرِ جایگاهی که به طرف‌های کمتری خدمت دهد یک هویتِ $A$ را
تنها در یک سو اشغال‌نشده می‌گذارد، که در آن سو $\rej$ پاسخ می‌گیرد و در
سوی دیگر پاسخی واقعی. میدانِ $\pars$ هرگز خوانده نمی‌شود --- مجموعهٔ
خموش‌شده همان $\{\id'' : \id'' \neq \id\}$ِ عینیِ سطرِ~\ref{ln:advsil}
است، که پوشانهٔ لازم آن را به‌عنوانِ کد حمل می‌کند --- و این تشریفات
نیست: شبیه‌سازی که روی \emph{کمتری} خموش شده باشد می‌توانست هویتِ یک طرف
را ادعا کند، از دروازهٔ فسادِ $\id'.F = A$ در سطرِ~\ref{ln:corrgate}
بگذرد، و واسطِ طرفِ درستکاری را براند، توانی که حریفی که او به‌جایش
می‌ایستد ندارد. میدانِ $\uses$ بی‌اثر است و تنها $\WF$ می‌خواندش، و آنجا
$\pi^{+}$ نامِ $\Adv$ِ رزرو را می‌برد و نه اشغال‌گرِ جایگاه را.

می‌ماند $\admits$، و همین است که چرا تعریفِ~\ref{def:sim} آن را رو
می‌خواهد. فراخ‌ترکردن ناممکن است، چون $\Adv.\admits$ از پیش هر شناسهٔ
فرایندی را در بر دارد که فراخوانی می‌گذارد؛ و تنگ‌ترکردن هیچ برای
شبیه‌ساز نمی‌آورد و هر دو قضیهٔ زیر را می‌گیرد. حذفِ $\Stdpid$ هر قلابِ
واسطه‌گری در طرفِ فاسد را به همان $\none$ِ محسوبِ
بخشِ~\ref{sec:execution} بدل می‌کند --- و قضیهٔ~\ref{thm:comp} روی هر
$\rho$ سور می‌زند، پس شبیه‌سازی که تنها شناسه‌های فرایندِ $\varphi$ را
بپذیرد در نخستین جهانِ بیرونی که شناسه‌های دیگری داشته باشد با آن‌ها
روبه‌رو می‌شود. حذفِ $\Zpid$ محیط را تنها در سوی ایده‌آل بیرون می‌گذارد، و
سطرِ~\ref{ln:caller} در جایگاه، درست همان دستورهایی را رد می‌کند که حریفِ
واقعی پاسخشان می‌دهد. و هیچ‌چیز در هیچ‌یک از این‌ها از زوجی که برایش
شبیه‌سازی می‌شود نام نمی‌برد، و همین است که می‌گذارد
قضیهٔ~\ref{thm:comp} شبیه‌سازِ فرض را بی‌تغییر تحویل دهد.
\end{remark}

\begin{definition}[محیط‌ها]\label{def:env}
یک \emph{محیط} ماشینی از اجراست که $\IDs(\Zenv)$ را اشغال می‌کند به همراهِ
هویت‌های تعدادِ متناهی کارکردِ استاندارد که \emph{میهمان} می‌کند --- برای
ماشینِ بخشِ~\ref{sec:env} هیچ، و برای محیط‌های جذب‌شدهٔ
تعریفِ~\ref{def:absorb}ِ زیر، رونوشت‌هایی از پروتکلِ پیرامون. این ماشین
واسطِ بخشِ~\ref{sec:env} را، با خموش‌کننده و پارامترِ $\pars$، در هر
هویتِ $Z$ حمل می‌کند، و واسطِ کاملِ خودِ هر کارکردِ میهمان را در
هویت‌های آن کارکرد؛ $\opl{Exec}$ در هر هویتی که او اشغال می‌کند به او
توزیع می‌کند، و خروجی‌اش همان است که بخشِ $\Zenv$ِ او برمی‌گرداند.
\end{definition}

\begin{definition}[محیط‌های مجاز]\label{def:admis}
برای سیستمِ $\pi$، بنویسیم $\ncl{\pi} := \{\, \id \;:\; \id.\PID.F = \F.F
\ \text{for some}\ \F \in \pi \,\}$ برای \emph{بستارِ نامیِ} آن: هویت‌های
هر شناسهٔ فرایندی که نامی را حمل می‌کند که $\pi$ به کار می‌برد، اشغال‌شده
یا نه، در هر نشست و هر نمونه. برای سیستم‌های $\pi$ و $\varphi$، محیط‌های
\emph{مجاز} برای این زوج این‌هایند:
\[
  \begin{gathered}
  \ZenvSet_{\pi,\varphi} \;:=\; \bigl\{\, \Zenv \;:\; \Zenv.\pars \supseteq \ncl{\pi} \cup \ncl{\varphi}
  \ \text{ و} \\[2pt]
  \Zenv \text{ در هیچ هویتی از } \IDs(\pi) \cup \IDs(\varphi) \text{ میهمان نمی‌کند} \,\bigr\} .
  \end{gathered}
\]
\end{definition}

شرطِ دوم بهداشتِ اشغال است: اجرا تنها برای اشغال‌گرانی با مجموعه‌های هویتِ
دوبه‌دو مجزا تعریف شده است (تعریفِ~\ref{def:occ})، پس محیطی که در هویتی
میهمان کند که یکی از سیستم‌ها اشغالش کرده نمی‌توانست کنارِ آن نشانده شود، و
مجازبودن آن را برای هر دو جهان یک‌جا ممنوع می‌کند.

شمول در زیر توجیه می‌شود. آن‌گاه یک حکمِ تقلید سه رده را تعیین می‌کند: ردهٔ
$\Aset$ از حریف‌ها (تعریفِ~\ref{def:advcls})، ردهٔ $\SimSet$ از
شبیه‌سازها (تعریفِ~\ref{def:sim})، و ردهٔ $\ZenvSet \subseteq
\ZenvSet_{\pi,\varphi}$ از محیط‌ها (تعریفِ~\ref{def:env}). گرفتنِ
$\ZenvSet = \ZenvSet_{\pi,\varphi}$ قوی‌ترین حکم از این سه را می‌دهد، و
خوانشِ پیش‌فرض است.

\begin{definition}[تقلیدِ UC]\label{def:uc}
بگذارید $\pi$ و $\varphi$ سیستم باشند، $\varepsilon \geq 0$، و $\ZenvSet
\subseteq \ZenvSet_{\pi,\varphi}$. می‌گوییم $\pi$ کارکردِ $\varphi$ را
\emph{در $\varepsilon$ و در برابرِ} $(\Aset,\SimSet,\ZenvSet)$
\emph{UC-تقلید می‌کند} اگر
\[
  \forall\, \Adv \in \Aset \ \ \exists\, \Sim \in \SimSet \ \ \forall\, \Zenv \in \ZenvSet \ : \qquad
  \advtg^{\op{uc}}_{\pi,\varphi}(\Zenv,\Adv,\Sim) \;\leq\; \varepsilon .
\]
\end{definition}

ترتیبِ سورها جانِ این تعریف است: شبیه‌ساز می‌تواند به حریف وابسته باشد اما
به محیط نه، پس یک $\Sim$ باید یک‌جا در برابرِ هر $\Zenv \in \ZenvSet$ کار
کند. جز این، آن دو اجرا یک آزمایش‌اند که یک سیستم در آن عوض شده، و $\Zenv$
عیناً همان یک ماشین در هر دو سو است. ضعیف‌کردنش به $\forall\, \Adv\
\forall\, \Zenv\ \exists\, \Sim$ --- شبیه‌سازی که برای محیطی که با آن
روبه‌روست ساخته شده --- مفهومی می‌دهد که تعریفِ~\ref{def:uc} آن را نتیجه
می‌دهد و عکسش را ما اثبات نمی‌کنیم، و آن دو در صورتِ عینیِ
بخشِ~\ref{sec:concrete} به‌صورتِ $\max_{\Adv} \min_{\Sim} \max_{\Zenv}$ و
$\max_{\Adv} \max_{\Zenv} \min_{\Sim}$ از یک کمیت خوانده می‌شوند. نتایجِ
ترکیب از این ضعیف‌سازی جان سالم می‌برند، چون هر اثباتِ زیر فرض را در
محیطی می‌نشاند که نتیجه تعیینش کرده؛ پس آنچه این ترتیب می‌خرد معناست و نه
ترکیب، چون تنها زیرِ آن است که $(\varphi,\Sim)$ یک سیستم است که مثلِ
$(\pi,\Adv)$ رفتار می‌کند، هر که تماشا کند --- جهانی ایده‌آل که می‌شود
نامش برد و تحویلش داد.

تعریفِ~\ref{def:admis} همان است که این مقایسه را معنادار می‌کند. چون
$\Zenv.\pars$ خودِ $\Zenv$ را خموش می‌کند، او هیچ هویتی از درونِ آن را
نمی‌تواند ادعا کند، و سه دسته باید دور از دسترس بمانند. هویتی در
$\IDs(\pi) \mathbin{\triangle} \IDs(\varphi)$ در یک جهان اشغال است و در
دیگری آزاد، پس به محیطی که اجازه یابد به‌عنوانِ آن سخن بگوید پیش از آنکه
چیزی رخ دهد گفته می‌شود در کدام جهان است، و هیچ شبیه‌سازی نمی‌تواند این
را تعمیر کند. هویتی که درونیِ \emph{هر دو} است چیزی را جدا نمی‌کند، اما
محیطی که آن را ادعا کند به‌عنوانِ جزئی از سیستم سخن می‌گوید و نه از
بیرون، و محیط برای این نیست. و هویتِ \emph{اشغال‌نشده}ٔ یک نامِ به‌کاررفته
به‌اندازهٔ اشغال‌شده کارساز است، چون $\opl{Guard}$ فراخواننده‌ها را با نام
می‌پذیرد: فراخواندهٔ محلی هر نمونه از نامِ والدش را می‌پذیرد و
سطرِ~\ref{ln:caller} بیش از این نمی‌خواهد، پس محیطی که نمونهٔ تازه‌ای از
آن نام را ادعا کند زیرروال‌های درونی را درست همان‌گونه می‌راند که
فراخواننده‌هایشان --- و در برابرِ زوجی مسدودشده
(فصلِ~\ref{sec:blockers}) رو جهان‌ها را از هم تشخیص می‌دهد، چون مسدودگر
فراخوانی را رد می‌کند که اشغال‌گرِ اصیل می‌پذیرد. خموش‌کردنِ بستار هر سه
را یک‌جا ممنوع می‌کند. آنچه ادعاکردنی می‌ماند نام‌های \emph{تازه} است، و
آن‌ها تنها به چیزی می‌رسند که کلِ $\Stdpid$ را می‌پذیرد: کارکردهای جهانی،
که در هر مقایسه‌ای که خودشان محلِ نزاع نباشند بینِ دو جهان مشترکند.

خموش‌سازی حکم می‌کند که $\Zenv$ چه چیزی را می‌تواند \emph{ادعا} کند، و
فراخواندن کانالِ دومی است که به آن دست نمی‌زند. هویتی از $\IDs(\pi)
\mathbin{\triangle} \IDs(\varphi)$ در هر دو جهان می‌تواند فرا خوانده شود:
آنجا که هیچ‌چیز اشغالش نکرده پاسخ $\rej$ است
(بخشِ~\ref{sec:execution})، و آنجا که چیزی اشغالش کرده پاسخ باز هم $\rej$
است --- به شرطی که آن چیز $\Zpid$ را در سطرِ~\ref{ln:caller} رد کند. این
شرط، شرطی است بر سیستم‌هایی که مقایسه می‌شوند و نه بر ردهٔ محیط، و زوجی
که آن را بشکند به‌سادگی تقلیدشدنی نیست، چون هیچ شبیه‌سازی نمی‌تواند پاسخی
از خود بسازد که جهانِ ایده‌آل ماشینی برای دادنش ندارد. همین است که چرا
زیرروال‌های محلیِ فصلِ~\ref{sec:realac} $Z$ را از $\admits$ خود بیرون
نگه می‌دارند.

چون مجموعه‌های کوچک‌تر کم‌محدودکننده‌ترند، کم‌محدودترین محیط‌های مجاز
آن‌هایی هستند که $\pars$شان درست همان بستار است --- و مسدودسازیِ
فصلِ~\ref{sec:blockers} آن‌ها را برای زوجی سازگار تحویل می‌دهد:
\[
  \IDs(\overline{\pi}) \;=\; \IDs(\overline{\varphi}) \;=\; \IDs(\pi) \cup \IDs(\varphi)
\]
بنا بر ادعای دومِ قضیهٔ~\ref{thm:blk}، و مسدودگر نامش را از عضوی که
به‌جایش می‌ایستد رونویسی می‌کند، پس
\[
  \ncl{\overline{\pi}} \;=\; \ncl{\overline{\varphi}} \;=\; \ncl{\pi} \cup \ncl{\varphi} .
\]
پس $\Zenv.\pars := \ncl{\overline{\pi}}$ تنها مقداری نیست که آن دو سیستمِ
مسدودشده می‌توانند مشترک داشته باشند، بلکه ضعیف‌ترین محدودیتی است که آن
زوج می‌پذیرد.\footnote{سور‌زدن روی همهٔ محیط‌ها بی‌اعتنا به $\pars$ سازگار
است اما برای به‌کارآمدن زیادی قوی است: شبیه‌ساز ناچار بود ترافیکِ درونیِ
$\varphi$ را فراخوان‌به‌فراخوان بازتولید کند. اگر می‌گذاشتیم شبیه‌ساز در
هویت‌هایی که $\varphi$ اشغال نمی‌کند \emph{کارکرد بزایاند}، این شرط شل
می‌شد و تنها $\IDs(\varphi) \setminus \IDs(\pi)$ برای دریغ‌کردن می‌ماند ---
که برای جفت‌شدنِ معمول تهی است؛ و اینکه زایاندن در ترکیب چه هزینه‌ای دارد
را به کارِ آینده وامی‌گذاریم.}

یک ویژگیِ این رابطه از پیش در دست است، و از گونه‌ای دیگر از هر چیزِ پس از
آن است. نتایجِ ترکیبِ زیر --- جانشانی، ترکیبِ موازی، تمامیتِ حریفِ بدل ---
نتیجه‌های یک لمِ جذبِ یگانه‌اند، که هر یک یک اجرا را بازگروه‌بندی می‌کند
و همان یک مزیت را دو بار می‌خواند. ترایایی هیچ بازگروه‌بندی نمی‌کند.
اثباتش نامساویِ مثلثی روی تعریفِ~\ref{def:adv} است، به همراهِ این آزمون
که ردهٔ نتیجه برای زوجِ بیرونی مجاز است، پس به متر تعلق دارد و نه به
ماشین‌آلات، و به هیچ‌یک از دستگاهِ باقیِ این بخش نیازی ندارد. به همین سبب
همین‌جا بیان می‌شود.

\begin{proposition}[ترایایی]\label{prop:trans}
بگذارید $\pi$، $\varphi$ و $\psi$ سیستم باشند و $\ZenvSet \subseteq
\ZenvSet_{\pi,\varphi} \cap \ZenvSet_{\varphi,\psi}$. اگر $\pi$ کارکردِ
$\varphi$ را در $\varepsilon_1$ و در برابرِ $(\Aset,\SimSet,\ZenvSet)$
UC-تقلید کند و $\varphi$ کارکردِ $\psi$ را در $\varepsilon_2$ و در برابرِ
$(\SimSet,\SimSet',\ZenvSet)$، آن‌گاه $\pi$ کارکردِ $\psi$ را در
$\varepsilon_1 + \varepsilon_2$ و در برابرِ $(\Aset,\SimSet',\ZenvSet)$
UC-تقلید می‌کند.
\end{proposition}

\begin{proof}
اول، نتیجه حکمِ مجازی از تقلید است. $\Zenv$ای که در هر دو رده باشد
$\Zenv.\pars \supseteq \ncl{\pi} \cup \ncl{\varphi}$ و $\Zenv.\pars
\supseteq \ncl{\varphi} \cup \ncl{\psi}$ دارد، پس $\Zenv.\pars \supseteq
\ncl{\pi} \cup \ncl{\psi}$؛ و در هیچ هویتی از $\IDs(\pi) \cup
\IDs(\varphi)$ و نه از $\IDs(\varphi) \cup \IDs(\psi)$ میهمان نمی‌کند، پس
در هیچ هویتی از $\IDs(\pi) \cup \IDs(\psi)$ هم نه --- پس هر دو بندِ
تعریفِ~\ref{def:admis} برقرارند و $\ZenvSet \subseteq
\ZenvSet_{\pi,\psi}$؛ و برای این، آن سه سیستم لازم نیست نسبتی با هم
داشته باشند. اکنون بگذارید $\Adv \in \Aset$. تقلیدِ اول $\Sim \in
\SimSet$ را فراهم می‌کند، و چون تقلیدِ دوم حریف‌هایش را روی $\SimSet$ سور
می‌زند، خواندنِ $\Sim$ به‌عنوانِ حریفی در برابرِ $\varphi$ به $\Sim' \in
\SimSet'$ می‌رسد. برای هر $\Zenv \in \ZenvSet$، نامساویِ مثلثی می‌دهد
\[
  \begin{aligned}
  \advtg^{\op{uc}}_{\pi,\psi}(\Zenv,\Adv,\Sim')
  &\;\leq\;
  \advtg^{\op{uc}}_{\pi,\varphi}(\Zenv,\Adv,\Sim) \;+\; \advtg^{\op{uc}}_{\varphi,\psi}(\Zenv,\Sim,\Sim') \\
  &\;\leq\; \varepsilon_1 + \varepsilon_2 ,
  \end{aligned}
\]
چون جملهٔ میانیِ $\Pr[\opl{Exec}(\varphi,\Zenv,\Sim,\Corr) = 1]$ حذف
می‌شود. و چون $\Sim'$ تنها به $\Adv$ وابسته است، همان شبیه‌سازِ لازم است.
\end{proof}

سه ساختِ زیر --- محیطِ جذب‌شدهٔ اینجا، و آن دوی
بخشِ~\ref{app:dummy} --- چند ماشین را در یک اشغال‌گرِ اجرا تا می‌کنند.
آن‌ها یک شکلِ مشترک دارند، که یک بار نامش می‌بریم.

\begin{definition}[دسته]\label{def:bundle}
یک \emph{دسته} چند ماشین را با هم به‌عنوانِ یک اشغال‌گرِ یگانهٔ اجرا
می‌دواند، و ژتون را بینشان همان‌گونه که $\opl{Exec}$ می‌کند دست‌به‌دست
می‌کند و هر یک را همان‌گونه که $\opl{Exec}$ می‌کرد مقدارِ آغازین می‌دهد.
دسته مجموعهٔ اعلام‌شده‌ای از هویت‌ها را اشغال می‌کند، که هویت‌های خودِ
اعضا باید بپوشانندش، و خروجی و پارامترهایش را از عضوی معین می‌گیرد.
فراخوانِ میانِ دو ماشینِ میهمان در درون خدمت می‌شود، با مسیریابیِ خودِ
دسته و \emph{بی‌آنکه روی اجرا گذاشته شود}؛ فراخوانی که از بیرون می‌رسد
به‌دستِ عضوی که هویتِ هدف را حمل می‌کند خدمت می‌شود، از راهِ واسطِ کاملِ
همان عضو، با نگهبان و همه چیز؛ و فراخوانی که ماشینِ میهمان به بیرون
می‌گذارد، هرگاه هدفش در مجموعهٔ اعلام‌شدهٔ \emph{بیرونی} باشد روی اجرا
گذاشته می‌شود و در غیرِ این صورت $\rej$ پاسخ می‌گیرد. یک نمونه می‌تواند
علاوه بر این، فراخوان‌های خاصی را میانِ اعضایش به‌دست مسیریابی کند --- یکی
را پیش از دیگری بنشاند، حتی وقتی هدف بیرونِ هویت‌های هر عضوی باشد --- و
چنین مسیریابیِ صریحی بر فراخوان‌هایی که نامشان می‌برد حکم می‌کند، و تنها
باقی را به قاعدهٔ هدفِ بالا وامی‌گذارد.

اینکه فراخوانِ بیرونیِ ماشینِ میهمان زیرِ چه ادعایی می‌رود بر اشغال
می‌گردد، چون خموش‌کننده هستهٔ یک واسط را می‌پوشاند و فراتر نمی‌رود
(فصلِ~\ref{sec:wrappers}). عضوی \emph{که دسته هویت‌هایش را اشغال می‌کند}
فراخوان‌هایش را زیرِ هویتِ خودش می‌فرستد و تنها با خموش‌کنندهٔ خودش
روبه‌رو می‌شود؛ عضوی \emph{که دسته هویت‌هایش را اشغال نمی‌کند} هیچ هویتی
برای ادعا ندارد، پس دسته هر یک از فراخوان‌های بیرونی‌اش را بازنویسی
می‌کند تا زیرِ هویتی برود که خودش داردش --- و اینکه کدام هویت، بخشی از
همان چیزی است که نمونه اعلامش می‌کند، در کنارِ مجموعهٔ بیرونی‌اش --- و
ادعای بازنویسی‌شده با خموش‌کنندهٔ آن واسط روبه‌رو می‌شود. اینکه عضوی در
کدام حالت است، همان است که دسته را به اجرایی که به‌جایش می‌ایستد وفادار
نگه می‌دارد.
\end{definition}

اینکه ترافیکِ درونی \emph{خدمت} می‌شود و نه \emph{گذاشته}، همان است که
می‌گذارد دسته جایگاهِ حریفِ سیستمی را اشغال کند که هسته‌هایش او را
پاسخگویانه فرا می‌خوانند، و ارزشِ گفتن دارد که چرا این تمایز مویِ دماغ
نیست. پوشانهٔ $\opl{Respond}$ (تعریفِ~\ref{def:respond}) هر فراخوانی را
که پاسخ‌دهنده \emph{می‌گذارد} رد می‌کند، پس دسته‌ای که ناچار بود
فراخوان‌های درونی‌اش را بگذارد در حالِ پاسخ‌دادن به هیچ‌یک از اعضای خودش
نمی‌رسید --- و شبیه‌سازِ فصلِ~\ref{sec:realac}، که در درونِ هر یک از دو
پاسخِ پاسخگو سه دورگهٔ میهمان را می‌راند، اصلاً نمی‌توانست بدود. خدمت‌کردنِ
آن‌ها در درون، ژتون را همان جا که بود می‌گذارد، و این کلِ چیزی است که
پاسخگویی می‌خواهد. آنچه دسته در درونِ یک پاسخِ پاسخگو باز هم نمی‌تواند
بکند، رسیدن به چیزی است که میهمانش نکرده: دفتر و کارکردهای مشترک بیرونِ
اویند، پس فراخوان‌هایشان گذاشته می‌شود، و رد می‌شود.

تعریفِ~\ref{def:bundle} پیش‌تر می‌گوید ماشینِ میهمان دو گونه دارد، بر حسبِ
اینکه دسته هویت‌هایش را اشغال می‌کند یا نه، و همین تفاوت کلِ آنچه در پی
می‌آید است. ماشینی که دسته هویت‌هایش را نگه می‌دارد از توزیعِ اجرا محو
می‌شود و در درون خدمت می‌شود؛ و ماشینی که دسته هویت‌هایش را واگذار می‌کند
باید آن‌ها را به کسی بسپارد، و فراخوان‌هایش را باید همان کس به‌نیابتِ او
بگذارد. هر دو حرکت پیش می‌آید --- اولی پروتکلی را جذب می‌کند و دومی حریفی
را --- پس با هم تعریفشان می‌کنیم و یک بار درستی‌شان را اثبات می‌کنیم.

\begin{definition}[جذب]\label{def:absorption}
بگذارید $\Zenv$ محیطی باشد و $H$ مجموعهٔ متناهی‌ای از ماشین‌های یک اجرا،
که به $H = H_{\op{keep}} \uplus H_{\op{cede}}$ شقّه می‌شود و اعضای
$H_{\op{cede}}$ اشغال‌گرانِ جایگاهِ حریف‌اند. \emph{جذبِ} $H$ در $\Zenv$
همان دستهٔ (تعریفِ~\ref{def:bundle}) $\Zenv^{H}$ از $\Zenv$ با یک رونوشت
از هر عضوِ $H$ است؛ این دسته $\IDs(\Zenv) \cup \IDs(H_{\op{keep}})$ را
اشغال می‌کند، خروجی و $\pars$ را از $\Zenv$ می‌گیرد، و مجموعهٔ بیرونی‌اش
همان هویت‌هایی است که اجرای پیرامون هنوز خدمتشان می‌کند. عضوی از
$H_{\op{keep}}$ هویت‌هایش را نگه می‌دارد، پس فراخوان‌هایش زیرِ ادعای خودش
می‌روند و فراخوان‌هایی که به او خطاب شوند در درونِ دسته خدمت می‌شوند.
عضوی از $H_{\op{cede}}$ هیچ هویتی نگه نمی‌دارد، پس دسته هر فراخوانِ
بیرونیِ $\id''.\fopl{Op}(x)$ را که او در حالِ دویدن در $(A,P)$ می‌گذارد
به این دستور بازنویسی می‌کند
\[
  (A,P).\fopl{A}\bigl(\id''.\Op,\ x\bigr) \ \text{ as } (Z,P) ,
\]
و هویت‌هایی که او واگذار می‌کند در اجرای تازه به‌دستِ یک \emph{بازپخش}
برای او اشغال می‌شوند، که سرتاسر همان بدلِ $\Dum$ از
بخشِ~\ref{app:dummy} است.
\end{definition}

ما حالتِ واگذاری را تنها برای جایگاهِ حریف بیان می‌کنیم. تنها همین یکی در
زیر به کار می‌رود، و تنها همین یکی است که واسطش شکلِ دستوری دارد که
فراخوان‌ها را به آن بازنویسی کنیم --- کارکردِ استاندارد آنچه کدِ خودش
نامش می‌برد را خدمت می‌کند و نمی‌توانست به‌جای فراخوان‌های ماشینِ دیگری
بایستد.

\begin{definition}[محیطِ جذب‌شده]\label{def:absorb}
بگذارید $\varphi \subseteq \rho$، و $\pi$ سیستمی باشد که به‌جای $\varphi$
نشانده می‌شود، و $\Zenv$ محیطی. محیطِ \emph{جذب‌شدهٔ}
$\Zenv^{\rho \setminus \varphi}$ همان جذبِ
(تعریفِ~\ref{def:absorption}) $\rho \setminus \varphi$ در $\Zenv$ است با
هر عضو نگه‌داشته و هیچ عضو واگذارشده، پس دستهٔ $\Zenv$ است با یک رونوشت
از هر $\F \in \rho \setminus \varphi$، که هر رونوشت هویت‌هایش را نگه
می‌دارد؛ این دسته $\IDs(\Zenv) \cup \IDs(\rho \setminus \varphi)$ را
اشغال می‌کند، خروجی و $\pars$ را از $\Zenv$ می‌گیرد، و مجموعهٔ بیرونی‌اش
$\IDs(\pi) \cup \IDs(\varphi) \cup \{\, \id : \id.F \in \{A,C\} \,\}$
است. این یک محیط است (تعریفِ~\ref{def:env}) و نه یک کارکرد، چون روی چند
شناسهٔ فرایند گسترده است. بالانویس تنها پوستهٔ میهمان‌شده را ثبت می‌کند ---
مجموعهٔ بیرونی $\pi$ را هم می‌خواند --- پس $\Zenv^{\rho \setminus
\varphi}$ نسبت به جانشانی‌ای که بافت می‌دهد تعیین می‌شود، و دو جانشانی که
$\rho \setminus \varphi$ مشترک دارند اما در $\pi$ متفاوتند نامِ دو ماشینِ
متفاوتند.
\end{definition}

هر دو نیمهٔ مجموعهٔ بیرونی را همان چیزی وادار می‌کند که اجرای جذب‌شده باید
بازتولیدش کند. اولی جایگاه است، که هویت‌هایش بر حسبِ اینکه محیطِ جذب‌شده
در کدام جهان نشسته به $\pi$ یا به $\varphi$ می‌رسند. ما اجتماع را
می‌گیریم و نه تنها $\IDs(\varphi)$ را، تا همان یک ماشین در هر دو خدمت
کند، و این حالا که آن دو لازم نیست همان هویت‌ها را اشغال کنند مهم است.
دومی آن دو شناسهٔ رزروِ $A$ و $C$ است، که به‌عنوانِ شناسه نامیده شده‌اند و
نه به‌عنوانِ ماشین، چون جایگاهِ حریف در سوی ایده‌آل $\Sim$ را در بر دارد و
گزارهٔ~\ref{prop:regroup} در هر دو سو یک و همان ماشینِ جذب‌شده را لازم
دارد. آن‌ها آنجایند چون ماشین‌های میهمان واسط‌های \emph{کامل}‌اند و نه
هسته‌های خالی: هر یک دفتر را در سطرهای~\ref{ln:stat}
و~\ref{ln:stattwo}ِ بخشِ~\ref{sec:fullinterface} می‌خواند، و هر یک
فراخوانی در طرفِ فاسد را از راهِ $\opl{Mediate}$ به حریف می‌سپارد، که به
$(A,\id.P)$ قلاب می‌شود. هیچ‌یک از این دو شناسه به هیچ سیستمی تعلق ندارد.
بی این نیمه، فراخوان به کارکردی میهمان در نخستین سطرش رد می‌شد، و قلابِ
طرفِ فاسد به‌جای رسیدن به حریف رد می‌شد: اجرای جذب‌شده درست همان جا که
$\Cs$ به $\rho \setminus \varphi$ می‌رسید از بازگروه‌بندیِ اجرای بیرونی
بودن دست می‌کشید. هر چیزِ دیگر با $\rej$ پاسخ می‌گیرد، چون در اجرای
بیرونی هم چیزی خدمتش نمی‌کرد.

پارامتر دست‌نخورده منتقل می‌شود، و باید چنین باشد. آن مجموعهٔ خموش‌شدهٔ
واسطِ خودِ $\Zenv$ است، پس هر تغییری در آن تغییر می‌دهد که $\Zenv$ چه
فراخوان‌هایی می‌تواند بگذارد --- و $\Zenv$ همان یک ماشینی است که باید در
هر دو سو یکسان رفتار کند. کسی ممکن است گمان کند کارکردهای میهمان نیاز
دارند $\IDs(\rho \setminus \varphi)$ از آن حذف شود تا بتوانند به‌عنوانِ
خودشان سخن بگویند، اما ندارند: هر یک واسطی کامل است، و
سطرِ~\ref{ln:call} از پوشانهٔ خودش همان است که ادعایش را مجاز می‌کند، و
پوشانهٔ $\Zenv$ فراتر از هستهٔ خودِ $\Zenv$ نمی‌رسد
(تعریفِ~\ref{def:bundle}). حذفِ آن هویت‌ها هیچ چیزِ تازه‌ای برایشان مجاز
نمی‌کرد، و $\Zenv$ را درست روی همان هویت‌هایی ناخموش می‌کرد که پیش از جذب
رویشان خموش بود.

پس جذب، محیطی برای زوجِ بیرونی را به محیطی برای زوجِ درونی بدل می‌کند و
مجازبودن را نگه می‌دارد. دفتر رفتارِ ویژه‌ای لازم ندارد: چون عضوِ هیچ
سیستمی نیست، نه در $\varphi$ است و نه در $\rho \setminus \varphi$، پس
محیطِ جذب‌شده میهمانش نمی‌کند و مسدودسازی هیچ مسدودگری برایش نمی‌سازد.
مثلِ $\Zenv$ و $\Adv$، $\opl{Exec}$ فراهمش می‌کند، و همان یک ماشینِ واحد
در هر دو سوی هر مقایسه‌ای است. رونوشت‌های میهمانِ آن هویت‌هایشان را نگه
می‌دارند، پس بنا بر تعریفِ~\ref{def:bundle} فراخوان‌هایشان را زیرِ
ادعاهای خودشان می‌فرستند و با خموش‌کننده‌های خودشان روبه‌رو می‌شوند ---
نه با خموش‌کنندهٔ $\Zenv$ --- و همین است که بازگروه‌بندیِ زیر لازم دارد؛
و ساخت‌های بخشِ~\ref{app:dummy} حالتِ مقابلند، چون حریفی را میهمان
می‌کنند که دسته هویت‌هایش را اشغال نمی‌کند.

\begin{proposition}[جذب مجازبودن را نگه می‌دارد]\label{prop:absorb}
بگذارید $\rho$ و $\pi$ سیستم باشند، $\varphi \subseteq \rho$، و جانشانیِ
$\varphi$ با $\pi$ در $\rho$ خوش‌ساخت باشد. آن‌گاه
\[
  \Zenv \in \ZenvSet_{\rho[\pi/\varphi],\,\rho}
  \qquad\Longrightarrow\qquad
  \Zenv^{\rho \setminus \varphi} \in \ZenvSet_{\pi,\varphi} .
\]
\end{proposition}

\begin{proof}
برای نیمهٔ $\pars$، جذب این میدان را دست‌نخورده منتقل می‌کند، پس همین بس
است که
\[
  \ncl{\rho[\pi/\varphi]} \cup \ncl{\rho} \;\supseteq\; \ncl{\pi} \cup \ncl{\varphi} .
\]
هر نامی که $\varphi$ به کار می‌برد $\rho$ هم به کار می‌برد، چون $\varphi
\subseteq \rho$، و هر نامی که $\pi$ به کار می‌برد $\rho[\pi/\varphi]$ هم،
چون $\pi \subseteq \rho[\pi/\varphi]$. برای نیمهٔ میهمانی، $\Zenv^{\rho
\setminus \varphi}$ فراتر از هر چه $\Zenv$ میهمان کرده بود، در $\IDs(\rho
\setminus \varphi)$ میهمان می‌کند: و آن دومی از $\IDs(\rho[\pi/\varphi])
\cup \IDs(\rho)$ می‌پرهیزد، که $\IDs(\pi) \cup \IDs(\varphi)$ را در بر
دارد؛ و $\IDs(\rho \setminus \varphi)$ هیچ کجا با $\IDs(\varphi)$ برخورد
نمی‌کند، چون آن دو زیرسیستم‌های مجزای $\rho$اند، و هیچ کجا با $\IDs(\pi)$
برخورد نمی‌کند، بنا بر خوش‌ساختیِ جانشانی --- تنها جایی که این فرض به کار
می‌رود، و تنها برای نوع‌بندی به کار می‌رود. هیچ نسبتی میانِ $\IDs(\pi)$ و
$\IDs(\varphi)$ وارد نمی‌شود. یک چیز که این حکم نمی‌گویدش: ماشینِ جذب‌شده
هویت‌هایی را در درونِ $\pars$ِ خودش \emph{ادعا} می‌کند، و هر رونوشتِ
میهمان هویتِ خودش را. این تناقضی نیست --- بنا بر تعریفِ~\ref{def:bundle}
ادعاهای رونوشتِ میهمان با خموش‌کنندهٔ خودش روبه‌رو می‌شوند و هرگز با
خموش‌کنندهٔ محیط، و مجازبودن شرطی است بر آن میدان، که هر جا آن میدان
اِعمال شود خوانده می‌شود.
\end{proof}

آنچه این قضیه را کارآمد می‌کند این است که جذب یک اجرا را بازگروه‌بندی
می‌کند بی‌آنکه تغییرش دهد. و چون همان تناظر به کارِ بخشِ~\ref{app:dummy}
هم می‌آید، پیش از لم آن را می‌نشانیم.

\heading{پیکربندی‌ها} مقصودمان از \emph{پیکربندیِ} یک اجرا حالتِ آن بینِ
فعال‌سازی‌هاست: حالتِ هر ماشین، بخشِ ناخواندهٔ هر نوارِ سکه، اینکه کدام
ماشین ژتون را در دست دارد و کجای کدِ خودش است، و فراخواننده‌هایی که معلق و
چشم‌به‌راهِ پاسخند. فراخوان‌ها تو‌درتو می‌شوند، و
بخشِ~\ref{sec:execution} می‌گذارد یک نمونه چند ادامهٔ معلق داشته باشد، پس
آخری از این‌ها یک زنجیره است و نه یک مجموعه: یک درایه به‌ازای هر
فراخوانی که گذاشته شده و هنوز پاسخ نگرفته، به همان ترتیبی که گذاشته شده،
که فراخواننده و فراخوانی که چشم‌به‌راهش است را ثبت می‌کند.

آن دو اجرا همان ماشین‌ها را نمی‌دوانند --- کارِ جذب همین است --- اما وقتی
$\Zenv^{H}$ را به‌جای یکی، همان ماشین‌هایی بشماریم که دسته‌شان می‌کند،
همان‌ها را می‌دوانند، جز بازپخش‌هایی که هویتِ واگذارشده لازم دارد.
تعریفِ~\ref{def:bundle} ژتون را بینِ آن‌ها همان‌گونه که $\opl{Exec}$
می‌کند دست‌به‌دست می‌کند، پس هر یک در هر دو سو به همان معنا آن را در دست
دارد. با این شمارش، ماشین‌های آن دو اجرا جز بازپخش‌ها یک‌به‌یک متناظرند، و
ما $\iota$ را برای این دوسویی می‌نویسیم: رونوشتِ میهمانِ هر عضوِ $H$ به
همان عضو، $\Zenv$ِ میهمان به $\Zenv$، و هر چه اجرا هنوز می‌دواندش،
$\Corr$ هم در آن، به خودش.

\begin{definition}[تناظرِ بازگروه‌بندی]\label{def:regcorr}
پیکربندیِ $c$ از اجرای چپ و پیکربندیِ $\hat c$ از اجرای راست
\emph{متناظرند}، و می‌نویسیم $c \approx \hat c$، هرگاه
\begin{enumerate}[label=(C\arabic*),leftmargin=3em]
  \item هر ماشینِ $c$ و تصویرش زیرِ $\iota$ در همان حالت باشند --- $\Corr$
        هم در آن، پس آن دو بر $\Cs$ توافق دارند؛
  \item هر ماشین و تصویرش همان نوارِ سکهٔ ناخوانده را داشته باشند؛
  \item همان ماشین در هر دو، زیرِ $\iota$، ژتون را در دست داشته باشد، در همان
        نقطه از همان عمل و با همان مقدارهای محلی؛ و
  \item دو زنجیرهٔ فراخوانندهٔ معلق زیرِ $\iota$ توافق داشته باشند: طولِ برابر، و
        به‌ازای هر $k$، درایه‌های $k$اُم نامِ ماشین‌هایی را ببرند که $\iota$
        یکی‌شان می‌کند و چشم‌به‌راهِ همان فراخوان باشند --- همان هویتِ هدف،
        همان عمل، همان ورودی و همان هویتِ ادعایی.
\end{enumerate}
هر جا خانوادهٔ راست بازپخشی حمل کند، $\iota$ دوسویی‌ای است میانِ ماشین‌های
چپ و ماشین‌های راستِ \emph{جز} بازپخش، و (C4) با حذفِ قاب‌های خودِ بازپخش
خوانده می‌شود: فراخوانِ بازپخش‌شده یک قاب می‌افزاید که چپ ندارد، و دقیقاً
یک بار با همان پاسخی که فراخوانِ چپ برگرداند از سر گرفته می‌شود --- جز
آنکه $\rej$ای که چپ می‌خواند می‌تواند در راست به‌صورتِ $\none$ برسد، و
این تنها جایی است که بی‌اعتنایی به رد خرج می‌شود.
\end{definition}

تنها (C4) از این کمانک‌گذاری در خطر است. در چپ، هر تعلیقی را $\opl{Exec}$
ثبت می‌کند؛ در راست، تعلیق‌ها بینِ $\opl{Exec}$ و دسته تقسیم می‌شوند، و
دسته آن‌هایی را ثبت می‌کند که به فراخوان‌های درونیِ خودش مربوطند. این بند
می‌خواهد که آن دو زنجیره، وقتی زنجیره‌های خودِ دسته همان جا که گذاشته
شده‌اند شمرده شوند، توافق داشته باشند. بازگروه‌بندی درست همین است: یک
زنجیرهٔ تعلیق، که دو توزیع‌گر نگهش می‌دارند.

\heading{نخستین فعال‌سازی} آن دو اجرا در تناظر آغاز می‌کنند، و ارزشِ گفتن
دارد که چرا، چون این همان یک فعال‌سازی است که تحلیلِ حالتِ لم به آن
نمی‌رسد: هیچ ماشینی نگذاشته است.
سطرِ~\ref{ln:execinit} $\Corr$ و هر چه اجرا در هر دو سو هنوز می‌دواندش را
مقدارِ آغازین می‌دهد، تعریفِ~\ref{def:absorption} هر رونوشتِ میهمان را
همان‌گونه که $\opl{Exec}$ می‌کرد مقدارِ آغازین می‌دهد، و بازپخش هیچ حالتی
حمل نمی‌کند، پس هر ماشین و تصویرش در همان حالت آغاز می‌کنند، که (C1) است.
هیچ نواری خوانده نشده، که (C2) را می‌دهد. عملِ $\op{Initialize}$ هیچ
فراخوانی نمی‌گذارد، پس هیچ‌چیز معلق نیست و هر دو زنجیره تهی‌اند، که (C4)
را می‌دهد --- و با آن، دلیلِ اینکه ترتیبی که دسته رونوشت‌هایش را مقدارِ
آغازین می‌دهد مشاهده‌شدنی نیست: آن ترتیب را از هیچ‌چیزِ پیکربندی نمی‌شود
خواند. سپس سطرِ~\ref{ln:execstart} $(Z,Z).\fopl{Z}()$ از $(Z,Z)$ را فعال
می‌کند. در چپ ژتون به $\Zenv$ می‌رود؛ در راست به $\Zenv^{H}$، که در آن
$(Z,Z)$ را $\Zenv$ِ میهمان حمل می‌کند، همان که
تعریفِ~\ref{def:absorption} دسته را وادار به اشغالش می‌کند. هر دو در همان
هویت و روی همان ورودی به $\fopl{Z}$ وارد می‌شوند، که (C3) است.

از آن پس هر فعال‌سازی یا فراخوانی است که گذاشته می‌شود یا مقداری که
برگردانده می‌شود، و لم به این می‌رسد که نشان دهد تناظر از یکی از آن‌ها در
هر دو سو جان سالم می‌برد، و پیکربندی‌های متناظر با هم و با همان خروجی
می‌ایستند. بندهایش آن فعال‌سازی‌ها را حالت‌به‌حالت می‌گیرند.

\begin{lemma}[جذب]\label{lem:absorb}
بگذارید اجرایی اشغال‌گرانِ $\{\Zenv\} \cup H \cup R$ را بدواند، با $\Corr
\in R$ و $\Zenv$ی محیط، و $H = H_{\op{keep}} \uplus H_{\op{cede}}$، و
$\Zenv^{H}$ همان جذبِ تعریفِ~\ref{def:absorption} باشد. از هر عضوِ
$H_{\op{cede}}$، و از هیچ‌چیزِ $H_{\op{keep}}$، بخواهید که
\begin{enumerate}[label=(\alph*),leftmargin=2.2em]
  \item بی‌اعتنا به رد باشد (تعریفِ~\ref{def:refobl})؛
  \item هیچ هستهٔ $R$ فراخوانِ پاسخگویی روی هویت‌هایی که او واگذار می‌کند
        نگذارد (تعریفِ~\ref{def:respcall})؛ و
  \item $\Zenv.\pars$ هیچ هویتِ $Z$ در بر نداشته باشد.
\end{enumerate}
آن‌گاه، به شرطِ آنکه آن دو خانواده به معنای تعریفِ~\ref{def:occ}
خانواده‌های مجازِ اشغال‌گر باشند،
\[
  \begin{gathered}
  \opl{Exec}\bigl(\{\Zenv\} \cup H \cup R\bigr)
  \qquad\text{and}\qquad \\[3pt]
  \opl{Exec}\bigl(\{\Zenv^{H}\} \cup R \cup \{\text{a relay at each ceded identity}\}\bigr)
  \end{gathered}
\]
همسان‌توزیعند --- به‌عنوانِ توزیع برابرند، و نه تنها نزدیک.
\end{lemma}

\begin{proof}
% BEGIN overview lem:absorb
این استدلال یک استقرا روی فعال‌سازی‌هاست، و پنج بندِ زیر گامش را انجام
می‌دهند. \emph{ماشین‌ها} دو خانوادهٔ اشغال‌گر را جفت می‌کند و می‌سنجد که
هر ماشین و تصویرش یکسان آغاز می‌کنند. \emph{مسیریابی} می‌پرسد هر فراخوان
کجا فرود می‌آید و چهار حالت می‌یابد --- درونیِ گروه، خروج از آن، ورود در
هویتی نگه‌داشته، ورود در هویتی واگذارشده --- که از آن‌ها تنها آخری میانِ دو
خانواده تفاوت دارد. \emph{ادعاها} بر حسبِ اینکه چه گونه‌ای از عضو فراخوان
را گذاشته شقّه می‌شود: عضوی که هویت‌هایش را نگه می‌دارد در هر دو سو با
همان خموش‌کننده‌ها روبه‌رو می‌شود، و عضوِ واگذارنده هیچ هویتی برای ادعا
ندارد، پس فراخوان‌های بیرونی‌اش بازنویسی می‌شوند، و این بازنویسی باید از
سه آزمون جان سالم ببرد --- دسته از کجا می‌تواند سخن بگوید، آیا آن ادعا
خموش شده، و آیا از نگهبانِ بازپخش می‌گذرد. \emph{ردها} آن یک مقداری را
جدا می‌کند که می‌تواند متفاوت باشد، یعنی $\rej$ای که چپ می‌خواند و در راست
به‌صورتِ $\none$ می‌رسد، و فرضِ (a) را رویش خرج می‌کند؛ و فرضِ (b) مسئلهٔ
آینه‌ایِ همان را سد می‌کند، یعنی فراخوانِ پاسخگو روی جایگاهی بازپخش‌شده.
\emph{ژتون} می‌سنجد که آن دو زنجیرهٔ تعلیق، وقتی زنجیره‌های خودِ دسته همان
جا که گذاشته شده‌اند شمرده شوند، توافق دارند. فرضِ (c) در \emph{ادعاها}
خرج می‌شود، بر حقِ عضوِ واگذارنده برای ادعای $(Z,P)$.
% END overview lem:absorb
دو قرارداد از فصل‌های~\ref{sec:functionalities} و~\ref{sec:execution} در
سرتاسر به کار می‌روند و ارزشِ بازگفتن دارند. ماشین‌های متمایز از سکه‌های
مستقل می‌کشند، و دسته‌کردنِ ماشین‌ها در یکی نوارهایشان را یکی نمی‌کند، پس
تصادفِ رونوشتِ میهمان همان است که پیش از جذب بود. و $\op{Initialize}$ هیچ
فراخوانی نمی‌گذارد، پس ترتیبی که دسته رونوشت‌های میهمانش را مقدارِ آغازین
می‌دهد مشاهده‌ناپذیر است. اثبات استقرایی است روی فعال‌سازی‌ها:
پیکربندی‌های آغازین متناظرند، چنان که پیش از لم نشان داده شد؛
پیکربندی‌های متناظر در هر دو سو همان فعال‌سازیِ بعدی را تولید می‌کنند و به
پیکربندی‌های متناظر گام می‌زنند، و هر ماشین و تصویرش همان سکه‌ها را به
همان ترتیب می‌خوانند، که همان است که (C2)ِ تعریفِ~\ref{def:regcorr} نگهش
می‌دارد؛ و پیکربندی‌های ایستای متناظر یکسان خروجی می‌دهند، چون خروجی مالِ
$\Zenv$ِ میهمان است، بنا بر تعریفِ~\ref{def:absorption}. بندهای زیر گامِ
استقرا را انجام می‌دهند: چه ماشین‌هایی هستند و یکسان آغاز می‌کنند، (C1)؛
فراخوان کجا می‌رود و چه ادعا می‌کند، (C4)؛ و ژتون چگونه حرکت می‌کند،
(C3). دو بندِ اول بینِ هر دو گونهٔ عضو مشترکند؛ و آن دو بندِ پس از آن‌ها
همان چیزی هستند که عضوِ واگذارنده هزینه می‌کند.

\emph{ماشین‌ها.} چپ $\{\Zenv\} \cup H \cup R$ را می‌دواند؛ راست
$\{\Zenv^{H}\} \cup R$ را به همراهِ یک بازپخش در هر هویتِ واگذارشده، و
$\Zenv^{H}$ همان $\Zenv$ است با یک رونوشت از هر عضوِ $H$. اگر به‌جای یکی،
همان ماشین‌هایی شمرده شود که دسته‌شان می‌کند، آن دو گردایه جز بازپخش‌ها بر
هم می‌افتند، پس $\iota$ هر عضوِ $H$ را با رونوشتش جفت می‌کند، هر عضوِ $R$
و $\Corr$ را با خودش، و بازپخش‌ها را باقی می‌گذارد، که همان خوانشی است که
تعریفِ~\ref{def:regcorr} به آن‌ها می‌دهد. سطرِ~\ref{ln:execinit} $R$ را در
هر دو سو مقدارِ آغازین می‌دهد، تعریفِ~\ref{def:absorption} رونوشت‌های
میهمان را همان‌گونه که $\opl{Exec}$ می‌کرد، و بازپخش بی‌حالت است، پس هر
ماشین و تصویرش یکسان آغاز می‌کنند. اشغال بنا بر فرض در هر دو سو مجزاست،
پس هر فراخوان هدفی یگانه دارد: در راست، دسته $\IDs(\Zenv) \cup
\IDs(H_{\op{keep}})$ را در بر دارد و بازپخش‌ها هر چه $H_{\op{cede}}$ رها
کرده است.

\emph{مسیریابی.} فراخوانی را بگیرید و بپرسید کجا فرود می‌آید. با هر دو سرِ
آن در $\{\Zenv\} \cup H$، در چپ $\opl{Exec}$ توزیعش می‌کند و در راست در
درونِ $\Zenv^{H}$ خدمت می‌شود، و در هر دو حال به‌دستِ همان ماشین. اگر از
آن گروه بیرون رود، در چپ $\opl{Exec}$ توزیعش می‌کند و در راست به بیرون
گذاشته می‌شود، چون هدفش در مجموعهٔ بیرونی است، که
تعریفِ~\ref{def:absorption} آن را همان هویت‌هایی می‌کند که اجرا هنوز
خدمتشان می‌کند. اگر در هویتی از $H_{\op{keep}}$ وارد گروه شود، در چپ
$\opl{Exec}$ توزیعش می‌کند و در راست رونوشتِ میهمانی که آن هویت را حمل
می‌کند خدمتش می‌کند، همان که دسته اشغالش کرده. اگر در هویتی از
$H_{\op{cede}}$ وارد شود، در چپ به عضوِ واگذارنده توزیع می‌شود و در راست
به بازپخش، و این همان یک جایی است که دو خانواده تفاوت دارند و موضوعِ دو
بندِ زیر است. اگر به هر جای دیگری خطاب شود، در هیچ‌یک از دو سو ماشینی
ندارد: $\opl{Exec}$ هدفی نمی‌یابد و $\Zenv^{H}$ $\rej$ پاسخ می‌دهد.
نگهبان‌ها در فراخوانده و از شناسهٔ ادعایی سنجیده می‌شوند، پس هرگاه ادعاها
توافق کنند آن‌ها هم توافق می‌کنند.

\emph{ادعاها، برای عضوی که هویت‌هایش را نگه می‌دارد.} هر چنین فراخوانی در
هر دو سو همان شناسهٔ ادعایی را حمل می‌کند، و آن دو خموش‌کننده‌ای که
می‌توانستند یکی را تغییر دهند در هر دو سو همان دو خموش‌کننده‌اند. واسطِ
میهمان بنا بر تعریفِ~\ref{def:bundle} با خموش‌کنندهٔ خودش حکم می‌شود و نه
با خموش‌کنندهٔ محیط: هویت‌هایش را نگه می‌دارد، پس $\opl{Silence}$ در
$\Zenv$ به او نمی‌رسد. و چون واسطی کامل است، در هر دو سو درست هویتِ خودش
را ادعا می‌کند، و سطرِ~\ref{ln:call} هسته‌اش را روی $\{\id'' \neq \id\}$
خموش می‌کند؛ و فراخوان‌هایی که پوشانه‌اش می‌گذارد، روی $\Corr$ و از راهِ
$\opl{Mediate}$، در هر دو سو حتی بیرونِ همان خموش‌کننده می‌ایستند. واسطِ
خودِ محیط با $\pars$ حکم می‌شود، که جذب آن را بی‌تغییر منتقل می‌کند، پس
$\overline{\Ids}$ از همان مجموعه و در همان هویت حساب می‌شود و درست همان
ادعاها را می‌پذیرد. هیچ‌یک از دو ماشین ادعایی به‌دست نمی‌آورد و از دست
نمی‌دهد.

\emph{ادعاها، برای عضوی که آن‌ها را واگذار می‌کند.} چنین عضوی هیچ هویتی از
خودش در درونِ دسته برای ادعا ندارد، پس تعریفِ~\ref{def:absorption}
فراخوان‌های بیرونی‌اش را بازنویسی می‌کند، و این بازنویسی باید از سه آزمون
جان سالم ببرد. اول، دسته از کجا می‌تواند سخن بگوید. هرگاه عضوِ میهمان
به‌نیابتِ طرفِ $P$ بدود --- یعنی آنجا که چپ او را در $(A,P)$ در حالِ
دویدن دارد --- دسته در $(Z,P)$ یا در $(Z,Z)$ می‌دود، یا وگرنه $P = Z$:
یا $\Zenv$ با فراخواندنِ یک هویتِ $A$ واردش می‌کند، که در درون خدمت
می‌شود و $\opl{Guard}_{\Adv}$ طرفِ ادعایی را وادار می‌کند طرفِ هدف باشد، و
سطرِ~\ref{ln:zparty} می‌گذارد واسطی در $(Z,Q)$ طرفِ $P$ را تنها وقتی
ادعا کند که $Q \in \{P,Z\}$ یا $P = Z$؛ یا بازپخش واردش می‌کند، که
فراخوانِ درونی‌اش درست $(Z,P)$ را هدف می‌گیرد. دوم، آیا ادعای $(Z,P)$
خموش شده است: نه با گرهِ طرف که همین حالا حل شد، نه با
سطرِ~\ref{ln:zadv} چون نامش $Z$ است، و نه با $\pars$، که فرضِ (c) آن را
از هویت‌های $Z$ تهی می‌کند. سوم، آیا از $\opl{Guard}$ در بازپخش می‌گذرد:
$P \in \Adv.\Ps$ در سطرِ~\ref{ln:interface}، $\Zpid \subseteq
\Adv.\admits$ با طرف‌های برابر در سطرِ~\ref{ln:caller}، و $\op{global}$ در
سطرِ~\ref{ln:session}. آن‌گاه بازپخش درست همان فراخوانی را می‌گذارد که چپ
گذاشت، زیرِ درست همان ادعایی که چپ کرد، پس هر دو سو با همان عطف‌ها در همان
فراخوانده روبه‌رو می‌شوند. فراخوانی که \emph{به} هویتِ واگذارشده می‌رسد
راهِ دیگر را با همان سه آزمون می‌پیماید: زیرِ نگهبانی که در چپ از آن گذشت
به بازپخش می‌رسد، با فراخواننده‌اش در بارِ پیام به $\Zenv$ در $(Z,P)$
سپرده می‌شود --- که خموش‌کنندهٔ خودِ بازپخش و $\opl{Guard}_{\Zenv}$ آن را
می‌پذیرند، با $\Apid \subseteq \Zenv.\admits$ و $\op{global}(\Zenv)$ ---
و از راهِ نگهبانِ همان واسط به واسطِ عضوِ میهمان در $(A,P)$ می‌رسد، همان
نگهبانی که در چپ از آن گذشت.

\emph{ردها.} یک مقدار می‌تواند متفاوت باشد. آنجا که فراخوانی که عضوِ
واگذارنده می‌گذارد رد می‌شود، چپ $\rej$ را از فراخوانِ خودش می‌خواند، و در
راست بازپخش همان $\rej$ را برمی‌گرداند و بخشِ~\ref{sec:execution} آن را
به‌صورتِ $\none$ به عضوِ میهمان تحویل می‌دهد. فرضِ (a) درست همان است که این
را می‌بندد: ماشینِ بی‌اعتنا به رد به آن دو یکسان پاسخ می‌دهد، پس توزیعش
تغییر نمی‌کند. فرضِ (b) همان است که جهتِ دیگرِ همین مسئله را می‌بندد،
یعنی فراخوانِ پاسخگو روی هویتی واگذارشده: $\opl{Respond}$ هر فراخوانی را
که پاسخ‌دهنده می‌گذارد رد می‌کند، و بازپخش باید یکی بگذارد تا اصلاً پاسخ
دهد، پس فراخوانِ پاسخگو روی جایگاهی بازپخش‌شده $\none$ پاسخ می‌گرفت آنجا
که چپ پاسخی با محتوا داد --- و همین است که چرا هیچ هستهٔ $R$ نمی‌تواند
یکی بگذارد. هیچ‌چیزِ دیگری از کارِ بازپخش دیدنی نیست: $\op{Corrupt}$ از
این دور جان سالم می‌برد، چون دروازه‌اش $\id'.F = A$ با ادعای بازپخش
برآورده می‌شود، و آن همان یک فراخوانی است که بازنویسی باید زیرِ هویتِ $A$
بگذاردش و دسته نمی‌تواند.

\emph{ژتون.} $\opl{Exec}$ یک عمل دارد، گذاشتنِ فراخوان، و ادامهٔ آن همان
جایی است که ژتون به آن بازمی‌گردد. در راست، فراخوان‌های درونیِ $\Zenv^{H}$
دیگر به آن هندلر نمی‌رسند، اما تعریفِ~\ref{def:bundle} ژتون را بینِ
ماشین‌های میهمان همان‌گونه که $\opl{Exec}$ می‌کند دست‌به‌دست می‌کند، پس هر
فراخوانندهٔ درونی روی ادامهٔ خودش معلق می‌شود و با همان پاسخ از سر گرفته
می‌شود؛ و هر فراخوانِ بازپخش‌شده در راست یک قاب می‌افزاید که دقیقاً یک بار
از سر گرفته می‌شود، و این همان حذفی است که تعریفِ~\ref{def:regcorr} اجازه
می‌دهد. دو ویژگیِ $\opl{Exec}$ به کار می‌رود، و
بخشِ~\ref{sec:execution} هر دو را به هر نمونه‌ای می‌دهد: ماشین می‌تواند
چند ادامهٔ معلق داشته باشد، و می‌تواند در حالِ تعلیق دوباره وارد شود. دومی
همان است که فراخوانی که از $R$ می‌رسد لازم دارد، در حالی که خودِ
$\Zenv^{H}$ چشم‌به‌راهِ فراخوانی بیرونی است. در هر لحظه و در هر دو سو،
دقیقاً یک ماشین ژتون را در دست دارد.

پس هر ماشین همان ورودی‌ها را به همان ترتیب می‌بیند و از همان نوار پاسخ
می‌دهد، پس تناظر در هر فعال‌سازی نگه داشته می‌شود و آن دو اجرا به‌عنوانِ
توزیع توافق دارند --- به‌ویژه بر احتمالِ اینکه $\Zenv$ مقدارِ $1$
برگرداند.
\end{proof}

از دید نموداری، این لم یک تساوی است با دو حالتش پهلوی هم: عضوِ
نگه‌داشته یک‌جا در دسته محو می‌شود، و عضوِ واگذارشده از راهِ بازپخشی به
بیرون پل می‌خورد.

\begin{center}
%% sizes itself, so it opts out of the fitting hook
\fitdiagramfalse
\resizebox{\linewidth}{!}{%
\begin{tikzpicture}[
  peerbox/.style={draw=accent, rounded corners=1pt, inner sep=3pt, font=\small, minimum height=6mm},
  absorbedbox/.style={draw=black!55, dashed, rounded corners=2pt, inner sep=2mm},
  wire/.style={draw=black!70},
  unionsign/.style={font=\small},
  rel/.style={font=\normalsize},
  taglabel/.style={font=\tiny, color=accentsoft},
]

%% ===== LEFT: before absorption, four flat peers =====
\node[peerbox] (Zl) {$\Zenv$};
\node[unionsign, right=2.5mm of Zl] (u1) {$\cup$};
\node[peerbox, right=2.5mm of u1] (Hk) {$H_{\mathrm{keep}}$};
\node[unionsign, right=2.5mm of Hk] (u2) {$\cup$};
\node[peerbox, right=2.5mm of u2] (Hc) {$H_{\mathrm{cede}}$};
\node[unionsign, right=2.5mm of Hc] (u3) {$\cup$};
\node[peerbox, right=2.5mm of u3] (Rl) {$R$};

%% ===== relation =====
\node[rel, right=14mm of Rl] (eqsign) {$=$};
\node[taglabel, above=0.5mm of eqsign] {(a)(b)(c)};

%% ===== RIGHT: after absorption, Z^H bundle plus a relay for each ceded id =====
\node[peerbox, right=14mm of eqsign] (Zr) {$\Zenv$};
\node[peerbox, right=2mm of Zr] (Hkr) {$H_{\mathrm{keep}}$};
\node[peerbox, right=2mm of Hkr] (Hcr) {$H_{\mathrm{cede}}$};
\node[absorbedbox, fit=(Zr)(Hkr)(Hcr)] (ZH) {};
\node[taglabel, above=0.5mm of ZH] {$\Zenv^{H}$};

\node[peerbox, right=16mm of ZH] (relay) {$\Dum$};
\draw[wire, <->] (Hcr.east) -- (relay);
\node[peerbox, right=8mm of relay] (Rr) {$R$};
\draw[wire, <->] (relay) -- (Rr);

\end{tikzpicture}%
}
\fitdiagramtrue
\end{center}

عضوی از $H_{\op{keep}}$ --- مثلِ کلِ $\rho \setminus \varphi$ در شکلِ
قضیهٔ~\ref{thm:comp}ِ بالا --- یک‌جا جذب می‌شود: $\Zenv^{H}$ هویتش را
اشغال می‌کند، پس $R$ تنها به‌عنوانِ بخشی از دسته به او می‌رسد. عضوی از
$H_{\op{cede}}$ هم میهمان می‌شود و روی همان نوار می‌دود، اما دسته هویتش را
اشغال نمی‌کند: بازپخشِ $\Dum$ (همان بدلِ بخشِ~\ref{app:dummy}) به‌جای او
می‌ایستد، فراخوان‌های $R$ را با برچسبِ فراخواننده‌شان به $\Zenv$ می‌سپارد و
فراخوان‌های بیرونیِ عضوِ میهمان را زیرِ ادعای اصلی‌اش از نو صادر می‌کند، پس
هیچ‌یک از دو سو نمی‌تواند بفهمد که بازپخشی میانشان نشسته است. فرض‌های
(a)--(c)، که تنها از $H_{\op{cede}}$ خواسته می‌شوند، همان‌هایی هستند که
این را تمایزناپذیر نگه می‌دارند: (a) درست همان بی‌اعتنایی به رد است، چون
$\Dum$ هر $\rej$ی را که بازپخش می‌کند به $\none$ بدل می‌کند؛ و (b) و (c)
فراخوانِ پاسخگو و برخوردِ ادعای خودِ $\Zenv$ را ممنوع می‌کنند، آن دو راهِ
دیگری که پرشِ اضافیِ بازپخش می‌توانست از آن‌ها پیدا شود. شکلِ
قضیهٔ~\ref{thm:comp} همان حالتِ $H_{\op{cede}} = \emptyset$ است؛ و حالتِ
آینه‌ای، یعنی یک اشغال‌گرِ واگذارشده که به‌جای کلِ جایگاهِ حریف می‌ایستد
($H_{\op{keep}} = \emptyset$)، همان است که ساختِ حریفِ بدلِ
بخشِ~\ref{app:dummy} رویش می‌دود.

دو نیمهٔ این لم جدا از هم به کار می‌روند. جذبِ یک پروتکل هیچ واگذار
نمی‌کند، پس هر سه فرض می‌افتند و آنچه می‌ماند یک کمانک‌گذاریِ محض است؛ و
این همان حالتی است که قضیهٔ~\ref{thm:comp} رویش می‌دود، و ارزشِ آن را
دارد که جداگانه بیان شود، برای همین تأکید که اصلاً هیچ هزینه‌ای ندارد.

\begin{proposition}[بازگروه‌بندی]\label{prop:regroup}
بگذارید $\rho$ و $\pi$ سیستم باشند، $\varphi \subseteq \rho$، جانشانیِ
$\varphi$ با $\pi$ در $\rho$ خوش‌ساخت باشد، و $\Zenv$ محیطی باشد که در هیچ
هویتی از $\IDs(\rho) \cup \IDs(\pi)$ میهمان نمی‌کند. آن‌گاه برای $\sigma$ی
برابر با $\pi$ یا $\varphi$، و برای هر اشغال‌گرِ $\mathcal{X}$ از جایگاهِ
حریف (تعریفِ~\ref{def:occ})،
\[
  \opl{Exec}\bigl(\rho[\sigma/\varphi],\,\Zenv,\,\mathcal{X},\,\Corr\bigr)
  \qquad\text{and}\qquad
  \opl{Exec}\bigl(\sigma,\,\Zenv^{\rho \setminus \varphi},\,\mathcal{X},\,\Corr\bigr)
\]
همسان‌توزیعند --- به‌عنوانِ توزیع برابرند، و نه تنها نزدیک.
\end{proposition}

\begin{proof}
لمِ~\ref{lem:absorb} در $H = \rho \setminus \varphi$ با $H_{\op{cede}} =
\emptyset$، به‌گونه‌ای که $\Zenv^{H}$ همان محیطِ جذب‌شدهٔ
تعریفِ~\ref{def:absorb} باشد و هیچ فرضی از آن لم به کار نیاید: (a)، (b) و
(c) از اعضای واگذارشده خواسته می‌شوند و هیچ عضوِ واگذارشده‌ای نیست. توجه
کنید که $\rho[\varphi/\varphi] = \rho$، پس حالتِ $\sigma = \varphi$ همان
سوی ایده‌آل است. آن دو خانوادهٔ اشغال‌گر مجازند: گردایهٔ سمتِ چپ در هر دو
حالت یک سیستم است، پس هیچ هویتی دو بار اشغال نمی‌شود --- برای $\sigma =
\pi$ چون خوش‌ساختی می‌گوید $\rho \setminus \varphi$ و $\pi$ هیچ شناسهٔ
فرایندِ مشترکی ندارند، و برای $\sigma = \varphi$ چون آن گردایه خودِ $\rho$
است --- و در سمتِ راست، فرضِ میهمانی روی $\Zenv$ دسته را از $\IDs(\sigma)$
برکنار نگه می‌دارد.
\end{proof}

اکنون قضیه می‌گوید که جانشین‌کردنِ یک زیرسیستم به‌جای دیگری در درونِ یک
پروتکل، هیچ مزیتی هزینه نمی‌کند.

به‌زبانِ غیررسمی، قضیهٔ~\ref{thm:comp} می‌گوید: اگر $\pi$ کارکردِ $\varphi$
را در $\varepsilon$ UC-تقلید کند، آن‌گاه $\rho[\pi/\varphi]$ کارکردِ $\rho$
را در همان $\varepsilon$ UC-تقلید می‌کند. رده‌های حریف و شبیه‌ساز
بی‌تغییر منتقل می‌شوند؛ و محیط‌های نتیجه آن‌هایی هستند که برای زوجِ
$(\rho[\pi/\varphi],\rho)$ مجازند و صورتِ جذب‌شده‌شان $\Zenv^{\rho
\setminus \varphi}$ در ردهٔ فرض می‌افتد.

\begin{theorem}[ترکیب]\label{thm:comp}
بگذارید $\rho$ و $\pi$ سیستم باشند، $\varphi \subseteq \rho$، و جانشانیِ
$\varphi$ با $\pi$ در $\rho$ به معنای تعریفِ~\ref{def:wf} خوش‌ساخت باشد.
فرض کنید
\[
  \pi \ \text{کارکردِ} \ \varphi \ \text{ را در } \varepsilon \text{ و در برابرِ } (\Aset,\SimSet,\ZenvSet) \text{ UC-تقلید می‌کند} .
\]
آن‌گاه
\[
  \rho[\pi/\varphi] \ \text{کارکردِ} \ \rho \ \text{ را در } \varepsilon \text{ و در برابرِ } (\Aset',\SimSet',\ZenvSet') \text{ UC-تقلید می‌کند} ,
\]
که در آن
\[
  \Aset' \;=\; \Aset, \qquad
  \SimSet' \;=\; \SimSet, \qquad
  \ZenvSet' \;=\; \bigl\{\, \Zenv \in \ZenvSet_{\rho[\pi/\varphi],\,\rho} \;:\; \Zenv^{\rho \setminus \varphi} \in \ZenvSet \,\bigr\} .
\]
و به‌صراحت،
\[
  \begin{gathered}
  \forall\, \Adv \in \Aset' \ \ \exists\, \Sim \in \SimSet' \ \ \forall\, \Zenv \in \ZenvSet' \ : \\[3pt]
  \advtg^{\op{uc}}_{\rho[\pi/\varphi],\,\rho}(\Zenv,\Adv,\Sim)
  \;\leq\;
  \advtg^{\op{uc}}_{\pi,\varphi}\bigl(\Zenv^{\rho \setminus \varphi},\Adv,\Sim\bigr)
  \;\leq\; \varepsilon .
  \end{gathered}
\]
\end{theorem}

\begin{proof}
% BEGIN overview thm:comp
سه نسبت، از سه گونه. رده‌های حریف و شبیه‌ساز با رونویسی منتقل می‌شوند:
حریف هر سیستمی را که حمله کند همان یک ماشین است، و هیچ‌چیز از آنچه
شبیه‌ساز را نوع‌بندی می‌کند از زوجی که برایش شبیه‌سازی می‌شود نام
نمی‌برد. ردهٔ محیط تنها چیزی است که حساب می‌شود و نه رونویسی، چون پیش‌نگارهٔ
ردهٔ فرض زیرِ جذب است. آن‌گاه کران دو گام است.
گزارهٔ~\ref{prop:regroup}، که یک بار در هر سو به کار می‌رود، مزیتِ بیرونی
را \emph{دقیقاً} به مزیتِ درونی بدل می‌کند --- چون همان یک ماشینِ جذب‌شده
در هر دو سو خدمت می‌کند، و همین است که مجموعهٔ بیرونی‌اش برای مجازکردنش
انتخاب شده بود --- و فرضِ تقلید کرانی بر مزیتِ درونی می‌گذارد. پس نسبتِ
اولِ آن نمایش یک تساوی است، که تنها برای یک‌دستی با نسبتِ دوم به‌صورتِ
نامساوی نشان داده شده.
% END overview thm:comp
آن سه نسبت از سه گونه‌اند. حریف‌ها بی‌تغییر منتقل می‌شوند چون آن فروکاست
همان یک ماشین را تحویل می‌دهد: $\Adv \in \Aset'$ که به $\rho[\pi/\varphi]$
حمله می‌کند، به‌عنوانِ $\Adv \in \Aset$ که به $\pi$ حمله می‌کند خوانده
می‌شود. هر $\Aset' \subseteq \Aset$ به کار می‌آمد، و تساوی بزرگ‌ترین چنین
انتخابی است. شبیه‌سازها به دلیلِ آینه‌ای منتقل می‌شوند: $\Sim \in \SimSet$
که فرض فراهمش می‌کند به‌عنوانِ شبیه‌سازِ نتیجه پس داده می‌شود، پس هر
$\SimSet' \supseteq \SimSet$ به کار می‌آمد، و تساوی کوچک‌ترین است. و اینکه
اصلاً می‌شود پسش داد، همان نکتهٔ~\ref{rem:advpars} است: هیچ‌چیز از آنچه
شبیه‌ساز را نوع‌بندی می‌کند از زوجی که برایش شبیه‌سازی می‌شود نام نمی‌برد،
پس شبیه‌سازی برای $(\pi,\varphi)$ از پیش در جایگاهِ
$(\rho[\pi/\varphi],\rho)$ می‌ایستد.

محیط‌ها همان یک جایی هستند که چیزی حساب می‌شود و نه رونویسی. اینجا
$\ZenvSet$ و $\ZenvSet'$ رده‌هایی برای دو زوجِ متفاوت از سیستم‌هایند، و
$\ZenvSet'$ پیش‌نگارهٔ $\ZenvSet$ زیرِ جذب است،
\[
  \ZenvSet' \;=\; \bigl((\cdot)^{\rho \setminus \varphi}\bigr)^{-1}(\ZenvSet) \;\cap\; \ZenvSet_{\rho[\pi/\varphi],\,\rho} ,
\]
پس محیطِ بیرونی درست وقتی به شمار می‌آید که آنچه در جذب به آن بدل می‌شود
یکی از آن‌هایی باشد که فرض می‌پوشاندش. این یک ردهٔ اصیل است، چون بنا بر
ساخت $\ZenvSet' \subseteq \ZenvSet_{\rho[\pi/\varphi],\rho}$. و وقتی
$\ZenvSet$ بزرگ‌ترین ردهٔ مجاز یعنی $\ZenvSet_{\pi,\varphi}$ باشد،
گزارهٔ~\ref{prop:absorb} شرطِ دوم را زائد می‌کند و پیش‌نگاره کلِ
$\ZenvSet_{\rho[\pi/\varphi],\rho}$ است: تقلید در برابرِ بزرگ‌ترین رده،
بزرگ‌ترین رده را باز می‌دهد.

گزارهٔ~\ref{prop:absorb} مزیتِ سمتِ راست را هم تعریف‌شده می‌کند، بی‌آنکه
فرضی از بستار بر محیط‌ها یا نسبتی میانِ $\IDs(\pi)$ و $\IDs(\varphi)$ لازم
باشد. نامساویِ اول همان گزارهٔ~\ref{prop:regroup} است که دو بار به کار
رفته، با $\sigma = \pi$ و $\mathcal{X} = \Adv$ در سوی واقعی و $\sigma =
\varphi$ و $\mathcal{X} = \Sim$ در سوی ایده‌آل. با تفریق،
\[
  \advtg^{\op{uc}}_{\rho[\pi/\varphi],\,\rho}(\Zenv,\Adv,\Sim)
  \;=\;
  \advtg^{\op{uc}}_{\pi,\varphi}\bigl(\Zenv^{\rho \setminus \varphi},\Adv,\Sim\bigr) ,
\]
پس در واقع یک تساوی است: آن دو آزمایش یک آزمایش‌اند که به‌گونهٔ دیگری
کمانک‌گذاری شده، و همان یک $\Sim$ در هر دو سو خدمت می‌کند چون اجرای
ایده‌آل با همان یک جذب بازگروه‌بندی می‌شود. نامساویِ دوم همان فرضِ تقلید
است که در $\Zenv^{\rho \setminus \varphi}$ به کار رفته، و آن در $\ZenvSet$
است چون $\Zenv$ در $\ZenvSet'$ است. قضیه نسبتِ اول را تنها برای یک‌دستی
با نسبتِ دوم به‌صورتِ نامساوی نشان می‌دهد.
\end{proof}

از دید نموداری، با تعیینِ یک $\Adv \in \Aset'$، یک $\Sim \in \SimSet'$ که
خدمت می‌کند و یک $\Zenv \in \ZenvSet'$، اثبات همان مربعِ جابه‌جاییِ زیر
است، که به سبکِ نمودارهای ریسمانیِ~\cite{FKKKW26} کشیده شده: جعبهٔ توپُر
یک سیستم است، سیم یک واسط، و جعبه‌ای که از پایینِ سیستم آویخته همان ماشینی
است که در آن لحظه جایگاهِ حریفانهٔ آن را اشغال کرده است.

\begin{center}
\compositionsquare
\figcaption{fig:comp}{جانشانی، به‌صورتِ یک مربعِ جابه‌جایی. خواندنِ هر ستون از بالا به پایین یک اجراست؛ جعبهٔ نقطه‌چین زوجی را نشان می‌دهد که مبادله می‌شود. هر دو یالِ افقی جذبند، پس هر دو تساویند و نه کران، و یالِ راست همان فرض است. آن‌گاه نتیجهٔ سمتِ چپ همان کمیت است، و نه تنها کمیتی با همان کران. (قضیهٔ~\ref{thm:comp}.)}
\end{center}

یال‌های بالا و پایین همان گزارهٔ~\ref{prop:regroup}اند، که به
$(\pi,\Adv)$ و به $(\varphi,\Sim)$ به کار رفته --- بازگروه‌بندی مفت است،
پس هر دو تساویند و نه کرانی صرف، و همین است که چرا جعبهٔ خط‌چینِ سمتِ راست
(همان محیطِ جذب‌شدهٔ $\Zenv^{\rho \setminus \varphi}$ از
گزارهٔ~\ref{prop:regroup}) تنها چیزی است که میانِ یک ستون و همسایه‌اش
تغییر می‌کند. یالِ راست همان فرضِ قضیه است که به آن محیطِ جذب‌شده به کار
رفته. آن‌گاه یالِ چپ --- نتیجهٔ قضیه --- وادار می‌شود: تعقیبِ مربع نشان
می‌دهد آن دو $\approx_\varepsilon$ نه تنها هر دو با همان $\varepsilon$
کران‌دار شده‌اند، بلکه بنا بر آن دو تساوی، عیناً همان یک کمیتند. و
ترکیب‌کردنِ $n$ مربعِ این‌چنین در یک زنجیره، یک دورگه در هر گام، درست همان
قضیهٔ~\ref{thm:parcomp}ِ زیر است.

اکنون جانشانی و ترایایی در صورتی ترکیب می‌شوند که ترکیب واقعاً در آن به
کار می‌رود. گزارهٔ~\ref{prop:trans} همراهِ متر در بالا بیان شد، چون
اثباتش به هیچ‌یک از این ماشین‌آلات نیاز نداشت؛ و اینجا همان جایی است که
کارش را می‌کند.

\begin{corollary}[ترکیب با یک مشخصه]\label{cor:comp}
بگذارید $\pi$، $\rho$ و $\psi$ سیستم باشند، $\varphi \subseteq \rho$،
جانشانیِ $\varphi$ با $\pi$ در $\rho$ خوش‌ساخت باشد، و $\pi$ کارکردِ
$\varphi$ را در $\varepsilon_1$ و در برابرِ $(\Aset,\SimSet,\ZenvSet)$
UC-تقلید کند. $\ZenvSet'$ را برای ردهٔ قضیهٔ~\ref{thm:comp} بنویسید. اگر
$\rho$ کارکردِ $\psi$ را در $\varepsilon_2$ و در برابرِ
$(\SimSet,\SimSet',\ZenvSet')$ UC-تقلید کند، آن‌گاه $\rho[\pi/\varphi]$
کارکردِ $\psi$ را در $\varepsilon_1 + \varepsilon_2$ و در برابرِ
$(\Aset,\SimSet',\ZenvSet')$ UC-تقلید می‌کند.
\end{corollary}

\begin{proof}
بنا بر قضیهٔ~\ref{thm:comp}، $\rho[\pi/\varphi]$ کارکردِ $\rho$ را در
$\varepsilon_1$ و در برابرِ $(\Aset,\SimSet,\ZenvSet')$ UC-تقلید می‌کند، و
$\ZenvSet' \subseteq \ZenvSet_{\rho[\pi/\varphi],\rho}$. فرضِ دوم حکمی از
تقلید در برابرِ همان رده است، پس $\ZenvSet' \subseteq \ZenvSet_{\rho,\psi}$
هم برقرار است. پس هر دو شمولی که گزارهٔ~\ref{prop:trans} لازم دارد
برقرارند، و آن گزاره به آن دو حکم با $\rho$ در میانه اِعمال می‌شود.
\end{proof}

این همان صورتی است که ترکیب در آن به کار می‌رود. آدم $\rho$ را با
زیرروالِ ایده‌آل‌شدهٔ $\varphi$ سرِ جایش تحلیل می‌کند و اثبات می‌کند که هر
مشخصهٔ $\psi$ی که بخواهد را تقلید می‌کند --- در برابرِ رده‌هایی که نتیجه
نامشان می‌برد و حریف‌هایشان شبیه‌سازهای گامِ اول‌اند؛ و آن‌گاه نتیجه، آن
حکم را به سیستمی منتقل می‌کند که واقعاً می‌دود. یک فرض به‌سادگی از چشم
می‌افتد: تقلید در برابرِ $\ZenvSet'$، به‌عنوانِ شرطِ جنبیِ
تعریفِ~\ref{def:uc}، ادعا می‌کند که $\ZenvSet' \subseteq
\ZenvSet_{\rho,\psi}$ --- پس نام‌های $\psi$ باید زیرِ $\pars$ِ هر محیطی
باشند که آن رده نگه می‌دارد، و هر جا نباشند آن نتیجه تُهی است.

%% ==================================================================
\section{ترکیبِ موازی}\label{app:parallel}

قضیهٔ~\ref{thm:comp} یک زیرسیستم را جانشین می‌کند. جانشین‌کردنِ چند
زیرسیستم یک‌جا استدلالِ تازه‌ای نیست، بلکه $n$ بار به‌کاربستنِ همان یکی در
امتدادِ زنجیره‌ای از دورگه‌هاست؛ و آنچه دقت می‌خواهد، حساب‌داریِ آن سه رده
در طولِ زنجیره است.

فرض‌های مجزابودنِ زیر تحت‌اللفظی‌اند، و تنها به‌سببِ قراردادِ
فصلِ~\ref{sec:functionalities} برآورده‌شدنی‌اند: دفتر ماشینی از اجراست و
عضوِ هیچ سیستمی نیست، و هرگز بیش از یکی از آن نیست. اگر عضو شمرده می‌شد،
یک‌جا در هر $\varphi_i$ و هر $\pi_k$ می‌افتاد، پس هیچ دو زیرسیستمی از
$\rho$ مجزا نمی‌بودند و هیچ دو $\pi_k$ شناسه‌های فرایندِ مجزا حمل
نمی‌کردند؛ آن فرض‌ها برای $n \geq 2$ برآورده‌نشدنی می‌شدند و باید به‌پیمانهٔ
$\Corr$ خوانده می‌شدند. بیرون نگه داشتنِ دفتر از سیستم‌ها همان است که
می‌گذارد آن‌ها همان‌گونه که نوشته شده‌اند خوانده شوند.

\begin{definition}[جانشانیِ هم‌زمان]\label{def:multirepl}
بگذارید $\rho$ سیستمی باشد و $\varphi_1,\dots,\varphi_n \subseteq \rho$
زیرسیستم‌هایی دوبه‌دو مجزا. برای سیستم‌های $\pi_1,\dots,\pi_n$،
\[
  \rho[\pi_1/\varphi_1,\dots,\pi_n/\varphi_n] \;:=\; \Bigl(\rho \setminus \bigcup_{i=1}^{n} \varphi_i\Bigr) \cup \bigcup_{i=1}^{n} \pi_i .
\]
\end{definition}

برای $0 \leq k \leq n$ بنویسیم $\rho_k :=
\rho[\pi_1/\varphi_1,\dots,\pi_k/\varphi_k]$، به‌گونه‌ای که $\rho_0 = \rho$
و $\rho_n = \rho'$. چون $\varphi_i$ها دوبه‌دو مجزایند، $\varphi_k$ از
$k-1$ جانشانیِ اول دست‌نخورده می‌ماند، پس $\varphi_k \subseteq \rho_{k-1}$
و
\[
  \rho_k \;=\; \rho_{k-1}[\pi_k/\varphi_k] .
\]
پس دورگه‌های متوالی تنها در یک جانشانی تفاوت دارند، و همین می‌گذارد
قضیهٔ~\ref{thm:comp} در هر گام کار کند؛ و اثباتِ زیر زنجیره را به‌صورتِ
$n$ جهشِ بازی می‌پیماید، یکی به‌ازای هر جانشانی. جفت‌کردنِ گام‌ها به این هم
نیاز دارد که رده‌ها تنظیم‌شدنی باشند، و هستند.

\begin{lemma}[یکنوایی]\label{lem:mono}
فرض کنید $\pi$ کارکردِ $\varphi$ را در $\varepsilon$ و در برابرِ
$(\Aset,\SimSet,\ZenvSet)$ UC-تقلید می‌کند. آن‌گاه همین کار را در برابرِ
$(\Aset^{-},\SimSet^{+},\ZenvSet^{-})$ هم می‌کند، برای هر $\Aset^{-}
\subseteq \Aset$، هر ردهٔ شبیه‌سازِ $\SimSet^{+} \supseteq \SimSet$ و هر
$\ZenvSet^{-} \subseteq \ZenvSet$.
\end{lemma}

\begin{proof}
تعریفِ~\ref{def:uc} حکمی از گونهٔ $\forall\exists\forall$ روی آن سه رده
است. تنگ‌کردنِ دامنهٔ هر یک از دو سورِ عمومی و فراخ‌کردنِ دامنهٔ سورِ وجودی،
هر دو آن را نگه می‌دارند، و $\ZenvSet^{-} \subseteq \ZenvSet \subseteq
\ZenvSet_{\pi,\varphi}$ نتیجه را حکمی مجاز از تقلید نگه می‌دارد.
\end{proof}

به‌زبانِ غیررسمی، قضیهٔ~\ref{thm:parcomp} می‌گوید جانشین‌کردنِ $n$
زیرسیستمِ دوبه‌دو مجزا یک‌جا، جمعِ خطاهای فردی را هزینه می‌کند: این همان
قضیهٔ~\ref{thm:comp} است که $n$ بار در امتدادِ زنجیره‌ای از دورگه‌ها به کار
رفته، که حریف‌ها حریف‌های آخرین جانشانی‌اند، شبیه‌سازها شبیه‌سازهای اولین،
و محیط‌ها آن‌هایی که یک‌جا برای هر گام مجازند.

\begin{theorem}[ترکیبِ موازی]\label{thm:parcomp}
بگذارید $\rho$ سیستمی باشد، $\varphi_1,\dots,\varphi_n \subseteq \rho$
دوبه‌دو مجزا، و $\pi_1,\dots,\pi_n$ سیستم‌هایی چنان که $\pi_k$ها
شناسه‌های فرایندِ دوبه‌دو مجزا حمل کنند و، به‌ازای هر $k$، $\pi_k$ هیچ
شناسهٔ فرایندی با $\rho \setminus \varphi_k$ مشترک نداشته باشد. فرض کنید
به‌ازای هر $k$
\[
  \pi_k \ \text{کارکردِ} \ \varphi_k \ \text{ را در } \varepsilon_k \text{ و در برابرِ } (\Aset_k,\SimSet_k,\ZenvSet_k) \text{ UC-تقلید می‌کند} ,
\]
و $\ZenvSet_k'$ را برای ردهٔ متناظری بنویسید که قضیهٔ~\ref{thm:comp} به
گامِ $k$اُم می‌چسباند،
\[
  \ZenvSet_k' \;=\; \bigl\{\, \Zenv \in \ZenvSet_{\rho_k,\,\rho_{k-1}} \;:\; \Zenv^{\rho_{k-1} \setminus \varphi_k} \in \ZenvSet_k \,\bigr\} .
\]
و سرانجام فرض کنید
\[
  \SimSet_{k+1} \subseteq \Aset_k \quad (1 \leq k < n),
  \qquad
  \ZenvSet^{\star} \;\subseteq\; \bigcap_{k=1}^{n} \ZenvSet_k' .
\]
آن‌گاه $\ZenvSet^{\star} \subseteq \ZenvSet_{\rho',\rho}$، و
\[
  \begin{gathered}
  \rho' \;:=\; \rho[\pi_1/\varphi_1,\dots,\pi_n/\varphi_n]
  \ \text{ کارکردِ } \ \rho \\[3pt]
  \text{ را در } \sum_{k=1}^{n} \varepsilon_k
  \text{ و در برابرِ } (\Aset_n,\SimSet_1,\ZenvSet^{\star}) \text{ UC-تقلید می‌کند} .
  \end{gathered}
\]
\end{theorem}

\begin{proof}
% BEGIN overview thm:parcomp
یک استدلالِ دورگه روی $n$ گام، و کار در این است که پیش از برداشتنِ هیچ‌یک
از آن‌ها، بسنجیم هر گام نمونه‌ای مجاز از قضیهٔ~\ref{thm:comp} است. سه
مقدمه این کار را می‌کنند. اولی نشان می‌دهد هر جانشانی خوش‌ساخت است و هر
$\rho_k$ یک سیستم، با استقرا روی $k$، تا قضیهٔ ترکیب در هر گام اِعمال شود.
دومی نشان می‌دهد ردهٔ $\ZenvSet^{\star}$ یک‌جا برای هر زوجِ میانی مجاز است،
و --- با گرفتنِ اجتماع روی $k$ --- برای زوجِ بیرونی هم، تا نتیجه خودش حکمی
مجاز از تقلید باشد. سومی زنجیرهٔ شبیه‌سازها را \emph{رو به پایین} می‌سازد،
با $\Sim_{n+1} := \Adv$ و هر $\Sim_k$ که گامِ $k$اُم برای حریفِ
$\Sim_{k+1}$ فراهم می‌کند، و همین است که $\Sim_1$ را تنها به $\Adv$ وابسته
می‌کند. تنها پس از آن است که دورگه برداشته می‌شود: بازی‌های متوالی در یک
زیرسیستمِ جانشین‌شده و یک حریف تفاوت دارند، آن تفاوت یک بار به‌کاربستنِ
قضیهٔ~\ref{thm:comp} است، و جمعِ $n$ جهش تلسکوپی به $\sum_k \varepsilon_k$
می‌رسد.
% END overview thm:parcomp
اول، هر گام خوش‌ساخت است. اعضای $\rho_{k-1} \setminus \varphi_k$ همان‌های
$\rho \setminus \varphi_k$ هستند که $k-1$ جانشانیِ اول سرِ جایشان
می‌گذارند، به همراهِ $\pi_1,\dots,\pi_{k-1}$؛ و بنا بر فرض $\pi_k$ با
هیچ‌یک از این دو گروه شناسهٔ فرایندی مشترک ندارد. پس جانشانیِ $\varphi_k$
با $\pi_k$ در $\rho_{k-1}$ به معنای تعریفِ~\ref{def:wf} خوش‌ساخت است. هر
$\rho_k$ یک سیستم است، با استقرا روی $k$: $\rho_0 = \rho$ یکی است، و
$\rho_k = \rho_{k-1}[\pi_k/\varphi_k]$ بنا بر گزارهٔ~\ref{prop:repl} از
همان جانشانیِ خوش‌ساخت یکی است. پس قضیهٔ~\ref{thm:comp} و
گزارهٔ~\ref{prop:absorb}، که $\rho_{k-1}$ و $\pi_k$ را سیستم می‌خواهند، به
گامِ $k$اُم اِعمال می‌شوند.

بعد، مجازبودن. بگذارید $\Zenv \in \ZenvSet^{\star}$ و $k$ را ثابت کنید. از
$\ZenvSet^{\star} \subseteq \ZenvSet_k' \subseteq
\ZenvSet_{\rho_k,\rho_{k-1}}$ به این می‌رسیم که
\[
  \Zenv.\pars \;\supseteq\; \ncl{\rho_k} \cup \ncl{\rho_{k-1}} ,
\]
پس $\Zenv$ برای زوجِ $k$اُمِ دورگه‌ها مجاز است و مزیتِ
$\advtg^{\op{uc}}_{\rho_k,\rho_{k-1}}$ِ زیر تعریف‌شده است.
گزارهٔ~\ref{prop:absorb}، که به جانشانیِ $\varphi_k$ با $\pi_k$ در
$\rho_{k-1}$ اِعمال می‌شود، آن‌گاه محیطِ جذب‌شده را در
$\ZenvSet_{\pi_k,\varphi_k}$ می‌نشاند: پارامترش پارامترِ $\Zenv$ است، که از
پیش هویت‌های $\pi_k$ و $\varphi_k$ را در بر دارد، پس مزیتِ درونیِ
$\advtg^{\op{uc}}_{\pi_k,\varphi_k}$ هم تعریف‌شده است. با گرفتنِ اجتماعِ
شمول‌های نمایش‌داده‌شده روی همهٔ $k$ها،
\[
  \Zenv.\pars \;\supseteq\; \bigcup_{k=0}^{n} \ncl{\rho_k} \;\supseteq\; \ncl{\rho'} \cup \ncl{\rho} ,
\]
و به همین شکل $\Zenv$ در هیچ هویتی از $\bigcup_k \IDs(\rho_k) \supseteq
\IDs(\rho') \cup \IDs(\rho)$ میهمان نمی‌کند --- بنا بر بندِ میهمانیِ هر
$\ZenvSet_{\rho_k,\rho_{k-1}}$، در $k = n$ و $k = 1$ --- پس هر دو بندِ
تعریفِ~\ref{def:admis} $\ZenvSet^{\star} \subseteq \ZenvSet_{\rho',\rho}$
را می‌دهند و نتیجه خودش حکمی مجاز از تقلید است.

$\Adv \in \Aset_n$ را ثابت کنید و $\Sim_{n+1} := \Adv$ بگذارید. برای $k =
n, n-1, \dots, 1$ بگذارید $\Sim_k \in \SimSet_k$ همان شبیه‌سازی باشد که
گامِ $k$اُم برای حریفِ $\Sim_{k+1}$ فراهم می‌کند. این در هر اندیسی مجاز
است: در $k = n$ چون $\Sim_{n+1} = \Adv \in \Aset_n$، و در $k < n$ چون
$\Sim_{k+1} \in \SimSet_{k+1} \subseteq \Aset_k$. زنجیرهٔ شبیه‌سازها رو به
پایین ساخته می‌شود، هر یک از یکیِ پیش از خودش، و $\Sim_1$ تنها به $\Adv$
وابسته است.

اکنون $\Zenv \in \ZenvSet^{\star}$ را ثابت کنید و این $n+1$ بازی را در
نظر بگیرید
\[
  G_k \;:=\; \opl{Exec}\bigl(\rho_k,\Zenv,\Sim_{k+1},\Corr\bigr), \qquad p_k \;:=\; \Pr[\,G_k = 1\,], \qquad 0 \leq k \leq n .
\]
اینجا $G_0$ اجرای $\rho$ زیرِ $\Sim_1$ است و $G_n$ اجرای $\rho'$ زیرِ
$\Adv$، پس این دنباله از $\rho_0$ آغاز می‌شود و به $\rho_n$ می‌رسد.
بازی‌های متوالی در یک زیرسیستمِ جانشین‌شده و در ماشینی که نقشِ حریف را
بازی می‌کند تفاوت دارند، و آن تفاوت درست یک بار به‌کاربستنِ
قضیهٔ~\ref{thm:comp} است: برای $1 \leq k \leq n$،
\[
  \begin{aligned}
  \lvert\,p_k - p_{k-1}\,\rvert
  \;&=\;
  \advtg^{\op{uc}}_{\rho_k,\,\rho_{k-1}}\bigl(\Zenv,\Sim_{k+1},\Sim_k\bigr) \\
  \;&\leq\;
  \advtg^{\op{uc}}_{\pi_k,\varphi_k}\bigl(\Zenv^{\rho_{k-1} \setminus \varphi_k},\Sim_{k+1},\Sim_k\bigr)
  \;\leq\;
  \varepsilon_k .
  \end{aligned}
\]
آن تساوی، تعریفِ~\ref{def:adv} است که برعکس خوانده شده، چون $G_k$ و
$G_{k-1}$ همان دو اجرایی هستند که مزیتِ
$\advtg^{\op{uc}}_{\rho_k,\rho_{k-1}}$ مقایسه‌شان می‌کند. نامساویِ اول
فروکاستِ قضیهٔ~\ref{thm:comp} است، که در دست است چون بنا بر فرض $\Zenv \in
\ZenvSet^{\star} \subseteq \ZenvSet_k'$، پس لمِ~\ref{lem:mono} گامِ
$k$اُم را حکمی در برابرِ $\ZenvSet^{\star}$ می‌کند. نامساویِ دوم همان
تقلیدِ $\varphi_k$ به‌دستِ $\pi_k$ است، که به حریفِ $\Sim_{k+1}$ و محیطِ
جذب‌شده اِعمال شده.

با جمعِ $n$ جهش و تلسکوپی‌شدن،
\[
  \advtg^{\op{uc}}_{\rho',\,\rho}\bigl(\Zenv,\Adv,\Sim_1\bigr)
  \;=\; \lvert\,p_n - p_0\,\rvert
  \;\leq\; \sum_{k=1}^{n} \lvert\,p_k - p_{k-1}\,\rvert
  \;\leq\; \sum_{k=1}^{n} \varepsilon_k .
\]
و چون $\Sim_1 \in \SimSet_1$ تنها از $\Adv$ انتخاب شد و $\Zenv \in
\ZenvSet^{\star}$ دلخواه بود، این همان حکمِ قضیه است.
\end{proof}

از دید نموداری، جهشِ $k$ همان مربعِ خودِ قضیهٔ~\ref{thm:comp} است (همان
شکلِ بالای اثباتش) با $\varphi \mapsto \varphi_k$، $\pi \mapsto \pi_k$،
$\Adv \mapsto \Sim_{k+1}$، $\Sim \mapsto \Sim_k$، و پوستهٔ $\rho \setminus
\varphi$ی که فراخ شده تا هر جایگاهِ دست‌نخورده را هم با خود ببرد. یک
تصویر، که جهشِ $1$ در آن کامل کشیده شده و جهش‌های $2$ و $3$ از همان جا که
تمام می‌شود جمع‌وجور ادامه می‌یابند، از $\rho_0 = \rho$ تا $\rho_3 =
\rho'$:

\clearpage
\begin{center}
%% Turned on its side, on a page of its own (Task 16) so nothing above or
%% below it competes for height. The chain is 2.9 times as wide as it is
%% tall; rotated, height binds before width does regardless --
%% \textheight:\textwidth is a shallower 1.63:1 than the chain's own
%% 2.9:1, so some margin is left either side no matter the scale, and
%% widening further would only push the top and bottom past \textheight.
%% 0.97 leaves a touch of white at top and bottom rather than bleeding to
%% the trim; it no longer needs to share the page with a paragraph, which
%% is what the previous 0.86 was leaving room for.
\fitdiagramfalse
%% The blank lines inside this picture are commented out, and they have to
%% be: bidi patches graphicx, and its \rotatebox reads its argument in a
%% mode where a blank line ends the box early -- "Paragraph ended before
%% \Grot@box@std was complete", and then sixty-four \node lines land
%% outside any picture. The English file may keep its paragraph breaks
%% here; this one may not.
\rotatebox{90}{\resizebox{0.97\textheight}{!}{%
\begin{tikzpicture}[
  slotbox/.style={draw=accent, rounded corners=1pt, inner sep=1.5pt, font=\scriptsize, minimum height=4mm},
  wire/.style={draw=black!70},
  shell/.style={draw=accent, rounded corners=2pt, inner sep=2mm},
  absorbed/.style={draw=black!55, dashed, rounded corners=2pt, inner sep=2mm},
  role/.style={font=\small},
  rel/.style={font=\normalsize},
  taglabel/.style={font=\tiny, color=accentsoft},
  hl/.style={draw=accentsoft, dotted, thick, rounded corners=1pt, inner sep=1pt},
]
%
%% ================= ROW 1: hop 1 in full =================
%% ----- cell: flat rho_0 -----
\node[slotbox] (r0-rest) {$\rho\setminus\bigcup_i\varphi_i$};
\node[slotbox, right=1mm of r0-rest] (r0-f1) {$\varphi_1$};
\node[slotbox, right=1mm of r0-f1] (r0-f2) {$\varphi_2$};
\node[slotbox, right=1mm of r0-f2] (r0-f3) {$\varphi_3$};
\node[shell, fit=(r0-rest)(r0-f1)(r0-f2)(r0-f3)] (shell0) {};
\node[role, above=6mm of shell0] (Z0) {$\Zenv$};
\draw[wire] (Z0) -- (shell0.north);
\node[role, below=6mm of shell0] (S0) {$\Sim_1$};
\draw[wire] (shell0.south) -- (S0);
\node[taglabel, below=0.5mm of S0] {$(\rho_0=\rho)$};
%
\node[rel, right=13mm of shell0] (eqA) {$=$};
%
%% ----- cell: absorbed, phi_1 exposed -----
\node[slotbox, right=13mm of eqA] (a0-rest) {$\rho\setminus\bigcup_i\varphi_i$};
\node[slotbox, right=1mm of a0-rest] (a0-f2) {$\varphi_2$};
\node[slotbox, right=1mm of a0-f2] (a0-f3) {$\varphi_3$};
\node[fit=(a0-rest)(a0-f2)(a0-f3), inner sep=0pt] (a0-mid) {};
\node[role, above=2mm of a0-mid] (Z0a) {$\Zenv$};
\draw[wire] (Z0a) -- (a0-mid);
\node[absorbed, fit=(Z0a)(a0-rest)(a0-f2)(a0-f3)] (env0) {};
\node[slotbox, below=6mm of env0] (a0-f1) {$\varphi_1$};
\draw[wire] (env0.south) -- (a0-f1);
\node[role, below=6mm of a0-f1] (Sa0) {$\Sim_1$};
\draw[wire] (a0-f1) -- (Sa0);
%
\node[rel, right=13mm of a0-f3] (eqB) {$\approx_{\varepsilon_1}$};
\node[taglabel, below=0.5mm of eqB] {(فرض روی $\pi_1,\varphi_1$)};
%
%% ----- cell: absorbed, pi_1 exposed -----
\node[slotbox, right=13mm of eqB] (a1-rest) {$\rho\setminus\bigcup_i\varphi_i$};
\node[slotbox, right=1mm of a1-rest] (a1-f2) {$\varphi_2$};
\node[slotbox, right=1mm of a1-f2] (a1-f3) {$\varphi_3$};
\node[fit=(a1-rest)(a1-f2)(a1-f3), inner sep=0pt] (a1-mid) {};
\node[role, above=2mm of a1-mid] (Z1a) {$\Zenv$};
\draw[wire] (Z1a) -- (a1-mid);
\node[absorbed, fit=(Z1a)(a1-rest)(a1-f2)(a1-f3)] (env1) {};
\node[slotbox, below=6mm of env1] (a1-p1) {$\pi_1$};
\draw[wire] (env1.south) -- (a1-p1);
\node[role, below=6mm of a1-p1] (Sa1) {$\Sim_2$};
\draw[wire] (a1-p1) -- (Sa1);
%
\node[rel, right=13mm of a1-f3] (eqC) {$=$};
%
%% ----- cell: flat rho_1 -----
\node[slotbox, right=13mm of eqC] (r1-rest) {$\rho\setminus\bigcup_i\varphi_i$};
\node[slotbox, right=1mm of r1-rest] (r1-p1) {$\pi_1$};
\node[slotbox, right=1mm of r1-p1] (r1-f2) {$\varphi_2$};
\node[slotbox, right=1mm of r1-f2] (r1-f3) {$\varphi_3$};
\node[shell, fit=(r1-rest)(r1-p1)(r1-f2)(r1-f3)] (shell1) {};
\node[role, above=6mm of shell1] (Z1) {$\Zenv$};
\draw[wire] (Z1) -- (shell1.north);
\node[role, below=6mm of shell1] (S1) {$\Sim_2$};
\draw[wire] (shell1.south) -- (S1);
%
%% ================= ROW 2: hops 2,3 compact, continuing from rho_1 =================
\node[slotbox, below=34mm of r0-rest] (r1r-rest) {$\rho\setminus\bigcup_i\varphi_i$};
\node[slotbox, right=1mm of r1r-rest] (r1r-p1) {$\pi_1$};
\node[slotbox, right=1mm of r1r-p1] (r1r-f2) {$\varphi_2$};
\node[slotbox, right=1mm of r1r-f2] (r1r-f3) {$\varphi_3$};
\node[shell, fit=(r1r-rest)(r1r-p1)(r1r-f2)(r1r-f3)] (shell1r) {};
\node[role, above=6mm of shell1r] (Z1r) {$\Zenv$};
\draw[wire] (Z1r) -- (shell1r.north);
\node[role, below=6mm of shell1r] (S1r) {$\Sim_2$};
\draw[wire] (shell1r.south) -- (S1r);
\node[hl, fit=(r1r-f2)] {};
\node[taglabel, below=0.5mm of S1r] {(ادامه از $\rho_1$ِ بالا)};
%
\node[rel, right=13mm of shell1r] (eqD) {$\approx_{\varepsilon_2}$};
%
\node[slotbox, right=13mm of eqD] (r2-rest) {$\rho\setminus\bigcup_i\varphi_i$};
\node[slotbox, right=1mm of r2-rest] (r2-p1) {$\pi_1$};
\node[slotbox, right=1mm of r2-p1] (r2-p2) {$\pi_2$};
\node[slotbox, right=1mm of r2-p2] (r2-f3) {$\varphi_3$};
\node[shell, fit=(r2-rest)(r2-p1)(r2-p2)(r2-f3)] (shell2) {};
\node[role, above=6mm of shell2] (Z2) {$\Zenv$};
\draw[wire] (Z2) -- (shell2.north);
\node[role, below=6mm of shell2] (S2) {$\Sim_3$};
\draw[wire] (shell2.south) -- (S2);
\node[hl, fit=(r2-p2)] {};
\node[hl, fit=(r2-f3)] {};
%
\node[rel, right=13mm of shell2] (eqE) {$\approx_{\varepsilon_3}$};
%
\node[slotbox, right=13mm of eqE] (r3-rest) {$\rho\setminus\bigcup_i\varphi_i$};
\node[slotbox, right=1mm of r3-rest] (r3-p1) {$\pi_1$};
\node[slotbox, right=1mm of r3-p1] (r3-p2) {$\pi_2$};
\node[slotbox, right=1mm of r3-p2] (r3-p3) {$\pi_3$};
\node[shell, fit=(r3-rest)(r3-p1)(r3-p2)(r3-p3)] (shell3) {};
\node[role, above=6mm of shell3] (Z3) {$\Zenv$};
\draw[wire] (Z3) -- (shell3.north);
\node[role, below=6mm of shell3] (S3) {$\Adv$};
\draw[wire] (shell3.south) -- (S3);
\node[hl, fit=(r3-p3)] {};
\node[taglabel, below=0.5mm of S3] {$(\rho_3=\rho')$};
%
\end{tikzpicture}%
}}
\fitdiagramtrue
\end{center}
\clearpage

از چپ به راست بخوانید: $\rho_0$ (با همهٔ $\varphi$ها) جذب می‌شود --- باقیِ
واقعیِ $\rho$ و $\varphi_2,\varphi_3$ِ دست‌نخورده همراهش در $\Zenv$ تا
می‌شوند و تنها $\varphi_1$ رو می‌ماند. فرضِ قضیهٔ~\ref{thm:comp} آن را در
$\varepsilon_1$ با $\pi_1$ عوض می‌کند؛ و پس‌دادنِ آن تای‌خورده $\rho_1$ را
از نو سوار می‌کند، و $\Sim$ برای جور‌شدن از $\Sim_1$ به $\Sim_2$ پیش
می‌رود. جهش‌های $2$ و $3$ همان مربعِ یکسانند، تنها فشرده به دو سرِ صافشان
--- کشیدنِ دوبارهٔ آن تکرار بود و نه افزودن --- و جعبه‌های نقطه‌چین در هر
جهش درست نشان می‌دهند کدام جایگاه $\varphi_k$ است (که در آستانهٔ برگشتن
است) و کدام $\pi_k$ (که همین حالا برگشته). جمعِ آن سه جهش و تلسکوپی‌شدن،
چنان که در سطرهای آخرِ اثباتِ بالا، سه $\varepsilon_k$ِ جدا را --- که
هیچ‌یک در خودِ حکمِ قضیه نمی‌آیند --- به آن یک کرانِ $\rho
\approx_{\varepsilon_1+\varepsilon_2+\varepsilon_3} \rho'$ بدل می‌کند.

\heading{فساد در طولِ زنجیره} فساد در طولِ زنجیره همان‌گونه رفتار می‌کند که
در یک گامِ تنها. در درونِ یک بازی، دفتر یک ماشین است در زیرِ $\rho
\setminus \varphi_k$ و $\pi_k$ و $\varphi_k$، به‌یک‌سان، و کلیدش طرف است،
پس یک جهش نمی‌تواند طرفی را نیمه‌کاره فاسد کند؛ و بینِ بازی‌ها دفترها
نمونه‌های جدایند و هیچ‌چیز از جهشی به جهشِ دیگر نمی‌برند. محیط در هر جهش
همان طرف‌ها را فاسد می‌کند، چون همان ماشین است که همان فراخوان‌ها را زیرِ
سطرِ~\ref{ln:corronlya} می‌گذارد؛ جایگاه ناچار نیست چنین کند، پس آن دو
ماشینی که در هر جهش تغییر می‌کنند --- زیرسیستمِ جانشین‌شده و اشغال‌گرِ
جایگاه $\Sim_k \to \Sim_{k+1}$ --- لازم نیست همان طرف‌ها را فاسد کنند؛ اما
محیط $\Cs$ را بی‌واسطه‌گری می‌خواند (نکتهٔ~\ref{rem:statusmed})، پس
$\Sim_k$ای که به‌گونهٔ دیگری فاسد کند با همان $\varepsilon_k$ گرفته
می‌شود، و هم‌اندازگی باز هم نتیجه است و نه فرض. طرف‌هایی که $\rho_{k-1}
\setminus \varphi_k$ خدمتشان می‌دهد یکسان گردانده می‌شوند: قلاب‌های
$\opl{Mediate}$شان محیطِ جذب‌شده را به‌قصدِ جایگاه ترک می‌کنند، و همین است
که چرا مجموعهٔ بیرونیِ تعریفِ~\ref{def:absorb} شناسه‌های $A$ و $C$ را حمل
می‌کند. پس قضیه تنها این را ادعا می‌کند که هیچ محیطی در $\ZenvSet^{\star}$
الگوهای فساد را با بیش از $\sum_k \varepsilon_k$ از هم تشخیص نمی‌دهد.


\begin{remark}[بزرگ‌ترین رده‌ها]\label{rem:parclasses}
اجتماعِ بستارهای دورگه صورتی بسته دارد. هر $\rho_k$ از بخشِ دست‌نخوردهٔ
$\rho$ به همراهِ $\pi_1,\dots,\pi_k$ تشکیل شده، پس
\[
  \bigcup_{k=0}^{n} \ncl{\rho_k} \;=\; \ncl{\rho} \cup \bigcup_{i=1}^{n} \ncl{\pi_i} \;=\; \ncl{\rho'} \cup \ncl{\rho} .
\]
اکنون فرض کنید هر $\ZenvSet_k$ بزرگ‌ترین ردهٔ مجاز یعنی
$\ZenvSet_{\pi_k,\varphi_k}$ باشد. آن‌گاه گزارهٔ~\ref{prop:absorb} شرطِ
دومی که $\ZenvSet_k'$ را تعریف می‌کند زائد می‌کند، پس $\ZenvSet_k' =
\ZenvSet_{\rho_k,\rho_{k-1}}$، و بزرگ‌ترین انتخابِ رواییِ
$\ZenvSet^{\star}$ این است
\[
  \bigcap_{k=1}^{n} \ZenvSet_{\rho_k,\rho_{k-1}}
  \;=\; \bigl\{\, \Zenv : \Zenv.\pars \supseteq \textstyle\bigcup_{k=0}^{n} \ncl{\rho_k} \,\bigr\}
  \;=\; \ZenvSet_{\rho',\rho} .
\]
پس ترکیبِ موازی در برابرِ بزرگ‌ترین رده‌ها، باز بزرگ‌ترین رده را تحویل
می‌دهد، همان‌گونه که قضیهٔ تک‌گامی می‌کند.
\end{remark}

آن سه رده همان‌گونه رفتار می‌کنند که زنجیره نشان می‌دهد. حریف‌ها حریف‌های
\emph{آخرین} جانشانی‌اند، چون استدلال از حریفی در برابرِ $\rho_n = \rho'$
آغاز می‌کند؛ و شبیه‌سازها شبیه‌سازهای \emph{اولین}، چون زنجیره با تولیدِ
شبیه‌سازی در برابرِ $\rho_0 = \rho$ تمام می‌شود. هر ردهٔ میانی دو بار خدمت
می‌کند، به‌عنوانِ شبیه‌سازهای یک گام و حریف‌های گامِ بعد، و همین است که
$\SimSet_{k+1} \subseteq \Aset_k$ ترتیبش می‌دهد. محیط‌ها باید یک‌جا برای
هر گام مجاز باشند، و اشتراک از همین‌جاست. و مزیت جمع‌شدنی است، چنان که
استدلالِ دورگه روی $n$ گام همیشه چنینش می‌کند.

%% ==================================================================
\section{تمامیتِ حریفِ بدل}\label{app:dummy}

حریفِ بدل هر چه محیط بگویدش انجام می‌دهد و هر چه به او گفته شود پس گزارش
می‌کند؛ هیچ راهبردی از خودش ندارد. هر چه حریفِ واقعی به‌دست می‌آورد، محیط
می‌تواند با راندنِ بدل به‌دست آورد، و همین است که چرا می‌شود سور‌زدن روی
حریف‌ها را با سور‌زدن روی محیط‌ها عوض کرد. این بخش او را تعریف می‌کند و
اثبات می‌کند که، برای سیستم‌ها و رده‌هایی که قضیه دقیقشان می‌کند، هیچ از
دست نمی‌رود --- و نکتهٔ~\ref{rem:respdummy} آن یک رژیمی را ثبت می‌کند که
این ادعا به آن نمی‌رسد: هسته‌هایی که فراخوانِ پاسخگو می‌گذارند.

\heading{آنچه چارچوب از پیش تعیین کرده} بدل گونهٔ تازه‌ای از ماشین نیست. او
همان حریفِ بخشِ~\ref{sec:adv} است با هسته‌ای خاص، پس میدان‌های $\Adv$ را
حمل می‌کند و فراخوان‌هایش از پوشانهٔ $\Adv$ می‌گذرند. سه پیامد، شکلِ کد را
پیش از نوشتنِ یک سطرِ آن تعیین می‌کنند.

اول، او تنها یک هویت را می‌تواند ادعا کند. سطرِ~\ref{ln:advsil} از جعبهٔ
$\Adv$ هسته را روی $\{\id'' \ne \id\}$ خموش می‌کند، پس هر فراخوانی که بدل
می‌گذارد $\id$ را ادعا می‌کند، همان هویتی که در آن می‌دود. پس دستوری از
شکلِ «فراخوانِ $\id''.\Op$ \emph{به‌عنوانِ} $\id'''$» را نمی‌شود برای
$\id'''$ِ دلخواه انجام داد: بدل به‌عنوانِ $(A,\id.P)$ سخن می‌گوید یا اصلاً
سخن نمی‌گوید. آنچه محیط انتخاب می‌کند شناسهٔ ادعایی نیست بلکه طرف است، با
دستوردادن به بدل در $(A,P)$ به‌جای $(A,Q)$؛ و ریشهٔ $\Zenv$ می‌تواند برای
هر طرفی عمل کند و پس می‌تواند هر یک را برگزیند.

دوم، طرفِ هدف وادار شده است. محمولِ $\opl{Guard}$ از فراخوانده $\id'.P =
\id.P$ را می‌پرسد، و ادعا $(A,\id.P)$ است، پس بدلِ نشسته در $(A,P)$ به
واسط‌های طرفِ $P$ می‌رسد و به هیچ طرفِ دیگر. هدفی که $\Apid$ را نپذیرد، یا
طرفی درستکار، فراخوان را با سطرِ~\ref{ln:caller} یا
سطرِ~\ref{ln:corrgate} رد می‌کند؛ بدل آن رد را پس می‌دهد، و قاعدهٔ
بخشِ~\ref{sec:execution} آن را به فراخوانندهٔ خودِ بدل به‌صورتِ $\none$
تحویل می‌دهد.

سوم، راهِ او به محیط هم تعیین شده است: $(A,P) \to (Z,P)$، که
$\opl{Guard}_{\Zenv}$ می‌پذیردش، چون $\Apid \subseteq \Zenv.\admits$ و
طرف‌ها جورند.

\heading{کد} هسته بر حسبِ اینکه چه کسی فرایش خوانده توزیع می‌کند. فراخوان
از سوی محیط دستوری است برای انجام‌دادن؛ و فراخوان از سوی هر چیزِ دیگر ---
یک قلابِ واسطه‌گری، یا یک $\Adv(\cdot)$ که در درونِ هسته‌ای گذاشته شده ---
پیامی است برای تحویل‌دادن. در هر دو جهت، بارِ پیام همان شکلی را دارد که
$\opl{Mediate}$ از پیش به کار می‌برد، یعنی یک عمل به همراهِ یک مقدار.

\begin{interface}{حریفِ بدلِ $\Dum$ \\
\params{$\PID := A$, \quad $\Ps := \bits \cup \{A,Z\}$, \quad
$\admits := \Stdpid \cup \Apid \cup \Zpid$, \quad $\uses := \{(\Zenv,\serves)\}$, \quad
$\pars := \none$}}
\opsig{$\id.\op{A}(in)$} \quad از $\id'$
\begin{algorithmic}[1]
\If{$\id'.F = Z$}
  \State $in$ را به‌صورتِ $\bigl(\id''.\Op,\ x\bigr)$ تجزیه کن \label{ln:dumparse}
    \Comment{تجزیه‌ناپذیر: $\none$ برگردان}
  \State \textbf{return} $\id''.\fopl{Op}(x) \text{ as } \id$ \label{ln:dumout}
    \Comment{دستور را انجام بده}
\Else
  \State \textbf{return} $(Z,\id.P).\fopl{Z}\bigl(\id',\ in\bigr) \text{ as } \id$ \label{ln:dumin}
    \Comment{به $\Zenv$ بسپارش، با فراخواننده و همه چیز}
\EndIf
\end{algorithmic}
\end{interface}

سطرِ~\ref{ln:dumout} جهتِ بیرونی است: محیط نامِ واسطِ هدف و یک ورودی را
می‌برد، و بدل درست همان فراخوان را می‌گذارد و درست همان پاسخ را
برمی‌گرداند. چیزی بازرسی نمی‌شود و چیزی به یاد نمی‌ماند --- بدل هیچ حالتی
نگه نمی‌دارد، پس $\op{Initialize}$ ندارد. سطرِ~\ref{ln:dumin} جهتِ درونی
است: هر چه به حریف می‌رسد، به همراهِ هویتی که از آن آمده، به محیط سپرده
می‌شود، و هر چه محیط پاسخ دهد بی‌تغییر به آن فراخواننده برگردانده می‌شود.
آن دو سطر همان یک بازپخش‌اند که در دو جهتِ مخالف خوانده شده، و همین است که
بدل را شفاف می‌کند.

دستوری که سطرِ~\ref{ln:dumparse} نتواند تجزیه‌اش کند با $\none$ پاسخ
می‌گیرد: نگهبان‌ها هر چه را بدل نتواند \emph{بگذارد} با $\rej$ پاسخ
می‌دهند، اما بارِ پیامی که نامِ هیچ فراخوانی را نبرد به مقداری از خودش
نیاز دارد، و $\none$ همان است که ماشینی با هیچ حرفی برمی‌گرداند.

یک جزء ارزشِ درنگ دارد. فراخوانندهٔ $\id'$ در سطرِ~\ref{ln:dumin} به
$\Zenv$ سپرده می‌شود و نه تنها بارِ پیام: بی آن، محیط نمی‌توانست بفهمد
کدام واسط دارد می‌پرسد، و بازپخشی که این را از دست بدهد نمی‌تواند شفاف
باشد.

\heading{نشاندنِ حریف در میانه} تمامیت حکمی است در بابِ بردنِ حریف از اجرا
به درونِ ماشین‌های پیرامونش، و برای هر حرکت نامی لازم دارد. هر دو، یک عمل
دیده‌شده از دو سویند: $\Adv$ میانِ ماشینی و واسطِ بدل نشانده می‌شود، پس
آنچه پیش‌تر به $\Adv$ می‌رسید حالا در درون به او می‌رسد، و آنچه $\Adv$
پیش‌تر می‌گذاشت حالا به‌صورتِ دستوری به $\Dum$ بیرون می‌رود.

\begin{definition}[حریفِ جذب‌شده در محیط]\label{def:zadv}
بگذارید $\Zenv$ محیطی باشد و $\Adv$ حریفی. محیطِ $\Zenv^{\Adv}$ همان دستهٔ
(تعریفِ~\ref{def:bundle}) $\Zenv$ با یک رونوشت از $\Adv$ است، که آن رونوشت
هویت‌هایش را نگه \emph{نمی‌دارد}؛ این دسته تنها $\IDs(\Zenv)$ را اشغال
می‌کند و خروجی و $\pars$ را از $\Zenv$ می‌گیرد. فراخوانی که $\Zenv$ روی
هویتی $A$ می‌گذارد در درون به‌دستِ واسطِ $\Adv$ در آنجا خدمت می‌شود، و
فراخوانی که $\Adv$ روی هویتی $Z$ می‌گذارد به‌دستِ واسطِ $\Zenv$. هر
فراخوانِ دیگری که $\Adv$ِ میهمان می‌گذارد، روی $\id''.\fopl{Op}(x)$ در
حالِ دویدن در $(A,P)$، به این دستور بازنویسی می‌شود
\[
  (A,P).\fopl{A}\bigl(\id''.\Op,\ x\bigr) \ \text{ as } (Z,P) ,
\]
که کلیدش طرفی است که $\Adv$ در آن می‌دود و هرگز طرفِ هدفش نیست، پس
فراخوانی که $\Adv$ِ جذب‌نشده تنها به‌عنوانِ $(A,P)$ می‌توانست بگذارد در
همان طرف بازپخش می‌شود و درست همان جا که بود پذیرفته یا رد می‌شود.
فراخوانی که از هویتی $A$ با بارِ پیامِ $(\id',in)$ می‌رسد، به واسطِ $\Adv$
در $(A,P)$ سپرده می‌شود، با نگهبان و همه چیز، به‌عنوانِ فراخوانی از
$\id'$.
\end{definition}

\begin{definition}[حریفِ نشانده پیش از شبیه‌ساز]\label{def:sadv}
بگذارید $\Sim$ شبیه‌سازی باشد و $\Adv$ حریفی. شبیه‌سازِ $\Sim[\Adv]$ همان
دستهٔ (تعریفِ~\ref{def:bundle}) $\Sim$ با یک رونوشت از $\Adv$ است؛ این دسته
$\IDs(\Adv)$ را اشغال می‌کند و باقیِ میدان‌های $\Sim$ را حمل می‌کند ---
چون $\Sim$ است و نه $\Adv$ که او به‌جایش می‌ایستد، پس $\Sim[\Adv]$ هر جا
که $\Sim$ شبیه‌ساز باشد شبیه‌ساز است. فراخوان از هویتی $Z$ را $\Adv$ِ
میهمان خدمت می‌کند؛ و فراخوان از هر شناسهٔ فرایندِ دیگری، به‌ویژه یک قلابِ
واسطه‌گری، را $\Sim$. فراخوان‌های $\Adv$ِ میهمان روی هویتی $Z$ --- راهِ
بازگشتش به محیط --- زیرِ $(A,P)$ از دسته بیرون می‌روند به‌سوی $\Zenv$ِ
بیرونی؛ فراخوان‌های بیرونیِ دیگرش به دستورِ تعریفِ~\ref{def:zadv} بدل
می‌شوند، که $\Sim$ آن را در $(A,P)$ به‌عنوانِ فراخوانی از $(Z,P)$ خدمت
می‌کند؛ و فراخوانی که $\Sim$ روی هویتی $Z$ می‌گذارد، پس از گذر از
$\opl{Guard}_{\Zenv}$، به $\Adv$ در $(A,P)$ سپرده می‌شود.
\end{definition}

هیچ‌یک از این دو ساخت گونهٔ تازه‌ای از ماشین نیست. دستهٔ $\Zenv^{\Adv}$ آنچه
محیط اشغال می‌کند را اشغال می‌کند و آنچه $\Zenv$ ادعا می‌کرد را ادعا
می‌کند، پس درست وقتی مجاز است که $\Zenv$ مجاز باشد. او هیچ هویتِ $A$ را
اشغال نمی‌کند، چون آن‌ها در اجرای سمتِ راست مالِ بازپخش‌اند، پس بنا بر
تعریفِ~\ref{def:bundle} $\Adv$ِ میهمان هیچ هویتی از خودش ندارد که فراخوانی
را زیرِ آن بگذارد: و از همین‌جاست آن بازنویسیِ بالا، و از همین‌جاست آن
تعهد --- که در بندِ «جانشینی» در زیر انجام می‌شود --- که ادعای
بازنویسی‌شده را در برابرِ خموش‌کنندهٔ $\Zenv$ بسنجیم. جذب در
فصلِ~\ref{sec:uc} آن حالتِ دیگرِ تعریفِ~\ref{def:bundle} است، و همین است که
چرا کارکردِ میهمان در آنجا به‌جایش ادعایش را نگه می‌دارد. و در همان حال،
$\Sim[\Adv]$ $\IDs(\Adv)$ را اشغال می‌کند، که همان معنای ایستادن در جایگاهِ
حریف است.

\heading{قضیه} تقلید \emph{تنها نسبت به بدل}، همان تعریفِ~\ref{def:uc} است
با ردهٔ حریف که به همان یک ماشینِ $\Dum$ تنگ شده. ادعا این است که هیچ از
دست نمی‌رود.

یک شرطِ ملایم بر حریف لازم است، و آن بهای بازپخش‌بودنِ بدل است.

\begin{definition}[بی‌اعتنا به رد]\label{def:refobl}
حریفی \emph{بی‌اعتنا به رد} است اگر رفتارش تغییر نکند وقتی هر $\rej$ی که
به‌عنوانِ پاسخِ فراخوانی که خودش گذاشته دریافت می‌کند، با $\none$ عوض
شود.
\end{definition}

بدل چنین پاسخی را بازپخش می‌کند، و بنا بر بخشِ~\ref{sec:execution}
\textbf{return}ِ خودش از یک $\rej$ به‌صورتِ $\none$ به فراخواننده می‌رسد؛
پس آنجا که حریفِ جذب‌نشده با گذاشتنِ فراخوانی ردشده $\rej$ را مستقیم
می‌خواند، همان حریف در حالِ راندنِ بدل $\none$ می‌خواند. حریفِ بی‌اعتنا به
رد نمی‌تواند آن دو را از هم تشخیص دهد، و همین است که بازگروه‌بندیِ زیر
لازم دارد. خودِ بدل بی‌اعتنا به رد است --- پاسخش را بازپخش می‌کند بی‌آنکه
رویش شاخه بزند --- پس این فرض برای ردهٔ بدل هیچ هزینه‌ای ندارد، و حریفِ
واقعی‌ای که سیگنالِ ردِ درون‌چارچوبی را بخواند در هر حال یک صنعتِ
مدل‌سازی است و نه یک حمله.

به‌زبانِ غیررسمی، قضیهٔ~\ref{thm:dummy} می‌گوید تقلید تنها نسبت به بدل هیچ
از دست نمی‌دهد: هر چه حریفِ واقعی به‌دست می‌آورد، محیط می‌تواند با راندنِ
بدل به‌دست آورد، و شبیه‌سازِ $\Adv$ همان $\Sim_{\Dum}[\Adv]$ است ---
همان یک شبیه‌سازی که فرض فراهم می‌کند، با $\Adv$ که پیش از آن نشانده شده
--- که تنها به $\Adv$ وابسته است و نه به $\Zenv$.

\begin{theorem}[تمامیتِ حریفِ بدل]\label{thm:dummy}
بگذارید $\pi$ و $\varphi$ سیستم‌هایی باشند که هسته‌هایشان هیچ فراخوانِ
پاسخگو نمی‌گذارند (تعریفِ~\ref{def:respcall})، $\varepsilon \geq 0$،
$\Aset$ ردهٔ حریف‌های بی‌اعتنا به رد باشد
(تعریفِ~\ref{def:refobl})، $\SimSet$ ردهٔ شبیه‌سازها، و $\ZenvSet \subseteq
\ZenvSet_{\pi,\varphi}$. فرض کنید
\[
  \pi \ \text{کارکردِ} \ \varphi \ \text{ را در } \varepsilon \text{ و در برابرِ } (\{\Dum\},\SimSet,\ZenvSet) \text{ UC-تقلید می‌کند} .
\]
آن‌گاه
\[
  \pi \ \text{کارکردِ} \ \varphi \ \text{ را در } \varepsilon \text{ و در برابرِ } (\Aset,\SimSet',\ZenvSet') \text{ UC-تقلید می‌کند} ,
\]
که در آن
\begin{gather*}
  \SimSet' \;=\; \bigl\{\, \Sim[\Adv] \;:\; \Sim \in \SimSet,\ \Adv \in \Aset \,\bigr\},
  \\
  \ZenvSet' \;=\; \bigl\{\, \Zenv \in \ZenvSet_{\pi,\varphi} \;:\;
  \Zenv.\pars \cap \{\, \id : \id.F = Z \,\} = \emptyset \\
  \wedge\ \Zenv^{\Adv} \in \ZenvSet \ \text{به‌ازای هر } \Adv \in \Aset \,\bigr\} .
\end{gather*}
و به‌صراحت، بگذارید $\Sim_{\Dum} \in \SimSet$ همان شبیه‌سازی باشد که فرض
برای $\Dum$ فراهم می‌کند. آن‌گاه
\[
  \begin{gathered}
  \forall\, \Adv \in \Aset \ \ \forall\, \Zenv \in \ZenvSet' \ : \\[3pt]
  \advtg^{\op{uc}}_{\pi,\varphi}\bigl(\Zenv,\Adv,\Sim_{\Dum}[\Adv]\bigr)
  \;\leq\;
  \advtg^{\op{uc}}_{\pi,\varphi}\bigl(\Zenv^{\Adv},\Dum,\Sim_{\Dum}\bigr)
  \;\leq\; \varepsilon ,
  \end{gathered}
\]
پس شبیه‌سازی که تعریفِ~\ref{def:uc} می‌خواهد همان $\Sim_{\Dum}[\Adv]$ است،
که تنها به $\Adv$ وابسته است و نه به $\Zenv$.
\end{theorem}

\begin{proof}
% BEGIN overview thm:dummy
کلِ این استدلال دو بازگروه‌بندی است، یکی به‌ازای هر جهان، که در زیر (R1) و
(R2) نامیده می‌شوند. با پذیرفتنِ آن‌ها، دو مزیتِ آن نمایش جمله‌به‌جمله
توافق می‌کنند و فرض در $\Zenv^{\Adv}$ کرانی بر دومی می‌گذارد، که همان
قضیه است. هیچ‌یک کمانک‌گذاریِ محض نیست --- اجراهای سمتِ راست ماشینی حمل
می‌کنند که اجراهای سمتِ چپ ندارند، یعنی بازپخشی که در هویت‌های واگذارشده
ایستاده --- پس هر دو نیمهٔ \emph{واگذارندهٔ} لمِ~\ref{lem:absorb}اند و نه
نیمهٔ گزارهٔ~\ref{prop:regroup}. (R1) همان لم است که مستقیم به کار رفته، و
سه فرضش همان سه فرضِ قضیه‌اند. (R2) همان استدلال است که دوباره با $\varphi$
به‌جای $\pi$ و $\Sim_{\Dum}$ به‌جای $\Dum$ دوانده شده، و استدلالی از خود
لازم دارد چون در آنجا حریف از دسته‌ای به دستهٔ دیگر می‌رود و نه از اجرا به
بیرون: آنچه اثباتِ آن لم از ماشینِ نشسته در جایگاه به کار برد بازگفته
می‌شود و نشان داده می‌شود که برای $\Sim_{\Dum}[\Adv]$ هم برقرار است. بندِ
پایانی ترتیبِ سورها را می‌سنجد --- $\Sim_{\Dum}[\Adv]$ تنها به $\Adv$
وابسته است، و فرض هیچ از $\Zenv$ نمی‌داند.
% END overview thm:dummy
$\Adv \in \Aset$ و $\Zenv \in \ZenvSet'$ را ثابت کنید، و بگذارید
$\Sim_{\Dum} \in \SimSet$ همان شبیه‌سازی باشد که فرض برای $\Dum$ فراهم
می‌کند. عضویت در $\ZenvSet'$ می‌دهد $\Zenv^{\Adv} \in \ZenvSet$، پس فرض به
آن اِعمال می‌شود، و $\Zenv^{\Adv}.\pars = \Zenv.\pars \supseteq \ncl{\pi}
\cup \ncl{\varphi}$، پس مزیتِ سمتِ راستِ آن نمایش تعریف‌شده است.
قراردادهای گزارهٔ~\ref{prop:regroup} --- سکه‌های مستقل، نوارهای یکی‌نشده،
$\op{Initialize}$ِ بی‌فراخوان --- برقرارند. کار دو بازگروه‌بندی است، یکی
به‌ازای هر جهان:
\begin{align*}
  \opl{Exec}\bigl(\pi,\Zenv,\Adv,\Corr\bigr) &\ \equiv\ \opl{Exec}\bigl(\pi,\Zenv^{\Adv},\Dum,\Corr\bigr) , \tag{R1}\\
  \opl{Exec}\bigl(\varphi,\Zenv,\Sim_{\Dum}[\Adv],\Corr\bigr) &\ \equiv\ \opl{Exec}\bigl(\varphi,\Zenv^{\Adv},\Sim_{\Dum},\Corr\bigr) , \tag{R2}
\end{align*}
که هر یک تساویِ توزیع‌هاست؛ و با پذیرفتنِ آن‌ها، دو مزیت جمله‌به‌جمله توافق
می‌کنند و فرض در $\Zenv^{\Adv}$ دومی را به $\varepsilon$ کران می‌زند، که
همان نمایش است. هیچ‌یک کمانک‌گذاریِ محض نیست --- اجراهای سمتِ راست ماشینی
حمل می‌کنند که اجراهای سمتِ چپ ندارند، یعنی بازپخشِ نشسته در جایگاهِ حریف
--- پس هر دو نیمهٔ واگذارندهٔ لمِ~\ref{lem:absorb}اند و نه نیمهٔ
گزارهٔ~\ref{prop:regroup}.

\emph{اثباتِ \textnormal{(R1)}.} لمِ~\ref{lem:absorb} در $H = \{\Adv\}$ با
$H_{\op{keep}} = \emptyset$، به‌گونه‌ای که $\Zenv^{H}$ همان
$\Zenv^{\Adv}$ِ تعریفِ~\ref{def:zadv} باشد و بازپخشِ نشسته در
$\IDs(\Adv)$ِ واگذارشده همان $\Dum$. سه فرضِ آن، سه فرضِ قضیه‌اند: $\Adv$
بی‌اعتنا به رد است چون $\Aset$ چنین است؛ هیچ هستهٔ $\pi$ فراخوانِ پاسخگو
نمی‌گذارد، که همان فرض بر سیستم‌هاست؛ و $\Zenv \in \ZenvSet'$ هیچ هویتِ
$Z$ را خموش نمی‌کند. خانواده‌های اشغال‌گر مجازند، چون $\Zenv^{\Adv}$ تنها
$\IDs(\Zenv)$ را اشغال می‌کند و $\Dum$ همان هویت‌هایی را که $\Adv$ رها
کرد.
\emph{اثباتِ \textnormal{(R2)}.} اینجا $\Adv$ از دسته‌ای به دستهٔ دیگر
می‌رود و نه از اجرا به بیرون، پس لمِ~\ref{lem:absorb} در ظاهر اِعمال
نمی‌شود و استدلال همان نیمهٔ واگذارندهٔ آن است که دوباره با $\varphi$
به‌جای $\pi$ و $\Sim_{\Dum}$ به‌جای $\Dum$ دوانده شده. و اینکه این منتقل
می‌شود، بسته به آن است که اثباتِ آن لم از ماشینِ نشسته در جایگاه چه به کار
برد، و آن چیزی بیش از این نیست که نگهبانش حکم‌های $\opl{Guard}_{\Adv}$ را
برمی‌گرداند، که بر حسبِ اینکه ادعا $\id'.F = Z$ دارد یا نه توزیع می‌کند،
و --- آنجا که $\Adv$ِ میهمان دستوری ردشده را می‌خواند، $\rej$ از
$\Sim_{\Dum}$ در درونِ دسته در سمتِ چپ در برابرِ $\none$ِ محسوب از آن‌سوی
جایگاه در سمتِ راست --- که $\Aset$ بی‌اعتنا به رد است، درست همان‌گونه که
در (R1). آنچه تفاوت دارد این است که کدام ماشین $\Adv$ را در بر دارد ---
$\Sim_{\Dum}[\Adv]$ در چپ، $\Zenv^{\Adv}$ در راست --- و
تعریف‌های~\ref{def:sadv} و~\ref{def:zadv} آن دو را وادار می‌کنند همان
تحویل‌ها را انجام دهند، که هر یک از راهِ واسطِ نگهبانی‌شدهٔ ماشینِ گیرنده
مسیریابی می‌شود، و آن دستور در یک سو در درونِ یک ماشین سفر می‌کند و در سوی
دیگر از آن‌سوی اجرا. زنجیرهٔ $\Zenv \to \Adv \to \Sim_{\Dum} \to \varphi$
در هر دو همان است، تنها در نقطهٔ دیگری بریده شده: $\Sim_{\Dum}$ با همان
دستورها به همان ترتیب و زیرِ همان ادعاها روبه‌رو می‌شود، بندِ «جانشینی»
مو‌به‌مو برقرار است، و اشغال‌گر در هر دو حال با همان میدان‌ها داوری
می‌شود --- $\PID$ و $\Ps$ با اشغال، و $\admits$ با تعریفِ~\ref{def:sim}،
که $\Sim_{\Dum}$ برآورده‌اش می‌کند و $\Sim_{\Dum}[\Adv]$ به ارث می‌بردش،
چون بنا بر تعریفِ~\ref{def:sadv} میدان‌های $\Sim_{\Dum}$ را حمل می‌کند؛ و
هیچ نگهبانی باقی را نمی‌خواند. تعریفِ~\ref{def:sadv} تحویلِ سمتِ چپ را از
راهِ همان واسطِ نگهبانی‌شده‌ای مسیریابی می‌کند که $\opl{Exec}$ در سمتِ
راست خدمتش می‌کند، و فراخوان‌های خودِ $\Adv$ِ میهمان به $\Zenv$ زیرِ
$(A,P)$ از دسته بیرون می‌روند، درست همان‌گونه که در سمتِ راست از
$\opl{Exec}$ می‌گذرند.

\emph{نتیجه.} بنا بر (R1) و (R2) دو مزیتِ آن نمایش برابرند، و فرض در
$(\Dum,\Zenv^{\Adv})$ دومی را به $\varepsilon$ کران می‌زند. شبیه‌سازِ
$\Sim_{\Dum}[\Adv]$ در $\SimSet'$ است و تنها به $\Adv$ وابسته است، چون فرض
هیچ از $\Zenv$ نمی‌داند، پس ترتیبِ سورهای تعریفِ~\ref{def:uc} برآورده
است؛ و $\Zenv \in \ZenvSet'$ دلخواه بود.
\end{proof}

از دید نموداری، اثبات همان مربعِ جابه‌جاییِ قضیهٔ~\ref{thm:comp} است (همان
شکلِ بالای اثباتش)، که این بار از نیمهٔ \emph{واگذارندهٔ}
لمِ~\ref{lem:absorb} ساخته شده و نه از نیمهٔ نگه‌دارنده: آنچه به درونِ
$\Zenv$ می‌رود بخشی از سیستم نیست بلکه خودِ حریف است، و بازپخشی همان جا که
او ایستاده بود می‌ایستد.

\begin{center}
%% sizes itself, so it opts out of the fitting hook
\fitdiagramfalse
\resizebox{\linewidth}{!}{%
\begin{tikzpicture}[
  sysnode/.style={draw=accent, rounded corners=1pt, inner sep=2pt, font=\scriptsize, minimum height=4.5mm},
  wire/.style={draw=black!70},
  absorbed/.style={draw=black!55, dashed, rounded corners=2pt, inner sep=2mm},
  role/.style={font=\small},
  rel/.style={font=\normalsize},
  taglabel/.style={font=\tiny, color=accentsoft},
]

%% ===== Row 1 (top): real world, R1 =====
\node[role] (Z0) {$\Zenv$};
\node[sysnode, below=6mm of Z0] (pi0) {$\pi$};
\draw[wire] (Z0) -- (pi0);
\node[sysnode, below=6mm of pi0] (Adv0) {$\Adv$};
\draw[wire] (pi0) -- (Adv0);

\node[rel, right=13mm of pi0] (eq1) {$=$};

\node[role, right=13mm of eq1] (Z0a) {$\Zenv$};
\node[sysnode, below=2mm of Z0a] (Adv0a) {$\Adv$};
\node[absorbed, fit=(Z0a)(Adv0a)] (env0) {};
\node[sysnode, below=6mm of env0] (Dum0) {$\Dum$};
\draw[wire] (env0.south) -- (Dum0);
\node[sysnode, below=6mm of Dum0] (pi0a) {$\pi$};
\draw[wire] (Dum0) -- (pi0a);

%% ===== Row 2 (bottom): ideal world, R2 =====
\node[role, below=44mm of Z0] (Z1) {$\Zenv$};
\node[sysnode, below=6mm of Z1] (phi0) {$\varphi$};
\draw[wire] (Z1) -- (phi0);
\node[sysnode, below=6mm of phi0] (SDA0) {$\Sim_{\Dum}[\Adv]$};
\draw[wire] (phi0) -- (SDA0);

\node[rel] at (eq1 |- phi0) (eq2) {$=$};

\node[role] at (Z0a |- Z1) (Z1a) {$\Zenv$};
\node[sysnode, below=2mm of Z1a] (Adv1a) {$\Adv$};
\node[absorbed, fit=(Z1a)(Adv1a)] (env1) {};
\node[sysnode, below=6mm of env1] (SD1) {$\Sim_{\Dum}$};
\draw[wire] (env1.south) -- (SD1);
\node[sysnode, below=6mm of SD1] (phi0a) {$\varphi$};
\draw[wire] (SD1) -- (phi0a);

%% ===== vertical relations =====
\draw[<->]
  (pi0.west) .. controls +(-16mm,-8mm) and +(-16mm,8mm) .. (phi0.west)
  node[midway, left=2mm, rel, align=center] {$\approx_{\varepsilon}$\\[1pt]{\tiny\color{accentsoft}(نتیجه)}};

\draw[<->]
  (env0.east) .. controls +(18mm,-12mm) and +(18mm,12mm) .. (env1.east)
  node[midway, right=2mm, rel, align=center] {$\approx_{\varepsilon}$\\[1pt]{\tiny\color{accentsoft}(فرض، روی $\Dum$)}};

\end{tikzpicture}%
}
\fitdiagramtrue
\end{center}

پایین‌چپ همان جایی است که قضیه از آن آغاز می‌کند: $\pi$ و $\varphi$ در
سرتاسر دست‌نخورده، $\Adv$ اشغال‌گرِ جهانِ واقعی و $\Sim_{\Dum}[\Adv]$ ---
یعنی $\Sim_{\Dum}$ با $\Adv$ که پیش از آن نشانده شده --- اشغال‌گرِ
مشتق‌شدهٔ جهانِ ایده‌آل. یال‌های بالا و پایین همان (R1) و (R2)ِ اثباتند:
جذبِ $\Adv$ در $\Zenv$، از هر سویی که خوانده شود، باز هم \emph{تساویِ}
توزیع‌هاست، درست همان‌گونه که در گزارهٔ~\ref{prop:regroup} بود --- اما این
بار نیمهٔ واگذارندهٔ لمِ~\ref{lem:absorb} است که بهایش را می‌دهد، با
فرض‌های (a)--(c) به‌جای هیچ، چون آنچه حرکت می‌کند خودِ حریف است و بازپخشی
($\Dum$، و سپس $\Sim_{\Dum}$) باید به‌جایش بایستد. \emph{همان} یک
$\Zenv^{\Adv}$ِ جذب‌شده در هر دو سو خدمت می‌کند، و همین است که چرا یالِ
راست اصلاً می‌تواند آن دو را بسنجد: آن یال همان فرضِ قضیه است، که تنها در
برابرِ $\Dum$ بیان شده. با تعقیبِ مربع، درست همان‌گونه که در
قضیهٔ~\ref{thm:comp}، یالِ چپ --- نتیجهٔ قضیه، که حالا برای هر $\Adv$ِ
بی‌اعتنا به رد است و نه برای یک ماشینِ ثابت --- وادار می‌شود: آن دو
$\approx_\varepsilon$، بنا بر آن دو تساوی، همان یک کمیتند.

آن سه رده همان‌گونه حرکت می‌کنند که در قضیهٔ~\ref{thm:comp}، و به همان
دلایل. حریف‌ها آن‌هایی هستند که نتیجه رویشان سور می‌زند، و فرض تنها $\Dum$
را داشت؛ شبیه‌سازها حساب می‌شوند، یکی به‌ازای هر حریف، با نشاندنِ آن حریف
پیش از شبیه‌سازی که فرض فراهم می‌کند؛ و محیط‌ها یک پیش‌نگاره‌اند، یعنی
$\Zenv$ِ بیرونی درست وقتی به شمار می‌آید که آنچه در جذبِ هر $\Adv$ به آن
بدل می‌شود یکی از آن‌هایی باشد که فرض می‌پوشاندش. گرفتنِ $\ZenvSet$
به‌عنوانِ بزرگ‌ترین ردهٔ مجاز، $\ZenvSet'$ را باز بزرگ‌ترین می‌کند ---
منهای تنها آن محیط‌هایی که هویت‌های $Z$ِ خودشان را خموش می‌کنند، شرطی که
هیچ مجازبودنی تحمیلش نمی‌کند و هیچ کاربردی نمی‌خواهدش، و کنار گذاشته
می‌شود چون دستورِ تعریفِ~\ref{def:zadv} زیرِ ادعایی $Z$ بیرون می‌رود و چنین
محیطی آن را تنها در یک سو رد می‌کرد --- چون جذب $\pars$ را دست‌نخورده
می‌گذارد.

اگر به معنای بخشِ~\ref{sec:concrete} عینی خوانده شود، آن دو ساخت همان
هزینه‌ای را دارند که حریف دارد: اگر $\Adv$ کراندارِ $\bd{a}$ باشد، آن‌گاه
$\Zenv^{\Adv}$ به $\bd{z} \oplus \bd{a}$ نیاز دارد آنجا که $\Zenv$ به
$\bd{z}$ نیاز داشت، و $\Sim_{\Dum}[\Adv]$ به $\bd{s} \oplus \bd{a}$ آنجا
که $\Sim_{\Dum}$ به $\bd{s}$. شمارِ فراخوانی‌ها جمع می‌شود و زمان‌های هر
فراخوانی ماکسیمم می‌گیرند، و هر دو مستقل از $\pi$ و $\varphi$اند.

%% ==================================================================
\section{امنیتِ عینی}\label{sec:concrete}

تعریفِ~\ref{def:uc} در برابرِ رده‌های دلخواه بیان شده و در بابِ آنچه در
آن‌هاست هیچ نمی‌گوید. پُرکردنشان با ماشین‌هایی با منابعِ کراندار، هر حکمِ
بخشِ پیشین را به حکمی عینی بدل می‌کند، و نسبت‌های \emph{میانِ} رده‌ها را به
حسابِ روی بودجه‌ها.

\begin{definition}[بودجه‌ها و رده‌های کراندار]\label{def:res}
یک \emph{بودجه} زوجی $\bd{a} = (t,r)$ از اعدادِ طبیعی است. ماشینی
\emph{کراندارِ $\bd{a}$} است اگر در هر فراخوانی دستِ‌بالا $t$ گام بردارد و
در کل دستِ‌بالا به $r$ فراخوان پاسخ دهد؛ پس در هر اجرایی دستِ‌بالا $t\,r$
گام می‌دود. ماشینی که در هر اجرایی هر دو کران را برآورد کراندارِ $\bd{a}$
شمرده می‌شود، خواه قاعدهٔ ایستادن در خودش نوشته شده باشد خواه نه --- و
دسته‌هایی که در زیر از ماشین‌های کراندار ساخته می‌شوند به همین معنا
کراندارند. بنویسیم $\Aset(\bd{a})$، $\SimSet(\bd{a})$ و
$\ZenvSet_{\pi,\varphi}(\bd{a})$ برای ردهٔ حریف‌های کراندارِ $\bd{a}$،
شبیه‌سازهای کراندارِ $\bd{a}$، و محیط‌های کراندارِ $\bd{a}$ که در
$\ZenvSet_{\pi,\varphi}$ می‌افتند. بودجه‌ها مؤلفه‌به‌مؤلفه مرتب می‌شوند؛ و
ماشین‌هایی که یک اجرا را شریکند --- دسته‌ای از چند تا، یا شبیه‌سازی که
حریفی را می‌دواند --- با این عمل ترکیب می‌شوند
\[
  (t,r) \oplus (t',r') \;:=\; \bigl(\max(t,t'),\ r + r'\bigr) ,
\]
چون در هر فعال‌سازی یک ماشین می‌دود، پس زمانِ هر فراخوانی کرانِ بزرگ‌تر را
می‌گیرد و شمارِ فراخوانی‌ها جمع می‌شود.
\end{definition}

بودجه بخشی از ماشین است، و نه وعده‌ای در بابِ او. کران‌گذاشتن بر آن سه
ماشینِ جایگاه هنوز برای کران‌گذاشتن بر یک اجرا بس نیست --- اعضای سیستم هم
ماشین‌اند، و هسته‌ای که روی خودش بازگشتی شود، یا دو عضو که تا ابد یکدیگر
را فرا بخوانند، ژتون را از هر ماشینِ بودجه‌داری دور می‌بردند و چیزی برای
پس‌گرفتنش نمی‌بود. پس حکم‌های عینی روی \emph{سیستم‌های بودجه‌دار} خوانده
می‌شوند: هر عضوِ هر سیستمی که مقایسه می‌شود بودجه‌ای از خودش حمل می‌کند، به
همان معنای این تعریف که مو‌به‌مو اِعمال شده. آن‌گاه اجرا تعدادِ متناهی
ماشین می‌دواند، هر یک تعدادِ متناهی گام برمی‌دارد، و در هر لحظه دقیقاً یک
ماشین گام می‌زند، پس اجرا در کلِ کرانِ گامِ $\sum_i t_i\,r_i$ِ ماشین‌هایش
می‌ایستد و $\opl{Exec}$ رو تعریف‌شده است. ایستادن بر سرِ بودجه مقداری
داده شده است: ماشینی که در میانِ فراخوان می‌ایستد هر فراخوانندهٔ معلق بر
خودش را با $\none$ از سر می‌گیرد، و محیطی که دور از بازگشتِ ریشه‌اش
می‌ایستد چنان خوانده می‌شود که $0$ برگردانده است. رده‌ها در بودجه
یکنوایند، یعنی $\Aset(\bd{a}) \subseteq \Aset(\bd{a}')$ هرگاه $\bd{a} \leq
\bd{a}'$، و همین برای $\SimSet$ و $\ZenvSet_{\pi,\varphi}$. و همین
لمِ~\ref{lem:mono} را حکمی در بابِ اعداد می‌کند: کوچک‌کردنِ بودجهٔ حریف،
بزرگ‌کردنِ بودجهٔ شبیه‌ساز، یا کوچک‌کردنِ بودجهٔ محیط، هر سه تقلید را نگه
می‌دارند.

\begin{definition}[مزیتِ عینیِ UC]\label{def:concadv}
بگذارید $\pi$ و $\varphi$ سیستم‌های بودجه‌دار باشند، $\bd{a}$ و $\bd{s}$
بودجه‌هایی برای حریف و شبیه‌ساز، و $\ZenvSet \subseteq
\ZenvSet_{\pi,\varphi}$ ردهٔ محیط‌ها. بگذارید
\[
  \advtg^{\op{uc}}_{\pi,\varphi}[\bd{a},\bd{s},\ZenvSet]
  \;:=\;
  \max_{\Adv \in \Aset(\bd{a})} \ \
  \min_{\Sim \in \SimSet(\bd{s})} \ \
  \max_{\Zenv \in \ZenvSet} \ \
  \advtg^{\op{uc}}_{\pi,\varphi}(\Zenv,\Adv,\Sim) ,
\]
و $\advtg^{\op{uc}}_{\pi,\varphi}[\bd{a},\bd{s},\bd{z}]$ را کوتاه‌نویسیِ
حالتِ $\ZenvSet = \ZenvSet_{\pi,\varphi}(\bd{z})$ بگیرید، یعنی بزرگ‌ترین
ردهٔ آن بودجه.
\end{definition}

آرگومانِ سوم یک رده است و نه یک بودجه، چون بزرگ‌ترین رده به آن زوجِ
سیستم‌ها وابسته است، و حکم‌های زیر مزیت‌هایی را می‌سنجند که روی زوج‌های
\emph{متفاوت} گرفته شده‌اند. هر جا ردهٔ مشترکی خواسته باشد، نمادگذاری باید
بتواند بگویدش.

آن سه عملگر، همان سه سورِ تعریفِ~\ref{def:uc}اند که به‌عنوانِ یک مقدار
خوانده شده‌اند، به همان ترتیب و به همان دلیل: شبیه‌ساز پس از حریف و پیش از
محیط انتخاب می‌شود. اکسترمم‌ها حاصل می‌شوند، هرچند دلیلش لحظه‌ای وقت
می‌خواهد، و آن دلیل متناهی‌بودن است و نه فشردگی. ماشینِ کراندارِ $\bd{a}$
در کل دستِ‌بالا $t\,r$ گام برمی‌دارد، پس دستِ‌بالا $t\,r$ نماد از هر
ورودی و دستِ‌بالا $t\,r$ سکه می‌خواند؛ و تا حدِ رفتار، او نگاشتی است از
آنچه خوانده به آنچه کرده، و چنین نگاشت‌ها به‌شمارِ متناهی‌اند. سکه‌هایش
آن‌ها را با احتمال‌های دوگانیِ مخرجِ $2^{t r}$ وزن می‌دهند، پس ردهٔ کراندار
در کل تعدادِ متناهی توزیعِ رفتار القا می‌کند. و با بودجه‌دار‌بودنِ اعضای
سیستم‌ها، $\Pr[\opl{Exec} = 1]$ در هر انتخابی از رفتارها عددی خوش‌تعریف
است، پس هر مزیتِ تعریفِ~\ref{def:adv} یکی از تعدادِ متناهی مقدار را
می‌گیرد، و ماکسیمم یا مینیممِ تعدادِ متناهی مقدار حاصل می‌شود. (روی ردهٔ
دلخواهِ $\ZenvSet$، آن $\max$ِ بیرونی به‌عنوانِ سوپریمم خوانده می‌شود.) پس
\[
  \pi \ \text{کارکردِ} \ \varphi \ \text{ را در } \varepsilon \ \text{ و در برابرِ}
  \ \bigl(\Aset(\bd{a}),\SimSet(\bd{s}),\ZenvSet_{\pi,\varphi}(\bd{z})\bigr) \text{ UC-تقلید می‌کند}
\]
درست وقتی برقرار است که آن کمیت دستِ‌بالا $\varepsilon$ باشد --- و آن
مینیممِ حاصل‌شده همان است که جهتِ عکس را در مرز دقیق می‌کند. آن‌گاه یک
ادعای امنیتِ عینی، یک نامساویِ تنهاست که سه بودجه و یک احتمال را به هم
مربوط می‌کند، بی هیچ پارامترِ امنیتی و هیچ مجانبی در چشم.




%% ==================================================================

%% ==================================================================
\chapter{ویژگی‌ها}\label{sec:properties}

تقلید دو سیستم را می‌سنجد. یک \emph{ویژگی} از یکی چیزی می‌خواهد: که هیچ
حریفی بازیِ بازی‌شده در برابرِ او را نبرد --- که
گزارهٔ~\ref{prop:propslotabs} آن را به نبردنِ بدل فرو می‌کاهد، با جایگاهی
که در محیط جذب می‌شود. برای هیچ‌یک ماشین‌آلاتِ تازهٔ چندانی نمی‌خواهیم. بازی
کارکردی است مثلِ هر کارکردِ دیگر، محیط همان حریفی است که بازی‌اش می‌کند، و
لوله‌کشیِ بخش‌های پیشین از پیش حل کرده که چه کسی چه چیزی را می‌تواند فرا
بخواند --- پس یک ویژگی عبارت است از یک سیستم، یک ردهٔ محیط، و یک پیشامد، و
میدان‌های تعریفِ~\ref{def:func} بیشترِ کاری را می‌کنند که قواعدِ یک بازی
روی کاغذ می‌کنند. آنچه باید افزوده شود کم است، و هر جا که بیاید نامش برده
می‌شود: یک شرطِ بهداشتی بر سیستمِ بازی، و آن سه شرطی که
بخشِ~\ref{app:dummy} از پیش از هر چه جابه‌جا می‌کند می‌خواهد.

\begin{definition}[بازی، سیستمِ بازی]\label{def:prop}
یک \emph{بازی} برای $\F$ کارکردی استاندارد به نامِ $\Gm$ است با
\[
  \begin{gathered}
  \Gm.\PID := (\Gm.F,\,0,\,0), \qquad
  \Gm.\Ps := \F.\Ps \cup \{Z\}, \\
  \Gm.\admits := \Zpid, \qquad
  \Gm.\uses := \{(\F,\serves_{\Gm})\} ,
  \end{gathered}
\]
که در کنارِ واسط‌های خودش یک یا چند \emph{نهایی‌سازی} حمل می‌کند: واسط‌هایی
که نامشان با $\op{Fin}$ آغاز می‌شود، پیشوندی که برای آن‌ها رزرو شده و هیچ
واسطِ دیگری آن را حمل نمی‌کند، که هر یک در ریشه خدمت می‌شود، هر یک یک بیت
برمی‌گرداند، و پس از آنکه یکی از آن‌ها خدمت شد هیچ واسطِ $\Gm$ ---
نهایی‌سازی یا غیرِ آن --- خدمت نمی‌شود. اینجا
\[
  \serves_{\Gm}(\Gm,\F) \;:\iff\; \F.\Ps \supseteq \Gm.\Ps \setminus \{Z\}
  \ \wedge\ \Gm.\PID \in \F.\admits
\]
همان $\serves$ است با ریشه که معاف شده. یک \emph{سیستمِ بازی} عبارت است از
$\gamma := \{\Gm\} \cup \sigma$ برای سیستمِ $\sigma$ که $\F$ و هر چه $\F$
به کار می‌برد را در بر دارد و $\Gm$ را در بر ندارد، با $\WF(\gamma)$، و با
$\gamma$ که به معنای نکتهٔ~\ref{rem:tight} در $\{\Gm\}$ \emph{تنگ} است.
\end{definition}

پس یک بازی می‌تواند چند ویژگی حمل کند: مثلاً $\op{FinCor}$ و $\op{FinUF}$
روی کارکردِ امضا، که اوراکل‌هایی را شریکند که ثبت می‌کنند چه چیزی امضا شده
و تنها در حکمی که در پایان از آن‌ها خوانده می‌شود تفاوت دارند. محیط با
انتخابِ اینکه کدام نهایی‌سازی را فرا بخواند انتخاب می‌کند که بر کدام داوری
شود، و آن انتخاب آخرین کاری است که می‌کند. پس بازی با چند نهایی‌سازی
عطفی از ویژگی‌ها نیست، بلکه پوسته‌ای است که چند ویژگی را در بر دارد و هر
یک مزیتِ خودش را در زیر دارد؛ محیطی که یکی را ماکسیمم کند لازم نیست دیگری
را ماکسیمم کند، و هیچ‌چیزِ اینجا از او این را نمی‌خواهد.

\begin{interface}[nofloat]{بازیِ $\Gm$ برای $\F$: پوسته \\
\params{$\PID := (\Gm.F,0,0)$, \quad $\Ps := \F.\Ps \cup \{Z\}$, \quad $\admits := \Zpid$, \quad
$\uses := \{(\F,\serves_{\Gm})\}$, \quad $\pars := \none$}}

\opsig{$\op{Initialize}()$}:
\begin{algorithmic}[1]
\setcounter{contline}{0}\algcont
\State $\done \gets \false$
\State \dots\ متغیرهای خودِ بازی
\algsave
\end{algorithmic}
\vspace{1.4ex}

\opsig{$\id.\op{Op}(in)$} \quad از $\id'$
\begin{algorithmic}[1]
\algcont
\Req{$\neg\done$} \label{ln:gmopen}
  \Comment{با نهایی‌شدن بسته می‌شود}
\State \dots\ کدِ خودِ بازی، که می‌تواند $\F$ را در طرفِ $\id.P$ فرا بخواند
\State \textbf{return} $out$
\algsave
\end{algorithmic}
\vspace{1.4ex}

\opsig{$\id.\op{Leak}()$} \quad از $\id'$
\begin{algorithmic}[1]
\algcont
\State \textbf{return} $\none$
  \Comment{بازی رازِ هیچ طرفی را نگه نمی‌دارد}
\algsave
\end{algorithmic}
\end{interface}

\begin{interface}[nofloat]{نهایی‌سازی‌های $\Gm$، یکی به‌ازای هر ویژگی؛ پوستهٔ $\F$ مشترک است \\
\params{هر یک به نامِ $\op{Fin}\dots$، هر یک خدمت‌شده در ریشه، هر یک بازی را می‌بندد}}

\opsig{$\id.\opdef{Fin}()$} \quad از $\id'$ \quad (شکلِ هر نهایی‌سازی)
\begin{algorithmic}[1]
\algcont
\Req{$\neg\done$} \label{ln:gmonce}
  \Comment{تنها یک نهایی‌سازی}
\Req{$\id.P = Z$} \label{ln:gmroot}
  \Comment{در ریشه، که درستکار می‌ماند}
\State $\done \gets \true$ \label{ln:gmdone}
\State $w \gets$ حکمِ این ویژگی، خوانده‌شده از حالتِ پوسته \label{ln:gmwin}
\State \textbf{return} $w$
\algsave
\end{algorithmic}
\vspace{1.4ex}

\opsig{$\id.\op{FinCor}()$, \ $\id.\op{FinUF}()$, \ \dots} \quad از $\id'$
\begin{algorithmic}[1]
\algcont
\State \dots\ همان سه سطر، و سپس حکمِ خودِ هر ویژگی
\algsave
\end{algorithmic}
\end{interface}

هر میدان وادار شده است، و ارزشِ گفتن دارد که با چه چیزی.

\heading{شناسهٔ فرایند} در هر آزمایش یک بازی هست، پس نشست و نمونه هیچ خبری
حمل نمی‌کنند و ما هر دو را در $0$ تعیین می‌کنیم. نشستِ $0$ قرارداد نیست
بلکه یک خواسته است: محیط با ادعای هویتِ $Z$ِ خودش به $\Gm$ می‌رسد، که
سطرِ~\ref{ln:caller} می‌پذیردش چون $\Zpid \subseteq \Gm.\admits$، و آن‌گاه
سطرِ~\ref{ln:session} می‌پرسد $\Gm.s = \Zenv.s = 0$ --- چون $\Gm$ جهانی
نیست. نکتهٔ~\ref{rem:sesszero} همان مشاهده است به‌طورِ کلی. بنا بر
تعریفِ~\ref{def:wfsys} یک $\F$ِ محلی نشستِ فراخواننده را شریک است، پس کلِ
آزمایش در نشستِ $0$ می‌نشیند، که همان جایی است که ماشین‌هایی که اجرا فراهم
می‌کند هم در آن زندگی می‌کنند.

\heading{طرف‌های خدمت‌گیرنده} هسته تنها هویتِ خودش را ادعا می‌کند
(سطرِ~\ref{ln:call}) و $\opl{Guard}$ می‌پرسد $\id'.P = \id.P$، پس $\Gm$ در
طرفِ $P$ به $\F$ در طرفِ $P$ می‌رسد و به هیچ طرفِ دیگر: بازی $\F$ را هر بار
در یک طرف تمرین می‌دهد، و محیط با انتخابِ اینکه کدام هویتِ $\Gm$ را خطاب
کند طرف را انتخاب می‌کند، و ریشه‌اش بنا بر سطرِ~\ref{ln:zparty} آزاد است
برای هر طرفی عمل کند. و $\F.\Ps$ از همین‌جاست. ریشه برای نهایی‌سازی‌ها
افزوده شده، و بندِ بعد می‌گوید چرا.

\heading{فراخواننده‌های پذیرفته} $\Zpid$ و هیچ‌چیزِ دیگر. نه $\Stdpid$، که
$\Gm$ را جهانی می‌کرد و می‌گذاشت هر کارکردِ استانداردی بازی را نهایی کند؛
و نه $\Apid$، که آن را به حریف می‌سپرد. با $\Zpid$ِ تنها،
سطرِ~\ref{ln:caller} هر فراخواننده‌ای جز محیط را رد می‌کند.

\heading{نام} دلخواه است، و نتیجهٔ~\ref{cor:callername} همین را می‌گوید:
تغییرنامِ بازی و $\F.\admits$ با هم هیچ احتمالی را تغییر نمی‌دهد، پس
ویژگی‌ای که در برابرِ بازی‌ای با یک نام اثبات شود، در برابرِ همان بازی زیرِ
هر نامِ دیگری هم برقرار است. و این مهم است چون سطرِ~\ref{ln:call} می‌گذارد
بازی تنها هویتِ خودش را ادعا کند، پس $\F.\admits$ باید نامِ بازی را ببرد،
در حالی که آن $\F$ی که در درونِ یک پروتکل عرضه می‌شود نامِ آن پروتکل را
می‌برد؛ و آن نتیجه همان است که این دو را یک ویژگی می‌کند. آنچه دلخواه
نیست، \emph{کدِ} فراخواننده است، که هیچ تغییرنامی به آن دست نمی‌زند.

\heading{وابستگی} $\serves_{\Gm}$ ریشه را معاف می‌کند چون ریشه هویتِ
حساب‌داریِ خودِ بازی است: $\F$ لازم نیست به طرفی خدمت دهد که بازی هرگز در
آن فرایش نمی‌خواند. بازی‌ای که حکمش ناچار است از $\F$ \emph{پرس‌وجو} کند
--- مثلاً برای تأییدِ جعلی که ادعا شده --- دو گزینه دارد، و همان گزینه‌های
معمولند: یا در طولِ فراخوان‌های اوراکل هر چه حکم لازم دارد را ثبت کند، یا
از $\F$ بخواهد به ریشه هم خدمت دهد، که در آن صورت $\serves_{\Gm}$ به
$\serves$ بدل می‌شود.

\heading{تنگی} میدان‌های بالا به‌تنهایی بس نیستند، و آن شکاف شکافِ
نکتهٔ~\ref{rem:tight} است. محیطِ مجاز روی $\ncl{\gamma}$ خموش است، پس هیچ
هویتی از نامی که $\gamma$ به کار می‌برد را نمی‌تواند \emph{ادعا} کند ---
اما تعریفِ~\ref{def:admis} تنها او را از \emph{میهمان‌کردن} در هویت‌هایی که
$\gamma$ اشغال می‌کند منع می‌کند، و کارکردِ میهمان واسطی کامل است که هویتِ
خودش را از راهِ خموش‌کنندهٔ خودش ادعا می‌کند
(تعریفِ~\ref{def:bundle}). پس محیط می‌تواند ماشینی را در شناسهٔ فرایندِ
اشغال‌نشده‌ای از نامی به‌کاررفته میهمان کند: مثلاً در $(\Gm.F,0,1)$، یعنی
نمونهٔ دومِ ساختگی از نامِ خودِ بازی. عضوی از $\sigma$ که فراخواننده‌اش را با
\emph{نام} می‌پذیرد --- همان الگوی محلیِ فصلِ~\ref{sec:functionalities} ---
آن ماشین را می‌پذیرد، و فراخوان‌هایش از هر سطرِ $\opl{Guard}$ می‌گذرند: نام
در $\admits$ است، طرف‌ها جورند، نشست‌ها جورند، و $\opl{Mediate}$ در طرفِ
درستکار شلیک نمی‌کند. آن‌گاه محیط اوراکلی روی $\F$ در دست داشت که بازی هیچ
از آن ثبت نمی‌کند، و $\Pr[\Win_{\op{Fin}}] = 1$ برای هر بازی‌ای که حکمش
این است که حریف چیزی تولید کرده که به او داده نشده بود.

تنگی در $\{\Gm\}$ این را رد می‌کند، و کاملاً ردش می‌کند. هر عضوی جز بازی،
کلِ شناسه‌های فرایندی را می‌پذیرد که سیستم اشغالشان کرده، پس نمونهٔ ساختگی
در سطرِ~\ref{ln:caller} شکست می‌خورد؛ و شرطِ طرفِ نکتهٔ~\ref{rem:tight}
هر هویتی را که فراخوانندهٔ سنجاق‌شده می‌توانست از آن سخن بگوید اشغال‌شده
می‌کند، همان جا که تعریفِ~\ref{def:admis} میهمانی را ممنوع می‌کند. آن یک
راهی که تنگی باز می‌گذارد --- به عضوی از خودِ $I$، که اینجا همان بازی است
--- با $\Gm.\admits = \Zpid$ بسته می‌شود: کارکردِ میهمان استاندارد است و
شناسهٔ فرایندِ استانداردی را ادعا می‌کند، که $\Zpid$ آن را کنار می‌گذارد، و
هویت‌های $Z$ِ خودِ محیط میهمان نیستند بلکه از خودِ اویند. پس بازی همان یک
عضوی است که می‌تواند ردهٔ خالی‌ای از شناسه‌های فرایند را بپذیرد، و درست به
این سبب در این کار ایمن است که $Z$ هیچ نامِ استانداردی نیست.

\begin{remark}[Why a finalization is served at the root]\label{rem:finroot}
Party ids in $\bits$ can be corrupted, and line~\ref{ln:medin} of the full interface hands a corrupt
party's call to the adversary \emph{before} the core runs. A finalization served at such a party
would therefore be answered by the adversary whenever that party is corrupt, and the answer to
``have I won?'' would be whatever the adversary says --- Remark~\ref{rem:statusmed}'s problem exactly,
and no check inside the core can repair it, mediation happening in the wrapper above.
The root is the fix, and the framework already supplies it: $\fopl{Corrupt}$ asks
$\id.P \in \bits$, so $Z$ never enters $\Cs$, so $\opl{Mediate}$ at $(\Gm.\PID,Z)$ can never fire and
the bit line~\ref{ln:gmwin} computes is always the game's own. The reserved party ids were introduced
to keep the three supplied machines honest; they keep the game honest at no extra cost.
Line~\ref{ln:gmroot} then confines every finalization to that identity, so no other is a place to
win.
\end{remark}

\begin{remark}[One finalization, once, and last]\label{rem:gmonce}
That exactly one finalization is answered, and that it is the last call the game answers, is
enforced rather than assumed --- and the one flag does both jobs. Line~\ref{ln:gmdone} sets a flag
the wrapper cannot bypass, line~\ref{ln:gmonce} refuses every later finalization, its own name or
another's, and line~\ref{ln:gmopen} refuses every other interface thereafter; a refused
\textbf{require} returns $\rej$ by Section~\ref{sec:execution}, so the environment learns that the
game is closed and nothing else --- at an honest party, at least. At a corrupt one line~\ref{ln:medin}
intercepts before the core is entered, so the caller is answered by the adversary rather than refused;
that changes nothing, the slot being unable to write the game's state and the winning event being
fixed already.
Nothing stops the environment calling a global functionality afterwards either, which is
right --- the game is over, and what it does then cannot change what it won. The order of the lines
matters for the reason $\Fsig$ re-tests its tables: line~\ref{ln:gmwin} may place calls, and
Section~\ref{sec:execution} lets a suspended instance be re-entered, so the flag is set before the
verdict is computed rather than after. Were it the other way about, two nested finalizations
could each find the game open.

That one flag governs them all is what makes several finalizations safe to hang off one shell.
Were each to keep a flag of its own, an environment could finalize twice and be judged on two
properties in one run --- reading, say, a correctness verdict off a state its unforgeability
verdict had already been computed from, with the oracles between them. Sharing $\done$ is what
confines a run to one property, and it is why the finalizations sit in a box of their own with the
shell's flag rather than each carrying its own bookkeeping.
\end{remark}

\begin{remark}[What corruption buys, and what it costs]\label{rem:gmcorr}
At a party $P \in \bits$ the game's own interfaces are wrapped like any other functionality's, so
once $P \in \Cs$ line~\ref{ln:medin} intercepts a call to $\Gm$ at $P$: the core does not run, the
game records nothing, and the caller is answered by the adversary. That gains an environment nothing
--- $\opl{Mediate}$ replaces an answer rather than running a core, so the game's state is never
written by the adversary --- and it costs it the oracle at that party. Since corruption is indexed by
party alone (Section~\ref{sec:corruption}), corrupting a party of $\F$ closes the game's oracle at
that party too. A game that wants an oracle at a corrupt party must therefore serve it somewhere the
corruption does not reach: at the root, or at a party of its own outside $\F.\Ps$, which it may do
whenever that oracle needs no call on $\F$.

What keeps the adversary out of the game's cores is not either of the two lines one would reach for.
Line~\ref{ln:corrgate} \emph{admits} $\Adv$ at a corrupt party, refusing it only at honest ones; and
$\opl{Mediate}$ tests $\id'.F \neq A$, so it does not intercept a call the adversary itself places.
A call from the slot at a corrupt party would therefore run $\Gm$'s core, exactly as
Section~\ref{sec:fullinterface} says of a corrupt party addressed by $\Adv$. What refuses it is
$\Gm.\admits = \Zpid$ at line~\ref{ln:caller}, and nothing else does. The narrowness of that field is
not tidiness but the only thing between the adversary and the game's own code.

And $\Gm$'s $\op{Leak}$ returns $\none$ because leakage
does not go through $\opl{Guard}$ (Section~\ref{sec:leakage}): the adversary can read it at a corrupt
party whatever $\Gm.\admits$ says, so a game must keep nothing there it is not willing to hand over.
\end{remark}

\begin{remark}[Several instances of $\F$]\label{rem:gminstances}
A game may address as many instances of $\F$ as it likes, and the two ways it can differ.
Instances in the game's own session need nothing extra, line~\ref{ln:session} comparing sessions and
ignoring the instance number, so $(\F.F,0,i)$ is reachable for every $i$. Instances in other sessions
need $\op{global}(\F)$, which is what waives that line. Either way the claim is $\Gm$'s own identity,
line~\ref{ln:call} silencing its cores on everything else, so every instance sees one and the same
caller and access control is decided once. And the instances really are copies: by
Chapter~\ref{sec:reloc} they run the same code and their fields differ in the process id alone, while
Section~\ref{sec:execution} gives each its own state and its own coins. So a multi-instance property
--- $n$ keys, $n$ sessions --- is written by having the game count them, with no extra convention
about what a second instance means.

What counting them does not do is make transfer cheaper. Proposition~\ref{prop:proptransfer} replaces
one subsystem, so a game over $n$ instances every one of which is to be replaced costs $n$
applications and $n\varepsilon$, by Theorem~\ref{thm:parcomp} along a chain of hybrids --- exactly as
Remark~\ref{rem:manyinstances} records on the relocation side. One application transfers the property
for the instances $\varphi$ contains and says nothing of the others.
\end{remark}

\begin{remark}[A game is not session-uniform, and need not be]\label{rem:gmsessunif}
$\Gm.\admits = \Zpid$ meets $\Zpid$ while $\Gm$ is not global, so Definition~\ref{def:sessunif} fails
of $\Gm$ and the results of Chapter~\ref{sec:reloc} do not apply to a game system. That is as it
should be. A game is pinned to session $0$ by the paragraph above, there is one of it per experiment,
and relocating it is not something anyone wants; what the chapter applies to is the $\F$ under test,
whose instances Remark~\ref{rem:gminstances} counts.
\end{remark}

Which environments may play is settled by the class Chapter~\ref{sec:uc} already defines, taken at a
single system.

\begin{definition}[Property environments]\label{def:propenv}
The environments \emph{admissible for} a game system $\gamma$ are $\ZenvSet_{\gamma,\gamma}$ of
Definition~\ref{def:admis}: those with $\Zenv.\pars \supseteq \ncl{\gamma}$ that host at no identity
of $\IDs(\gamma)$. Write $\ZenvSet^{\circ}_{\gamma}$ for those among them that silence no $Z$-identity,
the ones the absorption of Proposition~\ref{prop:propslotabs} can move.
\end{definition}

The two clauses say what a game needs them to say. Silenced on the name closure, an environment may
claim no identity of a name $\gamma$ uses, so it cannot speak as $\F$, as one of $\F$'s subroutines,
or as a further instance of any of them; what it may claim is its own $Z$-identities, which reach
$\Gm$, and fresh names, which reach only what admits all of $\Stdpid$. So the environment plays the
game through $\Gm$'s interfaces and touches $\F$ directly only when $\F$ is global --- and then it
touches it as anyone may, which is what being global means and what
Definition~\ref{def:admis} already grants in the emulation setting. A property of a global $\F$ is
therefore a property against an adversary that shares it, and the game must be written knowing that.
The hosting clause is the occupancy hygiene of Definition~\ref{def:occ}, and it also stops the
environment from installing a second game.

\begin{definition}[Win, property advantage]\label{def:propadv}
Let $\gamma = \{\Gm\} \cup \sigma$ be a game system, let $\op{Fin}$ be one of $\Gm$'s
finalizations, let $\Zenv$ be admissible for $\gamma$ (Definition~\ref{def:propenv}), and let
$\mathcal{X}$ occupy the adversary slot. Write $\Win_{\op{Fin}}$ for the event that the interface
$(\Gm.\PID,Z).\fopl{Fin}$ returns $1$ in $\opl{Exec}(\gamma,\Zenv,\mathcal{X},\Corr)$. The
\emph{advantage} of $\Zenv$ in the property $\op{Fin}$ of $\sigma$ is one of
\[
  \begin{gathered}
  \padv{\op{Fin}}{\Gm,\sigma,\Zenv,\mathcal{X}} \;:=\; 2\Pr[\,\Win_{\op{Fin}}\,] - 1 , \\[3pt]
  \padv{\op{Fin}}{\Gm,\sigma,\Zenv,\mathcal{X}} \;:=\; \Pr[\,\Win_{\op{Fin}}\,] ,
  \end{gathered}
\]
the \emph{decision} reading and the \emph{search} reading, the probability being over the coins of
every machine in the execution. Which of the two is meant is fixed for each finalization where the
game is given, and not by this definition; the paragraph on calibration below says how to choose.
We drop $\mathcal{X}$ from the subscript when it is $\Dum$, and write $\F$ for $\sigma$ when
$\sigma = \{\F\}$. For $\varepsilon \geq 0$, a
class $\ZenvSet \subseteq \ZenvSet_{\gamma,\gamma}$ of environments and a class $\Xset$ of slot
occupants, we say $\sigma$ has the property $\op{Fin}$ \emph{within $\varepsilon$ against}
$(\ZenvSet,\Xset)$ if
$\padv{\op{Fin}}{\Gm,\sigma,\Zenv,\mathcal{X}} \leq \varepsilon$ for every $\Zenv \in \ZenvSet$ and
every $\mathcal{X} \in \Xset$, and \emph{within $\varepsilon$ against $\ZenvSet$} when
$\Xset = \{\Dum\}$.
\end{definition}

Since one flag closes the game (Remark~\ref{rem:gmonce}), at most one $\Win_{\op{Fin}}$ occurs in
any run: the events of distinct finalizations are disjoint, and a game with several of them is
several properties measured in several executions rather than several numbers read off one.

Under the decision reading the advantage is signed, where that of Definition~\ref{def:adv} is not,
and the asymmetry is the point:
an environment cannot flip $\Win_{\op{Fin}}$ as it can flip a guess, the event being the game's
verdict rather than its own, so a game it cannot win scores $-1$ and there is nothing to take an
absolute value of.

$\Win_{\op{Fin}}$ is an event of the execution rather than its output, and that is deliberate: an
execution's
output is whatever the environment returns (Chapter~\ref{sec:uc}), and a property is not a guess to be
reported but a thing that happened. Nothing is lost by the distinction, because
Remark~\ref{rem:finroot} makes the two agree where it matters. The bit line~\ref{ln:gmwin} computes is
the bit the finalization hands the environment, unmediated, so we may name the environment that reports
it: for a fixed $\op{Fin}$, write $\Zenv^{\star}$ for $\Zenv$ with its return at the root replaced by
``return $1$ if any call I
placed on $(\Gm.\PID,Z).\fopl{Fin}$ was answered with $1$, and $0$ otherwise''. Every part of that
earns its place. The \emph{identity} is named because a finalization addressed at another party fails
line~\ref{ln:gmroot} where that party is honest and is answered by the \emph{adversary} where it is
corrupt, so the value coming back need not be the game's; at the root neither happens, by
Remark~\ref{rem:finroot}. The \emph{operation} is named because $Z \in \Gm.\Ps$, so
line~\ref{ln:interface} serves every interface of $\Gm$ at the root, the game's own oracles included
--- and because with several finalizations the reporting form must say which property it is reporting
on, one $\Zenv^{\star}$ per $\op{Fin}$.
The test is \emph{against $1$} rather than a reading of the answer as a bit, because a later
finalization comes back $\rej$ by Remark~\ref{rem:gmonce}, and $\rej$ is no bit
(Section~\ref{sec:execution}). And it is \emph{any} call rather than the first, because the first one
placed may be refused --- an environment claiming $(Z,P)$ for $P \neq Z$ fails line~\ref{ln:caller}'s
party test, leaving $\done$ unset --- while a later one is served; nothing is weakened by quantifying,
since at most one call is ever answered with $1$, the game serving one finalization and refusing the
rest. So $\opl{Exec}$ returns $1$ exactly when $\Win_{\op{Fin}}$ occurs, and $\Zenv^{\star}$ can compute
it, every
call on that interface being one it placed itself: $\Gm.\admits = \Zpid$ admits no other caller, and a
functionality the environment hosts claims a standard id, which $\Zpid$ excludes. Then
\[
  \Pr\bigl[\opl{Exec}(\gamma,\Zenv^{\star},\mathcal{X},\Corr) = 1\bigr]
  \;=\;
  \Pr\bigl[\Win_{\op{Fin}} \ \text{in} \ \opl{Exec}(\gamma,\Zenv,\mathcal{X},\Corr)\bigr] ,
\]
the two machines placing the same calls and differing only in the value they halt with. That also
puts $\Zenv^{\star}$ in every class $\Zenv$ is in whose membership turns on $\pars$ and hosting
alone --- $\ZenvSet_{\gamma,\gamma}$ among them, and the budgeted classes of
Definition~\ref{def:res} up to the one further step of reading a bit back.
Proposition~\ref{prop:proptransfer} is where all of this is used.

\heading{Why the dummy} The default slot occupant is the dummy of Section~\ref{app:dummy}, and the
reason is that a game has one adversary and it is the environment. The dummy relays in both
directions and keeps no strategy, so an environment driving it holds everything the slot can do; and
where $\F$'s own cores reach the adversary --- the $\Adv(\cdot)$ of Section~\ref{sec:adv}, which is how
an ideal functionality leaves a value unpinned --- it is then the environment that answers, which is
exactly the freedom a game means to give it. Dropping $\mathcal{X}$ from the subscript for the dummy's
case is thus no loss of generality but a choice of who plays; the longer form is kept because
Proposition~\ref{prop:proptransfer} needs a simulator in that slot. That the abbreviation costs nothing
is not a matter of taste, and the reason is the construction of Section~\ref{app:dummy}: the slot
absorbs into the environment.

\begin{proposition}[The slot absorbs into the environment]\label{prop:propslotabs}
Let $\gamma = \{\Gm\} \cup \sigma$ be a game system whose cores place no responsive calls
(Definition~\ref{def:respcall}), let $\mathcal{X}$ be a refusal-oblivious occupant of the adversary
slot (Definition~\ref{def:refobl}), and let $\Zenv \in \ZenvSet^{\circ}_{\gamma}$. Then
$\Zenv^{\mathcal{X}}$ of Definition~\ref{def:zadv} again lies in $\ZenvSet^{\circ}_{\gamma}$, and for
every finalization $\op{Fin}$ of $\Gm$
\[
  \padv{\op{Fin}}{\Gm,\sigma,\Zenv,\mathcal{X}}
  \;=\;
  \padv{\op{Fin}}{\Gm,\sigma,\Zenv^{\mathcal{X}}} .
\]
So if $\sigma$ has the property $\op{Fin}$ within $\varepsilon'$ against $\ZenvSet^{\circ}_{\gamma}$,
it has
it within $\varepsilon'$ against $\bigl(\ZenvSet^{\circ}_{\gamma},\Xset\bigr)$ for every class $\Xset$
of refusal-oblivious occupants. Read concretely, an $\bd{z}$-bounded environment and an
$\bd{x}$-bounded occupant give an absorbed environment that is $\bd{z} \oplus \bd{x}$-bounded.
\end{proposition}

\begin{proof}
Lemma~\ref{lem:absorb} at $H = \{\mathcal{X}\}$ with $H_{\op{keep}} = \emptyset$ gives
\[
  \opl{Exec}(\gamma,\Zenv,\mathcal{X},\Corr)
  \;\equiv\;
  \opl{Exec}(\gamma,\Zenv^{\mathcal{X}},\Dum,\Corr) ,
\]
its three hypotheses being the three assumed here, one for one.
Now $\Win_{\op{Fin}}$ is the event that the interface $\fopl{Fin}$ at $(\Gm.\PID,Z)$ returns $1$, and
absorption touches the
adversary slot alone --- $\Gm$ occupies that identity on both sides --- so the event has one
probability for each $\op{Fin}$ alike, and Definition~\ref{def:propadv} gives one number under
either reading. Admissibility is
Section~\ref{app:dummy}'s observation that $\Zenv^{\mathcal{X}}$ occupies what an environment occupies
and carries $\Zenv$'s $\pars$, so it is admissible exactly when $\Zenv$ is; the budget is the
$\bd{z} \oplus \bd{a}$ of Chapter~\ref{app:conccosts}. The final sentence is the display read at each
$\mathcal{X}$ in turn.
\end{proof}

The hypotheses are Theorem~\ref{thm:dummy}'s and no more: a game system placing responsive calls is
outside this, by Remark~\ref{rem:respdummy}, and so is an environment that silences its own
$Z$-identities, the instruction of Definition~\ref{def:zadv} leaving under a $Z$-claim. Neither is a
condition on what the occupant \emph{does}, which is the point --- a game need not care whether the
machine in the slot plays fairly, only that it can be moved.

One scope condition comes with all of this, and it is the same one throughout. The proposition above
rests on the regrouping of Theorem~\ref{thm:dummy}, which holds only for systems whose cores place no
responsive calls, Remark~\ref{rem:respdummy} showing where the relay fails otherwise; the same
condition governs Theorem~\ref{thm:dummy} itself, which is what supplies the emulation hypothesis
Proposition~\ref{prop:proptransfer} wants from a wider adversary class, and
Corollary~\ref{cor:proptransfer2} inherits it on both sides at once. A game system whose $\F$ reaches
the slot through $\Adv^{!}$ therefore sits outside all three and must have its dummy-only hypothesis
established directly. What a responsive theory of games would look like we leave where
Remark~\ref{rem:respdummy} leaves it.

\heading{The calibration} Which of Definition~\ref{def:propadv}'s two readings a finalization takes
is not a matter of taste, and the rule is this. The factor $2\Pr[\Win_{\op{Fin}}] - 1$ reads the
property as a \emph{decision}: an
environment that guesses wins with probability $\tfrac12$ and scores $0$, and one that always wins
scores $1$. That is right where guessing is a strategy --- where the verdict is a bit the
environment could have flipped a coin for. It is wrong for a \emph{search} property ---
unforgeability, say, where winning is meant to be hard outright --- because there the same quantity
reads $\Pr[\Win_{\op{Fin}}] = (1 + \padv{\op{Fin}}{\Gm,\sigma,\Zenv})/2$, so a bound
$\varepsilon$ on the advantage is a bound $(1+\varepsilon)/2$ on the winning probability, which for
small $\varepsilon$ is the wrong scale: an $\varepsilon$ of $0$ would permit a forger succeeding
half the time. Search properties therefore take $\Pr[\Win_{\op{Fin}}]$ itself. Declaring the
reading per finalization rather than fixing one globally is what lets a single shell carry both
kinds at once. Every property in this paper is of the second kind --- winning each of the games
below is meant to be impossible outright, not merely no better than a guess --- so all of them take
the search reading; the decision reading is there for properties where guessing is the strategy, an
indistinguishability written as a game among them.

The choice is not free, and Proposition~\ref{prop:proptransfer} is where it is paid for. Its two
displays are the same bound at the two scales --- $2\varepsilon$ for the decision reading,
$\varepsilon$ for the search reading --- so each finalization inherits transfer at the scale its own
reading fixes, and a game mixing the two inherits at both.

Finally, the reason to define properties this way rather than beside the framework: emulation carries
them.

\begin{proposition}[Properties transfer along emulation]\label{prop:proptransfer}
Let $\gamma = \{\Gm\} \cup \sigma$ be a game system, let $\varphi \subseteq \sigma$, let the
replacement of $\varphi$ by $\pi$ in $\gamma$ be well-formed in the sense of Definition~\ref{def:wf},
and let $\gamma[\pi/\varphi]$ be a game system as well. Suppose
\[
  \pi \ \text{UC-emulates} \ \varphi \ \text{ within } \varepsilon \text{ against }
  (\{\Dum\},\SimSet,\ZenvSet) .
\]
Then there is $\Sim \in \SimSet$ such that for every finalization $\op{Fin}$ of $\Gm$ and every
$\Zenv \in \ZenvSet_{\gamma[\pi/\varphi],\,\gamma}$
whose reporting form has $\Zenv^{\star\,\gamma \setminus \varphi} \in \ZenvSet$ --- a condition
Proposition~\ref{prop:absorb} makes vacuous whenever $\ZenvSet$ is the largest admissible class, though
not when it is the largest class at a \emph{budget}, absorption costing
$\bd{c}_{\gamma \setminus \varphi}$ and the reporting form one step more ---
\[
  \padv{\op{Fin}}{\Gm,\sigma[\pi/\varphi],\Zenv}
  \;\leq\;
  \padv{\op{Fin}}{\Gm,\sigma,\Zenv,\Sim} \;+\; 2\varepsilon
  \qquad \text{(decision reading),}
\]
and, writing $\Win_{\op{Fin}}$ and $\Win'_{\op{Fin}}$ for the winning events of the two executions
the display above compares --- the real one and the ideal one ---
\[
  \Pr[\,\Win_{\op{Fin}}\,] \;\leq\; \Pr[\,\Win'_{\op{Fin}}\,] \;+\; \varepsilon
  \qquad \text{(search reading).}
\]
\end{proposition}

\begin{proof}
% BEGIN overview prop:proptransfer
Three moves and one bound. First the statement is checked to be about the systems it names, which
is where the hypothesis on $\gamma[\pi/\varphi]$ is spent: emulation supplies neither the game's
party set at the new occupant nor its process id in that occupant's $\admits$, and tightness is
part of Definition~\ref{def:prop} besides. Then the reporting form of Definition~\ref{def:propadv}
turns the winning event into an output bit, rewriting what $\Zenv$ returns at the root and changing
no run. Then absorption moves the game itself inside the environment, so that an environment
\emph{playing} the game becomes one \emph{distinguishing} the two systems. What is left is
Theorem~\ref{thm:comp} applied at that absorbed environment. The search reading is the bound it
gives; the decision reading is that bound doubled, and nothing in the argument reads $\op{Fin}$, so
it holds of each finalization separately.
% END overview prop:proptransfer
First, the two advantages are defined and are about the systems the statement names. Since
$\Gm \notin \sigma$ we have $\Gm \notin \varphi$, so
$\gamma[\pi/\varphi] = \{\Gm\} \cup \sigma[\pi/\varphi]$; and
$\ZenvSet_{\gamma[\pi/\varphi],\gamma}$ lies in both $\ZenvSet_{\gamma,\gamma}$ and
$\ZenvSet_{\gamma[\pi/\varphi],\gamma[\pi/\varphi]}$, each clause of Definition~\ref{def:admis}
weakening when a union of two closures shrinks to one. That the left-hand side is a property advantage
at all is where the hypothesis on $\gamma[\pi/\varphi]$ is spent, and it is not free: emulation supplies
neither $\Gm.\Ps = \F'.\Ps \cup \{Z\}$ at the new occupant $\F'$ nor $\Gm.\PID \in \F'.\admits$, the
second being hypothesis~(ii) of Theorem~\ref{thm:wfrepl} and no part of Theorem~\ref{thm:comp}. Nor is
that all it buys: tightness is part of Definition~\ref{def:prop}, so the hypothesis asks in addition
that $\pi$'s members pin their callers by process id where $\varphi$'s did. Emulation supplies that
least of all, comparing behaviour at the boundary and saying nothing of the declared fields behind it.

Let $\Sim \in \SimSet$ be the simulator the hypothesis supplies for $\Dum$; it depends on $\Dum$
alone, so one $\Sim$ serves every $\Zenv$. Fix such a $\Zenv$, fix $\op{Fin}$, and take the
reporting form
$\Zenv^{\star}$ for that $\op{Fin}$ from the paragraph after Definition~\ref{def:propadv}, which gives
\[
  \Pr\bigl[\opl{Exec}(\delta,\Zenv^{\star},\mathcal{X},\Corr) = 1\bigr]
  \;=\;
  \Pr\bigl[\Win_{\op{Fin}} \ \text{in} \ \opl{Exec}(\delta,\Zenv,\mathcal{X},\Corr)\bigr]
\]
for either game system $\delta$ and either occupant $\mathcal{X}$; the hypothesis places
$\Zenv^{\star}$ in the class $\ZenvSet'$ that Theorem~\ref{thm:comp} attaches to this replacement, and
$\Zenv^{\star} \in \ZenvSet_{\gamma[\pi/\varphi],\gamma}$ because $\Zenv$ is and the two agree on
$\pars$ and on what they host. Theorem~\ref{thm:comp} applied to
the replacement of $\varphi$ by $\pi$ in $\gamma$, at the adversary $\Dum$ and the environment
$\Zenv^{\star}$, bounds the difference of the two left-hand probabilities --- with
$\delta = \gamma[\pi/\varphi]$ and $\mathcal{X} = \Dum$ against $\delta = \gamma$ and
$\mathcal{X} = \Sim$ --- by $\varepsilon$. Hence $\Pr[\Win_{\op{Fin}}]$ differs by at most
$\varepsilon$ between
the two, which is the second display; the advantage under the decision reading, being
$2\Pr[\Win_{\op{Fin}}]-1$, then differs by at most $2\varepsilon$, which is the first. The order
matters for use
rather than for proof: the probability bound is what the argument gives, and the decision-reading
bound is it doubled. Under the search reading the advantage \emph{is} the probability, so the
second display is already the transfer statement and the factor of two never arises. Nothing in the
argument reads $\op{Fin}$, so it holds of each finalization separately.
\end{proof}

Two things about that proof. The reduction is the textbook one, and absorption is what makes it so: the
absorbed environment $\Zenv^{\star\,\gamma \setminus \varphi}$ of Definition~\ref{def:absorb} hosts
the game itself along with the rest of $\gamma \setminus \varphi$, so ``an environment playing the game
against $\pi$'' becomes ``an environment distinguishing $\pi$ from $\varphi$'' by moving the game
inside it, which is what Theorem~\ref{thm:comp} was proved about. And the property on the right is
stated against $\Sim$ in the slot rather than against the dummy, which is why
Definition~\ref{def:propadv} keeps the occupant as an argument --- but it is not an obligation anyone
need meet in that form, because the slot absorbs into the environment as well.

The reduction is worth drawing, because it is one move and the move is easy to read past. It is
the commuting square of Figure~\ref{fig:comp} again, in the same conventions and with the dotted
box again marking the pair being exchanged --- but this time it is the \emph{game} that moves:

\begin{center}
\begin{tikzpicture}[
  sysnode/.style={draw=accent, rounded corners=1pt, inner sep=2pt, font=\scriptsize,
                  minimum height=4.5mm},
  wire/.style={draw=black!70},
  shell/.style={draw=accent, rounded corners=2pt, inner sep=1.8mm},
  absorbed/.style={draw=black!55, dashed, rounded corners=2pt, inner sep=1.8mm},
  role/.style={font=\small},
  rel/.style={font=\normalsize},
  tag/.style={font=\tiny, color=accentsoft},
  hl/.style={draw=accentsoft, dotted, thick, rounded corners=1pt, inner sep=1.2pt},
]

%% ================= ROW 1: the real world =================
%% ---- cell A: the property as stated ----
\node[sysnode] (A-rest) at (0mm, -29mm) {$\sigma\setminus\varphi$};
\node[sysnode, right=1mm of A-rest] (A-pi) {$\pi$};
\node[shell, fit=(A-rest)(A-pi)] (A-sys) {};
\node[sysnode, above=4mm of A-sys] (A-G) {$\Gm$};
\node[role, above=3.5mm of A-G] (A-Z) {$\Zenv$};
\node[role, below=4mm of A-sys] (A-D) {$\Dum$};
\draw[wire] (A-Z) -- (A-G);
\draw[wire] (A-G) -- (A-sys.north);
\draw[wire] (A-sys.south) -- (A-D);
\node[hl, fit=(A-pi)] {};
\node[tag, anchor=north] at ($(A-D.south)+(0,-1mm)$) {$\Win_{\op{Fin}}$ at $\Gm$};

\node[rel] at (26mm, -26mm) {$=$};
\node[tag, anchor=north, align=center] at (26mm, -28.5mm)
  {reporting form,\\ then absorption};

%% ---- cell B: the game absorbed ----
\node[sysnode] (B-rest) at (52mm, -25mm) {$\sigma\setminus\varphi$};
\node[sysnode, above=3mm of B-rest] (B-G) {$\Gm$};
\node[role, above=3mm of B-G] (B-Z) {$\Zenv^{\star}$};
\draw[wire] (B-Z) -- (B-G);
\draw[wire] (B-G) -- (B-rest);
\node[absorbed, fit=(B-Z)(B-G)(B-rest)] (B-env) {};
\node[sysnode, below=4mm of B-env] (B-pi) {$\pi$};
\node[role, below=4mm of B-pi] (B-D) {$\Dum$};
\draw[wire] (B-env.south) -- (B-pi);
\draw[wire] (B-pi) -- (B-D);
\node[hl, fit=(B-pi)] {};

%% ================= ROW 2: the ideal world =================
\def\dy{-54mm}

%% ---- cell D: the property as stated ----
\node[sysnode] (D-rest) at (0mm, -29mm+\dy) {$\sigma\setminus\varphi$};
\node[sysnode, right=1mm of D-rest] (D-phi) {$\varphi$};
\node[shell, fit=(D-rest)(D-phi)] (D-sys) {};
\node[sysnode, above=4mm of D-sys] (D-G) {$\Gm$};
\node[role, above=3.5mm of D-G] (D-Z) {$\Zenv$};
\node[role, below=4mm of D-sys] (D-S) {$\Sim$};
\draw[wire] (D-Z) -- (D-G);
\draw[wire] (D-G) -- (D-sys.north);
\draw[wire] (D-sys.south) -- (D-S);
\node[hl, fit=(D-phi)] {};
\node[tag, anchor=north] at ($(D-S.south)+(0,-1mm)$) {$\Win'_{\op{Fin}}$ at $\Gm$};

\node[rel] at (26mm, -26mm+\dy) {$=$};
\node[tag, anchor=north, align=center] at (26mm, -28.5mm+\dy)
  {absorption, then\\ reporting form};

%% ---- cell C: the game absorbed ----
\node[sysnode] (C-rest) at (52mm, -25mm+\dy) {$\sigma\setminus\varphi$};
\node[sysnode, above=3mm of C-rest] (C-G) {$\Gm$};
\node[role, above=3mm of C-G] (C-Z) {$\Zenv^{\star}$};
\draw[wire] (C-Z) -- (C-G);
\draw[wire] (C-G) -- (C-rest);
\node[absorbed, fit=(C-Z)(C-G)(C-rest)] (C-env) {};
\node[sysnode, below=4mm of C-env] (C-phi) {$\varphi$};
\node[role, below=4mm of C-phi] (C-S) {$\Sim$};
\draw[wire] (C-env.south) -- (C-phi);
\draw[wire] (C-phi) -- (C-S);
\node[hl, fit=(C-phi)] {};

%% ================= the two vertical relations =================
\draw[<->]
  (A-sys.west) .. controls +(-15mm,-7mm) and +(-15mm,7mm) .. (D-sys.west)
  node[midway, left=2mm, rel, align=center]
  {$\approx_{\varepsilon}$\\[1pt]{\tiny\color{accentsoft}(conclusion)}};

\draw[<->]
  (B-env.east) .. controls +(17mm,-10mm) and +(17mm,10mm) .. (C-env.east)
  node[midway, right=2mm, rel, align=center]
  {$\approx_{\varepsilon}$\\[1pt]{\tiny\color{accentsoft}(hypothesis, at $\Dum$)}};

\end{tikzpicture}
\figcaption{fig:proptransfer}{Property transfer, as the same square with the game moving instead of the protocol. In the left column $\Gm$ sits inside the system, where a game belongs; in the right it sits inside the dashed box, part of the machine doing the distinguishing. Crossing that boundary is the whole reduction. (Proposition~\ref{prop:proptransfer}.)}
\end{center}

Read it by watching $\Gm$. In the left column it sits inside the system, where a game belongs: the
environment plays it, and $\Win_{\op{Fin}}$ is the verdict it returns. In the right column it sits
inside the dashed box, part of the machine doing the distinguishing. Crossing that boundary is the
whole reduction --- an environment \emph{playing the game} against $\pi$ becomes an environment
\emph{telling $\pi$ from $\varphi$} --- and once the game is across, what is left is a bare
emulation question, which is what Theorem~\ref{thm:comp} was proved about.

The two horizontal edges are free. The reporting form only rewrites what $\Zenv$ returns at the
root, leaving every call it places untouched, so it changes no run; and absorption is
Lemma~\ref{lem:absorb}, an identity of executions rather than a bound. So the right edge is the
only thing paid for, and it is the hypothesis applied at $\Dum$ --- which is why the bottom row
carries $\Sim$ in the slot and not the dummy, and so why Definition~\ref{def:propadv} keeps the
occupant as an argument at all. The left edge is then forced, exactly as in
Figure~\ref{fig:comp}: chasing the square makes the two $\approx_\varepsilon$'s the same quantity
rather than two bounds that happen to share an $\varepsilon$.

One reading of that quantity is the proposition's, the other is its double. Under the search
reading the property advantage \emph{is} $\Pr[\Win_{\op{Fin}}]$, so the left edge is already the
transfer statement and no factor of two appears; under the decision reading the advantage is
$2\Pr[\Win_{\op{Fin}}]-1$, which doubles it. Every per-functionality property proof in the paper
is an instance of this one square, which is why it earns a diagram where the instances do not.

\begin{corollary}[Transfer, against every adversary]\label{cor:proptransfer2}
In the setting of Proposition~\ref{prop:proptransfer}, suppose in addition that the cores of $\gamma$
and of $\gamma[\pi/\varphi]$ place no responsive calls (Definition~\ref{def:respcall}), that the
simulator $\Sim$ that proposition supplies and the occupant $\Adv$ below are refusal-oblivious
(Definition~\ref{def:refobl}), and that the environments named lie in $\ZenvSet^{\circ}_{\gamma}$. Fix
a finalization $\op{Fin}$ of $\Gm$. If
$\sigma$ has the property $\op{Fin}$ within $\varepsilon'$ against $\ZenvSet_{\gamma,\gamma}$ in the
sense of Definition~\ref{def:propadv}, then
\[
  \padv{\op{Fin}}{\Gm,\sigma[\pi/\varphi],\Zenv,\Adv} \;\leq\; \varepsilon' + 2\varepsilon
\]
for every occupant $\Adv$ of the adversary slot and every $\Zenv$ the proposition admits with
$\Zenv^{\Adv}$ among them. If in place of the hypothesis on advantages the ideal side satisfies
$\Pr[\,\Win'_{\op{Fin}}\,] \leq p$ under the same quantification, then
$\Pr[\,\Win_{\op{Fin}}\,] \leq p + \varepsilon$, with
$\Win_{\op{Fin}}$ and $\Win'_{\op{Fin}}$ naming the real and ideal events as in
Proposition~\ref{prop:proptransfer} --- which is the form a search property inherits in, the first
display being the decision reading's.
\end{corollary}

\begin{proof}
Three steps, of which the outer two are absorption. By
Proposition~\ref{prop:propslotabs} at the real game system,
$\padv{\op{Fin}}{\Gm,\sigma[\pi/\varphi],\Zenv,\Adv} =
\padv{\op{Fin}}{\Gm,\sigma[\pi/\varphi],\Zenv^{\Adv}}$,
the right-hand side being the dummy's case. Proposition~\ref{prop:proptransfer} at $\Zenv^{\Adv}$
bounds it by $\padv{\op{Fin}}{\Gm,\sigma,\Zenv^{\Adv},\Sim} + 2\varepsilon$. And
Proposition~\ref{prop:propslotabs} again, now at the ideal game system, rewrites the first term as
$\padv{\op{Fin}}{\Gm,\sigma,(\Zenv^{\Adv})^{\Sim}}$, which the hypothesis on $\sigma$ bounds by
$\varepsilon'$, the doubly absorbed machine being an admissible environment by two applications of the
same proposition. The probability form is the same three steps read on the second display of
Proposition~\ref{prop:proptransfer} rather than the first, absorption changing no probability at all.
\end{proof}

Two things this buys. The hypothesis on the ideal side is now the property against the \emph{dummy}
alone --- one statement about $\sigma$, quantified over environments and nothing else --- where
Proposition~\ref{prop:proptransfer} on its own asked for the property against the particular simulator
the emulation supplied. Nothing need be assumed about how that simulator corrupts or leaks --- only that
it does not branch on a framework-internal refusal, which Section~\ref{app:dummy} asks of every
machine it moves: whatever else it does, absorbing it into the environment turns it into part of a
machine the ideal-side hypothesis already covers. And the conclusion is the property of the real system
against \emph{every} adversary, not only against the dummy, by the same move applied on the other side.
Read concretely, the two absorptions are the whole price: for an $\bd{a}$-bounded $\Adv$ and an
$\bd{s}$-bounded $\Sim$ the innermost statement is reached at environment
budget $\bd{z} \oplus \bd{a} \oplus \bd{s}$, one hosting cost per machine moved, and
Chapter~\ref{app:conccosts} charges nothing else.

\begin{remark}[What the occupant of the slot does on the ideal side]\label{rem:propslot}
Nothing here is needed for the proof above, which absorbs whatever the slot does into the
$\varepsilon$ Theorem~\ref{thm:comp} supplies. What is at stake is instead the cost of the
\emph{hypothesis} on the right: the two executions the proposition compares differ in the slot as well
as in the system. Corruption divides in two, and only half of it is at issue. The half that is not is
what the environment does for itself: line~\ref{ln:corronlya} admits $\id'.F \in \{A,Z\}$, so an
environment needs nobody's cooperation for the corruptions it wants, and claiming $(Z,P)$ from its root
it makes them identically on both sides. The half that is at issue is the occupant's own initiative:
$\Sim$ may corrupt parties nobody asked it to, or decline what it is asked through the slot, where
$\Dum$ would relay. That is why
the right-hand side of Proposition~\ref{prop:proptransfer} is the property against $\Sim$ and not
against the dummy, and why Definition~\ref{def:propadv} keeps the occupant as an argument. What an
ideal $\sigma$ has to satisfy \emph{in that form} is the game against every simulator the emulation may
produce; Corollary~\ref{cor:proptransfer2} reduces it to the dummy again, and does so by moving the
simulator into the environment rather than by asking anything of it.

Nothing escapes the $\varepsilon$, and the register is why. The environment reads $\Cs$ through
$\fopl{Status}$, which mediates nothing and returns the whole corrupted set
(Remark~\ref{rem:statusmed}), so it may read $\Cs$ at any point and compare it against the parties it
corrupted itself together with whatever it instructed the slot to add.
Let $\Zenv$ be an environment the class admits, and let $\Zenv'$ be $\Zenv$ with that comparison added
and a return of $0$ where it fails. On the real side it never fails, $\Dum$ corrupting what it is told
and nothing besides; on the ideal side it fails with whatever probability $\delta$ the simulator strays
with, in either direction --- declining what it is asked, or corrupting unbidden. So
$\Zenv'$ distinguishes with advantage $\delta$, and the emulation hypothesis gives $\delta \leq
\varepsilon$. One may therefore reason as though $\Sim$ echoes the corruptions it is instructed to make,
at a price already paid --- which is Section~\ref{sec:corruption}'s commensurateness, read for
properties. It is worth having, but it is not what discharges the hypothesis on the right of
Proposition~\ref{prop:proptransfer}; absorbing the slot does, and it needs nothing at all of the
simulator's conduct.
\end{remark}

\begin{remark}[Leakage on the ideal side]\label{rem:propleak}
Leakage divides between the sides as corruption does, and needs one check of its own because it does
not go through
$\opl{Guard}$. Only the slot can read a $\op{Leak}$: line~\ref{ln:leakadv} asks $\id'.F = A$, and the
environment cannot claim such an identity. So $\Gm.\admits$ does not protect the game there, and
something else must: line~\ref{ln:leakcorr}'s gate together with the party, since the game's leakage is
readable only at a corrupt party, the root is never corrupt, and at every other party $\Gm$'s
$\op{Leak}$ returns $\none$ by Remark~\ref{rem:gmcorr}. The game's state therefore reaches the
adversary by that route on neither side. What does differ between the sides is $\varphi$'s leakage
against $\pi$'s, which $\Sim$ must produce; that is an ordinary part of emulation, bounded by the same
$\varepsilon$, and a property inherits the bound with nothing further to check.
\end{remark}



%% ==================================================================
\chapter{مسدودگرها}\label{sec:blockers}

Definition~\ref{def:admis} of Chapter~\ref{sec:uc} silences an environment on one and the same
name closure in both worlds, so the claims open to it do not depend on which world it faces ---
and Proposition~\ref{prop:absorb} and Theorem~\ref{thm:comp} use no relation between the two
identity sets at all.
What blocking buys is narrower: it makes that closure attainable as an exact
silenced set, and so delivers the least restricted admissible class, as the discussion of
Definition~\ref{def:admis} in Chapter~\ref{sec:uc} records. A \emph{blocker} occupies identities and nothing else. It admits no
caller, $\admits := \emptyset$, so the conjunct on line~\ref{ln:caller} of $\opl{Guard}$ is false
for every call and the caller receives $\rej$. Its $\op{Block}$ core is therefore never entered and
its body is a formality. Leakage does not go through $\opl{Guard}$, so a corrupt party's
$\op{Leak}$ core is reachable; it returns $\none$, all a machine with no state has to give. Both
are cores, wrapped as Definition~\ref{def:sys} requires of any standard functionality, and it is
the wrapper that makes the first unreachable.

\begin{interface}{مسدودگرِ $\Blk[\F]$ \\
\params{$\PID := \F.\PID$, \quad $\Ps := \F.\Ps$, \quad $\admits := \emptyset$, \quad
$\uses := \emptyset$, \quad $\pars := \none$}}

\opsig{$\id.\op{Block}()$} \quad از $\id'$
\begin{algorithmic}[1]
\setcounter{contline}{0}\algcont
\State \textbf{return} $\none$
  \Comment{به‌سببِ $\opl{Guard}$ دست‌نیافتنی}
\algsave
\end{algorithmic}

\opsig{$\id.\op{Leak}()$} \quad از $\id'$
\begin{algorithmic}[1]
\algcont
\State \textbf{return} $\none$
  \Comment{در دسترس، اما حالتی نیست}
\algsave
\end{algorithmic}

\end{interface}

The blocker $\Blk[\F]$ is indexed by the functionality it stands in for, from which it copies the
two fields that fix an identity set and nothing else. By construction
$\IDs(\Blk[\F]) = \IDs(\F)$: it occupies precisely the identities of $\F$, session and instance
included, and nothing besides. Being unreachable, it stands in for $\F$ as an occupant of those
identities and in no other respect.

Write $I_\pi := \{\F.\PID : \F \in \pi\}$ for the process ids a system uses. By
Definition~\ref{def:sys} each $q \in I_\pi$ is carried by exactly one member, which we write
$\pi_q$, so $\pi = \{\pi_q : q \in I_\pi\}$ and $\IDs(\pi) = \bigcup_{q \in I_\pi} \IDs(\pi_q)$.
Blockers are indexed by \emph{ids} and not by members, and the distinction matters: a process id
may occur in both worlds carrying different code, in which case $\varphi_q$ is not a member of
$\pi$ although $\pi$ uses $q$, and a blocker for such a $q$ would put two machines at the same id.

Indexing by $I_\varphi \setminus I_\pi$ instead keeps a blocker clear of what is already there. No
member of $\pi$ carries such a $q$, hence no identity of $\IDs(\pi)$ has $q$ in its first
component and $\IDs(\varphi_q) \cap \IDs(\pi) = \emptyset$. A blocker built from $\varphi_q$ thus
lands neither on an identity $\pi$ already serves nor on the id of an existing member. Distinct ids give pairwise disjoint blockers, and all of them are new to $\pi$, so the
enlarged system again has distinct ids throughout.

The pair of systems being compared must in addition agree wherever they overlap.

\begin{definition}[Compatibility]\label{def:compat}
Systems $\pi$ and $\varphi$ are \emph{compatible} if, whenever a process id occurs in both, the
two functionalities sitting at it agree on their served parties and on whom they admit; the
remaining fields --- what they use, their parameters, and the code they run --- may differ.
\end{definition}

Compatibility is used at exactly one point: the ids common to both worlds. No blocker is added
for such an id, so whatever it contributes must already agree on the two sides. Without it a
shared id could serve different parties in $\pi$ and in $\varphi$, and the difference would go
uncovered. Requiring $\admits$ to agree as well is more than the argument needs, since only $\Ps$ enters an
identity set once the process id is fixed; but it costs nothing, a shared id being carried by the
same functionality on both sides in every case we use.


\begin{definition}[Blocker system, blocked systems]\label{def:blksys}
For systems $\pi$ and $\varphi$, the \emph{blocker system} $\opdef{Blockers}(\varphi \setminus \pi)$ is
\[
  \opl{Blockers}(\varphi \setminus \pi) \;:=\; \bigl\{\, \Blk[\varphi_q] \;:\; q \in I_\varphi \setminus I_\pi \,\bigr\},
\]
and the \emph{blocked systems} of the pair are
\[
  \overline{\pi} \;:=\; \pi \cup \opl{Blockers}(\varphi \setminus \pi),
  \qquad
  \overline{\varphi} \;:=\; \varphi \cup \opl{Blockers}(\pi \setminus \varphi).
\]
\end{definition}

The argument of $\opl{Blockers}$ is notation rather than an input: the construction reads the ids
$I_\varphi \setminus I_\pi$ and $\varphi$'s occupants there, not the member set the difference
denotes --- at a shared id carrying different code, $\varphi_q \in \varphi \setminus \pi$ and yet
no blocker is built.
Each id in $I_\varphi \setminus I_\pi$ is thus covered by the blocker of the functionality living
there, with that functionality's party set. Each world gains exactly the process ids the other has
and it lacks, all blocking has to achieve. Ids common to both worlds get no blocker, which is
where Definition~\ref{def:compat} does its work.

Informally, Theorem~\ref{thm:blk} says blocking equalises the two worlds: each blocked system is
again a system, and for a compatible pair the two occupy one and the same identity set,
$\IDs(\pi) \cup \IDs(\varphi)$ --- each world having gained exactly the process ids the other
has and it lacks.

\begin{theorem}[Equalisation]\label{thm:blk}
Let $\pi$ and $\varphi$ be systems. Then the blocked systems $\overline{\pi}$ and
$\overline{\varphi}$ of Definition~\ref{def:blksys} are systems. If
$\pi$ and $\varphi$ are moreover compatible, then
\[
  \IDs(\overline{\pi}) \;=\; \IDs(\overline{\varphi}) \;=\; \IDs(\pi) \cup \IDs(\varphi).
\]
\end{theorem}

\begin{proof}
% BEGIN overview thm:blk
Two claims, and only the second needs compatibility. The first asks the three things
Definition~\ref{def:sys} asks of a system --- finiteness, standardness, distinct process ids --- of
each blocked system, and gets them from the indexing: blockers are indexed by the ids the other
world has and this one lacks, so they land on no existing member. The second compares the two
identity sets index by index over $I_\pi \cup I_\varphi$. At an id only one world uses, both sides
contribute the same set, a blocker copying the two fields that fix an identity set. At a shared id
no blocker is built, so the two contribute $\IDs(\pi_q)$ and $\IDs(\varphi_q)$ as they stand --- and
compatibility supplies the one thing the indexing does not, that the served party sets agree. That
is the only place it is used.
% END overview thm:blk
For the first claim, being a system asks three things of $\overline{\pi}$: that it be finite, that
its members be standard, and that they carry distinct process ids. Finiteness is immediate, the
added members being indexed by $I_\varphi \setminus I_\pi \subseteq I_\varphi$. Standardness is
inherited: a blocker carries the wrappers Definition~\ref{def:sys} asks of a standard
functionality, $\op{Block}$ as in Section~\ref{sec:fullinterface} and $\op{Leak}$ as in
Section~\ref{sec:leakage}, and its process id is copied from a member of $\varphi$, so its name is
none of $A$, $Z$, $C$. For distinctness, the members of
$\opl{Blockers}(\varphi \setminus \pi)$ are the $\Blk[\varphi_q]$ with
$q \in I_\varphi \setminus I_\pi$, and $\Blk[\varphi_q].\PID = \varphi_q.\PID = q$. These ids are
pairwise distinct and none lies in $I_\pi$, so the members of $\overline{\pi}$ carry distinct
process ids and $I_{\overline{\pi}} = I_\pi \cup I_\varphi$. The same argument applies to
$\overline{\varphi}$.

For the second, $\IDs(\Blk[\F]) = \IDs(\F)$ holds for every $\F$, since a blocker copies the
two fields that determine an identity set. Hence
\[
  \begin{gathered}
  \IDs(\overline{\pi}) \;=\; \bigcup_{q \in I_\pi} \IDs(\pi_q) \;\cup \bigcup_{q \in I_\varphi \setminus I_\pi} \IDs(\varphi_q), \\
  \IDs(\overline{\varphi}) \;=\; \bigcup_{q \in I_\varphi} \IDs(\varphi_q) \;\cup \bigcup_{q \in I_\pi \setminus I_\varphi} \IDs(\pi_q),
  \end{gathered}
\]
and both unions are indexed by $I_\pi \cup I_\varphi$. Compare them at each $q$ in that set. If
$q \in I_\pi \setminus I_\varphi$ then both contribute $\IDs(\pi_q)$, and if $q \in I_\varphi \setminus I_\pi$
then both contribute $\IDs(\varphi_q)$. If $q$ lies in both, the first contributes
$\IDs(\pi_q)$ and the second $\IDs(\varphi_q)$. Both are carried at the index $q$, so
$\pi_q.\PID = \varphi_q.\PID = q$ outright, and compatibility supplies the one thing that does not
follow from the indexing, $\pi_q.\Ps = \varphi_q.\Ps$. An identity set is fixed by those two
fields, so $\IDs(\pi_q) = \IDs(\varphi_q)$. The two unions therefore agree term by
term, and each equals
\[
  \bigcup_{q \in I_\pi} \IDs(\pi_q) \;\cup\; \bigcup_{q \in I_\varphi} \IDs(\varphi_q)
  \;=\; \IDs(\pi) \cup \IDs(\varphi) .
\]
\end{proof}

The second claim is a table. A process id falls in one of three places, and the theorem is what
sits at each:

\begin{center}
\begin{tikzpicture}[
  cell/.style={draw=accent!60, minimum width=31mm, minimum height=7mm,
               inner sep=1pt, font=\small, anchor=center},
  live/.style={cell, fill=accent!5},
  blkc/.style={cell, fill=accentsoft!25, draw=accentsoft, dashed},
  colhead/.style={font=\small},
  rowhead/.style={font=\small},
  tag/.style={font=\tiny, color=accentsoft},
]

\def\cw{31mm}
\def\rh{7mm}

\node[colhead] at (0*\cw, 7mm) {$I_\pi \setminus I_\varphi$};
\node[colhead] at (1*\cw, 7mm) {$I_\pi \cap I_\varphi$};
\node[colhead] at (2*\cw, 7mm) {$I_\varphi \setminus I_\pi$};

\draw[decorate, decoration={brace, amplitude=4pt, mirror}, accentsoft]
  (-0.5*\cw, 12mm) -- (2.5*\cw, 12mm);
\node[tag, anchor=south] at (1*\cw, 13mm)
  {one column per process id $q \in I_\pi \cup I_\varphi$};

\node[rowhead] at (-0.72*\cw, 0mm) {$\overline{\pi}$};
\node[live] at (0*\cw, 0mm) {$\pi_q$};
\node[live] at (1*\cw, 0mm) {$\pi_q$};
\node[blkc] at (2*\cw, 0mm) {$\Blk[\varphi_q]$};

\node[rowhead] at (-0.72*\cw, -\rh) {$\overline{\varphi}$};
\node[blkc] at (0*\cw, -\rh) {$\Blk[\pi_q]$};
\node[live] at (1*\cw, -\rh) {$\varphi_q$};
\node[live] at (2*\cw, -\rh) {$\varphi_q$};

\node[tag, anchor=north, align=center] at (0*\cw, -\rh-5mm)
  {$\IDs(\Blk[\pi_q]) = \IDs(\pi_q)$\\ a blocker copies $\PID$ and $\Ps$};
\node[tag, anchor=north, align=center] at (1*\cw, -\rh-5mm)
  {\textbf{compatibility}: $\pi_q.\Ps = \varphi_q.\Ps$\\ the one place it is used};
\node[tag, anchor=north, align=center] at (2*\cw, -\rh-5mm)
  {$\IDs(\Blk[\varphi_q]) = \IDs(\varphi_q)$\\ a blocker copies $\PID$ and $\Ps$};

\draw[decorate, decoration={brace, amplitude=4pt}, accentsoft]
  (2.5*\cw, -\rh-16mm) -- (-0.5*\cw, -\rh-16mm);
\node[font=\small, anchor=north] at (1*\cw, -\rh-19mm)
  {so both rows occupy $\IDs(\pi) \cup \IDs(\varphi)$};

\node[tag, anchor=west, align=left] at (2.62*\cw, -0.5*\rh)
  {shaded: a blocker ---\\ occupies the identities,\\ admits no caller};

\end{tikzpicture}
\figcaption{fig:blk}{Where a process id can fall, and what blocking does at each place. The three columns are the three cases of the second claim of Theorem~\ref{thm:blk}.}
\end{center}

In the outer columns the world that lacks the id gains a blocker, and a blocker occupies exactly
the identities of the functionality it stands in for, having copied the only two fields that fix
an identity set. So the two rows contribute the same identities there --- although one contributes
a working functionality and the other a machine that answers nothing, which is the point of
shading them differently. Equalising identity sets is all blocking has to achieve, and it had
better be all it achieves: a blocker that answered anything would change what an environment could
learn, and the comparison it was built to enable would be worthless.

The middle column is where the hypothesis lives. No blocker is built for a shared id, so the two
rows contribute $\IDs(\pi_q)$ and $\IDs(\varphi_q)$ as they stand, and nothing in the indexing
makes those equal --- both are carried at $q$, so the process ids agree outright, but the served
party sets need not. Compatibility is exactly the assumption that they do, and this column is the
only place it is used. Read the other way, it is also what the assumption costs: a shared id
serving different parties in the two worlds is a difference blocking cannot cover, because the one
device it has --- adding an occupant --- is unavailable where an occupant is already present.

Only the second claim needs any of this. That $\overline{\pi}$ and $\overline{\varphi}$ are
systems at all asks nothing of compatibility, which is why the theorem states the two separately
and why Proposition~\ref{prop:wfblk} below may lean on the first alone.

Blocking enlarges a system, and enlarging a system is one of the two ways
Definition~\ref{def:wfsys} could come undone.

\begin{proposition}[Blocking preserves well-formedness]\label{prop:wfblk}
Let $\pi$ and $\varphi$ be systems with $\WF(\pi)$ and $\WF(\varphi)$. Then $\WF(\overline{\pi})$
and $\WF(\overline{\varphi})$.
\end{proposition}

\begin{proof}
$\overline{\pi}$ is a system by the first claim of Theorem~\ref{thm:blk}, which asks no
compatibility, so $\WF(\overline{\pi})$ is a legal question. We argue for $\overline{\pi}$; the
argument for $\overline{\varphi}$ is the same with $\WF(\varphi)$ in place of $\WF(\pi)$. Take
$\F \in \overline{\pi}^{+}$ and $(\F',R) \in \F.\uses$.

If $\F$ is one of the added blockers there is nothing to check: a blocker has
$\uses = \emptyset$ by construction, so it contributes no entry. Otherwise $\F$ lies in $\pi^{+}$,
and $\WF(\pi)$ supplies a witness $\G \in \pi^{+}$ for that entry. That witness is still present,
since $\pi^{+} \subseteq \overline{\pi}^{+}$, and blocking alters no member's fields, so $\G.F$,
$\G.s$ and everything $R$ reads are what they were. The three conditions therefore still hold of
the same $\G$.

The argument uses nothing about blockers beyond their empty $\uses$: enlarging a system to another
system preserves well-formedness whenever the members added carry no entries of their own, because
a witness once present is never removed and the fields it is judged on never change. Neither
compatibility nor Theorem~\ref{thm:blk}'s second claim is used.
\end{proof}

\begin{remark}[Tight realizations]\label{rem:tight}
The closure keeps the environment's own claims off the used names; hosting is the one route
around it. An admissible environment may host a machine at an \emph{unoccupied} identity of a
used name --- an invented second instance of a caller, say --- and a hosted machine is a full
interface, claiming its own identity through its own silencer rather than through $\pars$
(Definition~\ref{def:bundle}). Against a pair whose members admit that name, the two worlds can then
answer differently at an identity only one of them occupies, and no condition on the class can
refuse the intruder without refusing absorption: the absorbed environment of
Definition~\ref{def:absorb} hosts machines of exactly the same shape --- outside occupants at
admitted names --- and differs from the intruder only in the code it runs, which admissibility
rightly does not read. Part of the repair belongs to the realization. Call a system \emph{tight at} a
subset $I$ of its members if every member outside $I$ admits whole process ids occupied by the
system, rather than names, and serves only parties each of its admitted callers serves --- with
$\serves$ this makes a private member's party set exactly its callers', which is what a private
subroutine is. Tightness closes the host route \emph{to the members outside $I$}: an unoccupied
process id lies outside every pinned $\admits$ and fails line~\ref{ln:caller}, and the party
condition makes every identity a pinned caller could speak from occupied, where
Definition~\ref{def:admis} forbids hosting; so a call an intruder addresses to a member outside
$I$ is refused, on the ideal side too, by the blocker's empty $\admits$ or by $\opl{Exec}$
finding no occupant. What tightness does \emph{not} close is the route to a member of $I$ itself:
an interface member that admits a bare used name --- the local pattern of
Chapter~\ref{sec:functionalities} --- is reachable by a hosted fake instance exactly as a private
member would be, and tightness constrains no member of $I$. Pinning those too, or admitting only
$\Zpid \cup \Apid$ where the environment reaches them directly, would close it; we do not pursue a
general sufficient condition here. Tightness is a hygiene condition in the spirit of a
subroutine-respecting protocol, necessary for the private members not to leak; a full
characterization of when comparisons against blocked pairs are faithful we leave open.
\end{remark}


%% ==================================================================
\chapter{جابه‌جایی و تغییرنام}\label{sec:reloc}

Chapter~\ref{sec:functionalities} records that a functionality's code hardly reads its own session
and never its instance number: $s$ occurs at line~\ref{ln:session} of $\opl{Guard}$, in the
wrappers that invoke it, and in Definition~\ref{def:wfsys}, and $i$ occurs nowhere. This chapter
turns that observation into a statement. Renaming sessions and instances is an isomorphism of
executions, so an emulation proved at one process id holds at every other, within the same
$\varepsilon$ and at the same budgets. One hygiene condition is needed and no more --- that a
functionality reachable from outside every session be reachable from every session, which every
machine the paper specifies satisfies --- and isolating it is half the work. That is what lets a realization be
written for a single session, as $\Fsig$ of Chapter~\ref{sec:fsig} is, and read as covering the rest.

\begin{definition}[Relocation]\label{def:reloc}
A \emph{relocation} is a permutation $\theta$ of the process ids such that
\begin{enumerate}[label=(\roman*),leftmargin=2.2em]
  \item $\theta(q).F = q.F$ for every $q$;
  \item $q.s = q'.s \iff \theta(q).s = \theta(q').s$ for all $q,q'$ with $q.F, q'.F \in \bits$;
  \item $\theta(q) = q$ for each of the three ids $q \in \{(A,0,0),(Z,0,0),(C,0,0)\}$ the execution
        supplies.
\end{enumerate}
\end{definition}

Conditions (i) and (ii) say exactly that on the standard names
\[
  \theta(F,s,i) \;=\; \bigl(F,\ \hat\theta(s),\ \theta_{F,s}(i)\bigr)
\]
for a permutation $\hat\theta$ of session ids and, for each name and session, a permutation
$\theta_{F,s}$ of instance numbers --- we write the two components of a relocation this way rather
than with letters of their own, $\sigma$ and $\iota$ being spoken for by
Proposition~\ref{prop:regroup} and Definition~\ref{def:regcorr}, whose notation the proofs below
import. Names are kept, because $\opl{Guard}$ and $\WF$ match by name and
nothing may become a different functionality; sessions are permuted \emph{as blocks}, because
line~\ref{ln:session} asks whether two ids share a session and not which; and instances are
permuted freely within a block, that component being read nowhere.

Condition (iii) pins the three process ids
that occur as \emph{literals} in code: $A$, in the hook $(A,\id.P)$ of $\opl{Mediate}$; $C$, in the
register read $(C,\id.P)$ of lines~\ref{ln:stat} and~\ref{ln:stattwo}; and $Z$, at
line~\ref{ln:execstart} and in the relay of line~\ref{ln:dumin}. Every identity named outright in the
paper is one of these three paired with a party in scope, so pinning the three pins them all --- and
it pins those three and no more. It must not be read as pinning every id of a reserved name, and the difference is
not cosmetic: by Chapter~\ref{sec:functionalities} the set $\Apid$ holds a triple $(A,s,i)$ for
every $s$ and $i$, only $(A,0,0)$ ever occurring, so a condition fixing all of them and read
together with (ii) unrestricted would give $\hat\theta(s) = s$ at every $s$ and leave no relocation
that moves a session at all. Restricting (ii) to the standard names is the other half of the same
care: $\hat\theta$ is determined there, and the ids of a reserved name that never occur are left to move
harmlessly inside their own class, which Proposition~\ref{prop:relocbasic} shows is where they stay.

Nothing here pins session $0$, and $\hat\theta$ is free to move it. That freedom has a price, and it
is worth seeing exactly where, since this is the only point at which the framework tells one session
from another. Line~\ref{ln:session} compares the \emph{caller's} session while $\op{global}$ is tested
of the callee, and the three machines the execution supplies sit in session $0$; so a call from
$\Adv$ or $\Zenv$ to a callee that is not global passes that line exactly when the callee's own
session is $0$. A callee of that shape would tell a relocation moving session $0$ from one leaving it
alone. Nothing else can: for every other callee the supplied machines are either waived by
$\op{global}$ or refused at line~\ref{ln:caller} before the session is looked at. So one condition on
the callee closes the gap, and it is a hygiene condition rather than a restriction on $\theta$.

\begin{definition}[Session-uniform]\label{def:sessunif}
A functionality $\F$ is \emph{session-uniform} if
\[
  \F.\admits \cap (\Apid \cup \Zpid) \neq \emptyset
  \qquad\Longrightarrow\qquad
  \op{global}(\F) ,
\]
so that it admits a caller from outside every session only if it is reachable from every session. A
system is session-uniform if each of its members is, and an environment if each functionality it
hosts is.
\end{definition}

The sessions a session-uniform functionality can be reached from do not depend on the session it sits
in, which is what relocation needs of it. The three supplied machines meet the condition outright,
all three being global --- $\Adv$ and $\Corr$ by the first disjunct of $\op{global}$, $\Zenv$ by the
second --- and so does every occupant of the adversary slot, global by that same second disjunct,
and so does a blocker, whose $\admits$ is empty and which line~\ref{ln:caller} therefore refuses
whatever the session. Those are every $\admits$ the paper declares; $\Fsig$ takes its as a parameter,
where the condition is one line to check. What the condition excludes is the shape
Chapter~\ref{sec:functionalities} leaves open in saying that a local $\F$ may itself declare whether
$\Apid$ and $\Zpid$ may call it: a local functionality admitting the adversary is reachable from the
slot in session $0$ and in no other, and that is the one thing a relocation cannot carry.
Remark~\ref{rem:sesszero} records the alternative, a one-line widening of line~\ref{ln:session} that
would dispense with the condition altogether.

\begin{definition}[Action of a relocation]\label{def:relocact}
A relocation $\theta$ acts on identities by $\theta(\PID,P) := (\theta(\PID),P)$, leaving the party
component alone, and on identity sets pointwise. It sends a functionality $\F$ to $\theta\F$ with
\[
  \begin{gathered}
  \theta\F.\PID := \theta(\F.\PID), \quad
  \theta\F.\Ps := \F.\Ps, \quad
  \theta\F.\admits := \theta(\F.\admits), \\[2pt]
  \theta\F.\uses := \theta(\F.\uses), \quad
  \theta\F.\pars := \theta(\F.\pars) ,
  \end{gathered}
\]
where $\theta$ acts on a $\uses$ entry by $\theta(\F',R) := (\theta\F',R)$ and on $\pars$ through
whatever identities and process ids it contains; the code of $\theta\F$ is the code of $\F$
with every process id occurring in it --- in a call target, a claimed identity, a test, or a stored
value --- replaced by its image; and every identity set a machine \emph{declares} rather than
computes is replaced by its image too, which for a bundle (Definition~\ref{def:bundle}) means its
occupied set, its outward set and any routing it fixes by hand. It acts on a system memberwise, and
on any other machine of the execution the same way, a bundle's hosted members included, so that
$\theta\Zenv$ hosts the images of what $\Zenv$ hosts. We require the requirement predicates to be
equivariant, $R(\F,\G) \iff R(\theta\F,\theta\G)$, as $\serves$ is, it reading only $\Ps$ and a
membership of $\admits$. Every clause above relabels pointwise, so the action is one:
$\op{id}$ acts as the identity and $(\theta'\theta)\F = \theta'(\theta\F)$. In particular
$\theta^{-1}$ undoes $\theta$ on machines as it does on process ids, which is what lets the results
below argue in both directions.
\end{definition}

Relabelling the process ids a program \emph{names} is not by itself relabelling what it does, and
the gap between the two is a condition on the machine model. We state it, because
Lemma~\ref{lem:reloc} is false without it.

\begin{definition}[Identity-atomic code]\label{def:atomic}
Code is \emph{identity-atomic} if it handles identities and process ids as opaque tokens. It may
project one to its name or party component; test two of them, or two of their name, party, session or
instance components, for equality; test one for membership in a set of them the code has been given,
as line~\ref{ln:caller} tests $\F.\admits$ and $\opl{Silence}$ tests $\Ids$; form a set of them by
those tests, as lines~\ref{ln:zadv} and~\ref{ln:zparty} form $\overline{\Ids}$ and
line~\ref{ln:call} forms $\{\id'' \neq \id\}$; name one outright; and
destructure a value that carries one. It may not compute with a session or an instance number: no
arithmetic, no ordering, and no session or instance written down on its own, apart from inside an
identity or process id named outright. Each of these counts as one step, whatever the identities
involved, as Definition~\ref{def:res} counts a call as one answer.
\end{definition}

Every machine in this paper is identity-atomic. The predicate $\opl{Guard}$ compares sessions for equality and
nothing else, no core reads $i$ at all, and the only process ids named outright are the three of
condition (iii). We take the model to be this throughout, and it is what makes
Definition~\ref{def:relocact} sound: on identity-atomic code the substitution there does not merely
rewrite literals but computes the relabelled behaviour, each permitted operation being equivariant
--- projections because $\theta$ leaves names and parties alone, equality and membership because
$\theta$ is injective and the sets tested and formed are relabelled with the code, a named identity
because substitution replaces it by its image. It is also what makes a
machine and its image take the same steps: they perform the same operations in the same order, the
control flow being the same on both, and the last clause of the definition charges one step for each
however the ids inside it have been renamed. So the budgets of Definition~\ref{def:res} carry over
unchanged --- which a cost model charging by the symbol would not grant, a relocation being free to
lengthen the encoding of a session.

Without the restriction both fail. A core returning $\id.i \bmod 2$ names no process id, so
substitution leaves it untouched; carried to another instance number it answers differently, and the
relabelled machine is not the original relabelled. A core calling
$(F,\id.s,\id.i + 1)$ misses its target in the same way. Neither is a machine anyone would write ---
Definition~\ref{def:advcls} permits an \emph{arbitrary} core, which is why the exclusion has to be
stated rather than assumed.

\begin{proposition}[Basic invariances]\label{prop:relocbasic}
Let $\theta$ be a relocation, $\F$ a functionality and $\pi$ a system. Then
\begin{enumerate}[label=(\roman*),leftmargin=2.2em]
  \item $\IDs(\theta\F) = \theta\bigl(\IDs(\F)\bigr)$, and occupants carrying pairwise disjoint
        identity sets keep them;
  \item $\theta$ carries each name class onto itself, so $\theta(\Stdpid) = \Stdpid$,
        $\theta(\Apid) = \Apid$ and $\theta(\Zpid) = \Zpid$, and every $\admits$ declared by names
        --- a local $\{\PID : \PID.F = F^{\parent}\}$ among them --- is its own image;
  \item $\op{global}(\theta\F) \iff \op{global}(\F)$;
  \item $\ncl{\theta\pi} = \theta\bigl(\ncl{\pi}\bigr) = \ncl{\pi}$;
  \item $\theta\pi$ is a system, and $\WF(\pi) \iff \WF(\theta\pi)$; and
  \item $\theta^{-1}$ is again a relocation.
\end{enumerate}
\end{proposition}

\begin{proof}
% BEGIN overview prop:relocbasic
Six items, and the useful thing to see is which of Definition~\ref{def:reloc}'s three conditions
each one spends. (i) is conditions (i) and injectivity: an identity set is fixed by two fields, one
relabelled and one untouched. (ii) is condition (i) alone --- name classes map into themselves, and
being a bijection on a partition, onto them as well, which is what makes the three reserved id sets
invariant. (iii) then follows from (ii), $\op{global}$ reading only a name and a subset of $\Stdpid$.
(iv) is condition (i) again, name closures being specified by names. (v) is the substantial one:
well-formedness transports because names match, the session conjunct is either waived by
$\op{global}$ or read between two standard ids where condition (ii) preserves it, and the relation
$R$ is equivariant by Definition~\ref{def:relocact}; the converse runs the same argument on
$\theta^{-1}$. (vi) is why that is legitimate --- conditions (i) and (ii) are an identity and an
equivalence, so they hold of the inverse, and (iii) holds because $\theta$ fixes those three ids.
% END overview prop:relocbasic
(i) An identity set is fixed by $\PID$ and $\Ps$, and a relocation relabels the first while leaving
the second alone; disjointness survives because $\theta$ is injective on identities.

(ii) By (i) of Definition~\ref{def:reloc} each name class maps into itself. The classes partition the
process ids, so for $x$ in the class of $F$ the point $\theta^{-1}(x)$ lies in some class, which maps
into itself and therefore is the class of $F$; so the map is onto that class as well. The three
displayed sets are unions of name classes, and a set of ids specified by a name is a union of them.

(iii) $\op{global}(\F)$ reads $\Stdpid \subseteq \F.\admits$ or $\F.F \in \{A,Z,C\}$. The second is
untouched by (i) of Definition~\ref{def:reloc}; for the first,
\[
  \Stdpid \subseteq \theta(\F.\admits)
  \quad\iff\quad
  \theta^{-1}(\Stdpid) \subseteq \F.\admits ,
\]
and $\theta^{-1}(\Stdpid) = \Stdpid$ by item (ii) above, $\theta$ being a bijection.

(iv) A name closure is specified by the names its system uses; $\theta$ keeps those, so $\theta\pi$
uses the names $\pi$ does and $\ncl{\theta\pi} = \ncl{\pi}$. And $\ncl{\pi}$ is a union of name
classes crossed with party ids, so item (ii) gives $\theta(\ncl{\pi}) = \ncl{\pi}$ as well.

(v) The members of $\theta\pi$ carry distinct process ids, $\theta$ being injective, and are
standard, their names being unchanged and their wrappers untouched by relabelling; the family is
finite. For $\WF$, take $\F \in \pi^{+}$ and an entry $(\F',R)$ with witness $\G$. Names match on
both sides. For the session conjunct $\G.s = \F.s$, either $\op{global}(\G)$ and item (iii) above
waives the conjunct on both sides, or $\G$ is standard --- and then so is $\F$, the entries of the
three supplied machines naming only each other and $\Corr.\uses$ being empty --- so that both ids are
standard and (ii) of Definition~\ref{def:reloc} preserves the test. Finally
$R(\F,\G) \iff R(\theta\F,\theta\G)$ by the equivariance
Definition~\ref{def:relocact} requires. The entries of the
three supplied machines are settled identically on both sides, every field $\serves$ reads of them
being fixed by items (ii) above and (iii) of Definition~\ref{def:reloc}. For the converse, run the
same argument on $\theta^{-1}$, a relocation by item (vi), using
$\theta^{-1}(\theta\pi) = \pi$ from Definition~\ref{def:relocact}.

(vi) Conditions (i) and (ii) of Definition~\ref{def:reloc} are stated as an identity of components and
as an equivalence, and so hold of the inverse; (iii) does because $\theta$ fixes those three ids.
\end{proof}

Two of these deserve a word. By (ii), only the pinned $\admits$ declarations of
Remark~\ref{rem:tight} move at all, and they move with the members they pin. And by (iv) name
closures are invariant, which is what makes Proposition~\ref{prop:relocemul} come out against an
environment class constrained exactly as before rather than one merely corresponding to it.

\begin{definition}[Pid-blind core]\label{def:pidblind}
A core is \emph{pid-blind} if $\theta\Op = \Op$ for every relocation $\theta$: it names no process
id that a relocation moves. It may compare the sessions of two identities in scope, as $\opl{Guard}$
does, and it may name one of the three process ids that (iii) of Definition~\ref{def:reloc} fixes ---
but not some other id of a reserved name, which (iii) leaves free to move.
\end{definition}

Every core in this paper is pid-blind. The wrappers of Chapter~\ref{sec:wrappers} read names and
parties only; $\opl{Mediate}$ hooks to $(A,\id.P)$ and lines~\ref{ln:stat} and~\ref{ln:stattwo} to
$(C,\id.P)$, both reserved; the $\Zenv$ box computes $\overline{\Ids}$ from $\pars$ by tests on
names and parties; $\Dum$ places the call its payload names; and $\Fsig$ reads $\id.P$ and its own
$\Ps$ and nothing else of an identity.

That is a stronger property than Lemma~\ref{lem:reloc} needs, and a different one. The lemma needs
Definition~\ref{def:atomic}, which every machine of an execution must meet, the environment and the
occupant of the adversary slot included --- and neither of those can be pid-blind, each having to
name the process ids it drives. What pid-blindness buys is the reading of the conclusion: for a
system of pid-blind cores, $\theta\pi$ is $\pi$ with five fields relabelled and not a line of code
rewritten, so a relocated protocol is the protocol, run elsewhere.

Before the lemma, the other half of the question the definition answers: which fields an execution
reads at all.

\begin{proposition}[Inert fields]\label{prop:inert}
Let $\F$ and $\F'$ be standard functionalities agreeing on $\PID$, $\Ps$, $\pars$ and code.
\begin{enumerate}[label=(\roman*),leftmargin=2.2em]
  \item If they agree on $\admits$ as well, so that they differ in $\uses$ alone, then no execution
        distinguishes them: in the two systems got by putting one or the other at that process id,
        every activation produces the same configuration.
  \item If they differ in $\admits$ alone, then from the same state and on the same call
        $\id.\fopl{Op}(in)$ from $\id'$, either both wrappers refuse at line~\ref{ln:guard} with
        $\rej$, or both admit and enter the same core at the same point in the same state.
\end{enumerate}
\end{proposition}

\begin{proof}
By inspection of where the two fields are read during an execution. The field $\uses$ is read only by the
static check of Definition~\ref{def:wfsys} and the results about it, never by $\opl{Exec}$ or by any
interface, which gives (i). The readers of $\admits$ inside an execution are two: line~\ref{ln:caller}
of $\opl{Guard}$, and the predicate $\op{global}$ and hence line~\ref{ln:session}. Both sit inside
$\opl{Guard}_{\F}$, which line~\ref{ln:guard} of the full interface \textbf{require}s before
anything else. (The field is read outside an execution too --- by Definition~\ref{def:wfsys} through
those two, by Definition~\ref{def:compat}, and by the typing conditions of
Definition~\ref{def:sim} and Remark~\ref{rem:tight} --- but those are conditions on a system, not
steps of one.) So the field decides that verdict and is read nowhere else in the wrapper: the lines
between line~\ref{ln:guard} and the core --- the register read, the corruption gate, the first
$\opl{Mediate}$ hook, the silenced set $\{\id'' \neq \id\}$ --- are functions of $\PID$, $\Ps$,
$\pars$, the code and the call, so two wrappers that both admit reach line~\ref{ln:call} together
and enter the same core at the same point in the same state. That is (ii).
\end{proof}

So $\admits$ is access control and nothing else, and $\uses$ is inert in every execution, as
Remark~\ref{rem:advpars} already records for the adversary slot. Inertness past the guard is a
statement about one step, and it cannot be strengthened to one about the value returned. The
obvious reason is that a call one admits and the other refuses sends the two states apart. The less
obvious one is that even a call both admit can end differently, because the core entered at
line~\ref{ln:call} may re-enter this very wrapper --- Section~\ref{sec:execution} permits an
instance to be re-entered, and line~\ref{ln:call} lets a core claim its own identity --- and the
inner call meets $\opl{Guard}_{\F}$ again, where the field differs. A core that places
$\id.\fopl{Op}'()$ as $\id$ and returns $1$ just if the answer is not $\rej$ separates an $\F$
admitting its own process id from an $\F'$ that does not, on a call both admitted. What survives is
the exact claim: the two differ in guard verdicts alone, and a single differing verdict, outer or
inner, is enough to separate them.

\begin{remark}[The two fields that are read]\label{rem:fields}
Neither $\Ps$ nor $\pars$ admits a statement of this kind, and for different reasons. The field $\Ps$ fixes
$\IDs(\F)$, hence which calls $\opl{Exec}$ dispatches to $\F$ at all and the second conjunct of
line~\ref{ln:interface}; and cores read it directly --- $\op{Initialize}$ of $\Fsig$ types
$\V{VK} : \Fsig.\Ps \to \Keys \cup \{\unset\}$, and the test
$\exists\, P' \notin \Cs$ of $\op{Verify}$ quantifies over that domain. So
enlarging or shrinking a party set is not a symmetry. What \emph{is} one is \emph{renaming}: a
permutation $\beta$ of $\bits$ extended to fix $A$ and $Z$, acting on identities, on $\Ps$, on
$\Cs$ and on the party-indexed state, gives an isomorphism of executions by the argument of
Lemma~\ref{lem:reloc}, provided cores read party ids only through equality and
indexing --- as $\Fsig$ does, its $\exists P' \notin \Cs : \vk = \V{VK}[P']$ being equivariant. The
two places a party id enters otherwise survive it as well: line~\ref{ln:zparty} names the party $Z$,
which $\beta$ fixes, and $\fopl{Corrupt}$ tests $\id.P \in \bits$, a set $\beta$ permutes.
The field $\pars$ is read by the code by definition of the field, and $\Zenv.\pars$ is read by
Definition~\ref{def:admis} besides; there is nothing to be invariant in. A relocation moves it,
and moves it to itself whenever it is declared by names, which
$\Zenv.\pars := \ncl{\overline{\pi}}$ is.
\end{remark}

Informally, Lemma~\ref{lem:reloc} says that relabelling sessions and instances throughout an
execution changes nothing observable: it is Proposition~\ref{prop:regroup}'s induction with $\theta$ in
place of the absorption bijection, and every check is the equivariance of one line of one wrapper.

\begin{lemma}[Relocation]\label{lem:reloc}
Let $\theta$ be a relocation, $\pi$ a session-uniform system, $\mathcal{X}$ an occupant of the
adversary slot (Definition~\ref{def:occ}), and $\Zenv$ a session-uniform environment hosting at no
identity of $\IDs(\pi)$, so that $\opl{Exec}$ is defined at the four. Then
\[
  \opl{Exec}\bigl(\theta\pi,\,\theta\Zenv,\,\theta\mathcal{X},\,\Corr\bigr)
  \qquad\text{and}\qquad
  \opl{Exec}\bigl(\pi,\,\Zenv,\,\mathcal{X},\,\Corr\bigr)
\]
are identically distributed.
\end{lemma}

\begin{proof}
% BEGIN overview lem:reloc
The correspondence of Definition~\ref{def:regcorr} at $\iota := \theta$, read up to $\theta$: two
states, program points or caller chains correspond when one is the other's image. Since $\theta$ is
a bijection the relation is symmetric, and the proof is again an induction on activations, whose
step the four paragraphs below discharge. \emph{Machines} checks that occupancy is relabelled
bijectively, so both executions are defined and start in correspondence. \emph{Guards} takes the
three conjuncts of $\opl{Guard}$ in turn; the first two are injectivity, and the third is the one
paragraph that does real work, since line~\ref{ln:session} is the only place the framework tells one
session from another --- this is where session-uniformity is spent, and without it a relocation
moving session $0$ would change the verdict. \emph{Silencers, mediation, the register} checks the
wrappers, and observes that the register needs no image at all, corruption being indexed by party
alone. \emph{Routing, claims, token} closes the step, identity-atomic code taking the same branches
on both sides.
% END overview lem:reloc
The correspondence is that of Definition~\ref{def:regcorr} with $\iota := \theta$ --- each machine
paired with its image, $\Corr$ with itself --- and with (C1), (C3) and (C4) read up to $\theta$: two
states, two program points, two chains of suspended callers correspond when one is the $\theta$-image
of the other. Since $\theta$ is a bijection on identities the relation is symmetric, and as before
the proof is an induction on activations. The paragraphs below discharge the step.

\emph{Machines.} $\theta\pi$ is a system and occupancy is relabelled bijectively, by (i) and (v) of
Proposition~\ref{prop:relocbasic}; since $\Zenv$ hosts clear of $\IDs(\pi)$ and $\theta\Zenv$ hosts
at the images of what $\Zenv$ hosts, the occupants are pairwise disjoint on the right exactly where
they are on the left, so both executions are defined. Line~\ref{ln:execinit} initialises $\Corr$ and
every member on both sides, and a member's initial state is its image's relabelled, giving (C1) at
the outset.

\emph{Guards.} Take a call $\id.\fopl{Op}(in)$ from $\id'$ at $\F$ and its image at $\theta\F$.
Line~\ref{ln:interface}: $\id.\PID = \F.\PID$ iff $\theta(\id.\PID) = \theta(\F.\PID)$ by
injectivity, and $\id.P \in \Ps$ is untouched. Line~\ref{ln:caller}: $\id'.\PID \in \F.\admits$ iff
$\theta(\id'.\PID) \in \theta(\F.\admits)$, again by injectivity, and the party test reads
components $\theta$ leaves alone. Line~\ref{ln:session} is a disjunction whose second disjunct agrees
on the two sides by (iii) of Proposition~\ref{prop:relocbasic}, so suppose $\op{global}(\F)$ fails ---
it fails on both sides together or on neither --- and consider the first. Every machine of the
execution is session-uniform: $\pi$ and $\Zenv$'s hosted copies by hypothesis, and $\mathcal{X}$,
$\Corr$ and $\Zenv$'s own interfaces because their process ids are $A$, $C$ and $Z$, which
$\op{global}$ admits by its second disjunct. So $\F.\admits$ meets neither $\Apid$ nor $\Zpid$, and a
claim carrying the id of $\Adv$ or of $\Zenv$ is refused at line~\ref{ln:caller} on both sides --- (iii)
fixing those ids, and $\theta(\F.\admits)$ missing $\Apid$ and $\Zpid$ just as $\F.\admits$ does,
those two sets being $\theta$-invariant by (ii) of Proposition~\ref{prop:relocbasic}. The third
supplied id cannot arrive at all, $\Corr$ placing no calls. The claims that reach
line~\ref{ln:session} therefore carry standard process ids, and so does $\id$, the name of $\F$ lying
in $\bits$ once $\op{global}(\F)$ fails; so $\id'.s = \id.s$
holds of the images exactly when it holds of the originals, by (ii) of Definition~\ref{def:reloc}.
This is the one paragraph session-uniformity is for; without it a relocation moving session $0$ would
change the verdict here. So the
verdict is the same on both sides, at every callee, including the leakage interface of
Section~\ref{sec:leakage}, whose three requirements read only what these lines read.

\emph{Silencers, mediation, the register.} The silenced set of line~\ref{ln:call} is
$\{\id'' \neq \id\}$, whose image is $\{\id'' \neq \theta\id\}$, so a core claims its own identity
on both sides and nothing else. The set $\overline{\Ids}$ of the $\Zenv$ box is computed from
$\theta(\pars)$ by tests on names and parties, so it is the image of the set computed on the left,
and admits exactly the image of what that admits. The wrapper $\opl{Mediate}$ reads $\Cs$, $\id.P$ and
$\id'.F$ and hooks to $(A,\id.P)$; lines~\ref{ln:stat} and~\ref{ln:stattwo} address $(C,\id.P)$.
Every id named outright here is one of the three that (iii) of Definition~\ref{def:reloc} fixes. And
the register itself needs no image:
$\Cs$ collects \emph{party} ids, which a relocation does not touch, so $\Corr$ is literally the
same machine holding the same set at corresponding configurations. Corruption being indexed by
party alone (Section~\ref{sec:corruption}) is what buys this.

\emph{Routing, claims, token.} $\opl{Exec}$ dispatches a call to the occupant of its target
identity, and targets correspond under $\theta$ while occupancy does by the first paragraph; a call
addressed to an unoccupied identity is answered $\rej$ on the left exactly when its image is on the
right. Claims correspond by the previous paragraph. The code being identity-atomic
(Definition~\ref{def:atomic}), a machine and its image take the same branches and place calls that
are $\theta$-images of one another, so they reach the same program point with values related by
$\theta$; and each reads the same coins in the same order, which is (C2). The
token discipline is unchanged, $\opl{Exec}$ having one operation on both sides and the same chain of
suspensions arising, which is (C3) and (C4).

Corresponding halting configurations output alike, the output being $\Zenv$'s and
$\theta\Zenv$'s return at $(Z,Z)$, an identity (iii) of Definition~\ref{def:reloc} fixes. So the two
executions agree as distributions, in particular on the probability of returning $1$.
\end{proof}

The lemma is easier to see than to state. One execution, and the same execution with every process
id renamed, side by side:

\begin{center}
\begin{tikzpicture}[
  sysnode/.style={draw=accent, rounded corners=1pt, inner sep=2pt, font=\scriptsize,
                  minimum height=4.5mm},
  wire/.style={draw=black!70},
  shell/.style={draw=accent, rounded corners=2pt, inner sep=1.8mm},
  role/.style={font=\small},
  rel/.style={font=\normalsize},
  tag/.style={font=\tiny, color=accentsoft},
]

%% ---------------- left: the execution as given ----------------
\node[sysnode] (L-m1) at (-9mm, -12mm) {$(\op{Sig},3,0)$};
\node[sysnode, right=1mm of L-m1] (L-m2) {$(\op{Net},3,1)$};
\node[shell, fit=(L-m1)(L-m2)] (L-sys) {};
\node[role, above=5mm of L-sys] (L-Z) {$\Zenv$};
\node[role, below=5mm of L-sys] (L-X) {$\mathcal{X}$};
\draw[wire] (L-Z) -- (L-sys.north);
\draw[wire] (L-sys.south) -- (L-X);
\node[role, left=1.5mm of L-sys] {$\pi$};
\node[sysnode] (L-C) at (26mm, -31mm) {$\Corr$};

%% ---------------- the relation ----------------
\node[rel] at (44mm, -12mm) {$\equiv$};
\node[tag, anchor=south, align=center] at (44mm, -9mm)
  {identically\\ distributed};

%% ---------------- right: its image ----------------
\node[sysnode] (R-m1) at (71mm, -12mm) {$(\op{Sig},7,1)$};
\node[sysnode, right=1mm of R-m1] (R-m2) {$(\op{Net},7,0)$};
\node[shell, fit=(R-m1)(R-m2)] (R-sys) {};
\node[role, above=5mm of R-sys] (R-Z) {$\theta\Zenv$};
\node[role, below=5mm of R-sys] (R-X) {$\theta\mathcal{X}$};
\draw[wire] (R-Z) -- (R-sys.north);
\draw[wire] (R-sys.south) -- (R-X);
\node[role, left=1.5mm of R-sys] {$\theta\pi$};
\node[sysnode] (R-C) at (106mm, -31mm) {$\Corr$};

%% ---------------- annotations ----------------
\draw[accentsoft, dotted, thick] (L-C) -- (R-C);
\node[tag, anchor=north, align=center] at (66mm, -32.5mm)
  {the same machine, not an image: $\Cs$ holds party ids, which $\theta$ leaves alone};

\node[tag, anchor=north, align=center] at (48mm, -40mm)
  {$\theta : (F,s,i) \mapsto \bigl(F,\ \hat\theta(s),\ \theta_{F,s}(i)\bigr)$ --- the name kept,
   the session block moved, the instance free within it;\\
   $(Z,0,0)$, $(A,0,0)$ and $(C,0,0)$ pinned, so $\Zenv$, $\mathcal{X}$ and $\Corr$ keep their
   identities and only the process ids \emph{inside} their code move};

\end{tikzpicture}
\figcaption{fig:reloc}{One execution beside the same execution with every process id renamed. Lemma~\ref{lem:reloc} is that the two agree as distributions.}
\end{center}

Everything a relocation may do is legible in the two systems. The \emph{name} is kept, because
$\opl{Guard}$ and $\WF$ match by name and nothing may quietly become a different functionality.
The \emph{session} moves as a block --- both members go from $3$ to $7$ together --- because
line~\ref{ln:session} asks whether caller and callee \emph{share} a session and never which one
they are in. And the \emph{instance} is free to move within that block, $0$ and $1$ here changing
places, because no line of any code reads it. The three ids that occur as literals are pinned, so
$\Zenv$, $\mathcal{X}$ and $\Corr$ sit at the same identities on both sides and only the process
ids \emph{inside} their code are rewritten.

$\Corr$ is not even that: it is the same machine holding the same set, because $\Cs$ collects
party ids and a relocation does not touch a party. Indexing corruption by party alone
(Section~\ref{sec:corruption}) is what buys this, and it is worth noticing how much it buys --- a
register indexed by process id would have to be carried across by $\theta$ like everything else,
and the corruption pattern would then be a thing a relocation could disturb.

One warning the picture cannot give. Drawn this way the lemma looks free, and it nearly is: every
guard conjunct, every silencer, every mediation hook reads only names, parties and the three
pinned ids, all of which either survive $\theta$ or are fixed by it. The exception is the single
session test, and it is the reason session-uniformity is a hypothesis rather than a remark --- a
callee reachable from outside every session but sitting in one of them would answer the supplied
machines differently before and after a $\theta$ that moves session $0$, and the two executions
would come apart at that one line. That paragraph of the proof is the whole of the work; the rest
is bookkeeping the picture already shows.

\begin{proposition}[Emulation transports]\label{prop:relocemul}
Let $\theta$ be a relocation and $\pi$, $\varphi$ systems; for (ii) and (iii) let the two of them and
every environment of the classes named be session-uniform, which (i) does not need. Then
\begin{enumerate}[label=(\roman*),leftmargin=2.2em]
  \item $\ZenvSet_{\theta\pi,\theta\varphi} = \theta\bigl(\ZenvSet_{\pi,\varphi}\bigr)$, and
        $\theta$ carries adversaries to adversaries and simulators to simulators, fixing the dummy
        of Section~\ref{app:dummy};
  \item if $\pi$ UC-emulates $\varphi$ within $\varepsilon$ against
        $(\Aset,\SimSet,\ZenvSet)$, then $\theta\pi$ UC-emulates $\theta\varphi$ within
        $\varepsilon$ against $(\theta\Aset,\theta\SimSet,\theta\ZenvSet)$; and
  \item for budgeted systems and every $\bd{a},\bd{s},\bd{z}$,
        $\advtg^{\op{uc}}_{\theta\pi,\theta\varphi}[\bd{a},\bd{s},\bd{z}]
         = \advtg^{\op{uc}}_{\pi,\varphi}[\bd{a},\bd{s},\bd{z}]$.
\end{enumerate}
\end{proposition}

\begin{proof}
% BEGIN overview prop:relocemul
Three claims, and each rests on a different part of Proposition~\ref{prop:relocbasic}. (i) is
bookkeeping about the classes: both clauses of Definition~\ref{def:admis} transport, the $\pars$
clause because name closures are $\theta$-invariant and the hosting clause by injectivity, with the
containment going both ways because $\theta^{-1}$ is a relocation too; adversaries go to adversaries
and simulators to simulators because the fields those definitions read are the ones $\theta$ fixes,
and the dummy is fixed outright, its core being pid-blind. (ii) is then immediate from
Lemma~\ref{lem:reloc} applied to each world, the returned simulator being $\theta\Sim$, which
depends on $\theta\Adv$ alone. (iii) is the concrete form: $\theta$ is a bijection of each of the
three budgeted classes onto the class the conclusion quantifies over, and reindexing a function
along bijections changes neither a supremum nor an infimum. Session-uniformity is preserved
throughout, which is checked once rather than in each part.
% END overview prop:relocemul
For (i), Definition~\ref{def:admis} asks two things of $\Zenv$. The $\pars$ clause is
$\theta$-invariant outright: $\theta\Zenv.\pars = \theta(\Zenv.\pars)$ contains
$\theta(\ncl{\pi} \cup \ncl{\varphi}) = \ncl{\theta\pi} \cup \ncl{\theta\varphi}$ exactly when
$\Zenv.\pars$ contains $\ncl{\pi} \cup \ncl{\varphi}$, name closures being invariant by (iv) of
Proposition~\ref{prop:relocbasic}. The hosting clause transports by injectivity, $\theta\Zenv$
hosting at the images of what $\Zenv$ hosts and $\IDs(\theta\pi) \cup \IDs(\theta\varphi)$ being the
image of $\IDs(\pi) \cup \IDs(\varphi)$ by (i) of that proposition. The containment goes both ways,
$\theta^{-1}$ being a relocation by (vi). For
the slot, an adversary occupies $\IDs(\Adv)$, whose identities carry the process id (iii) of
Definition~\ref{def:reloc} fixes and party components $\theta$ leaves alone, and it carries $\Adv$'s
wrapper, which is pid-blind; so $\theta$ relabels the core alone and $\theta\Adv$ is an adversary of
the same shape. Definition~\ref{def:sim} asks $\admits \supseteq \Stdpid \cup \Zpid$, a set $\theta$
fixes by (ii), so simulators go to simulators. The dummy's core is pid-blind --- line~\ref{ln:dumout}
places the call its payload names --- so $\theta\Dum = \Dum$. Nothing of the kind is claimed of the
$\Adv$ of Section~\ref{sec:adv}, whose core Definition~\ref{def:advcls} leaves arbitrary, and nothing
below needs it.

For (ii), fix $\Adv \in \Aset$ and let $\Sim \in \SimSet$ be the simulator the hypothesis supplies
for it. Given $\theta\Adv \in \theta\Aset$, return $\theta\Sim \in \theta\SimSet$; it depends on
$\theta\Adv$ alone, so the quantifier order of Definition~\ref{def:uc} is met. For any
$\theta\Zenv \in \theta\ZenvSet$, Lemma~\ref{lem:reloc} applied to each world gives
\[
  \advtg^{\op{uc}}_{\theta\pi,\theta\varphi}\bigl(\theta\Zenv,\theta\Adv,\theta\Sim\bigr)
  \;=\;
  \advtg^{\op{uc}}_{\pi,\varphi}\bigl(\Zenv,\Adv,\Sim\bigr)
  \;\leq\; \varepsilon ,
\]
the two executions on the left being the $\theta$-images of the two on the right. By (i) the
conclusion is a legal statement of emulation, $\theta\ZenvSet \subseteq \ZenvSet_{\theta\pi,\theta\varphi}$.

Session-uniformity is preserved throughout: $\theta$ fixes $\Apid$ and $\Zpid$ by (ii) of
Proposition~\ref{prop:relocbasic} and $\op{global}$ by its (iii), so $\theta\F$ is session-uniform
exactly when $\F$ is, and the same of a system and of an environment memberwise. Every class below is
therefore closed under $\theta$ in this respect as well.

For (iii), a machine and its image take the same steps and answer the same calls, the code being
identity-atomic and its control flow therefore the same on both
(Definition~\ref{def:atomic}); a member of a budgeted system and its image carry the same budget for
the same reason. So $\theta$ is a bijection of each of the three classes onto the class the
conclusion quantifies over: $\Aset(\bd{a})$ and $\SimSet(\bd{s})$ onto themselves, and
$\ZenvSet_{\pi,\varphi}(\bd{z})$ onto $\ZenvSet_{\theta\pi,\theta\varphi}(\bd{z})$ by (i), each with
inverse induced by $\theta^{-1}$. Lemma~\ref{lem:reloc} makes the advantage of
Definition~\ref{def:adv} equal at corresponding triples, so the two nested extrema of
Definition~\ref{def:concadv} are the same three operators applied to one function reindexed along
bijections, and reindexing changes neither a supremum nor an infimum. No finiteness is needed here;
what the finiteness of Definition~\ref{def:concadv} settles is that the extrema are \emph{attained},
which is a separate matter and holds on both sides alike.
\end{proof}

Emulation is thus a property of a system up to relocation, and a proof may fix the session and the
instance at the outset. To say so for a family of instances, one needs to know when two lists of
process ids are related by a relocation at all, and the answer is three conditions.

\begin{proposition}[When one tuple relocates to another]\label{prop:orbit}
Let $\vec q = (q_1,\dots,q_m)$ and $\vec q\,' = (q_1',\dots,q_m')$ be tuples of process ids of
standard name. There is a relocation $\theta$ with $\theta(q_j) = q_j'$ for every $j$ if and only if,
for all $j$ and $k$,
\[
  \begin{gathered}
  q_j.F = q_j'.F, \qquad
  \bigl(q_j.s = q_k.s \iff q_j'.s = q_k'.s\bigr), \\[3pt]
  \bigl(q_j = q_k \iff q_j' = q_k'\bigr) .
  \end{gathered}
\]
In particular any two process ids of one name are related by a relocation.
\end{proposition}

\begin{proof}
% BEGIN overview prop:orbit
Necessity is the three conditions of Definition~\ref{def:reloc} read off in order. Sufficiency is a
construction: build $\hat\theta$ from the session assignment, which the second condition makes a
well-defined bijection between the sessions occurring on the two sides, and extend it to a
permutation because what is left over is of one size on each side; then, at each name and session,
build $\theta_{F,s}$ from the instance assignment, which the third and fourth conditions make
injective and hence again extensible. Taking the identity everywhere else assembles a relocation
carrying $\vec q$ to $\vec q\,'$, and (iii) holds because the reserved ids were never touched.
% END overview prop:orbit
Necessity is the conditions of Definition~\ref{def:reloc} read off in order: names by (i), sessions
shared or not by (ii), and equality by $\theta$ being injective.

For sufficiency, let $S$ and $S'$ be the finitely many sessions occurring in $\vec q$ and in
$\vec q\,'$. The second condition makes $q_j.s \mapsto q_j'.s$ a well-defined bijection $S \to S'$ ---
onto, every $q_j'.s$ being an image --- so what is left over on the two sides is of one size, finite or
infinite alike, and it extends to a permutation $\hat\theta$ of session ids. Now fix a name $F$
and a session $s \in S$ and consider the indices $j$ with $q_j.F = F$ and $q_j.s = s$. Their
instance assignment $q_j.i \mapsto q_j'.i$ is well defined and injective: if $q_j.i = q_k.i$ then
$q_j = q_k$, so $q_j' = q_k'$ by the fourth condition and the instances agree; and if
$q_j.i \neq q_k.i$ then $q_j \neq q_k$, so $q_j' \neq q_k'$ while their names and sessions agree, and
the instances differ. Each is again a bijection between the instance numbers occurring on the two
sides, so it extends to a permutation $\theta_{F,s}$ as $\hat\theta$ did; take the
identity at every other name and session, and let $\theta$ act as
$(F,s,i) \mapsto (F,\hat\theta(s),\theta_{F,s}(i))$ on the standard names and as the identity on the
reserved ones. Then $\theta$ meets Definition~\ref{def:reloc}, (iii) because the reserved ids are
fixed outright, and it carries $\vec q$ to $\vec q\,'$. For the last claim, a single pair $q$, $q'$ of
one name meets the three conditions vacuously.
\end{proof}

\begin{corollary}[One instance suffices]\label{cor:onepid}
Call a pair of maps $\vec q \mapsto \pi(\vec q)$ and $\vec q \mapsto \varphi(\vec q)$ an
\emph{instance family} if $\pi(\theta\vec q) = \theta\bigl(\pi(\vec q)\bigr)$ and
$\varphi(\theta\vec q) = \theta\bigl(\varphi(\vec q)\bigr)$ for every relocation $\theta$, where
$\vec q$ lists the process ids occupied by $\pi(\vec q)$ and $\varphi(\vec q)$ together --- a
protocol, its private subroutines and the specification it is measured against indexed by one tuple
and moving with it. Suppose $\pi(\vec q)$ and $\varphi(\vec q)$ are session-uniform and that
$\pi(\vec q)$ UC-emulates $\varphi(\vec q)$ within $\varepsilon$ against
$\bigl(\Aset(\bd{a}),\SimSet(\bd{s}),\ZenvSet_{\pi(\vec q),\varphi(\vec q)}(\bd{z})\bigr)$ cut down to
the session-uniform environments, for a single $\vec q$. Then at every $\vec q\,'$ meeting the three
conditions of Proposition~\ref{prop:orbit}, $\pi(\vec q\,')$ UC-emulates $\varphi(\vec q\,')$ within
the same $\varepsilon$ against $\Aset(\bd{a})$, $\SimSet(\bd{s})$ and the same cut-down class for that
pair. In particular, for a family occupying one process id, emulation at one id gives emulation at
every id of that name.
\end{corollary}

\begin{proof}
Proposition~\ref{prop:orbit} supplies a relocation carrying $\vec q$ to $\vec q\,'$, equivariance
turns it into one carrying the pair of systems to the pair at $\vec q\,'$, and
Proposition~\ref{prop:relocemul} transports the statement --- session-uniformity of the relocated
pair coming from that of the original, as the proof of that proposition records. The last claim is the
final clause of Proposition~\ref{prop:orbit}: a one-id tuple meets the three conditions against any
other id of its name.
\end{proof}

Indexing by the whole tuple is what makes the condition satisfiable, and the point is worth
recording, since indexing by the principal process id alone --- the first thing one writes --- does
not. A relocation may fix a given $q$ and still move a private member sitting at another id, so
equivariance in $q$ alone would force every member of $\pi(q)$ to have an id fixed by the whole
stabiliser of $q$; and the only such ids are $q$ itself and the reserved three, every other being
movable by some relocation that fixes $q$. Equivariance in $q$ alone therefore admits only systems
with no private members at all, which is not the case anyone wants.

Written for a pair of systems instead, and asking what else has to move, the statement splits in
two. A new \emph{instance} of a name can be taken up by itself; a new \emph{session} cannot, and
the reason is condition~(ii) of Definition~\ref{def:reloc}: a relocation permutes sessions as
blocks, so nothing can carry one member out of session $s$ while leaving another behind in it.

\begin{corollary}[Moving one instance, moving one session]\label{cor:moveone}
Let $\pi$ and $\varphi$ be session-uniform systems each carrying a member at one and the same process
id $q_0$, and let $q$ be a process id with $q.F = q_0.F$ occupied by neither system. Suppose $\pi$
UC-emulates $\varphi$ within $\varepsilon$ against
$\bigl(\Aset(\bd{a}),\SimSet(\bd{s}),\ZenvSet_{\pi,\varphi}(\bd{z})\bigr)$. In each case below the
conclusion is that the moved pair UC-emulates within the same $\varepsilon$ against $\Aset(\bd{a})$,
$\SimSet(\bd{s})$ and the largest admissible class of that pair at $\bd{z}$.
\begin{enumerate}[label=(\roman*),leftmargin=2.2em]
  \item If $q.s = q_0.s$, the transposition $\theta := (q_0\ q)$ of instance numbers at the name
        $q_0.F$ in that session is a relocation, and it fixes the process id of every other member
        of either system. So $\theta\pi$ and $\theta\varphi$ are $\pi$ and $\varphi$ with the
        instance at $q_0$ carried to $q$ and nothing else disturbed --- save that a member naming
        the moved instance in a field, by pinning it among its callers in the manner of
        Remark~\ref{rem:tight} or otherwise, has that mention carried along with it.
  \item If $q.s \neq q_0.s$, no relocation carries $q_0$ to $q$
        while fixing a process id of standard name in session $q_0.s$ --- so no member of a system
        lying there is fixed, and the block moves entire. One carrying $q_0$
        to $q$ exists; and if in addition neither system has a member in session $q.s$, one may be
        chosen that carries each member of $\pi$ and of $\varphi$ lying in session $q_0.s$ into
        session $q.s$ and fixes the process id of every member lying elsewhere. Without that
        further hypothesis the members already in session $q.s$ move too --- by the same argument
        applied at $q$, no relocation carrying $q_0$ to $q$ can fix one of them --- so the
        conclusion is then about a pair with two blocks moved rather than one.
\end{enumerate}
\end{corollary}

\begin{proof}
% BEGIN overview cor:moveone
Both parts are Proposition~\ref{prop:relocemul} once a relocation has been exhibited, so the work is
exhibiting one. For (i) the transposition of two ids of a common session does it, and the check is
that it moves nothing else occupied. For (ii) the negative half comes first --- a relocation carrying
$q_0$ to an id of a different session moves \emph{every} standard id of that session, by condition
(ii) of Definition~\ref{def:reloc} --- and the positive half then constructs the relocation that
moves session $q_0.s$ and fixes every third session. The last sentence applies the negative half to
$\theta^{-1}$, which is a relocation by (vi) of Proposition~\ref{prop:relocbasic}.
% END overview cor:moveone
Both parts are Proposition~\ref{prop:relocemul} once the relocation is exhibited, its part~(iii)
supplying the classes. For (i), the transposition has $\hat\theta = \op{id}$, so conditions (i) and (ii)
of Definition~\ref{def:reloc} hold, and (iii) does because $q_0.F$ is the name of a member of a
system and so none of $A$, $Z$, $C$. It moves only the two ids $q_0$ and $q$, and $q$ is occupied by
neither system, so every other occupied id is fixed; what is not fixed is a field naming $q_0$, which
Definition~\ref{def:relocact} relabels.

For (ii), suppose $\theta(q_0) = q$ and let $p$ be any
process id of standard name with $p.s = q_0.s$; by condition (ii) of Definition~\ref{def:reloc},
$\theta(p).s = \theta(q_0).s = q.s \neq p.s$, so $\theta(p) \neq p$. For existence, take $\hat\theta$ the
transposition of the two sessions and instance permutations
carrying $q_0.i$ to $q.i$ at the name $q_0.F$ in session $q_0.s$, the identity elsewhere. That
$\theta$ carries session $q_0.s$ into session $q.s$ and fixes every id of a third session; if no
member lies in session $q.s$, the members it moves are exactly those of session $q_0.s$. For the
final sentence, let $\theta$ carry $q_0$ to $q$ and let $m$ be a member lying in session $q.s$. Then
$\theta^{-1}$ is a relocation by (vi) of Proposition~\ref{prop:relocbasic} and carries $q$ to $q_0$
with $q_0.s \neq q.s$, so the negative claim applied to it moves every standard id of session $q.s$,
$m$ among them; and $\theta(m) = m$ would give $\theta^{-1}(m) = m$, so $\theta$ moves $m$ too.
\end{proof}

Sessions and instances are not the only components a permutation can move, and the third is what
settles a question Chapter~\ref{sec:properties} raises: a game exercises $\F$ under a claim carrying
the game's own name (line~\ref{ln:call}), so $\F.\admits$ must name the game, while the $\F$ that ships
inside a protocol admits that protocol instead. The two differ in one field, and it is the field
$\opl{Guard}$ reads.

\begin{definition}[Renaming]\label{def:rename}
A \emph{renaming} is a permutation $\nu$ of the process ids of the form
\[
  \nu(F,s,i) \;=\; \bigl(\bar\nu(F),\ s,\ i\bigr)
\]
for a permutation $\bar\nu$ of the names fixing $A$, $Z$ and $C$. It acts on identities, fields, code
and declared sets as Definition~\ref{def:relocact} prescribes, and in addition replaces every name
occurring outright in code by its image.
\end{definition}

The extra clause is needed and was not before. Definition~\ref{def:relocact} substitutes process ids,
which for a relocation is enough, names being what a relocation keeps; a renaming moves them, so a
test against a bare name --- $\id'.F \neq A$ in $\opl{Mediate}$, or the $\id'.F = Z$ the dummy of
Section~\ref{app:dummy} dispatches on --- has to travel too. For the machines of this paper it travels nowhere: the
only names they mention outright are the three $\bar\nu$ fixes, so a renaming changes no line of code
at all, only fields.

\begin{proposition}[Renaming]\label{prop:rename}
Let $\nu$ be a renaming. Then
\begin{enumerate}[label=(\roman*),leftmargin=2.2em]
  \item $\IDs(\nu\F) = \nu(\IDs(\F))$; $\nu(\Stdpid) = \Stdpid$, $\nu(\Apid) = \Apid$ and
        $\nu(\Zpid) = \Zpid$; $\op{global}(\nu\F) \iff \op{global}(\F)$; $\nu\F$ is session-uniform
        exactly when $\F$ is; $\ncl{\nu\pi} = \nu(\ncl{\pi})$; $\nu\pi$ is a system with
        $\WF(\pi) \iff \WF(\nu\pi)$; and $\nu^{-1}$ is a renaming;
  \item for every system $\pi$, every occupant $\mathcal{X}$ of the adversary slot and every
        environment $\Zenv$ hosting at no identity of $\IDs(\pi)$ --- with no session-uniformity
        asked of any of them ---
        \[
          \opl{Exec}\bigl(\nu\pi,\,\nu\Zenv,\,\nu\mathcal{X},\,\Corr\bigr)
          \quad\text{and}\quad
          \opl{Exec}\bigl(\pi,\,\Zenv,\,\mathcal{X},\,\Corr\bigr)
        \]
        are identically distributed; and
  \item $\ZenvSet_{\nu\pi,\nu\varphi} = \nu(\ZenvSet_{\pi,\varphi})$; if $\pi$ UC-emulates $\varphi$
        within $\varepsilon$ against $(\Aset,\SimSet,\ZenvSet)$ then $\nu\pi$ UC-emulates $\nu\varphi$
        within $\varepsilon$ against $(\nu\Aset,\nu\SimSet,\nu\ZenvSet)$; and the concrete advantages
        of Definition~\ref{def:concadv} agree at every budget.
\end{enumerate}
\end{proposition}

\begin{proof}
% BEGIN overview prop:rename
All three parts are the corresponding relocation proofs re-run, and the overview is what changes.
Names are now \emph{permuted} rather than carried onto themselves, which is the point of the
definition; the reserved classes and $\Stdpid$ stay fixed, and with them $\op{global}$ and
session-uniformity, while a name closure is carried rather than fixed. In (ii) the one hard paragraph
of Lemma~\ref{lem:reloc} --- the session conjunct --- becomes trivial, a renaming moving no session,
which is why no session-uniformity hypothesis appears here at all. In (iii) what is lost is the
reading, not the result: the environment class of the conclusion corresponds to the hypothesis's
rather than being it.
% END overview prop:rename
(i) is the proof of Proposition~\ref{prop:relocbasic} with two changes of bookkeeping. Name classes are
now permuted rather than carried onto themselves, so an $\admits$ declared by names becomes the one
declared by the images --- which is the whole point of the definition --- while $\Stdpid$, being the union
of every standard class, and $\Apid$ and $\Zpid$, being single reserved classes, are fixed, and with
them $\op{global}$, Definition~\ref{def:sessunif} and the reserved ids of (iii) of
Definition~\ref{def:reloc}. And a name closure is carried rather than fixed, $\ncl{}$ being specified by
the names a system uses. For $\WF$ the session conjunct is untouched, $\nu$ moving no session, and the
name match holds on one side exactly when on the other.

(ii) is the proof of Lemma~\ref{lem:reloc} verbatim but for its one hard paragraph, which becomes
trivial: line~\ref{ln:session}'s first disjunct is literally unchanged, $\nu$ leaving every session
alone, and its second is preserved by (i). Nothing needs excluding, which is why no session-uniformity
appears here. Every other check reads a name for equality or a membership in a set that travels with
the code, and Definition~\ref{def:atomic} together with the extra clause of
Definition~\ref{def:rename} makes the substitution compute the relabelled behaviour as before.

(iii) is the proof of Proposition~\ref{prop:relocemul}, which used of the closures only that
$\nu(\ncl{\pi} \cup \ncl{\varphi}) = \ncl{\nu\pi} \cup \ncl{\nu\varphi}$ and never that they stay put.
What is lost is the reading recorded after Proposition~\ref{prop:relocbasic}: the environment class of
the conclusion corresponds to the one of the hypothesis rather than being constrained identically,
since $\pars$ must now name the images.
\end{proof}

\begin{corollary}[The caller's name is immaterial]\label{cor:callername}
Let $\nu$ be a renaming and let $\gamma = \{\Gm\} \cup \sigma$ be a game system for $\F$
(Definition~\ref{def:prop}). Then $\nu\gamma = \{\nu\Gm\} \cup \nu\sigma$ is a game system, $\nu\Gm$ is
a game for $\nu\F$, and for every finalization $\op{Fin}$ of $\Gm$
\[
  \padv{\op{Fin}}{\nu\Gm,\nu\sigma,\nu\Zenv,\nu\mathcal{X}}
  \;=\;
  \padv{\op{Fin}}{\Gm,\sigma,\Zenv,\mathcal{X}}
\]
for every $\Zenv \in \ZenvSet_{\gamma,\gamma}$ and every occupant $\mathcal{X}$ of the adversary slot.
Since $\nu\Dum = \Dum$, the abbreviated form reads
\[
  \padv{\op{Fin}}{\nu\Gm,\nu\sigma,\nu\Zenv} \;=\; \padv{\op{Fin}}{\Gm,\sigma,\Zenv} .
\]
So $\sigma$ has the property $\op{Fin}$ within $\varepsilon$ against a class of environments exactly
when $\nu\sigma$ has it within $\varepsilon$ against the image of that class.
\end{corollary}

\begin{proof}
That $\nu\gamma$ is a system with $\WF(\nu\gamma)$ is (i). For $\nu\Gm$ being a game,
check the four fields of Definition~\ref{def:prop}: its process id is
$(\bar\nu(\Gm.F),0,0)$, $\nu$ moving neither session nor instance; its $\Ps$ is untouched and
$\nu\F.\Ps = \F.\Ps$, so it is $\nu\F.\Ps \cup \{Z\}$ as required; its $\admits$ is
$\nu(\Zpid) = \Zpid$ by (i); its entry is $(\nu\F,\serves_{\Gm})$, and $\serves_{\Gm}$ is equivariant,
reading $\Ps$ and a membership of $\admits$. Line~\ref{ln:gmroot} still serves every finalization at
the
root, party ids being no part of a process id, and a renaming moves no interface name, so $\op{Fin}$
names the same one on both sides. For the advantage, $\Win_{\op{Fin}}$ is the event that the
interface $\fopl{Fin}$
at $(\Gm.\PID,Z)$ returns $1$, and $\nu$ carries that identity to $(\nu\Gm.\PID,Z)$; by (ii)
the two executions are identically distributed under the correspondence, so the event has the same
probability on both sides and Definition~\ref{def:propadv} gives the same number. That
$\nu\Zenv$ is admissible for $\nu\gamma$ is (iii). For the abbreviated form, $\nu\Dum = \Dum$: the
dummy's fields are $\PID := A$ and sets of ids that (i) fixes, and its core dispatches on
$\id'.F = Z$ and places the call its payload names, so no name it mentions moves.
\end{proof}

Read at the transposition of $\Gm.F$ with a name $\gamma$ does not use, and with $\Gm.F$ carried by
$\Gm$ alone, this says what one wants of Definition~\ref{def:prop}. The game may be given any name and
$\F.\admits$ renamed to match; no other field moves, and by Definition~\ref{def:rename} no line of
code does either. A property is therefore not tied to what its game is called, and the $\F$ whose
property one proves may admit whatever caller name the deployment will use.

What no renaming gives is a different caller \emph{programme}. It carries a game to a game and a
protocol to a protocol, never a game to the protocol that will use $\F$ in earnest --- and a protocol
that hands $\F$'s answers to the adversary satisfies no property of $\F$, so nothing could.
Proposition~\ref{prop:inert} narrows the gap from the other side, $\admits$ being read nowhere but
inside $\opl{Guard}_{\F}$: a property is a statement about $\F$'s code and the access control around
it, and about nothing else that field could carry. Crossing the rest is what emulation is for, and
Proposition~\ref{prop:proptransfer} is the crossing. Relocations and renamings compose, each being an
isomorphism of executions, so a system may be moved in name, session and instance at once.

\begin{remark}[Session $0$ is distinguished]\label{rem:sesszero}
Session-uniformity is not an artefact of the proofs, and the functionality it excludes is a real one.
Line~\ref{ln:session} compares the \emph{caller's} session, and $\Adv$ and $\Zenv$ sit in session $0$,
so a functionality that is \emph{local} in the sense of Chapter~\ref{sec:functionalities} yet admits
$\Apid$ is addressable from the adversary slot exactly when its own session is $0$: with
$\F.s \neq 0$ the conjunct $\id'.s = \id.s$ fails and $\op{global}(\F)$ does not save it. The
difference is observable --- it is how $\Dum$ carries out an instruction of Section~\ref{app:dummy}
at a corrupt party --- so for such a pair no relocation moving session $0$ is an isomorphism, and
Definition~\ref{def:sessunif} is necessary rather than convenient. It is also the least one can
exclude: for every other callee the supplied machines are waived by $\op{global}$ or refused at
line~\ref{ln:caller}, so no smaller condition would do and no larger one is needed.

The shape it excludes is not exotic, and one place the paper itself reaches for it should be said
plainly. Remark~\ref{rem:tight} offers, as a way of closing the hosting route to an
\emph{interface} member, that such a member admit only $\Zpid \cup \Apid$ --- the environment and the
adversary and no one else. A member of that shape is not global, $\Stdpid$ lying outside its $\admits$,
so it is not session-uniform. That is no defect of Definition~\ref{def:sessunif} but the artefact
showing through at the place it costs most: line~\ref{ln:session} lets the environment reach such a
member in session $0$ and in no other, so a protocol whose interface members admit exactly
$\Zpid \cup \Apid$ is drivable only in session $0$, whatever relocation has to say about it. The two
routes an admissible environment has into a protocol divide accordingly. It may claim a fresh name,
which line~\ref{ln:caller} admits only at a member admitting all of $\Stdpid$ --- global, hence shared
across sessions, hence not one run of one protocol. Or it may claim its own $Z$-identity, which
line~\ref{ln:caller} admits at a member admitting $\Zpid$ --- and then line~\ref{ln:session} pins that
member to session $0$. Neither route reaches a session-local protocol interface in a session of its
own choosing.

That the environment can claim its own $Z$-identity at all is worth recalling, since it is what makes
the second route available: the name $Z$ is used by no system, so $Z$-identities lie outside
$\ncl{\pi} \cup \ncl{\varphi}$ and Definition~\ref{def:admis} leaves them claimable. A member whose
$\admits$ is only its parent's name class is reachable by neither route, and there the artefact shows
through the adversary's slot alone.

Reading line~\ref{ln:session} as
\[
  \bigl(\id'.s = \id.s \ \vee\ \op{global}(\F) \ \vee\ \id'.F \in \{A,Z\}\bigr)
\]
would dispense with the condition. Session $0$ would then confer nothing, a local functionality
admitting the adversary would be reachable from the slot in every session alike, and
Definition~\ref{def:sessunif} could be dropped from every statement above: the results would hold of
all systems rather than the session-uniform ones. Nothing else in the chapter would change ---
condition (iii) stays, the reserved ids being named outright in code whatever the session rule is,
and the tuple conditions of Proposition~\ref{prop:orbit} are already free of session $0$. The
widening is
small and in the spirit of the rest: the two machines it names are global by $\op{global}$'s second
disjunct wherever $\op{global}$ is tested of them, the corruption gate of line~\ref{ln:corrgate}
still confines $\Adv$ to corrupt parties, and $\Zenv$ stays out by $\admits$ and by its own $\pars$.
What it would buy is more than relocation: it is what would let a protocol be local to its session and
driven by the environment at the same time, which is the pairing the two routes above deny and which
Remark~\ref{rem:tight}'s suggestion silently assumes. We keep the narrower rule, which is
Canetti's~\cite{Canetti20},
and record both the option and its price here.
\end{remark}

\begin{remark}[Relocation is not a many-instance theorem]\label{rem:manyinstances}
Proposition~\ref{prop:relocemul} moves one copy and charges nothing. It does not follow, and is not
true, that $n$ copies at $n$ indices cost $\varepsilon$ rather than $n\varepsilon$: the executions
of $\pi(\vec q_1),\dots,\pi(\vec q_n)$ and of $\varphi(\vec q_1),\dots,\varphi(\vec q_n)$ are not
each other's relabellings --- no relocation carries a system to $n$ copies of itself --- and the only
route between them is the hybrid chain of
Theorem~\ref{thm:parcomp}, which adds. What the two results give together is the statement one
wants from a single proof: prove $\pi(\vec q)$ emulates $\varphi(\vec q)$ at one $\vec q$;
Corollary~\ref{cor:onepid} gives it at every other within $\varepsilon$; and
Theorem~\ref{thm:parcomp}, whose disjointness hypotheses come to the $\vec q_k$ being pairwise
disjoint and clear of the surrounding shell, replaces all $n$ at once within $n\varepsilon$, at the
budgets Chapter~\ref{app:conccosts} charges. Relocation is also functorial in the constructions of
Section~\ref{sec:repl} and Chapter~\ref{sec:blockers}, each of which reads only names, ids and
identity sets: $\theta\bigl(\rho[\pi/\varphi]\bigr) = (\theta\rho)[\theta\pi/\theta\varphi]$ since an
injective action commutes with the set difference, $\theta\overline{\pi} = \overline{\theta\pi}$ ---
the right-hand side read, as Definition~\ref{def:blksys} requires, as the blocked system of the pair
$(\theta\pi,\theta\varphi)$ --- since a blocker copies the two fields $\theta$ relabels and
$I_{\theta\varphi} \setminus I_{\theta\pi} = \theta(I_\varphi \setminus I_\pi)$, and
$\theta\bigl(\Zenv^{\rho \setminus \varphi}\bigr) = (\theta\Zenv)^{\theta\rho \setminus \theta\varphi}$
since Definition~\ref{def:relocact} relabels a bundle's declared outward set, whose reserved half
$\{\id : \id.F \in \{A,C\}\}$ is fixed by (ii) of Proposition~\ref{prop:relocbasic}. So a relocated
proof of composition is the composition of the relocated proof.
\end{remark}




%% ==================================================================

%% ==================================================================
\chapter{هزینه‌های عینی}\label{app:conccosts}

Two threads of concrete cost sit outside the main line and are collected here: the cost of a
system in itself, read off its members, and the cost the composition results of
Section~\ref{sec:concrete} charge, arithmetic on the budgets of Definition~\ref{def:concadv}.

Budgeted member by member, a system also has a cost as a whole, and it is a functionality's cost
read one level up.

\begin{definition}[Cost of a system]\label{def:syscost}
An \emph{invocation} of a system $\pi$ is a call reaching one of its members from outside $\pi$;
$\pi$ services it by a cascade of internal activations before control returns outside. The
\emph{cost} of $\pi$ is the pair $(T_\pi, R_\pi)$: the \emph{per-invocation time} $T_\pi$, the
largest total number of steps in one such cascade --- summed over the activations it triggers,
one machine running at a time --- and the \emph{invocation count} $R_\pi$, the number of
invocations $\pi$ answers. So $\pi$ runs in at most $T_\pi R_\pi$ steps.
\end{definition}

\begin{proposition}[System cost from member costs]\label{prop:syscost}
Let $\pi = \{\F_1,\dots,\F_n\}$ with each $\F_i$ of cost $(t_i, r_i)$, and let $k_i$ bound the
activations of $\F_i$ within one invocation of $\pi$. Then
\[
  T_\pi \;\le\; \sum_{i} k_i\, t_i , \qquad R_\pi \;\le\; \sum_{i} r_i ,
\]
and $\pi$ runs in at most $\sum_i t_i r_i$ steps. If every invocation activates each member at most
once --- $\pi$ subroutine-respecting, $k_i \le 1$ --- then $T_\pi \le \sum_i t_i$; if moreover each
invocation drives all members, $R_\pi \le \min_i r_i$.
\end{proposition}

\begin{proof}
An invocation activates $\F_i$ at most $k_i$ times, each for at most $t_i$ steps, so it spends at
most $\sum_i k_i t_i$. The activations of $\F_i$ over the whole execution number at most $r_i$, and
every invocation is at least one activation, so the invocations number at most $\sum_i r_i$; the
total steps are $\sum_i (\text{activations of }\F_i)\, t_i \le \sum_i r_i t_i$. When each
invocation drives all members, $R_\pi$ invocations activate each $\F_i$ exactly $R_\pi$ times, so
$R_\pi \le r_i$ for every $i$.
\end{proof}

\begin{corollary}[Cost of a union]\label{cor:syscost}
Let $\pi_1$ and $\pi_2$ be systems with disjoint process ids and costs $(T_1,R_1)$ and
$(T_2,R_2)$. If no member of one invokes a member of the other, then
\[
  \begin{aligned}
  \bigl(T_{\pi_1 \cup \pi_2},\ R_{\pi_1 \cup \pi_2}\bigr) &\;=\; (T_1,R_1) \oplus (T_2,R_2) \\
  &\;=\; \bigl(\max(T_1,T_2),\ R_1 + R_2\bigr) .
  \end{aligned}
\]
If instead each invocation of $\pi_1$ makes at most $k$ calls into $\pi_2$ and the cross-call graph
is acyclic, then $T_{\pi_1 \cup \pi_2} \le T_1 + k\,T_2$ from a $\pi_1$ entry, and symmetrically,
with $R_{\pi_1 \cup \pi_2} \le R_1 + R_2$; with the cross-calls unbounded the summary costs no
longer determine the union's.
\end{corollary}

\begin{proof}
With no cross-calls a call from outside $\pi_1 \cup \pi_2$ enters $\pi_1$ or $\pi_2$ and its cascade
stays there, costing at most $T_1$ or $T_2$ and belonging to that subsystem's invocation count; so
the per-invocation time is $\max(T_1,T_2)$ and the counts add. A cross-call from $\pi_1$ into
$\pi_2$ is internal to the union and extends the cascade by one invocation of $\pi_2$, at most $T_2$
steps; at most $k$ of them, acyclically, give $T_1 + k T_2$, while a call once external to $\pi_2$
but placed from $\pi_1$ is no longer counted, so the invocations still number at most $R_1 + R_2$.
\end{proof}

The subroutine-respecting proviso is not idle: without it one call can be driven around the system
without end, and $T_\pi$ degenerates towards the whole of $\sum_i t_i r_i$ --- the same discipline
that keeps a private member's callers pinned in Remark~\ref{rem:tight}. This system cost is dual to
the $\oplus$ of Definition~\ref{def:res}: hosting a system as a bundle, one member per activation,
combines its members by $\oplus$ into $(\max_i t_i,\, \sum_i r_i)$, the cost
$\bd{c}_{\rho \setminus \varphi}$ pays in the composition costs below; calling it as a subsystem
sums the times over a cascade and counts external calls, $(\sum_i t_i,\, \min_i r_i)$ in the
synchronous case. Both bound the same total $\sum_i t_i r_i$, the second more tightly when every
invocation is a full pass. And $\oplus$ recurs at the system level itself: by
Corollary~\ref{cor:syscost} non-interacting subsystems compose by it, one operator at both scales.

When the members of $\pi$ call one another the per-invocation time turns on the shape of that
calling, and one case is clean.

\begin{proposition}[DAG systems]\label{prop:dagcost}
Suppose the call graph of $\pi = \{\F_1,\dots,\F_n\}$ is acyclic and every invocation activates
each internal member $\F_i$ at most $\max_{j:\,\F_j \to \F_i} t_j$ times --- the largest
per-invocation time among the members that call it. Then
\[
  T_\pi \;\le\; \sum_i \Bigl(\max_{j:\,\F_j \to \F_i} t_j\Bigr) t_i \;\le\; \Bigl(\max_i t_i\Bigr) \sum_i t_i ,
\]
a per-invocation time independent of the depth of the graph.
\end{proposition}

\begin{proof}
An invocation entering a source activates each $\F_i$ some $a_i$ times and spends $\sum_i a_i t_i$
steps. A source runs once and each internal $\F_i$ at most $\max_{j \to i} t_j$ times by hypothesis,
so $a_i \le \max_i t_i$ throughout, whence
$T_\pi \le \sum_i (\max_{j\to i} t_j)\, t_i \le (\max_i t_i) \sum_i t_i$.
\end{proof}

The hypothesis earns its place, for acyclicity alone does not bound $T_\pi$. Writing $n_{ji}$ for
the calls $\F_j$ makes to $\F_i$ in one activation, the per-invocation activation counts solve
$a_i = \sum_{j \to i} n_{ji}\, a_j$, so $a_i$ counts the weighted source-to-$\F_i$ paths, and a
reconvergent graph multiplies them: a chain of $N$ diamonds ---
$s \to \{A_k, B_k\} \to C_k \to \{A_{k+1}, B_{k+1}\} \to \dotsb$ --- doubles the count at each
level, so $a$ at the foot is $2^N$ and the per-invocation time is exponential in the depth. Capping
each node by the \emph{maximum} over its callers rather than by their sum is exactly what stops the
paths accumulating --- the $\oplus$-over-summation theme once more, and the DAG counterpart of the
subroutine-respecting case $k_i \le 1$ of Proposition~\ref{prop:syscost}, one activation relaxed to
$\max_{j\to i} t_j$. The invocation caps bound $a_i \le r_i$ in any case; the point here is a bound
in the times alone, independent of the caps and of the depth.

That last remark is the whole of the general statement: the caps alone already close the cost
class, with no condition on the call structure.

\begin{corollary}[Cost composition]\label{cor:costcomp}
A system $\pi = \{\F_1,\dots,\F_n\}$ whose members have costs $(t_i, r_i)$ runs in at most
$\sum_i t_i r_i$ steps over at most $\sum_i r_i$ invocations, so $T_\pi \le \sum_i t_i r_i$ and
$R_\pi \le \sum_i r_i$ whatever its call structure. Hence if each $t_i$ and $r_i$ is polynomial in
a parameter $\lambda$ and $n$ is polynomially bounded, $(T_\pi, R_\pi)$ is polynomial: the
polynomially bounded systems are closed under formation.
\end{corollary}

\begin{proof}
The step and invocation bounds are Proposition~\ref{prop:syscost}, and one invocation's cascade is
part of the whole execution, so $T_\pi \le \sum_i t_i r_i$ too. A sum of polynomially many
polynomials is a polynomial.
\end{proof}

No condition on the call graph is needed here, unlike the depth-free \emph{time} bound of
Proposition~\ref{prop:dagcost}, because the caps carry it. A reconvergent cascade that would
activate a member $2^N$ times can do so only if that member answers $2^N$ calls, so $r_i \ge 2^N$
--- a member already outside the polynomial class; within the class the cap refuses the excess and
truncates the cascade at $\sum_i r_i$ activations. This is the efficiency counterpart of
Theorem~\ref{thm:comp}: where composition preserves \emph{security} --- $\rho[\pi/\varphi]$
emulates $\rho$ once $\pi$ emulates $\varphi$ --- cost composition preserves \emph{efficiency}, so
a protocol assembled from polynomial-time functionalities is itself polynomial-time. Together the
two are what the concrete apparatus is for: polynomial-time UC-secure protocols compose into
polynomial-time UC-secure protocols, which is what lets the asymptotic reading below be a corollary
rather than a theory of its own.

The same closure survives the operation the composition theorem runs.

\begin{proposition}[Replacement preserves polynomial boundedness]\label{prop:replpoly}
Let $\rho$ and $\pi$ be systems whose members are polynomial and polynomially many, and let the
replacement of $\varphi$ by $\pi$ in $\rho$ be well-formed. Then $\rho[\pi/\varphi]$ is
polynomially bounded.
\end{proposition}

\begin{proof}
The members of $\rho[\pi/\varphi] = (\rho \setminus \varphi) \cup \pi$ (Definition~\ref{def:repl})
are those of $\rho$ outside $\varphi$, each polynomial as a member of $\rho$, together with those
of $\pi$, again polynomial; they number at most $\lvert\rho\rvert + \lvert\pi\rvert$, polynomially many.
Corollary~\ref{cor:costcomp} applies.
\end{proof}

So the system that runs after a replacement is efficient whenever the specification $\rho$ and the
realization $\pi$ are --- Theorem~\ref{thm:comp} read for cost rather than security, and on the
same operation.

\heading{Composition, numerically} In Theorem~\ref{thm:comp} the adversaries and the simulators
carry over unchanged, so their budgets do. The environments are the one place something is
computed, and what absorption costs is the cost of hosting. Systems are budgeted member by member, so the cost can be
looked up: write $\bd{c}_{\rho \setminus \varphi}$ for the $\oplus$-combination of the budgets of
the members of $\rho \setminus \varphi$ --- the largest of their per-invocation times, the sum of
their invocation counts, the hosted-bundle cost dual to the called-subsystem cost of
Proposition~\ref{prop:syscost} --- which bounds them in any execution, whoever is driving them. The
absorbed environment of Definition~\ref{def:absorb} runs
$\Zenv$ and those functionalities and nothing else, so
\[
  \Zenv \in \ZenvSet_{\rho[\pi/\varphi],\,\rho}(\bd{z})
  \qquad\Longrightarrow\qquad
  \Zenv^{\rho \setminus \varphi} \in \ZenvSet_{\pi,\varphi}\bigl(\bd{z} \oplus \bd{c}_{\rho \setminus \varphi}\bigr) ,
\]
the containment in $\ZenvSet_{\pi,\varphi}$ being Proposition~\ref{prop:absorb}. Feeding this
into the theorem, the conclusion is
\[
  \advtg^{\op{uc}}_{\rho[\pi/\varphi],\,\rho}[\bd{a},\bd{s},\bd{z}]
  \;\leq\;
  \advtg^{\op{uc}}_{\pi,\varphi}\bigl[\bd{a},\,\bd{s},\,\bd{z} \oplus \bd{c}_{\rho \setminus \varphi}\bigr] .
\]
Composition is free in the adversary and in the simulator, and costs the surrounding system
once, on the environment side.

\heading{Transitivity, numerically} Proposition~\ref{prop:trans} asks that the first step's
simulator class be the second step's adversary class, which becomes an inequality between
budgets: $\SimSet(\bd{s}) \subseteq \Aset(\bd{a}')$ holds as soon as $\bd{s} \leq \bd{a}'$. A
simulator one is willing to build must be an adversary the next step is willing to face. The
proposition also asks for one class admissible for both steps, and that hypothesis must be
carried here: for $\ZenvSet \subseteq \ZenvSet_{\pi,\varphi} \cap \ZenvSet_{\varphi,\psi}$, which
by the proof of Proposition~\ref{prop:trans} lies in $\ZenvSet_{\pi,\psi}$ as well,
\[
  \advtg^{\op{uc}}_{\pi,\psi}[\bd{a},\bd{s}',\ZenvSet]
  \;\leq\;
  \advtg^{\op{uc}}_{\pi,\varphi}[\bd{a},\bd{s},\ZenvSet]
  \;+\;
  \advtg^{\op{uc}}_{\varphi,\psi}[\bd{s},\bd{s}',\ZenvSet] .
\]
Budgets alone will not do here. Written with $\bd{z}$ in all three places the brackets would
range over $\ZenvSet_{\pi,\psi}(\bd{z})$, $\ZenvSet_{\pi,\varphi}(\bd{z})$ and
$\ZenvSet_{\varphi,\psi}(\bd{z})$, and the first is contained in neither of the others --- an
environment admissible for the outer pair is told nothing about $\IDs(\varphi)$ --- so no
hypothesis would cover the middle probability that the triangle inequality cancels. The
composition display above has no such trouble: absorption maps its one class into the other.
Corollary~\ref{cor:comp} is the two displays put together. Along the chain of
Theorem~\ref{thm:parcomp} the same two effects accumulate: the errors add, each hop widens the
environment budget by the cost of hosting what it absorbs at that step, and the class relations
of Remark~\ref{rem:parclasses} become the inequalities $\bd{s}_{k+1} \leq \bd{a}_k$.

\heading{Nesting, and whether the simulator blows up} Say that an emulation statement has
\emph{overhead} $f$ if it holds against $\bigl(\Aset(\bd{a}),\SimSet(f(\bd{a})),\cdot\bigr)$ for
every adversary budget $\bd{a}$: whatever the adversary is given, the simulator one builds costs
$f$ of it. In the per-invocation budget the two components behave differently,
\[
  f(t,r) \;=\; \bigl(\max(t,\tau),\ \ \gamma r + \delta\bigr) ,
\]
the query count affine --- $\gamma$ the calls the simulator places for each call of the adversary
it runs, $\delta$ its own --- and the per-invocation time a maximum, $\tau$ the simulator's own
per-activation work, since one machine runs per activation and the simulator adds only bounded
work to each.

Theorem~\ref{thm:comp} does not touch $f$ at all. Its adversary and simulator classes are those
of the hypothesis, so a single replacement returns the very machine the hypothesis supplied and
the whole price is paid on the environment side. Overheads compose in one place only: where a
simulator is re-read as an \emph{adversary}, which is Proposition~\ref{prop:trans}. There $\Sim'$
is built from $\Sim$, itself built from $\Adv$, so the composite has overhead $f_2 \circ f_1$.
Applying the results $k$ times, the query budget obeys $r_j = \gamma\, r_{j-1} + \delta$, that is
\[
  r_k \;=\; \gamma^{k} r_0 \;+\; \delta\,\frac{\gamma^{k}-1}{\gamma-1}
  \quad (\gamma \neq 1),
  \qquad
  r_k \;=\; r_0 + k\delta \quad (\gamma = 1) ,
\]
while the per-invocation time saturates, $t_k = \max(t_0, \tau_1, \dots, \tau_k)$, bounded by the
largest level's own work and constant in depth.

The blow-up therefore lives entirely in the query count, a dichotomy turning on $\gamma$ alone. A
simulator that must place more than one call for some call of the adversary it runs has
$\gamma > 1$, and $k$ applications cost $\gamma^k$: genuinely exponential, and at $\gamma = 2$ a
depth of $20$ already costs a factor of $2^{20}$, about a million, on the query count. A simulator
that relays each call once and adds bounded traffic of its own has $\gamma = 1$, and $k$
applications cost $r_0 + k\delta$: linear, $\delta$ paid once per level. This is the property to
check when a simulator is built; nothing proved here implies it. It is also why straight-line
simulators matter beyond the usual reasons, rewinding being the standard way to acquire
$\gamma > 1$: one rewinding simulator anywhere in the tower puts a factor $\gamma$ on every level
below it.

The same formula governs breadth as well as depth. Theorem~\ref{thm:parcomp} builds
$\Sim_n,\dots,\Sim_1$ by handing each simulator to the next step as its adversary, so $n$
parallel instances compose $n$ overheads exactly as an $n$-deep tower does, and
Remark~\ref{rem:parclasses} read with budgets asks $\bd{a}_k \geq f_{k+1}(\bd{a}_{k+1})$ at every
link. Composing $n$ instances of a $\gamma = 1$ protocol costs $n\delta$ extra queries; composing
$n$ instances of a $\gamma = 2$ one is hopeless for even moderate $n$.

Nothing comparable happens on the environment side, or in the per-invocation time. Each level adds
the cost of hosting the shell it absorbs, and a tower's shells are disjoint, so a depth-$k$ nesting
reaches the innermost statement at budget $\bd{z} \oplus \bigoplus_{i \leq k} \bd{c}_i$ --- its
invocation count that of running the ambient system once, its per-invocation time the largest
shell's, however deep the tower. The errors behave likewise, $k$ applications of
Proposition~\ref{prop:trans} giving $\sum_i \varepsilon_i$. Where a security claim degrades
exponentially under nesting, the simulator's query overhead $\gamma > 1$ is the only possible
source: one machine runs per activation, so neither the per-invocation time nor the environment
budget can multiply.

\heading{Recovering the asymptotic reading} Nothing above mentions a security parameter. To get
the familiar statement, index everything by $\lambda$, let the budgets be polynomials in
$\lambda$, and ask that the advantage be negligible in it; the classes then become the
polynomially bounded machines and Definition~\ref{def:uc} the usual definition of UC emulation.
The concrete form is the one the composition results are proved in, and the asymptotic form is a
corollary of it, not the other way round --- the ordering practice-oriented provable security has
long argued for~\cite{Bellare97}.

%% ==================================================================
\chapter{فراخوان‌های پاسخگو}\label{app:responsive}

Chapter~\ref{sec:fsig} needs a promise the framework has so far had no way to write down. When
$\op{Sign}$ asks the adversary for a signature it must get one back, and get it back \emph{now}:
were the adversary free to go off and drive the rest of the execution before answering, the party
would be left suspended mid-signature and locality lost --- the same immediacy synchronous UC asks
of a responsive environment~\cite{KMTZ13}. The text there settles for restricting
the simulator class. This chapter says the thing in code instead.

\heading{What has to be forbidden} Not that the adversary answer helpfully --- it may return
anything, and Chapter~\ref{sec:sanitization} deals with a bad answer. What has to be forbidden is
the \emph{token} leaving: between receiving the call and returning, the responder must place no
call at all. That is a restriction on the calls a machine makes, and the paper already has a
wrapper for exactly that.

\begin{definition}[Responsive answering]\label{def:respond}
$\opl{Respond}$ is $\opl{Silence}$ at its maximum: with $\Ids_{\op{all}}$ the set of all
identities,
\[
  \begin{gathered}
  \opl{Respond}\bigl(\id.\Op(in) \text{ from } \id'\bigr) \\[3pt]
  \;:=\;
  \opl{Silence}\bigl(\Ids_{\op{all}},\ \id.\Op(in) \text{ from } \id'\bigr) .
  \end{gathered}
\]
\end{definition}

Every claim lies in $\Ids_{\op{all}}$, so by Chapter~\ref{sec:wrappers} every call the wrapped
programme issues is refused with $\rej$ \emph{without ever being placed}. The programme may compute
as long as it likes and must still return a value; what it cannot do is hand the token on.

\begin{procedure}{Wrapper $\opdef{Respond}$ \\
\params{programme $\id.\Op(in)$ from $\id'$}}
\opsig{$\opl{Respond}\bigl(\id.\Op(in) \text{ from } \id' \bigr)$}:
\settowidth{\clausewd}{$\id''.\fopl{Op}(in'')\ \text{as}\ \id''' \,;\ k$}
\settowidth{\reswd}{$k\bigl(\id''.\fopl{Op}(in'')\ \text{as}\ \id'''\bigr)$}
\begin{algorithmic}[1]
\State \textbf{handle}\ \ $\id.\Op(in)$ from $\id'$\ \ \textbf{with}
\State \quad \makebox[\clausewd][l]{$\textbf{return}\ out$}
       $\ \mapsto\ $ \makebox[\reswd][l]{$out$}
\State \quad \makebox[\clausewd][l]{$\id''.\fopl{Op}(in'')\ \text{as}\ \id''' \,;\ k$}
       $\ \mapsto\ $ \makebox[\reswd][l]{$k(\rej)$}
\end{algorithmic}
\end{procedure}

\heading{Read as a handler} Written out, $\opl{Respond}$ is $\opl{Silence}$ with its last clause
deleted. The wrapper $\opl{Silence}$ forwards a call whose claim it permits, so its operation clause is partly
a re-raise; $\opl{Respond}$ has no re-raise, only a constant. In the vocabulary of
Chapter~\ref{sec:wrappers} it \emph{discharges} the ambient call effect rather than forwarding it:
inside its scope the effect is handled locally and completely, so the wrapped programme is pure
with respect to it. It is also the definition of \emph{locality}, the property
Chapter~\ref{sec:fsig} wanted: a call on the slot is served within the responder's own
activation, the token never leaving it. That is the crisp statement of responsiveness --- \emph{a responsive core is
one whose call effect does not escape to $\opl{Exec}$} --- and the operational reading follows,
since $\opl{Exec}$ dispatches exactly the calls that do escape. A responsive call is one
activation of the responder and no more.

\begin{definition}[Responsive interface, responsive call]\label{def:respcall}
The adversary carries, beside $\fopl{A}$, a second interface $\fopl{A}^{!}$ with the same guard and
the same core but $\opl{Respond}$ in place of $\opl{Silence}$:
\[
  \begin{gathered}
  \id.\fopl{A}^{!}(in) \text{ from } \id'
  :\quad
  \textbf{require}\ \opl{Guard}_{\Adv}(\id,\id') ;\quad \\
  \textbf{return}\ \opl{Respond}\bigl(\id.\op{A}(in) \text{ from } \id'\bigr) .
  \end{gathered}
\]
A \emph{responsive call} is a call on $\fopl{A}^{!}$. Inside an interface with identity $\id$ we
write $\Adv^{!}(in)$ for $(A,\id.P).\fopl{A}^{!}(in)$ as $\id$, as $\Adv(in)$ abbreviates the
ordinary one.
\end{definition}

Three things follow, and they are the price.

The responder cannot read the register. The set $\Cs$ is reached by calling $\fopl{Status}$, and that call
is refused, so a responsive answer is computed from what the responder already knows. It cannot
delegate: no subroutine, no other party, no environment. And it cannot corrupt, $\op{Corrupt}$
being a call like any other. The last deserves stating on its own, because
Section~\ref{sec:fullinterface} reads the register twice and the gap between the two reads is
exactly where a corruption can slip in.

\begin{proposition}[A responsive call cannot corrupt]\label{prop:respcorr}
In any execution, the register's state $\Cs$ when a responsive call returns is its state when
the call was placed.
\end{proposition}

\begin{proof}
By Section~\ref{sec:corruption}, $\Cs$ changes only on line~\ref{ln:corrwrite} of $\fopl{Corrupt}$,
and only when that interface is reached by a call. Between the placing and the return of a
responsive call the responder holds the token throughout, and by Definition~\ref{def:respond}
every call it issues is refused before being placed, so $\fopl{Corrupt}$ is not reached. No other
machine holds the token in that interval, so none can reach it either.
\end{proof}

So if a core reaches the adversary only through $\Adv^{!}$, the party it serves is corrupt at
line~\ref{ln:medout} exactly when it was corrupt at line~\ref{ln:corrgate}, and the second
mediation hook cannot fire on account of that call. Remark~\ref{rem:midcall} describes what happens
when it does; responsiveness is the condition under which it does not.

This is what Chapter~\ref{sec:fsig} wanted. Writing $\Adv^{!}$ in place of $\Adv$ in $\op{Gen}$,
$\op{Sign}$ and $\op{Verify}$ makes locality a property of the code rather than a condition on the
simulator class, and the wrapper enforces it, so it holds of \emph{every} occupant of the
adversary slot --- the simulator included, Definition~\ref{def:sim} giving a simulator the
adversary's interfaces, $\fopl{A}^{!}$ among them. What it does not do is make the answer good: a responsive adversary may still return
$\none$, and $\San$ repairs that. Responsiveness buys the token back; sanitization buys the
value.

\begin{remark}[Responsiveness and the dummy]\label{rem:respdummy}
The two chapters pull against each other, and it is worth being explicit about where. The dummy
of Section~\ref{app:dummy} answers a call from a functionality by relaying it to the environment
on line~\ref{ln:dumin} --- and that is a call, which $\opl{Respond}$ refuses. So $\Dum$ cannot
serve a responsive call: the relay is refused, and the rule of Section~\ref{sec:execution}
delivers its answer as $\none$ where a real adversary would answer with substance. The regrouping
(R1) would therefore fail for a system whose cores use
$\Adv^{!}$, at the first responsive call --- which is why Theorem~\ref{thm:dummy} hypothesises
systems whose cores place none. Recovering completeness for responsive systems needs a matching
notion on the environment's
side, so that the relayed question can be answered without the token moving on; we leave that
open.
\end{remark}

%% ==================================================================
%% The one chapter that is not itself a functionality: it names the
%% language every core from here to the end of the worked examples is
%% written in, in the tradition of Bellare-Rogaway's L (their Appendix
%% B), so that "this code is safe to compile" is a checkable claim
%% rather than a hope.
%% ==================================================================
\chapter{Core Language}\label{app:codelang}

Chapters~\ref{sec:frand} to~\ref{sec:sigdeep} give the code of this book's functionalities,
protocols and games, and that code is meant to be read as a program, not as an illustration of
one --- a promise that has to be kept by a stated convention rather than by the reader's good
will. \cite{BR06} keep exactly this promise for game-based proofs with a small, strongly typed,
terminating language they call $\mathcal{L}$, given by a context-free grammar (their Figure~15)
and a short named list of the sugars they allow on top of it. This chapter takes the same two
steps here: a closed grammar for what every chapter from here on calls a \emph{core} --- an
operation's body, already this book's own word for it --- and this book's own short list of
departures from that grammar, each named once so that every core still to come can be checked
against it rather than taken on faith.

Everything before this chapter is exempt, and deliberately so. $\opl{Guard}$, $\opl{Silence}$ and
$\opl{Mediate}$ (Chapter~\ref{sec:wrappers}), $\opl{Exec}$ (Section~\ref{sec:execution}), the
dummy adversary (Section~\ref{app:dummy}) and $\opl{Respond}$ (Chapter~\ref{app:responsive}) do
not run inside a core: they are what runs a \emph{system} of cores, dispatching calls, evaluating
guards, handing the token from one activation to the next. A definition of execution is not
itself a program in the language it executes --- exactly as \cite{BR06}'s own account of a game
meeting an adversary sits in their informal Section~3, outside the grammar of their Figure~15 ---
and the one notation that could be mistaken for a core's own, the continuation form
$\id''.\fopl{Op}(in'')\ \text{as}\ \id'''\,;\,k$ carrying an explicit, unbound resumption $k$,
belongs entirely to that outer layer: it never once occurs inside an operation's own code from
Chapter~\ref{sec:frand} on. A plainer relative of it does, constantly --- an ordinary call to
another functionality's operation, resumed under the caller's own identity rather than handed a
fresh continuation --- and that one is a core-language matter in its own right, below.

\heading{The core grammar} A core is a sequence of statements, and every operation this book
writes is exactly one core. Scoped to what this book's cores actually use, rather than to
$\mathcal{L}$'s full grammar:

\begin{center}
\footnotesize
\[
\begin{array}{r@{\ }c@{\ }l}
\mathit{core} & ::= & \varepsilon \ \mid\ \mathit{stmt} \,;\, \mathit{core} \\[2pt]
\mathit{stmt} & ::= & \mathit{lv} \gets e
                  \ \mid\ \mathit{lv} \gets_{\$} S \\
              &     & \mid\ \ \textbf{if}\ e\ \textbf{then}\ \mathit{core}\ [\,\textbf{else}\ \mathit{core}\,] \\
              &     & \mid\ \ \textbf{for}\ x \in S\ \textbf{do}\ \mathit{core} \\
              &     & \mid\ \ \textbf{return}\ e \\[2pt]
\mathit{lv}   & ::= & \V{x} \ \mid\ \V{x}[e] \\[2pt]
e             & ::= & \dots \ \mid\ \mathit{call} \\[2pt]
\mathit{call} & ::= & \id''.\Op(e_1,\ldots,e_k)\ \textbf{as}\ \id
\end{array}
\]
\figcaption{fig:codelang}{The core grammar, scoped to this book's own usage. $\mathcal{L}$'s
grammar (\cite{BR06}, their Figure~15) with the sub-grammar of expressions left to ordinary
mathematical notation, since every expression this book writes is already standard over the six
types below; $\mathit{call}$ is $\mathcal{L}$'s own, generalized past their Adversary-only
restriction (below). (Chapter~\ref{app:codelang}.)}
\end{center}

Six base types, not $\mathcal{L}$'s five: integer, boolean, string, set, table --- $\mathcal{L}$'s
array, kept as a function's type throughout, $\V{L} : \mathbb{N} \to \{0,1\}^{n} \cup \{\unset\}$
reading exactly as $A : \{0,1\}^{*} \to \{0,1\}^{*} \cup \{\textsf{undefined}\}$ does there,
$\unset$ this book's $\textsf{undefined}$ --- and identity, this book's own addition, below. Only
one $\textbf{for}$-form is attested, $\textbf{for}\ x \in S\ \textbf{do}$, and the one instance of
it binds a pair directly, $\textbf{for}\ (m,y) \in N\ \textbf{do}$ (Chapter~\ref{sec:realac},
lines~\ref{ln:rloop} and~\ref{ln:simloop}) --- an obvious sugar for binding one variable and
destructuring it. The bounded integer-range form $\mathcal{L}$ also allows is in reserve, unused
so far, and licensed the moment a core needs it.

\heading{The caller is implicit} Every operation's core is headed $\id.\Op(in)$
\textbf{from}~$\id'$, and $\id'$ --- the caller's identity --- is readable anywhere in the core
though it names no argument of $\Op$. A plain oracle in $\mathcal{L}$ carries nothing like it; this
book's cores need it because so many begin by asking exactly who is calling before anything else
runs, $\Gclock$'s operations being representative (the guard at line~\ref{ln:clockaz} tests
$\id'.F$ and $\id.P$'s corruption before the core does anything else).

\heading{A core may call another core} $\mathcal{L}$ reserves the $\mathit{call}$ expression for
the Adversary alone; an oracle inside one of their games may not place one. Every functionality,
protocol and game in this book does --- the protocol of Chapter~\ref{sec:realac} reads $\Fpki$
and $\Fsig$ mid-operation (lines~\ref{ln:rverify} and neighbours), and $\Frand$'s own
unpredictability game calls back into $\Frand$ itself (line~\ref{ln:grnd}) --- because this book's
functionalities and protocols are built from others by declaration, $\uses$ naming exactly what a
core may call (Chapter~\ref{sec:functionalities}), and a game's shell is a functionality like any
other. The call itself is $\mathcal{L}$'s own, unrelaxed: $\id''.\Op(e_1,\ldots)\ \textbf{as}\ \id$
evaluates like any other expression and resumes the calling core at the same line, under its own
identity. What stays strictly outside a core is the \emph{general} form, with an unbound
continuation $k$ standing for an arbitrary resumption point rather than ``the next line of this
same core'' --- that is the execution model's own notation for what a call means in general,
never a core's, and Section~\ref{sec:execution} is where it is defined.

\heading{Require is sugar} A core's own $\textbf{require}\ e$ line, used throughout the
execution-model chapters and inside plenty of the cores to come, is sugar for
$\textbf{if not}\ e\ \textbf{then}\ \textbf{return}\ \none$ --- not literally $\rej$, even though a
failed guard elsewhere returns exactly that. Section~\ref{sec:execution} makes $\rej$ the one
value no core may return as its own output, laundering a core-level $\rej$ into $\none$ before
delivery; a core's own $\textbf{require}$ merely names the common case compactly. $\Gclock$'s operations spell the same
guard out by hand, with an ordinary $\textbf{if}\dots\textbf{then}\ \textbf{return}$
(line~\ref{ln:clockaz} again): both are the one core-language statement, sugared or not.

\heading{Table indices need no new rule} $\V{reg}[\id.P]$, $\V{L}[\V{ctr}]$: neither is
string-indexed, and $\mathcal{L}$ already has the answer, its own third enhancement
(\cite{BR06}) --- a non-string index is passed through $\mathrm{encode}(\cdot)$ first, silently.
Every table this book writes that is not string-indexed already reads this way; there is nothing
to add here beyond citing them for it.

\heading{Identity is a sixth type} One departure is not sugar. Definition~\ref{def:func} makes an
identity a pair $\id = (\PID,P)$ of a process id and a party, a process id itself a triple
$(F,s,i)$ --- and a core reads through both levels at once, $\id.F$, $\id.s$, $\id.P$ named
projections, $\id.i$ the one field no core ever reads (Chapter~\ref{sec:functionalities} says so
outright). None of $\mathcal{L}$'s five types carries a projection, and $\mathrm{encode}(\cdot)$
gives only equality-preservation, by \cite{BR06}'s own account of it --- an identity is read
apart far too often for that to serve. So identity joins the type list as a sixth base type,
structured rather than atomic, with projection and componentwise equality as its only operations
--- the one place this book's cores ask $\mathcal{L}$ for something it does not already have.

\heading{Reading the code} The grammar above is what a core may compute; how it is set on the
page is a separate promise, and one already mostly kept. Keywords are bold, algorithmicx's
default: \textbf{if}, \textbf{then}, \textbf{else}, \textbf{for}, \textbf{do}, \textbf{return},
and \textbf{require}. An operation's name carries a link wherever it is declared and wherever it
is called back to, both rendered in \textcolor{accentsoft}{accentsoft} small caps --- one macro
anchors the declaration, two more point back to it from every call site. State lives in $\V{x}$,
typewriter, reserved for a core's
own mutable fields; $\id$, $in$, $out$ stay italic, bound parameters rather than state, and the
two must not be run together --- a future reader deciding what a core's own record type looks
like needs exactly this distinction and no other. Comments print in
\textcolor{commenttint}{commenttint}, italic, off to the side, never load-bearing. That
copy-editing pass over the chapters already written is done, not new notation but a uniformity
sweep: a first prose mention of an already-declared operation is linked back to its declaration
throughout, and the constant atoms --- $\ok$, $\rej$, $\true$, $\false$, $\done$, together with
every trace tag and storage key elsewhere in this book --- carry
\textcolor{commenttint}{commenttint} rather than an operation name's own colour, so bare colour
now does separate ``this is a name'' from ``this is a value.''

What the grammar and the departures above buy, together, is a single test: from
Chapter~\ref{sec:frand} on, a line of this book's math is a core's own code exactly when it
type-checks against the table above plus the departures named here, and it is a specification ---
a game's verdict, a property's quantifier, $\San$'s own $\mathrm{pick}$ --- exactly when it does
not. Closing the gap to an actual language, OCaml or otherwise, is then a matter of compiling six
types and a handful of statement forms once, not of first deciding, chapter by chapter, what a
given display of symbols was ever claiming to compute.

%% ==================================================================
%% ==================================================================
\chapter{همسان تا خرابی}\label{app:iub}

Code-based proofs advance by \emph{hops}~\cite{BR06,Shoup04}: two games that are the same program
except where a flag has been raised, and a bound on the probability of the raising. The
fundamental lemma of game playing is the licence --- if the games are identical until $\V{bad}$,
their outcomes differ with probability at most $\Pr[\Bad]$ --- and its proof is a coupling: run
both games on the same coins and they cannot part company before the flag goes up. In this paper a
game is not one program. An execution runs an environment, a slot occupant, a register and a
system; the token wanders among them; the occupant may be handed the token mid-operation and
drive the rest of the execution before answering --- the freedom
Chapter~\ref{app:responsive} exists to discipline --- and a hop must survive all of it. This
chapter replays the lemma over executions, and the finding is that the wandering costs nothing:
the lemma below has no hypothesis at all.

\heading{Where the flag lives} $\V{bad}$ is an ordinary state variable of an ordinary machine:
declared $\V{bad} \gets \false$ in $\op{Initialize}$, never assigned any value but
$\true$ afterwards, owned like the rest of a machine's state by the machine that declares
it. The framework already keeps every other hand off it --- the one thing a machine can do to
state it does not own is provoke calls that read or extend it, and at a corrupt party
$\opl{Mediate}$ replaces \emph{answers}, never state (Remark~\ref{rem:gmcorr}) --- so a flag is
raised by its owner's own code or not at all. Whether anyone may \emph{read} it is immaterial,
through an interface or through $\op{Leak}$, which bypasses $\opl{Guard}$: until the flag is
raised the two machines below answer any read alike, and after it is raised the lemma has
already paid. Several machines may carry flags at once --- members sharing a name included, each
owning its own --- and $\Bad$ below is the event that \emph{some} machine raises \emph{some}
flag.

\begin{definition}[Identical until bad]\label{def:iub}
Standard functionalities $\F$ and $\F'$ are \emph{identical until bad} if their five fields
agree; their $\op{Initialize}$s agree line for line, declarations common to both sides and
$\V{bad} \gets \false$ among them; and their remaining code agrees except inside
\emph{bad-guarded} blocks, of two shapes: the continuation of an assignment
$\V{bad} \gets \true$, from just past the assignment to the end of its \emph{enclosing
block} --- the branch, loop or operation body the assignment stands in, so that everything in
the continuation is reached only across the assignment --- and the body of a conditional whose
test is $\V{bad}$. Erasing the bad-guarded blocks from both
sides leaves one functionality, operation for operation and line for line, $\op{Leak}$ included;
in particular the assignments themselves are common code, each heading its own divergent
continuation, and neither side assigns $\V{bad}$ any value but $\true$. Equal
functionalities are identical until bad, with nothing to erase. Two occupants of the adversary
slot are identical until bad when their $\op{Initialize}$s and cores are, read the same way. Systems $\pi$ and $\pi'$
are identical until bad if some bijection pairs every member of one with a member of the other
identical until bad with it.
\end{definition}

Two things about the shape, both~\cite{BR06}'s. Membership is checked by inspection --- it is a
comparison of code, not a claim about behaviour --- which is what keeps hops cheap to audit. And
it is built so that a divergent statement cannot run early: the only ways into a bad-guarded
block are its head, an assignment that raises the flag as it is crossed, and a test of the flag
itself, and the flag is monotone --- so nothing bad-guarded runs while its owner's flag is down.
Because paired members carry equal fields, the two systems occupy the same identities and use
the same names, $\ncl{\pi} = \ncl{\pi'}$, and one class of environments serves both:
$\ZenvSet_{\pi,\pi'} = \ZenvSet_{\pi,\pi} = \ZenvSet_{\pi',\pi'}$.

\begin{lemma}[The fundamental lemma, over systems]\label{lem:iub}
Let systems $\pi$ and $\pi'$ be identical until bad, let occupants $\mathcal{X}$ and
$\mathcal{X}'$ of the adversary slot be identical until bad, and let
$\Zenv \in \ZenvSet_{\pi,\pi'}$. Write $\Bad$ for the event that some machine executes an
assignment $\V{bad} \gets \true$, in whichever execution is named. Then
\[
  \Pr\bigl[\Bad \text{ in } \opl{Exec}(\pi,\Zenv,\mathcal{X},\Corr)\bigr]
  \;=\;
  \Pr\bigl[\Bad \text{ in } \opl{Exec}(\pi',\Zenv,\mathcal{X}',\Corr)\bigr] ,
\]
and, writing $\Pr[\Bad]$ for the common value,
\[
  \begin{gathered}
  \bigl|\, \Pr\bigl[\opl{Exec}(\pi,\Zenv,\mathcal{X},\Corr) = 1\bigr] \\[3pt]
  \;-\; \Pr\bigl[\opl{Exec}(\pi',\Zenv,\mathcal{X}',\Corr) = 1\bigr] \,\bigr|
  \;\leq\; \Pr[\Bad] .
  \end{gathered}
\]
If moreover the two systems are game systems --- the pairing then matches game to game, fields
being preserved --- the same holds of every finalization $\op{Fin}$, which the two games share:
\[
  \begin{gathered}
  \bigl|\, \Pr\bigl[\Win_{\op{Fin}} \text{ in } \opl{Exec}(\pi,\Zenv,\mathcal{X},\Corr)\bigr] \\[3pt]
  \;-\; \Pr\bigl[\Win_{\op{Fin}} \text{ in } \opl{Exec}(\pi',\Zenv,\mathcal{X}',\Corr)\bigr]
  \,\bigr| \\[3pt]
  \;\leq\; \Pr[\Bad] .
  \end{gathered}
\]
\end{lemma}

\begin{proof}
% BEGIN overview lem:iub
One coupling, read pointwise. Fixing every tape fixes both executions outright, so the two may be
taken from one sample and compared step by step. The induction shows that while no flag is up the
two correspond strictly and the token stands on \emph{common} code --- which is exactly what
Definition~\ref{def:iub}'s two bad-guarded shapes buy, neither being reachable except across the
assignment that raises the flag or through a test that fails while it is down. The first raising is
therefore the same step on both sides, and that is the first display: $\Pr[\Bad]$ is one number
rather than two that happen to share a bound. Off $\Bad$ the two runs are equal outright, so every
fact of the run has a single indicator there, and each difference of probabilities is at most the
measure of the set where indicators may differ.
% END overview lem:iub
Two conventions of Chapters~\ref{sec:functionalities} and~\ref{sec:execution} are used, the ones
the absorption lemma restates: each machine draws from its own coin tape, uniform and
independent of every other's, read only when its code samples --- an environment's hosted copies
keeping tapes of their own, bundling merging nothing --- and $\op{Initialize}$ places no calls.
Fixing every tape therefore fixes an execution outright: which machine holds the token, at which
line of which operation, with what state, and what its next statement does are functions of the
configuration, coins entering only as the next unread positions of the holder's own tape.

Couple the two executions. The pairing matches machine to machine --- member to member, occupant
to occupant, the environment and the register each to themselves --- and paired machines occupy
the same identities, so give paired machines equal tapes. Each execution separately keeps its
distribution, tapes on either side being uniform and independent, so both may be read off one
sample of tapes and compared pointwise. Say the two \emph{correspond strictly} at a
configuration if (C1)--(C4) of Definition~\ref{def:regcorr} hold under the pairing --- equal
states, equal unread tapes, the token in paired hands at the same line with the same local
values, suspended callers matching frame for frame --- with no relay anywhere and nothing
elided.

At the start they correspond strictly. Line~\ref{ln:execinit} initialises the register and every
member on both sides, in whatever order --- $\op{Initialize}$ places no calls and touches no
state but its owner's, so the order leaves no trace in the configuration; paired members run
$\op{Initialize}$s that agree line for line,
$\V{bad} \gets \false$ among them, placing no calls and drawing, where they draw, the same
positions of equal tapes; and line~\ref{ln:execstart} hands both tokens to the one environment
at the same identity.

Suppose the two correspond strictly, no flag yet raised on either side, and let one statement
run. The token stands in paired machines at the same line, and that line is \emph{common} code:
a bad-guarded statement is reached only through its head or through a test of the flag, the head
raises the flag as it is crossed, the test fails while the flag is down, and no flag is up. So
the statement is the same statement, read over the same state and the same next coins, histories
having consumed tape for tape alike. If it writes, both sides write alike; if it returns, both
resume the same frame with the same value; if it places a call, both place the same call ---
target, payload and claimed identity --- met by the same wrappers over equal fields and an equal
register state, landing at paired machines, or at no machine and answered $\rej$ alike; and
every delivery Section~\ref{sec:execution} fixes --- $\rej$ for a refused \textbf{require},
$\none$ where a core would return $\rej$ or a mediated slot answers one --- is a rule, not a
choice. Strict correspondence survives the step. And if the statement is an assignment
$\V{bad} \gets \true$ --- common code, standing at the head of its divergent continuation
--- both sides execute it, and the step is the first raising in both executions at once.

Pointwise on the coupled space, then: either no flag is ever raised and the two executions
correspond strictly at every step --- so one halts exactly when the other does, with the same
output, the environment being one machine resumed alike --- or there is a first raising, it is
simultaneous, and the executions agree up to and through it. The event $\Bad$ therefore has one
indicator on the coupled space, which is the first display, with equality. Off $\Bad$ the runs
are equal outright, so any fact of the run has one indicator there too: the output, and, for
game systems, the event $\Win_{\op{Fin}}$ that the interface $(\Gm.\PID,Z).\fopl{Fin}$ returns
$1$ (Definition~\ref{def:propadv}). Each difference of probabilities is at most the measure of
the set where its indicators may differ, which is $\Pr[\Bad]$.
\end{proof}

The definition and the proof have a picture each, and they are worth seeing together, since one is
exactly what the other consumes. Above is what Definition~\ref{def:iub} asks of the code; below is
what the coupling then does with it. Neither panel uses the box-and-wire vocabulary of the earlier
figures, and that is the point: nothing here is moved from one machine to another, so there is
nothing to draw a wire between.

\begin{center}
\begin{tikzpicture}[
  band/.style={draw=accent!55, minimum width=27mm, minimum height=5.4mm,
               inner sep=1pt, font=\scriptsize, anchor=north west},
  common/.style={band, fill=accent!4},
  head/.style={band, fill=accent!16},
  badblk/.style={band, fill=accentsoft!28, draw=accentsoft, dashed},
  colhead/.style={font=\small},
  tag/.style={font=\tiny, color=accentsoft},
  runline/.style={draw=accent, thick},
  dot/.style={circle, fill=accent, inner sep=1.1pt},
  outnode/.style={draw=accent, rounded corners=1pt, inner sep=2pt, font=\scriptsize},
]

%% ================= PANEL (a): the code shape =================
\def\bw{27mm}
\def\bh{5.4mm}

\foreach \r/\sty/\txt in {%
  0/common/{common code},
  1/head/{$\V{bad} \gets \true$},
  2/badblk/{\itshape divergent},
  3/common/{common code},
  4/head/{\textbf{if} $\V{bad}$},
  5/badblk/{\itshape divergent},
  6/common/{common code}%
} {
  \node[\sty] (L\r) at (0, -\r*\bh) {\txt};
  \node[\sty] (R\r) at (34mm, -\r*\bh) {\txt};
}

\node[colhead] at (13.5mm, 5mm) {$\F$};
\node[colhead] at (47.5mm, 5mm) {$\F'$};

\node[tag, anchor=west] at (62mm, -2*\bh+2.7mm) {continuation of the};
\node[tag, anchor=west] at (62mm, -2*\bh-0.5mm) {assignment, to the end};
\node[tag, anchor=west] at (62mm, -2*\bh-3.7mm) {of its enclosing block};

\node[tag, anchor=west] at (62mm, -5*\bh+1.2mm) {body of a test};
\node[tag, anchor=west] at (62mm, -5*\bh-2mm) {of the flag};

\draw[decorate, decoration={brace, amplitude=4pt}, accentsoft]
  (61mm, -2*\bh) -- (61mm, -3*\bh);
\draw[decorate, decoration={brace, amplitude=4pt}, accentsoft]
  (61mm, -5*\bh) -- (61mm, -6*\bh);

\node[font=\scriptsize, align=center, anchor=north] at (30mm, -7.4*\bh)
  {erase the two shaded blocks and one functionality remains,\\
   operation for operation and line for line};

%% ================= PANEL (b): the coupled run, below panel (a) =================
\def\px{0mm}
\def\py{-64mm}

%% --- off Bad ---
\node[font=\scriptsize, anchor=west] at (\px, \py+3mm) {off $\Bad$};
\draw[runline] (\px+16mm, \py) -- (\px+40mm, \py);
\node[dot] at (\px+16mm, \py) {};
\node[outnode, anchor=west] at (\px+41mm, \py) {one run, one output};
\node[tag, anchor=north west] at (\px+16mm, \py-1.5mm)
  {strict correspondence at every step};

%% --- on Bad ---
\node[font=\scriptsize, anchor=west] at (\px, \py-22mm) {on $\Bad$};
\draw[runline] (\px+16mm, \py-22mm) -- (\px+34mm, \py-22mm);
\node[dot] at (\px+16mm, \py-22mm) {};
\node[dot] (fork) at (\px+34mm, \py-22mm) {};
\draw[runline] (fork) -- (\px+40mm, \py-17mm);
\draw[runline] (fork) -- (\px+40mm, \py-27mm);
\node[outnode, anchor=west] (o1) at (\px+41mm, \py-17mm) {$\opl{Exec}(\pi,\dots)$};
\node[outnode, anchor=west] (o2) at (\px+41mm, \py-27mm) {$\opl{Exec}(\pi',\dots)$};

\draw[decorate, decoration={brace, amplitude=4pt}, accentsoft]
  ($(o1.north east)+(1.5mm,0)$) -- ($(o2.south east)+(1.5mm,0)$);
\node[tag, anchor=west, align=left] at ($(o2.east)+(3mm,5mm)$)
  {may differ here,\\ on a set of measure $\Pr[\Bad]$};

\node[tag, anchor=north, align=center] at (\px+25mm, \py-29mm)
  {first $\V{bad} \gets \true$ --- common code,\\ so one moment in both runs};

\end{tikzpicture}
\figcaption{fig:iub}{Identical until bad, in two panels. Above, what Definition~\ref{def:iub} asks of the code; below, what the coupling does with it. Neither panel uses the box-and-wire vocabulary of the earlier figures, because nothing here moves from one machine to another. (Definition~\ref{def:iub} and Lemma~\ref{lem:iub}.)}
\end{center}

The upper panel is checkable by inspection, which is what keeps a hop cheap to audit: it is a
comparison of code, not a claim about behaviour. The two shaded shapes are the only ones allowed,
and between them they are why nothing shaded can run while the flag is down --- the only ways in
are the head, which raises the flag as it is crossed, and a test of the flag, which fails while it
is down. So the shading above is exactly what runs beyond the fork below, and nowhere else.

The fork is a single point, and that is the whole of the argument. The statement heading each
divergent block is itself common code, so both runs execute it at the same step of the same
coupled sample --- which gives the lemma's \emph{first} display, that $\Pr[\Bad]$ is the same
number in the two executions rather than two numbers with one bound. Off $\Bad$ there is no fork
at all: the two runs are one run, so every fact of it has a single indicator there --- the output,
and, for game systems, $\Win_{\op{Fin}}$ --- and each difference of probabilities is at most the
measure of the set where indicators may differ. That set is the lower branch, and its measure is
$\Pr[\Bad]$.

\begin{remark}[Nothing is assumed]\label{rem:iubfree}
The lemma asks nothing of the systems or of the slot: not well-formedness, not tightness, not
refusal-obliviousness, and --- worth savouring --- not the absence of responsive calls. Every
result that moves a machine pays a hypothesis for the move: Theorem~\ref{thm:dummy} and
Proposition~\ref{prop:propslotabs} exclude responsive cores, Corollary~\ref{cor:proptransfer2}
inherits the exclusion on both sides at once, and the notifying clock, the network and the
channel sit outside all three for exactly that reason. The lemma above is indifferent --- their
responsive notifications hop as cheaply as anything else --- because nothing is moved. Like the
regrouping it is an identity about runs, resting on nothing but the determinism of an execution
once its tapes are fixed, not a comparison of what two runs achieve. No budget appears for the
same reason: exactness needs none, and Definition~\ref{def:res} enters only when an application
sets about bounding $\Pr[\Bad]$.
\end{remark}

\begin{remark}[Why not by absorption]\label{rem:iubabs}
There is a tempting shorter route: absorb the whole execution into the environment ---
Lemma~\ref{lem:absorb} exists to move machines --- until each side is a single program, and
quote~\cite{BR06} on the pair. It works where it works, and it is instructive to see why it is
the worse lemma. Absorption's hypotheses come along: the slot occupant moved must be
refusal-oblivious, no core may place a responsive call on a ceded identity
(Remark~\ref{rem:respdummy}), and the environment's $\pars$ must hold no $Z$-identity --- so the
reduction proves the lemma only for the systems the dummy theorem already covers, and the
notifying clock, for one, falls outside it. The direct coupling is no longer than the reduction
and carries nothing: the machinery of Definition~\ref{def:regcorr} was built for one execution
rebracketed, and two executions coupled on equal tapes is the same induction with an easier case
analysis --- no relay, no rewriting of claims, no refusal to smooth over.
\end{remark}

What the lemma is for, concretely. Emulation arguments hop between neighbouring \emph{ideal}
objects: a functionality and a lazier variant, a simulator and the same simulator with one check
removed. Each hop is a pair identical until bad, the flags raised where the two part company,
and the lemma prices the hop at $\Pr[\Bad]$ --- in the same execution, against the same
environment, with the occupants hopping alongside where a proof needs it. What remains is to
bound $\Pr[\Bad]$, and that is generally a \emph{property}: a finalization whose verdict is
precisely that the raising condition occurred, priced by Definition~\ref{def:propadv} and
inherited, where transfer applies, through Corollary~\ref{cor:proptransfer2}. The canonical
instance sits at the signature functionality of Chapter~\ref{sec:fsig}. Instrument $\Fsig$ to
raise $\V{bad}$ where $\op{Verify}$'s honest-key test fires --- an addition no caller can see,
the flag being written and never read --- and set beside it the variant that then lets the
forgery stand. The two are identical until bad, and the hop between them costs the probability
that a forgery is ever submitted --- the quantity $\op{FinUF}$ exists to police.

%% ==================================================================
\chapter{پاک‌سازی}\label{sec:sanitization}

We will often need to \emph{sanitize} simulator outputs, so attacks do not affect ideal
functionalities too adversely. In honest-majority settings, for instance, we typically sanitize a
$\none$ output to a non-$\none$ one to preserve liveness. We write such a procedure
$\San[\Clean](v; c)$, where $\Clean$ is a predicate on a value together with whatever context it
consults, and replace a simulator output $v$ by $v' \gets \San[\Clean](v; c)$, which operates as
below unless stated otherwise. The context $c$ after the semicolon lists the functionality state
the predicate reads; it is read-only, and different predicates take different contexts.

\begin{procedure}{Sanitizer $\opdef{San}[\Clean]$ \\
\params{predicate $\Clean$, \quad value $v$, \quad context $c$}}
\opsig{$\San[\Clean](v; c)$}:
\begin{algorithmic}[1]
\If{$\Clean(v; c)$}
  \State \textbf{return} $v$
\EndIf
\State pick $v' \in \{\, v' : \Clean(v'; c) \,\}$ \label{ln:sanpick}
\State \textbf{return} $v'$
\end{algorithmic}
\end{procedure}

So sanitization acts only when $\Clean(v; c)$ fails, and what it returns then satisfies
$\Clean$.\footnote{A good value exists whenever the space $\Clean$ draws from is infinite and
everything it excludes beyond non-membership is a value named in the context $c$. Both halves are
needed for the predicates of Chapter~\ref{sec:fsig}: membership in an infinite $\Keys$ or $\Sigs$
is the first conjunct of each, and everything else they exclude is an entry of a table passed in
$c$. Under Definition~\ref{def:res} the context is finitely supported in any execution --- a
bounded run records finitely many entries --- so the excluded set is finite inside an infinite
space.} Our sanitizers need no more
than \emph{read-only} access to the functionality state. How line~\ref{ln:sanpick} picks is part
of the specification --- the picked value reaches the caller, so different rules are observably
different functionalities --- and we fix the lexicographically first good value throughout. The
choice matters less than it looks: the pick surfaces only when the slot's own answer fails
$\Clean$, and an adversary preferring some particular good value may simply submit it.

%% ==================================================================
%% ==================================================================
\chapter{تصادف}\label{sec:frand}

\section{کارکرد}\label{sec:frandfunc}

$\Frand$ is the smallest functionality in this paper and the right one to read first. It hands out
uniform $n$-bit strings, remembers what it handed out, lets a party forget an entry, and leaks
what it still remembers. Chapter~\ref{sec:functionalities} named it and owed no code; here is the
code.

Two things make it worth its own chapter rather than a footnote. It places no call at all --- not
on the clock, not on the adversary --- so it is the one example where nothing at all is left to
the slot, and Chapter~\ref{app:responsive} has nothing to price. And its property, unlike every
other in this paper, is not attained at $0$: randomness is unpredictable up to a probability the
bound has to name, which is where the $\varepsilon$ of Definition~\ref{def:propadv} stops being
decorative.

\heading{The fields} $\Fsig$'s pattern. $\PID$, $\Ps$ and $\admits$ are parameters, one instance
per session, $\admits$ naming the protocol served in the local pattern of
Chapter~\ref{sec:functionalities}. $\uses := \emptyset$, there being no call to make.
$\pars := n$, the width of the strings, which $\opl{Rnd}$ reads at line~\ref{ln:rndsample}.

\heading{Erasure is what the leak is for} $\op{Leak}$ returns $\V{L}$ entire, so a corrupt party's
whole history of draws is the adversary's. That is not a defect but the point of having
$\opl{Erase}$ beside it: a protocol that draws randomness it must not retain erases the entry, and
what has been erased cannot leak. Read that way $\Frand$ is a model of secure erasure as much as
of sampling, and the pair is why the store of Chapter~\ref{sec:fstore} is shaped the same way.
One consequence is worth stating plainly, since it is easy to state wrongly: while nothing is
erased, the number of entries in $\V{L}$ is the number of times $\op{Rnd}$ was called, so the leak
reports the count as well as the values; erasing is exactly what takes an entry out of that count.

\begin{interface}[nofloat]{کارکردِ $\Frand$ \\
\params{$\PID$, \quad $\Ps$, \quad $\admits$, \quad $\uses := \emptyset$, \quad $\pars := n$}}

{\scriptsize
\setlength{\columnsep}{1.4em}
\begin{multicols}{2}
\opsig{$\op{Initialize}()$}:
\begin{algorithmic}[1]
\setcounter{contline}{0}\algcont
\State $\V{ctr} \gets 0$
\State $\V{L} : \mathbb{N} \to \{0,1\}^{n} \cup \{\unset\}$
\State $\V{L}[*] \gets \unset$
\algsave
\end{algorithmic}
\vspace{1.4ex}

\opsig{$\id.\op{Rnd}()$} \quad از $\id'$
\begin{algorithmic}[1]
\algcont
\State $\V{ctr} \gets \V{ctr} + 1$
\State $r \gets_{\$} \{0,1\}^{n}$ \label{ln:rndsample}
  \Comment{تنها نمونه‌برداریِ این کتاب}
\State $\V{L}[\V{ctr}] \gets r$
\State \textbf{return} $(\V{ctr},\, r)$ \label{ln:rndreturn}
  \Comment{اندیس، تا فراخواننده بتواند پاک کند}
\algsave
\end{algorithmic}
\vspace{1.4ex}

\columnbreak

\opsig{$\id.\op{Erase}(i)$} \quad از $\id'$
\begin{algorithmic}[1]
\algcont
\State $\V{L}[i] \gets \unset$ \label{ln:rnderase}
\State \textbf{return} $\ok$
\algsave
\end{algorithmic}
\vspace{1.4ex}

\opsig{$\id.\op{Leak}()$} \quad از $\id'$
\begin{algorithmic}[1]
\algcont
\State \textbf{return} $\V{L}$
\algsave
\end{algorithmic}
\end{multicols}}
\end{interface}

Three remarks on the code. The operation $\op{Rnd}$ answers with the index it wrote as well as the string. A
caller that means to erase a draw must be able to name it, and $\V{ctr}$ is the functionality's
state and not the caller's, so without line~\ref{ln:rndreturn} the $\op{Erase}$ beside it could
only be used by a caller willing to guess --- the protocol of Section~\ref{sec:psig} erases every
coin it draws, and this is what lets it. Entries are indexed by a counter and never rewritten, so $\op{Erase}$
needs no companion to $\Fstore$'s distinction between never-written and erased: an index
$\op{Rnd}$ has passed is never reused, and $\unset$ at an index below $\V{ctr}$ says the entry was
erased while $\unset$ above it says the counter has not reached there. And the counter is shared
across parties, as any functionality's state is: two parties drawing from one instance draw
distinct entries, which is what makes $\Frand$ a source rather than a per-party generator.

\section{ویژگی‌ها}\label{sec:randprops}

One property, $\opl{FinUnp}$, \emph{unpredictability}: no environment names in advance a value a
later draw returns. It takes the search reading of Definition~\ref{def:propadv}. The game is
played over $\gamma := \{\Gm\} \cup \{\Frand\}$ with $\Frand.\admits := \{\Gm.\PID\}$ pinned and
$\Frand.\Ps \subseteq \Gm.\Ps$, Section~\ref{sec:sigprops}'s tightness argument once more, and
for its reason: a draw the game does not see is a draw the verdict cannot rule on.

The shell carries one interface that is not a relay. The operation $\opl{Guess}$ places no call and touches
$\Frand$ not at all; it exists so that the environment can go on record, and the trace being a
sequence is what lets the verdict ask whether it went on record \emph{first}. That is
Section~\ref{sec:clockprops}'s use of order, borrowed for a different purpose: steadiness read
the order of answers, unpredictability reads the order of a claim against an answer.

\begin{interface}[nofloat]{پوستهٔ بازی روی $\Frand$؛ نهایی‌سازی در پی می‌آید \\
\params{$\PID := (\Gm.F,0,0)$, \quad $\Ps := \Frand.\Ps \cup \{Z\}$, \quad $\admits := \Zpid$, \quad
$\uses := \{(\Frand,\serves_{\Gm})\}$, \quad $\pars := \none$}}

{\scriptsize
\setlength{\columnsep}{1.4em}
\begin{multicols}{2}
\opsig{$\op{Initialize}()$}:
\begin{algorithmic}[1]
\setcounter{contline}{0}\algcont
\State $\done \gets \false$
\State $\V{tr} \gets (\,)$
  \Comment{رد، یک دنباله}
\algsave
\end{algorithmic}
\vspace{1.4ex}

\opsig{$\id.\opdef{Rnd}()$} \quad از $\id'$
\begin{algorithmic}[1]
\algcont
\Req{$\neg\done$}
\Req{$\id.P \neq Z$}
  \Comment{ریشه تنها نهایی می‌کند}
\State $(i,r) \gets \id_{\Frand}.\fopl{Rnd}() \text{ as } \id$ \label{ln:grnd}
  \Comment{اندیس ثبت نمی‌شود}
\State $\V{tr} \gets \V{tr} \cdot (\atom{rnd},\id.P,r)$
\State \textbf{return} $r$
\algsave
\end{algorithmic}
\vspace{1.4ex}

\opsig{$\id.\opdef{Erase}(i)$} \quad از $\id'$
\begin{algorithmic}[1]
\algcont
\Req{$\neg\done$}
\Req{$\id.P \neq Z$}
\State $\id_{\Frand}.\fopl{Erase}(i) \text{ as } \id$
\State $\V{tr} \gets \V{tr} \cdot (\atom{erase},\id.P,i)$
\State \textbf{return} $\ok$
\algsave
\end{algorithmic}
\vspace{1.4ex}

\columnbreak

\opsig{$\id.\opdef{Guess}(y)$} \quad از $\id'$
\begin{algorithmic}[1]
\algcont
\Req{$\neg\done$}
\State $\V{tr} \gets \V{tr} \cdot (\atom{guess},\id.P,y)$ \label{ln:gguess}
  \Comment{بی‌فراخوان؛ ریشه هم می‌تواند حدس بزند}
\State \textbf{return} $\ok$
\algsave
\end{algorithmic}
\vspace{1.4ex}

\opsig{$\id.\op{Leak}()$} \quad از $\id'$
\begin{algorithmic}[1]
\algcont
\State \textbf{return} $\none$
\algsave
\end{algorithmic}
\end{multicols}}
\end{interface}

\begin{interface}[nofloat]{The finalization of $\Gm$; the shell over $\Frand$ is shared}
\opsig{$\id.\opdef{FinUnp}()$} \quad from $\id'$ \quad (unpredictability)
\begin{algorithmic}[1]
\algcont
\Req{$\neg\done$}
\Req{$\id.P = Z$}
\State $\done \gets \true$
\State parse $\V{tr}$ as $\bigl(e_j\bigr)_{j=1}^{m}$
\State $w \gets 1$ if $\exists\, j < k \ \ \exists\, y :\ \
       e_j = (\atom{guess},\cdot,y) \ \wedge\ e_k = (\atom{rnd},\cdot,y)$ \label{ln:gunp}
\State else $w \gets 0$
\State \textbf{return} $w$
\end{algorithmic}
\end{interface}

Two things about the code. The relays refuse the root and $\op{Guess}$ does not, and the asymmetry
is exactly Section~\ref{sec:clockprops}'s: a relay at the root would address a party $\Frand$
does not serve, be answered $\rej$ at line~\ref{ln:interface}, and append $(\atom{rnd},Z,\rej)$ to
the trace --- whereupon an environment that guessed $\rej$ would win outright, having predicted
nothing. Refused at the relay, nothing is appended, and every $\atom{rnd}$ entry carries a string
$\Frand$ actually drew. The operation $\op{Guess}$ places no call, so it has nothing to be refused for and is
served everywhere, the root included. And the verdict asks honesty of nobody: at a corrupt party
the game's own wrapper intercepts before the core runs (Remark~\ref{rem:gmcorr}), so every entry
in $\V{tr}$ is an honest core's already, exactly as in Section~\ref{sec:clockprops}.

\begin{proposition}[$\Frand$ is unpredictable]\label{prop:randunp}
Let $\Zenv$ place at most $q_{\op{g}}$ calls on $\op{Guess}$ and at most $q_{\op{r}}$ on
$\op{Rnd}$. Then for every occupant $\mathcal{X}$ of the adversary slot,
\[
  \Pr[\Win_{\op{FinUnp}}] \;\leq\; q_{\op{g}} \, q_{\op{r}} \, 2^{-n} .
\]
So $\{\Frand\}$ has unpredictability within $q_{\op{g}} q_{\op{r}} 2^{-n}$ against
$\bigl(\ZenvSet,\Xset\bigr)$ for $\ZenvSet$ any class meeting those counts and $\Xset$ the class
of all occupants.
\end{proposition}

\begin{proof}
Fix $k$ with $e_k = (\atom{rnd},\cdot,r)$. By tightness no caller but $\Gm$
reaches $\Frand$, and $\Gm$ reaches it only through the shell's $\op{Rnd}$ relay
(line~\ref{ln:grnd}), so $r$ is the value $\Frand$'s own $\op{Rnd}$ drew at
line~\ref{ln:rndsample}: uniform on $\{0,1\}^{n}$ and independent of everything
the execution did before it, $\Frand$ placing no call and reading no state to draw it.

Condition on the execution up to the moment before that draw. Every $\atom{guess}$ entry with $j <
k$ is determined by that history, and there are at most $q_{\op{g}}$ of them, so the probability
that $r$ equals any of their values is at most $q_{\op{g}} 2^{-n}$. There are at most $q_{\op{r}}$
choices of $k$, and the verdict at line~\ref{ln:gunp} returns $1$ only if some pair $(j,k)$
matches, so a union bound over $k$ gives $q_{\op{g}} q_{\op{r}} 2^{-n}$.

Two details the bound rests on. The occupant of the slot does not appear: $\Frand$'s cores place
no call, so nothing $\mathcal{X}$ does reaches the sampling, and the same bound holds against
every occupant without an absorption argument. And a guess made \emph{after} a draw is worthless
by construction, the verdict asking $j < k$; an environment that reads $r$ and then guesses
it has predicted nothing, which is what makes the trace's order the content of this verdict.
\end{proof}

The bound is the first in this paper that is not $0$, and it is worth saying what that changes.
Corollary~\ref{cor:proptransfer2} reads the same either way: a system $\pi$ put in place of
$\{\Frand\}$ has unpredictability within $q_{\op{g}} q_{\op{r}} 2^{-n} + \varepsilon$ under the
search reading, the ideal bound and the emulation error simply adding. What it changes is that the
ideal side is no longer free --- a realization inherits a bound that was already nontrivial, and
the two terms are of different kinds: one is the sampling's own, and no protocol does better than
it, while the other is what the realization costs. Transfer itself is as cheap as it gets:
$\Frand$'s cores place no call, responsive or otherwise, so Remark~\ref{rem:respdummy}'s caveat
never arises and Theorem~\ref{thm:dummy} is untouched, exactly as for $\Gpki$
(Section~\ref{sec:pkiprops}).

%% ==================================================================
%% ==================================================================
\chapter{ذخیره‌سازی}\label{sec:fstore}

\section{کارکرد}\label{sec:fstorefunc}

$\Fstore$ is $\Frand$'s companion: where one produces values and lets them be forgotten, the other
takes values and holds them until they are. A party puts a value at an index, gets it back, and
erases it; the leak reports what is still held. Together they are what
Chapter~\ref{sec:functionalities} means by a protocol keeping no state on a machine tape --- all
randomness through $\Frand$, all persistence through $\Fstore$, each exposing what it holds
through its own $\op{Leak}$.

\heading{The fields} $\Frand$'s exactly, save the parameter: $\PID$, $\Ps$ and $\admits$
parameters, $\uses := \emptyset$, and $\pars := \none$, a store reading none. Like $\Frand$'s, its
state is shared across the parties it serves --- one party may get what another put, which is what
makes it a store rather than a per-party tape, and what
Chapter~\ref{sec:functionalities} means by sharing being the point.

\heading{Three states, not two} An index is in one of three conditions, and the code keeps them
apart on purpose. It is $\unset$ if nothing was ever put there; it holds a value if something was
put and not erased; and it is $\none$ if it was put and then erased. The operation $\opl{Put}$ accepts an index
in either empty condition, so an erased index is reusable; $\opl{Get}$ refuses only $\unset$, so a
get on an erased index is answered $\none$ rather than refused. The distinction earns its place
in the difference between those two answers: a caller can tell \emph{nothing was ever here} from
\emph{something was here and is gone}, which is what makes erasure observable to the protocol that
performed it rather than a silent hole. A store that hid the difference would refuse both alike
and could not be asked whether an erasure had happened.

\begin{interface}[nofloat]{کارکردِ $\Fstore$ \\
\params{$\PID$, \quad $\Ps$, \quad $\admits$, \quad $\uses := \emptyset$, \quad $\pars := \none$}}

{\scriptsize
\setlength{\columnsep}{1.4em}
\begin{multicols}{2}
\opsig{$\op{Initialize}()$}:
\begin{algorithmic}[1]
\setcounter{contline}{0}\algcont
\State $\V{L} : \mathbb{N} \to \bits \cup \{\none,\unset\}$
\State $\V{L}[*] \gets \unset$
\algsave
\end{algorithmic}
\vspace{1.4ex}

\opsig{$\id.\op{Put}(i,x)$} \quad از $\id'$
\begin{algorithmic}[1]
\algcont
\Req{$\V{L}[i] \in \{\none,\unset\}$} \label{ln:putfree}
  \Comment{اندیسِ اشغال‌شده رد می‌کند}
\State $\V{L}[i] \gets x$ \label{ln:putwrite}
\State \textbf{return} $\ok$
\algsave
\end{algorithmic}
\vspace{1.4ex}

\opsig{$\id.\op{Erase}(i)$} \quad از $\id'$
\begin{algorithmic}[1]
\algcont
\State $\V{L}[i] \gets \none$ \label{ln:storeerase}
\State \textbf{return} $\ok$
\algsave
\end{algorithmic}
\vspace{1.4ex}

\columnbreak

\opsig{$\id.\op{Get}(i)$} \quad از $\id'$
\begin{algorithmic}[1]
\algcont
\Req{$\V{L}[i] \neq \unset$} \label{ln:getset}
  \Comment{هرگز نوشته‌نشده رد می‌کند}
\State \textbf{return} $\V{L}[i]$ \label{ln:getread}
  \Comment{$\none$ اگر پاک شده باشد}
\algsave
\end{algorithmic}
\vspace{1.4ex}

\opsig{$\id.\op{Leak}()$} \quad از $\id'$
\begin{algorithmic}[1]
\algcont
\State \textbf{return} $\V{L}$
\algsave
\end{algorithmic}
\end{multicols}}
\end{interface}

Two remarks on the code. The refusals are written as \textbf{require}s rather than as returns of
$\rej$, which is the same thing said in the framework's own words: Section~\ref{sec:execution}
answers a refused \textbf{require} with $\rej$, so lines~\ref{ln:putfree} and~\ref{ln:getset} give
exactly the two refusals a store owes, and no core of this paper returns $\rej$ by writing it. And
$\op{Erase}$ refuses nothing: erasing an index that holds nothing is a no-op that answers $\ok$,
because a protocol erasing defensively should not have to know whether there was anything there.

\section{ویژگی‌ها}\label{sec:storeprops}

Two properties, and they are the two halves of what a store is for. $\opl{FinKeep}$,
\emph{persistence}: a value put and not erased is the value got back. $\opl{FinGone}$,
\emph{erasure}: a value erased is not got back. Both take the search reading, and both are two
finalizations of one game --- they read the same trace of the same three relays and differ in the
verdict alone, which is the test Section~\ref{sec:clockprops} gives for sharing a shell.

They are played over $\gamma := \{\Gm\} \cup \{\Fstore\}$ with $\Fstore.\admits := \{\Gm.\PID\}$
pinned and $\Fstore.\Ps \subseteq \Gm.\Ps$, tightness as before and for the same reason: a put the
game does not see is a value the verdict cannot account for. The relays refuse the root, for
Section~\ref{sec:randprops}'s reason: both verdicts compare a got value against a put one, so a
$\rej$ of the game's own making in the trace would be read as a value and could win either.

\begin{interface}[nofloat]{پوستهٔ بازی روی $\Fstore$؛ نهایی‌سازی‌ها در پی می‌آیند \\
\params{$\PID := (\Gm.F,0,0)$, \quad $\Ps := \Fstore.\Ps \cup \{Z\}$, \quad $\admits := \Zpid$, \quad
$\uses := \{(\Fstore,\serves_{\Gm})\}$, \quad $\pars := \none$}}

{\scriptsize
\setlength{\columnsep}{1.4em}
\begin{multicols}{2}
\opsig{$\op{Initialize}()$}:
\begin{algorithmic}[1]
\setcounter{contline}{0}\algcont
\State $\done \gets \false$
\State $\V{tr} \gets (\,)$
\algsave
\end{algorithmic}
\vspace{1.4ex}

\opsig{$\id.\opdef{Put}(i,x)$} \quad از $\id'$
\begin{algorithmic}[1]
\algcont
\Req{$\neg\done$}
\Req{$\id.P \neq Z$}
\State $a \gets \id_{\Fstore}.\fopl{Put}(i,x) \text{ as } \id$ \label{ln:gput}
\State $\V{tr} \gets \V{tr} \cdot (\atom{put},\id.P,i,x,a)$
\State \textbf{return} $a$
\algsave
\end{algorithmic}
\vspace{1.4ex}

\opsig{$\id.\op{Erase}(i)$} \quad از $\id'$
\begin{algorithmic}[1]
\algcont
\Req{$\neg\done$}
\Req{$\id.P \neq Z$}
\State $\id_{\Fstore}.\fopl{Erase}(i) \text{ as } \id$
\State $\V{tr} \gets \V{tr} \cdot (\atom{erase},\id.P,i)$ \label{ln:gsterase}
\State \textbf{return} $\ok$
\algsave
\end{algorithmic}
\vspace{1.4ex}

\columnbreak

\opsig{$\id.\opdef{Get}(i)$} \quad از $\id'$
\begin{algorithmic}[1]
\algcont
\Req{$\neg\done$}
\Req{$\id.P \neq Z$}
\State $y \gets \id_{\Fstore}.\fopl{Get}(i) \text{ as } \id$ \label{ln:gget}
\State $\V{tr} \gets \V{tr} \cdot (\atom{get},\id.P,i,y)$
\State \textbf{return} $y$
\algsave
\end{algorithmic}
\vspace{1.4ex}

\opsig{$\id.\op{Leak}()$} \quad از $\id'$
\begin{algorithmic}[1]
\algcont
\State \textbf{return} $\none$
\algsave
\end{algorithmic}
\end{multicols}}
\end{interface}

\begin{interface}[nofloat]{The finalizations of $\Gm$; the shell over $\Fstore$ is shared}
\opsig{$\id.\opdef{FinKeep}()$} \quad from $\id'$ \quad (persistence)
\begin{algorithmic}[1]
\algcont
\Req{$\neg\done$}
\Req{$\id.P = Z$}
\State $\done \gets \true$; \ parse $\V{tr}$ as $\bigl(e_j\bigr)_{j=1}^{m}$
\State $w \gets 1$ if $\exists\, j < k \ \ \exists\, i,x,y :\ \
       e_j = (\atom{put},\cdot,i,x,\ok) \ \wedge\ e_k = (\atom{get},\cdot,i,y)$ \label{ln:gkeep}
\State \quad ${}\wedge\ y \neq x
       \ \wedge\ \neg\exists\, l \in (j,k) : e_l = (\atom{erase},\cdot,i)$;
\State else $w \gets 0$
\State \textbf{return} $w$
\algsave
\end{algorithmic}
\vspace{1.4ex}

\opsig{$\id.\opdef{FinGone}()$} \quad from $\id'$ \quad (erasure)
\begin{algorithmic}[1]
\algcont
\Req{$\neg\done$}
\Req{$\id.P = Z$}
\State $\done \gets \true$; \ parse $\V{tr}$ as $\bigl(e_j\bigr)_{j=1}^{m}$
\State $w \gets 1$ if $\exists\, j < k < l \ \ \exists\, i,x :\ \
       e_j = (\atom{put},\cdot,i,x,\ok)$ \label{ln:ggone}
\State \quad ${}\wedge\ e_k = (\atom{erase},\cdot,i)
       \ \wedge\ e_l = (\atom{get},\cdot,i,x)$
\State \quad ${}\wedge\ \neg\exists\, l' \in (k,l) : e_{l'} = (\atom{put},\cdot,i,\cdot,\ok)$;
\State else $w \gets 0$
\State \textbf{return} $w$
\algsave
\end{algorithmic}
\end{interface}

The verdicts, in words. The finalization $\op{FinKeep}$ returns $1$ if some get answered other than the value the
last successful put at that index left there, no erasure intervening. The finalization $\op{FinGone}$ returns $1$
if some get answered exactly the value an erasure was supposed to have removed, no put having
restored it in between. Both read the trace's order, as Section~\ref{sec:clockprops}'s steadiness
does and for the same reason: a store is a claim about what happens between one call and the next,
and between is an ordering.

\begin{proposition}[$\Fstore$ has both properties]\label{prop:storeprops}
For $\op{Fin}$ either of $\op{FinKeep}$ and $\op{FinGone}$, for every
$\Zenv \in \ZenvSet_{\gamma,\gamma}$ and every occupant $\mathcal{X}$ of the adversary slot,
$\Pr[\Win_{\op{Fin}}] = 0$. So $\{\Fstore\}$ has each property within $0$ against
$\bigl(\ZenvSet_{\gamma,\gamma},\Xset\bigr)$ in the sense of Definition~\ref{def:propadv}, for
$\Xset$ the class of all occupants.
\end{proposition}

\begin{proof}
% BEGIN overview prop:storeprops
One invariant carries both properties: $\V{L}[i]$ changes only inside a call the trace records, and
to the value that call names. It comes from tightness (no caller but the game passes the guard), from
the game's own wrapper intercepting at a corrupt party before any core runs, and from $\Fstore$'s
cores placing no call, so nothing runs between a relay's call and its return and entries append in
execution order. Both properties are then read off it by the same move: assume the verdict's
conjuncts, and track which lines could have written $\V{L}[i]$ between the entries named. Persistence
turns on the put refusing an occupied index, so no later put can intervene unrecorded; erasure turns
on $\bits$ containing neither $\none$ nor $\unset$, so the value a get returns after an erase cannot
be the value put before it.
% END overview prop:storeprops
Everything rests on one invariant: $\V{L}[i]$ changes only inside a call the trace records, and to
the value that call names. By tightness no caller but $\Gm$ passes $\opl{Guard}$'s
line~\ref{ln:caller} at
$\Fstore$, and $\Gm$ reaches it only through the shell's three relays (lines~\ref{ln:gput},
\ref{ln:gsterase} and~\ref{ln:gget});
at a corrupt party the game's own wrapper intercepts a step earlier and no core
of $\gamma$ runs at all (Remark~\ref{rem:gmcorr}). $\Fstore$'s cores place no call, so no other
machine of $\gamma$ runs between a relay's call and its return, and entries append in execution
order. Line numbers below name $\Fstore$'s own cores unless said otherwise.

For persistence: suppose $e_j = (\atom{put},\cdot,i,x,\ok)$ and $e_k = (\atom{get},\cdot,i,y)$ with
$j < k$ and no $\atom{erase}$ at $i$ between. The put answered $\ok$, so $\op{Put}$'s
line~\ref{ln:putfree}
passed and its line~\ref{ln:putwrite} set $\V{L}[i] \gets x$. Only that line
and $\op{Erase}$'s line~\ref{ln:storeerase} write $\V{L}[i]$, and by the invariant each is
recorded: no $\atom{erase}$
entry at $i$ lies between, and a further $\atom{put}$ at $i$ answering $\ok$ cannot lie between
either, line~\ref{ln:putfree} refusing an occupied index and $\V{L}[i]$ being occupied from $j$
on. So $\V{L}[i] = x \neq \unset$ still at $k$, $\op{Get}$'s line~\ref{ln:getset} passes, and its
line~\ref{ln:getread} returns $x$; hence $y = x$ and the verdict at line~\ref{ln:gkeep} gives $0$.

For erasure: suppose $e_j = (\atom{put},\cdot,i,x,\ok)$, $e_k = (\atom{erase},\cdot,i)$ and
$e_l = (\atom{get},\cdot,i,x)$ with $j < k < l$ and no successful $\atom{put}$ at $i$ between $k$ and
$l$. $\op{Erase}$'s line~\ref{ln:storeerase} set $\V{L}[i] \gets \none$ at $k$, and by the
invariant the only
line that could restore a value is $\op{Put}$'s line~\ref{ln:putwrite} inside a recorded
$\atom{put}$ at $i$
answering $\ok$, which the hypothesis excludes. So $\V{L}[i] = \none$ at $l$, and $\op{Get}$'s
line~\ref{ln:getread} returns $\none$. Now $x \neq \none$: $x$ was written by
line~\ref{ln:putwrite} from a $\op{Put}$'s argument, and $\bits$ contains neither $\none$ nor
$\unset$. Hence the get did not answer $x$, contrary to $e_l$, and the verdict at
line~\ref{ln:ggone} gives $0$.
\end{proof}

Both transfer by Corollary~\ref{cor:proptransfer2} with nothing to discharge beyond its four
standing conditions: $\Fstore$'s cores place no call at all, so the responsive-call hypothesis is
met the cheap way and dummy completeness is untouched, exactly as for $\Frand$ and $\Gpki$. The
three functionalities that ask the adversary nothing are the three that pay nothing, and the
pattern is not a coincidence --- what Section~\ref{sec:clockprops} charges the clock is the
notification, and a functionality with no notification has no bill.

%% ==================================================================
\chapter{امضاهای دیجیتال}\label{sec:fsig}

\section{کارکرد}\label{sec:fsigfunc}

$\Fsig$ below idealizes digital signatures, in the tradition of Canetti's own ideal signature
functionality~\cite{Canetti04}. Its process id is a parameter and not a constant. The
name component is fixed --- that is what the subscript records, by the naming convention of
Chapter~\ref{sec:functionalities} --- but the session and instance are not, so one system may hold
a copy per session and several within a session, the situation composition exists to handle, and
one Definition~\ref{def:sys} would forbid outright if the whole process id were pinned. The code
below reads neither component, so Chapter~\ref{sec:reloc} lets it be read at any other --- for an
$\admits$ meeting the condition of Definition~\ref{def:sessunif}, which any global or purely local one
does. Recall
from Section~\ref{sec:adv} that $\Adv(x)$ abbreviates a call to the adversary on behalf of the
party being served; in the ideal execution the simulator answers those calls, and every one of
them here is sanitized in the sense of Chapter~\ref{sec:sanitization} before use.

\heading{How we write ideal functionalities} The style is the same throughout. Wherever a value
has to be produced and the specification does not pin it down, we let the simulator choose it and
sanitize the choice on the fly. The simulator is then meant to have as much freedom as the
security goal permits and no more --- a design intention, not a theorem --- and the goal stays
visible in the predicates rather than buried in the code.

In $\Fsig$ this happens three times, and the first two protect correctness. The operation $\opl{Gen}$ takes the
verification key from the adversary, and $\Clean_{\vk}$ makes it an actual key, distinct from
every key already issued and free of any signature already recorded valid under it, so generation
always succeeds and keys neither collide nor arrive pre-forged. The operation $\opl{Sign}$ takes the signature,
and $\Clean_{\sigma}$ makes it an actual signature not already recorded invalid, so signing
succeeds whenever a key has been generated, and what it produces verifies. Locality asks for that,
and asks in addition that the simulator answer without passing the token elsewhere, a condition on
the simulator class rather than a property of the code. The operation $\opl{Verify}$ is sanitized for
unforgeability instead: a verdict outside $\{0,1\}$ becomes $0$, and so does any verdict of $1$ on
a fresh triple under a key an honest party generated, which makes signatures \emph{strongly}
unforgeable. There the good value is unique, so the sanitization is written inline rather
than through $\San$. Verification keys are not assumed
authenticated, so $\op{Verify}$ takes a $\vk$ of its own.

$\op{Leak}$, which Section~\ref{sec:leakage} requires of an ideal functionality, hands the
adversary the party's verification key together with the signatures recorded valid under it, and
nothing else of the state. The operation $\op{Initialize}$ fixes the two tables and their types: $\V{VK}$ holds
one key per served party and $\V{Ver}$ the verdict on each triple of key, message and signature,
both starting everywhere at $\unset$, the value that separates ``not yet decided'' from a
recorded $0$ or $1$. Both $\Keys$ and $\Sigs$ are infinite and contain neither $\none$ nor
$\unset$: infinite so that the sanitizers never run out of good values, and clean of the two
markers so that membership already refuses them.

Where locality is wanted, the simulator class must be restricted to the \emph{responsive} ones:
those answering each simulator call without placing any call of their own, so the token returns
at once and simulator
calls are subroutine rather than coroutine calls --- the discipline
Chapter~\ref{app:responsive} enforces in code. For non-interactive, stateless signatures this
costs nothing, the usual simulators being responsive already.

One further line in $\op{Gen}$ and $\op{Verify}$ guards against re-entrance.
Section~\ref{sec:execution} lets a suspended instance be re-entered, and both operations suspend
at an adversary call between testing a table entry and writing it; each therefore re-tests the
entry on resumption, and a value recorded meanwhile wins. Without the re-test two nested calls
could issue two keys to one party, or return different verdicts on one triple --- the very
consistency $\V{Ver}$ exists to record.


\begin{interface}[nofloat]{کارکردِ $\Fsig$ \\
\params{$\PID$, \quad $\Ps$, \quad $\admits$, \quad
$\uses := \{(\Adv,\serves)\}$, \quad $\pars := \none$}}

{\scriptsize
\setlength{\columnsep}{1.4em}
\begin{multicols}{2}
\opsig{$\op{Initialize}()$}:
\begin{algorithmic}[1]
\setcounter{contline}{0}\algcont
\State $\V{VK} : \Fsig.\Ps \to \Keys \cup \{\unset\}$
\State $\V{VK}[*] \gets \unset$
\State $\V{Ver} : \Keys \times \Msgs \times \Sigs \to \{0,1\} \cup \{\unset\}$
\State $\V{Ver}[*,*,*] \gets \unset$
\algsave
\end{algorithmic}
\vspace{1.4ex}

\opsig{$\id.\op{Gen}()$} \quad از $\id'$
\begin{algorithmic}[1]
\algcont
\If{$\V{VK}[\id.P] \neq \unset$}
  \State \textbf{return} $\V{VK}[\id.P]$
\EndIf
\State $\vk \gets \Adv\bigl(\id.\op{Gen}\bigr)$
\If{$\V{VK}[\id.P] \neq \unset$}
  \State \textbf{return} $\V{VK}[\id.P]$
    \Comment{در حالِ تعلیق ثبت شد}
\EndIf
\State $\vk \gets \San[\Clean_{\vk}](\vk; \V{VK}, \V{Ver})$
\State $\V{VK}[\id.P] \gets \vk$
\State \textbf{return} $\vk$
\algsave
\end{algorithmic}
\vspace{1.4ex}

\opsig{$\id.\op{Sign}(\msg)$} \quad از $\id'$
\begin{algorithmic}[1]
\algcont
\State $\vk \gets \V{VK}[\id.P]$
\If{$\vk = \unset$}
  \State \textbf{return} $\none$
\EndIf
\State $\sigma \gets \Adv\bigl(\id.\op{Sign}, \msg\bigr)$
\State $\sigma \gets \San[\Clean_{\sigma}](\sigma; \msg, \vk, \V{Ver})$
\State $\V{Ver}[\vk, \msg, \sigma] \gets 1$
\State \textbf{return} $\sigma$
\algsave
\end{algorithmic}
\vspace{1.4ex}

\columnbreak

\opsig{$\id.\op{Verify}(\vk, \msg, \sigma)$} \quad از $\id'$
\begin{algorithmic}[1]
\algcont
\If{$\V{Ver}[\vk, \msg, \sigma] \neq \unset$}
  \State \textbf{return} $\V{Ver}[\vk, \msg, \sigma]$
\EndIf
\State $b \gets \Adv\bigl(\id.\op{Verify}, \vk, \msg, \sigma\bigr)$
\If{$\V{Ver}[\vk, \msg, \sigma] \neq \unset$}
  \State \textbf{return} $\V{Ver}[\vk, \msg, \sigma]$
    \Comment{در حالِ تعلیق ثبت شد}
\EndIf
\If{$b \not\in \{0,1\}$}
  \State $b \gets 0$
\EndIf
\If{$\exists\, P' \not\in \Cs \, : \, \vk = \V{VK}[P']$}
  \State $b \gets 0$
    \Comment{هیچ جعلی زیرِ کلیدِ درستکار}
\EndIf
\State $\V{Ver}[\vk, \msg, \sigma] \gets b$
\State \textbf{return} $b$
\algsave
\end{algorithmic}
\vspace{1.4ex}

\opsig{$\id.\op{Leak}()$} \quad از $\id'$
\begin{algorithmic}[1]
\algcont
\If{$\V{VK}[\id.P] = \unset$}
  \State \textbf{return} $(\none, \emptyset)$
\EndIf
\State \textbf{return} $\bigl(\V{VK}[\id.P],$
\Statex \quad\ $\{\, (\msg,\sigma) : \V{Ver}[\V{VK}[\id.P],\msg,\sigma] = 1 \,\}\bigr)$
\algsave
\end{algorithmic}
\vspace{1.4ex}

\opsig{$\Clean_{\vk}(\vk; \V{VK}, \V{Ver})$}:
\begin{algorithmic}[1]
\algcont
\State \textbf{return} $\vk \in \Keys \ \wedge\ \neg\exists\, P : \V{VK}[P] = \vk$
\State \algand\ $\neg\exists\, \msg,\sigma : \V{Ver}[\vk,\msg,\sigma] = 1$
\algsave
\end{algorithmic}
\vspace{1.4ex}

\opsig{$\Clean_{\sigma}(\sigma; \msg, \vk, \V{Ver})$}:
\begin{algorithmic}[1]
\algcont
\State \textbf{return} $\sigma \in \Sigs \wedge \V{Ver}[\vk,\msg,\sigma] \neq 0$
\algsave
\end{algorithmic}
\end{multicols}}
\end{interface}

\section{ویژگی‌ها}\label{sec:sigprops}

Definition~\ref{def:prop} is a shape, and a shape is worth an instance. We give one game $\Gm$ over
$\Fsig$ carrying two properties: $\opl{FinCor}$, saying that what the signing interface produces
verifies, and $\opl{FinUF}$, saying that nothing else does. They are two finalizations of one shell
rather than two games, which is what the shape now allows and what these two ask for --- they read
the same trace, of the same oracles, and differ in the verdict alone. The environment chooses which
it is judged on by choosing which to call, and calling one closes the game against the other.

Both take the \emph{search} reading of Definition~\ref{def:propadv},
\[
  \begin{gathered}
  \padv{\op{FinCor}}{\Gm,\Fsig,\Zenv,\mathcal{X}} \;:=\; \Pr[\,\Win_{\op{FinCor}}\,] , \\[3pt]
  \padv{\op{FinUF}}{\Gm,\Fsig,\Zenv,\mathcal{X}} \;:=\; \Pr[\,\Win_{\op{FinUF}}\,] ,
  \end{gathered}
\]
because winning either is meant to be impossible and not merely no better than a coin: a signature
scheme whose forgeries succeed half the time is broken, not average.

They are played over the game system $\gamma := \{\Gm\} \cup \{\Fsig\}$, with
$\Fsig.\admits := \{\Gm.\PID\}$ --- the game's process id, pinned, not its name --- and
$\Fsig.\Ps \subseteq \Gm.\Ps$, so the relays exist at every party $\Fsig$ serves and, at the
root, address a party it does not serve and are answered $\rej$ at line~\ref{ln:interface} --- the
root being where the game finalizes, not where it plays. Inside an interface at $\id$ we write
$\id_{\Fsig}$ for $(\Fsig.\PID,\id.P)$, the interface of $\Fsig$ at the party the game was called at.
That is tightness at $\{\Gm\}$, which Definition~\ref{def:prop}
requires, and here it earns its keep twice over. It makes $\gamma$ a legal game system; and it makes
the trace \emph{complete}. No caller but $\Gm$ passes line~\ref{ln:caller} at $\Fsig$: not the
environment, whose $Z$-identities and fresh names lie outside a pinned $\admits$; not the adversary,
whose claim is an $A$-identity; and not a machine the environment hosts at an invented instance of the
game's name, which is what pinning the id rather than the name refuses. So every signature $\Fsig$
ever issues passes through line~\ref{ln:gsign} below and is recorded. A game whose verdict is
``this was not obtained from the signing interface'' has no other way to know.

\begin{interface}[nofloat]{پوستهٔ بازی روی $\Fsig$؛ نهایی‌سازی‌ها در پی می‌آیند \\
\params{$\PID := (\Gm.F,0,0)$, \quad $\Ps := \Fsig.\Ps \cup \{Z\}$, \quad $\admits := \Zpid$, \quad
$\uses := \{(\Fsig,\serves_{\Gm})\}$, \quad $\pars := \none$}}

{\scriptsize
\setlength{\columnsep}{1.4em}
\begin{multicols}{2}
\opsig{$\op{Initialize}()$}:
\begin{algorithmic}[1]
\setcounter{contline}{0}\algcont
\State $\done \gets \false$
\State $\V{tr} \gets (\,)$
  \Comment{رد، یک دنباله}
\algsave
\end{algorithmic}
\vspace{1.4ex}

\opsig{$\id.\opdef{Gen}()$} \quad از $\id'$
\begin{algorithmic}[1]
\algcont
\Req{$\neg\done$}
\State $\vk \gets \id_{\Fsig}.\fopl{Gen}() \text{ as } \id$
\State $\V{tr} \gets \V{tr} \cdot (\atom{gen},\id.P,\vk)$
\State \textbf{return} $\vk$
\algsave
\end{algorithmic}
\vspace{1.4ex}

\opsig{$\id.\opdef{Sign}(\msg)$} \quad از $\id'$
\begin{algorithmic}[1]
\algcont
\Req{$\neg\done$}
\State $\sigma \gets \id_{\Fsig}.\fopl{Sign}(\msg) \text{ as } \id$ \label{ln:gsign}
\State $\V{tr} \gets \V{tr} \cdot (\atom{sign},\id.P,\msg,\sigma)$
\State \textbf{return} $\sigma$
\algsave
\end{algorithmic}
\vspace{1.4ex}

\columnbreak

\opsig{$\id.\opdef{Verify}(\vk,\msg,\sigma)$} \quad از $\id'$
\begin{algorithmic}[1]
\algcont
\Req{$\neg\done$}
\State $b \gets \id_{\Fsig}.\fopl{Verify}(\vk,\msg,\sigma)$
\Statex \quad $\text{as } \id$
\State $\V{tr} \gets \V{tr} \cdot (\atom{ver},\id.P,\vk,\msg,\sigma,b)$
\State \textbf{return} $b$
\algsave
\end{algorithmic}
\vspace{1.4ex}

\opsig{$\id.\op{Leak}()$} \quad از $\id'$
\begin{algorithmic}[1]
\algcont
\State \textbf{return} $\none$
\algsave
\end{algorithmic}
\vspace{1.4ex}

\end{multicols}}
\end{interface}

\begin{interface}[nofloat]{The finalizations of $\Gm$; the shell over $\Fsig$ is shared}
\opsig{$\id.\opdef{FinCor}()$} \quad from $\id'$ \quad (correctness)
\begin{algorithmic}[1]
\algcont
\Req{$\neg\done$}
\Req{$\id.P = Z$}
\State $\done \gets \true$
\State $\Cs \gets (C,Z).\fopl{Status}() \text{ as } \id$
\State $w \gets 1$ if $\exists\, P, Q, \vk, \msg, \sigma :\ \
       P \notin \Cs \ \wedge\ Q \notin \Cs \ \wedge\ \sigma \in \Sigs$ \label{ln:gcor}
\State \quad ${}\wedge\ (\atom{gen},P,\vk) \in \V{tr}
       \ \wedge\ (\atom{sign},P,\msg,\sigma) \in \V{tr}
       \ \wedge\ (\atom{ver},Q,\vk,\msg,\sigma,0) \in \V{tr}$;
\State else $w \gets 0$
\State \textbf{return} $w$
\algsave
\end{algorithmic}
\vspace{1.4ex}

\opsig{$\id.\opdef{FinUF}()$} \quad from $\id'$ \quad (unforgeability)
\begin{algorithmic}[1]
\algcont
\Req{$\neg\done$}
\Req{$\id.P = Z$}
\State $\done \gets \true$
\State $\Cs \gets (C,Z).\fopl{Status}() \text{ as } \id$
\State $w \gets 1$ if $\exists\, P, Q, \vk, \msg, \sigma :\ \
       P \notin \Cs \ \wedge\ Q \notin \Cs$ \label{ln:guf}
\State \quad ${}\wedge\ (\atom{gen},P,\vk) \in \V{tr}
       \ \wedge\ (\atom{ver},Q,\vk,\msg,\sigma,1) \in \V{tr}
       \ \wedge\ (\atom{sign},P,\msg,\sigma) \notin \V{tr}$;
\State else $w \gets 0$
\State \textbf{return} $w$
\algsave
\end{algorithmic}
\end{interface}

Four things about the code. The trace is a sequence, $\in$ means occurrence in it, and no verdict
reads the order. The verdicts read $\V{tr}$ and the register and call nothing, which is why
the root exemption in $\serves_{\Gm}$ suffices and $\Fsig$ need not serve $Z$. Both ask
$P \notin \Cs$ at the end rather than at the time of signing, and monotonicity
(Section~\ref{sec:fullinterface}) makes that the stronger reading: a party honest when the game closes
was honest throughout. Correctness asks $\sigma \in \Sigs$ because $\op{Sign}$ answers $\none$ where no
key has been generated, and $\none$ is no signature to hold against the functionality. And
unforgeability is \emph{strong}: line~\ref{ln:guf} excludes the triple, not the message, so a second
signature on a signed message counts as a forgery --- which is the notion
Chapter~\ref{sec:fsig}'s sanitizer was written to give.

\begin{proposition}[$\Fsig$ has both properties]\label{prop:sigprops}
For $\op{Fin}$ either of $\op{FinCor}$ and $\op{FinUF}$, for every
$\Zenv \in \ZenvSet_{\gamma,\gamma}$
and every occupant $\mathcal{X}$ of the adversary slot, $\Pr[\Win_{\op{Fin}}] = 0$. So $\{\Fsig\}$ has
each
property within $0$ against $\bigl(\ZenvSet_{\gamma,\gamma},\Xset\bigr)$ in the sense of
Definition~\ref{def:propadv}, for $\Xset$ the class of all occupants.
\end{proposition}

\begin{proof}
% BEGIN overview prop:sigprops
One preliminary and two verdicts. The preliminary is that both parties named by a verdict must be
honest, and it is not bookkeeping: at a corrupt party the core does not run and the answer is the
adversary's, so a mediated $\op{Verify}$ could report a $0$ on a good signature or a $1$ on a
forgery, and either finalization would fire with $\Fsig$ having done nothing wrong. Correctness then
follows from $\op{Sign}$ recording $\V{Ver}[\vk,\msg,\sigma] \gets 1$ and $\op{Verify}$ returning a
recorded entry at its first line. Unforgeability splits on whether $\op{Verify}$ read an entry or
computed one: it cannot have computed a $1$ under an honest party's key, so the entry was recorded,
and only $\op{Sign}$ records one.
% END overview prop:sigprops
Take parties honest at the close. Monotonicity makes them honest throughout, so no call the game
places at them is mediated, line~\ref{ln:medin} firing only at a corrupt party and the game's own name
lying in $\bits$. Both verdicts ask this of the signer $P$ \emph{and} of the verifier $Q$, and both
need it of both: at a corrupt party the core does not run and the answer is the adversary's, so a
mediated $\op{Verify}$ reports whatever it likes --- a $0$ on a good signature, or a $1$ on a
forgery --- and either finalization would return $1$ with $\Fsig$ having done nothing wrong.

For correctness, suppose $(\atom{sign},P,\msg,\sigma) \in \V{tr}$ with $\sigma \in \Sigs$. The core of
$\op{Sign}$ ran, so it sanitized $\sigma$ against $\Clean_{\sigma}$ and then set
$\V{Ver}[\vk,\msg,\sigma] \gets 1$ for $\vk = \V{VK}[P]$, which $\op{Gen}$ fixed once and never
changes. Any later $\op{Verify}(\vk,\msg,\sigma)$ finds that entry set and returns it at its first
line, so it returns $1$; thus $(\atom{ver},Q,\vk,\msg,\sigma,0) \notin \V{tr}$ for every $Q$, and
line~\ref{ln:gcor} gives $0$.

For unforgeability, suppose $(\atom{ver},Q,\vk,\msg,\sigma,1) \in \V{tr}$ with
$\vk = \V{VK}[P]$ for an honest $P$. Either $\op{Verify}$ read a recorded entry or it computed one.
It cannot have computed a $1$: the last test before writing $\V{Ver}$ sets $b \gets 0$ whenever
$\vk = \V{VK}[P']$ for some $P' \notin \Cs$, which $P$ witnesses. So the entry was recorded, and only
$\op{Sign}$ records a $1$ --- at the party whose key is $\vk$, that party being $P$ since
$\Clean_{\vk}$ issues no key twice. Hence $(\atom{sign},P,\msg,\sigma) \in \V{tr}$ and
line~\ref{ln:guf} gives $0$.
\end{proof}

The bound is not the interesting part; that it is $0$ merely says the ideal functionality was written
to have these properties, as an ideal functionality should be. The content is what
Corollary~\ref{cor:proptransfer2} then does with it, and the search reading is what makes that
useful: a system $\pi$ put in place of $\{\Fsig\}$ wins either finalization with probability at most
$\varepsilon$, against every adversary and not only the dummy, and under the search reading that
probability \emph{is} the advantage, so $\varepsilon$ bounds the property itself. Had the decision
reading been taken the same conclusion would read $2\varepsilon$ on $2\Pr[\Win_{\op{Fin}}]-1$, which
for a game meant to be unwinnable is no constraint at all --- a forger
succeeding half the time meets it. That is the calibration rule of Chapter~\ref{sec:properties}
applied where it bites, and why Proposition~\ref{prop:proptransfer} carries both displays: each
finalization inherits at the scale its own reading fixes, and these two inherit at $\varepsilon$.

Four conditions come with the corollary and none is discharged by emulation alone. The replacement of
$\{\Fsig\}$ by $\pi$ must be well-formed; $\{\Gm\} \cup \pi$ must again be a game system, which asks
$\pi$'s members to pin $\Gm.\PID$ where $\Fsig$ did, since tightness is part of
Definition~\ref{def:prop}; the simulator and the adversary must be refusal-oblivious; and the
environments must silence no $Z$-identity. A realization is thus never separately argued to be correct
or unforgeable --- it inherits both from the one emulation proof, together with those four checks, and
the two finalizations above are what the inheritance is \emph{of}.

One caveat travels with it, and this is where it bites. Corollary~\ref{cor:proptransfer2} asks that
the cores of the game system place no responsive call, and $\Fsig$ as written reaches the slot through
$\Adv(\cdot)$, so the corollary applies. Should one adopt Chapter~\ref{app:responsive}'s advice and
write $\Adv^{!}$ in $\op{Gen}$, $\op{Sign}$ and $\op{Verify}$ --- which buys locality --- the transfer
must be argued directly instead, for the reason Remark~\ref{rem:respdummy} gives. Locality and
inheritance are, as things stand, priced against each other.

%% ==================================================================
%% ==================================================================
\chapter{زیرساختِ کلید عمومی}\label{sec:pki}

\section{کارکرد}\label{sec:pkifunc}

$\Fsig$ said plainly that it does not authenticate the key it hands out --- $\opl{Verify}$ takes a
$\vk$ of its own, unconnected to any caller, because verification keys are not assumed
authenticated (Section~\ref{sec:fsigfunc}). The functionality $\Gpki$ below is the setup that closes the gap, in the shape of
Canetti's own certification functionality~\cite{Canetti04}: a
single, shared directory binding a key to the party that claims it, so that a $\vk$ handed to
$\op{Verify}$ can be trusted because it came from here rather than because it was simply believed.

What it does is quickly said. Each party may register one key, and only one: $\opl{Register}$
takes the key as \emph{input} --- unlike $\opl{Gen}$, which conjures one from the slot, this is a
directory and not a factory --- and keeps the first key it is given, sanitized against every key
already on file so that no two parties ever hold the same one. The operation $\opl{Retrieve}$ answers anyone with
whatever is on file for the party they name, or nothing if no one has registered.

\heading{The fields} They are $\Gclock$'s, and for the same reason. The process id is a constant,
$(\op{G}\op{PKI},0,0)$, so Definition~\ref{def:sys} caps the system at one copy: a second
directory would be a second notion of who owns which key, exactly what a PKI exists to rule out.
$\admits := \Stdpid \cup \Apid \cup \Zpid$ makes $\op{global}(\Gpki)$ hold by the same disjunct as
the clock's, and $\Ps$ is again the universe the directory covers rather than who has used it ---
registration is state, in $\V{VK}$, not membership. One field is simpler than the clock's:
$\uses := \emptyset$. The operations $\op{Register}$ and $\op{Retrieve}$ place no call at all, the key arriving
with the request rather than from the slot, so there is nothing $\Gpki$ depends on beyond the
ambient machinery every functionality already has.

\heading{Two borrowed guards} $\op{Register}$ takes one line from each of the two examples before
it. The first is $\Gclock$'s: a caller named $A$ or $Z$ acting for an \emph{honest} party is
answered that party's own key and changes nothing, exactly as it is answered the time rather than
allowed to tick on an honest party's behalf --- an outsider may ask what is on file, not decide
what goes there. At a corrupt party the check does not fire, so the call proceeds --- the adversary
registers on a corrupt party's behalf as freely as it ticks one, and whatever key it submits is
what the directory will show, mediation letting exactly this call through unmediated as it does
for $\Fac$'s corrupt sends (Section~\ref{sec:facfunc}). The second is $\Fsig$'s: the first call to
succeed wins, later ones returning the key already on file rather than overwriting it, exactly as
$\op{Gen}$ never re-issues a key once one exists. Sanitizing here differs from both signature
predicates in what it repairs: $\Clean_{\vk}$ and $\Clean_{\sigma}$ each fix a value the \emph{slot}
proposed; $\Clean_{\op{reg}}$ fixes one the \emph{caller} supplied, which is the more exposed
position --- an honest party's own key is presumably already good, but a corrupt one's is whatever
the adversary wrote in, and the one thing a directory cannot let through is a second party
claiming a key the first already holds. A good value always exists here for the same reason it
does for $\Clean_{\vk}$: $\Keys$ is infinite and the keys excluded are exactly those a bounded run
can have recorded, finitely many (footnote of Chapter~\ref{sec:sanitization}).

\begin{interface}[nofloat]{کارکردِ $\Gpki$ \\
\params{$\PID := (\op{G}\op{PKI},0,0)$, \quad $\Ps$, \quad
$\admits := \Stdpid \cup \Apid \cup \Zpid$, \quad $\uses := \emptyset$, \quad $\pars := \none$}}

{\scriptsize
\setlength{\columnsep}{1.4em}
\begin{multicols}{2}
\opsig{$\op{Initialize}()$}:
\begin{algorithmic}[1]
\setcounter{contline}{0}\algcont
\State $\V{VK} : \Gpki.\Ps \to \Keys \cup \{\unset\}$
\State $\V{VK}[*] \gets \unset$
\algsave
\end{algorithmic}
\vspace{1.4ex}

\opsig{$\id.\op{Register}(\vk)$} \quad از $\id'$
\begin{algorithmic}[1]
\algcont
\If{$\id'.F \in \{A,Z\} \ \wedge\ \id.P \notin \Cs$}
  \State \textbf{return} $\V{VK}[\id.P]$ \label{ln:pkipeek}
    \Comment{یک نگاه، نه یک نوشتن}
\EndIf
\If{$\V{VK}[\id.P] \neq \unset$}
  \State \textbf{return} $\V{VK}[\id.P]$ \label{ln:pkifixed}
\EndIf
\State $\vk \gets \San[\Clean_{\op{reg}}](\vk; \V{VK})$
\State $\V{VK}[\id.P] \gets \vk$ \label{ln:pkiwrite}
\State \textbf{return} $\vk$
\algsave
\end{algorithmic}
\vspace{1.4ex}

\columnbreak

\opsig{$\id.\op{Retrieve}(P)$} \quad از $\id'$
\begin{algorithmic}[1]
\algcont
\State \textbf{return} $\V{VK}[P]$
\algsave
\end{algorithmic}
\vspace{1.4ex}

\opsig{$\id.\op{Leak}()$} \quad از $\id'$
\begin{algorithmic}[1]
\algcont
\State \textbf{return} $\V{VK}$
\algsave
\end{algorithmic}
\vspace{1.4ex}

\opsig{$\Clean_{\op{reg}}(\vk; \V{VK})$}:
\begin{algorithmic}[1]
\algcont
\State \textbf{return} $\vk \in \Keys \ \wedge\ \neg\exists\, P : \V{VK}[P] = \vk$
\algsave
\end{algorithmic}
\end{multicols}}
\end{interface}

Three remarks on the code. Registration is recorded against the party, not the caller, for the
same reason ticks are: $\opl{Guard}$ has already tied the call to $\id.P$ by line~\ref{ln:caller},
so whichever machine holds the token when $P$ registers, the key lands on $P$'s own entry. There
is no $\op{Deregister}$, and none is added: a real PKI can support revocation, but a key once
bound is the simplest directory there is, and revocation is a genuine addition left for elsewhere
rather than a gap papered over. And $\op{Leak}$ answers with the whole table, which is no more
than $\op{Retrieve}$ already gives out one entry at a time --- the directory is public at every
one of its interfaces, not only this one, so $\op{Leak}$ tells the adversary nothing $\op{Retrieve}$,
asked enough times, would not.

Set against both examples before it, $\Gpki$ needs less than either. No value is the simulator's
to choose out of nothing, as in $\op{Gen}$ --- the key arrives with the call --- so sanitizing here
repairs a caller's input rather than a slot's output, the one place this differs from every $\San$
before it. And no call leaves the core at all, so there is no notification to be responsive about
and no price for Chapter~\ref{app:responsive} to settle: the two properties below transfer as
cheaply as a property can.

\section{ویژگی‌ها}\label{sec:pkiprops}

$\Gclock$ carried two properties for having two kinds of promise --- a local one, that the counter
moves as coded, and a global one, that it moves only when everyone honest has agreed. The directory $\Gpki$
splits the same way, and into two finalizations of one game. $\opl{FinFix}$, \emph{fixity}, says a
registered key never changes --- the local
promise, about one party's own entry. $\opl{FinBind}$, \emph{binding}, says no two parties ever
hold the same key --- the global one, quantified over all of $\Gpki.\Ps$ at once. Both take the
search reading of Definition~\ref{def:propadv}.

Both are played over the game system $\gamma := \{\Gm\} \cup \{\Gpki\}$, with $\Gpki$ taken as in
Section~\ref{sec:pkifunc} save one field: $\admits := \{\Gm.\PID\}$, pinned --- Remark~\ref{rem:tight}'s
forcing, exactly as for the clock (Section~\ref{sec:clockprops}), and for the identical reason: a
property of a global $\F$ is a property against an adversary that shares it, and fixity and
binding are both claims about \emph{every} key there is, certifiable only where the game's trace
is all there is.

Pinned, the directory is no longer global, so line~\ref{ln:session} asks for one session and the
experiment's own session $0$ meets it. No caller but $\Gm$ passes line~\ref{ln:caller}: not the
environment, not a machine it hosts, and not the adversary's $A$-identities, refused before
line~\ref{ln:pkipeek} could even apply. So every key $\Gpki$ ever holds enters through the relay
below and is recorded --- a game whose verdict is ``two parties share a key'' has no other way to
know it. Inside an interface at $\id$ we write $\id_{\Gpki}$ for $(\Gpki.\PID,\id.P)$.

\begin{interface}[nofloat]{پوستهٔ بازی روی $\Gpki$؛ نهایی‌سازی‌ها در پی می‌آیند \\
\params{$\PID := (\Gm.F,0,0)$, \quad $\Ps := \Gpki.\Ps \cup \{Z\}$, \quad $\admits := \Zpid$, \quad
$\uses := \{(\Gpki,\serves_{\Gm})\}$, \quad $\pars := \none$}}

{\scriptsize
\setlength{\columnsep}{1.4em}
\begin{multicols}{2}
\opsig{$\op{Initialize}()$}:
\begin{algorithmic}[1]
\setcounter{contline}{0}\algcont
\State $\done \gets \false$
\State $\V{tr} \gets (\,)$
  \Comment{رد، یک دنباله}
\algsave
\end{algorithmic}
\vspace{1.4ex}

\opsig{$\id.\opdef{Register}(\vk)$} \quad از $\id'$
\begin{algorithmic}[1]
\algcont
\Req{$\neg\done$}
\Req{$\id.P \neq Z$}
\State $\vk' \gets \id_{\Gpki}.\fopl{Register}(\vk) \text{ as } \id$ \label{ln:gpkireg}
\State $\V{tr} \gets \V{tr} \cdot (\atom{reg},\id.P,\vk')$
\State \textbf{return} $\vk'$
\algsave
\end{algorithmic}
\vspace{1.4ex}

\opsig{$\id.\op{Leak}()$} \quad از $\id'$
\begin{algorithmic}[1]
\algcont
\State \textbf{return} $\none$
\algsave
\end{algorithmic}

\columnbreak

\opsig{$\id.\opdef{Retrieve}(P)$} \quad از $\id'$
\begin{algorithmic}[1]
\algcont
\Req{$\neg\done$}
\Req{$\id.P \neq Z$}
\State $\vk' \gets \id_{\Gpki}.\fopl{Retrieve}(P) \text{ as } \id$
\State \textbf{return} $\vk'$
\algsave
\end{algorithmic}
\vspace{1.4ex}

\end{multicols}}
\end{interface}

\begin{interface}[nofloat]{The finalizations of $\Gm$; the shell over $\Gpki$ is shared}
\opsig{$\id.\opdef{FinFix}()$} \quad from $\id'$ \quad (fixity)
\begin{algorithmic}[1]
\algcont
\Req{$\neg\done$}
\Req{$\id.P = Z$}
\State $\done \gets \true$
\State $\Cs \gets (C,Z).\fopl{Status}() \text{ as } \id$
\State $w \gets 1$ if $\exists\, P,\vk_1,\vk_2 :\ \
       P \notin \Cs \ \wedge\ \vk_1 \neq \vk_2$ \label{ln:gpkifix}
\State \quad ${}\wedge\ (\atom{reg},P,\vk_1) \in \V{tr} \ \wedge\ (\atom{reg},P,\vk_2) \in \V{tr}$;
\State else $w \gets 0$
\State \textbf{return} $w$
\algsave
\end{algorithmic}
\vspace{1.4ex}

\opsig{$\id.\opdef{FinBind}()$} \quad from $\id'$ \quad (binding)
\begin{algorithmic}[1]
\algcont
\Req{$\neg\done$}
\Req{$\id.P = Z$}
\State $\done \gets \true$
\State $\Cs \gets (C,Z).\fopl{Status}() \text{ as } \id$
\State $w \gets 1$ if $\exists\, P,Q,\vk :\ \
       P \notin \Cs \ \wedge\ Q \notin \Cs \ \wedge\ P \neq Q$ \label{ln:gpkibind}
\State \quad ${}\wedge\ (\atom{reg},P,\vk) \in \V{tr} \ \wedge\ (\atom{reg},Q,\vk) \in \V{tr}$;
\State else $w \gets 0$
\State \textbf{return} $w$
\algsave
\end{algorithmic}
\end{interface}

Two things about the code. The operation $\op{Retrieve}$ is relayed and logged though neither verdict reads its
trace: a game that could not look anyone's key up would not be exercising $\Gpki$ at all, and
leaving the interface whole costs nothing, exactly as $\Gclock$'s shell relayed $\op{Read}$ for
games that read only the mutators. And $\op{FinFix}$'s verdict does not need $\vk_1$ and $\vk_2$ ordered,
unlike steadiness's reading of its trace: idempotence means at most one of the two calls that
produced them can be the write, the other only a later read of it, and either order loses the game
identically, a changed key being a changed key whichever one came first.

\begin{proposition}[$\Gpki$ has fixity]\label{prop:pkifix}
For every $\Zenv \in \ZenvSet_{\gamma,\gamma}$ and every occupant $\mathcal{X}$ of the adversary
slot, $\Pr[\Win_{\op{FinFix}}] = 0$. So $\{\Gpki\}$ has fixity within $0$ against
$\bigl(\ZenvSet_{\gamma,\gamma},\Xset\bigr)$ in the sense of Definition~\ref{def:propadv}, for
$\Xset$ the class of all occupants.
\end{proposition}

\begin{proof}
Suppose $(\atom{reg},P,\vk_1),(\atom{reg},P,\vk_2) \in \V{tr}$ with $P \notin \Cs$. By tightness both
were logged at line~\ref{ln:gpkireg}, each running $\Gpki.\op{Register}$'s own core at $\id.P = P$:
the relay's own name is $\Gm$'s, not $A$ or $Z$, so line~\ref{ln:pkipeek} never intercepts it. The
core's only write is line~\ref{ln:pkiwrite}, reached only while $\V{VK}[P] = \unset$, and once
reached it fixes $\V{VK}[P]$ to some sanitized value for good --- every later call at $P$ finds
$\V{VK}[P] \neq \unset$ and returns line~\ref{ln:pkifixed}'s value, the same one, unchanged. So the
two calls producing $\vk_1$ and $\vk_2$ agree: whichever ran first wrote the value, and the other
only read it back. The verdict at line~\ref{ln:gpkifix} gives $0$.
\end{proof}

\begin{proposition}[$\Gpki$ has binding]\label{prop:pkibind}
For every $\Zenv \in \ZenvSet_{\gamma,\gamma}$ and every occupant $\mathcal{X}$ of the adversary
slot, $\Pr[\Win_{\op{FinBind}}] = 0$. So $\{\Gpki\}$ has binding within $0$ against
$\bigl(\ZenvSet_{\gamma,\gamma},\Xset\bigr)$ in the sense of Definition~\ref{def:propadv}, for
$\Xset$ the class of all occupants.
\end{proposition}

\begin{proof}
Suppose $(\atom{reg},P,\vk),(\atom{reg},Q,\vk) \in \V{tr}$ with $P,Q \notin \Cs$ and $P \neq Q$. By
Proposition~\ref{prop:pkifix}'s argument $\vk = \V{VK}[P]$ and $\vk = \V{VK}[Q]$ from the moment
each was first written onward. Say $Q$'s write came second, the case $P$'s symmetric; at that
moment $\op{Register}$'s line~\ref{ln:pkiwrite} ran $\vk \gets \San[\Clean_{\op{reg}}](\cdot;\V{VK})$ with $\V{VK}[P]$
already set, and $\San$'s guarantee (Chapter~\ref{sec:sanitization}) is that its output satisfies
$\Clean_{\op{reg}}$ whatever the input --- membership in $\Keys$ and freshness against every key
already on file, $P$'s among them. So $Q$'s write cannot equal $\V{VK}[P]$, contradicting
$\vk = \V{VK}[P] = \V{VK}[Q]$. The verdict at line~\ref{ln:gpkibind} gives $0$.
\end{proof}

Both transfer by Corollary~\ref{cor:proptransfer2} with nothing extra to argue. The hypothesis it
asks --- that the game system's cores place no responsive call --- is met the cheap way: $\Gpki$'s
cores place no call of any kind, responsive or not, so Remark~\ref{rem:respdummy}'s caveat never
arises and dummy completeness (Section~\ref{app:dummy}) is untouched. Unlike $\Gclock$, whose
notification bought locality at inheritance's expense, and unlike $\Fnet$ and $\Fac$, whose slot
calls priced their own transfer, $\Gpki$ pays nothing: a system $\pi$ put in place of $\{\Gpki\}$
inherits fixity and binding directly, the four conditions of Section~\ref{sec:sigprops}'s closing
paragraph being the only ones left to discharge. A directory that asks the adversary nothing has
nothing to spend.

%% ==================================================================
%% ==================================================================
\chapter{ساعتِ جهانی}\label{sec:gclock}

\section{کارکرد}\label{sec:gclockfunc}

$\Fsig$ is local and randomized, and most of what was said about it concerned the simulator. The
second example is its opposite on both counts. A clock functionality was introduced by Katz,
Maurer, Tackmann and Zikas to model synchrony in UC~\cite{KMTZ13}; recast by Badertscher, Maurer,
Tschudi and Zikas as a \emph{shared} setup around a monotone counter~\cite{BMTZ17}, it has become
the canonical global functionality, the one the composable analyses of Bitcoin and its successors
keep time by. The clock $\Gclock$ below is that functionality in the language of this paper.

What it does is quickly said. The counter $\V{now}$ starts at $0$ and only ever steps forward.
Anyone the clock serves may read it. Parties enter and leave the round structure through
$\op{Register}$ and $\opl{Deregister}$, and the counter steps when every honest registered party
has asked it to: $\opl{Update}$ records a tick against the calling party, and once the honest
registered ticks are all in, the test on line~\ref{ln:advance} advances the counter and clears
them. A round ends exactly when the last honest party in it is done, and the adversary can
neither stall the round nor close it early --- how the code refuses each is below. Each recorded
tick is reported to the adversary before $\op{Update}$ returns, on line~\ref{ln:clocknotify},
\emph{responsively} (Chapter~\ref{app:responsive}) and with the answer discarded; what the
notification buys and what it costs is settled at the end of Section~\ref{sec:clockprops}.
Every operation answers the counter as it stands on return, so the tick that closes a round
learns the new time and the ticks before it the old; a party that must know where it stands
reads again.

\heading{The fields} The process id is a constant, $(\op{G}\op{Clock},0,0)$, where
$\Fsig$'s was a parameter. Definition~\ref{def:sys} then puts at most one copy in any system,
which is the point: global means one instance, shared, and a second clock would be a second time.
$\admits := \Stdpid \cup \Apid \cup \Zpid$ makes $\op{global}(\Gclock)$ hold by its first disjunct,
so line~\ref{ln:session} of $\opl{Guard}$ waives the session check and every session reaches the
one instance --- the very separation the session component exists to make, deliberately not made
here. A protocol that keeps time declares $(\Gclock,\serves)$ in its $\uses$, and the global
disjunct of Definition~\ref{def:wfsys} is what lets members of every session meet the entry at
the shared copy. And $\uses := \{(\Adv,\serves)\}$ for the one call the code places, the notification
of line~\ref{ln:clocknotify}. And $\Ps$ stays a parameter but as the universe rather than the
membership: it fixes the parties an interface exists for, while who is \emph{in} the round
structure at a given moment is state, entered and left through $\op{Register}$ and
$\op{Deregister}$ as in~\cite{BMTZ17} --- dynamic membership is one more table, not a new kind
of field.

\heading{Reading and ticking} $\admits$ is a field of the functionality, not of an operation, so
one set governs every core, and the set above is the right one for $\opl{Read}$: the adversary and
the environment may ask the time --- \cite{BMTZ17} lets both, and the environment reading the
shared clock is much of what makes a global setup global. It is too wide for the three mutators,
where an adversary free to act for a party could close a round the party had not finished, or
drop it from the round structure altogether. The refinement sits as the first line of each core:
a caller named $A$ or $Z$ acting for an \emph{honest} party is answered the time and changes
nothing, while at a corrupt party it passes --- the adversary registers, deregisters and ticks
for the parties it owns, as in~\cite{BMTZ17}, and those ticks are inert where it matters,
line~\ref{ln:advance} never reading a corrupt party's flag. One consequence deserves naming:
while no honest party is registered the test of line~\ref{ln:advance} is vacuous, so the
adversary may register a corrupt party and tick it freely --- until the first honest
registration, time is the adversary's to spend, there too as in~\cite{BMTZ17}.

The quantifier of line~\ref{ln:advance} reads ``every honest registered party'', and it is worth
seeing how little of it $\Gclock$ enforces itself. At a corrupt party a standard caller's
$\op{Update}$ is handed to the adversary by line~\ref{ln:medin} of the full interface before the
core runs, and the adversary's own call at the same party is admitted by line~\ref{ln:clockaz};
either way a corrupt party's ticks are the adversary's to give or withhold, and the $\forall$ ---
which reads the register by the convention of Section~\ref{sec:corruption} --- ignores them
regardless. Corrupt parties therefore neither hold a round open nor help close one. At an honest
party both routes shut: $\opl{Mediate}$ does not fire, and line~\ref{ln:clockaz} turns the
observers' ticks into reads. Monotonicity adds the last piece: a party corrupted after ticking
leaves a stale flag behind, but it is outside the quantifier from then on, and the next advance
clears it.

\begin{interface}[nofloat]{کارکردِ $\Gclock$ \\
\params{$\PID := (\op{G}\op{Clock},0,0)$, \quad $\Ps$, \quad
$\admits := \Stdpid \cup \Apid \cup \Zpid$, \quad $\uses := \{(\Adv,\serves)\}$, \quad $\pars := \none$}}

{\scriptsize
\setlength{\columnsep}{1.4em}
\begin{multicols}{2}
\opsig{$\op{Initialize}()$}:
\begin{algorithmic}[1]
\setcounter{contline}{0}\algcont
\State $\V{now} \gets 0$
\State $\V{reg},\, \V{tick} : \Gclock.\Ps \to \{0,1\}$
\State $\V{reg}[*] \gets 0$; \ $\V{tick}[*] \gets 0$
\algsave
\end{algorithmic}
\vspace{1.4ex}

\opsig{$\id.\op{Register}()$} \quad از $\id'$
\begin{algorithmic}[1]
\algcont
\If{$\id'.F \in \{A,Z\} \ \wedge\ \id.P \notin \Cs$}
  \State \textbf{return} $\V{now}$
\EndIf
\State $\V{reg}[\id.P] \gets 1$; \ $\V{tick}[\id.P] \gets 0$
  \Comment{عضوِ تازه یک تیک بدهکار است}
\State \textbf{return} $\V{now}$
\algsave
\end{algorithmic}
\vspace{1.4ex}

\opsig{$\id.\op{Deregister}()$} \quad از $\id'$
\begin{algorithmic}[1]
\algcont
\If{$\id'.F \in \{A,Z\} \ \wedge\ \id.P \notin \Cs$}
  \State \textbf{return} $\V{now}$
\EndIf
\State $\V{reg}[\id.P] \gets 0$
\State \textbf{return} $\V{now}$
\algsave
\end{algorithmic}
\vspace{1.4ex}

\columnbreak

\opsig{$\id.\opdef{Update}()$} \quad از $\id'$
\begin{algorithmic}[1]
\algcont
\If{$\id'.F \in \{A,Z\} \ \wedge\ \id.P \notin \Cs$}
  \State \textbf{return} $\V{now}$ \label{ln:clockaz}
    \Comment{خواندن، نه تیک‌زدن}
\EndIf
\If{$\V{reg}[\id.P] = 0$}
  \State \textbf{return} $\V{now}$
    \Comment{در ساختارِ دور نیست}
\EndIf
\State $\V{tick}[\id.P] \gets 1$
\If{$\forall\, P \in \Gclock.\Ps \setminus \Cs \, : \, \V{reg}[P] = 1 \Rightarrow \V{tick}[P] = 1$} \label{ln:advance}
  \State $\V{now} \gets \V{now} + 1$
  \State $\V{tick}[*] \gets 0$
    \Comment{دور تمام شد}
\EndIf
\State $\Adv^{!}\bigl(\id.\op{Update},\ \V{now}\bigr)$ \label{ln:clocknotify}
  \Comment{خبر بده؛ پاسخ ریخته می‌شود}
\State \textbf{return} $\V{now}$
\algsave
\end{algorithmic}
\vspace{1.4ex}

\opsig{$\id.\op{Read}()$} \quad از $\id'$
\begin{algorithmic}[1]
\algcont
\State \textbf{return} $\V{now}$
\algsave
\end{algorithmic}
\vspace{1.4ex}

\opsig{$\id.\op{Leak}()$} \quad از $\id'$
\begin{algorithmic}[1]
\algcont
\State \textbf{return} $(\V{now},\ \V{reg},\ \V{tick})$
\algsave
\end{algorithmic}
\end{multicols}}
\end{interface}

Three remarks on the code. Registrations and ticks are recorded against the \emph{party}, not the
caller: any machine acting for $P$ may register or tick $P$, $\opl{Guard}$ having already tied the
call to the party on line~\ref{ln:caller}. \cite{BMTZ17} instead has every registered party and
functionality tick separately before the round closes; that discipline, if wanted, is tables
indexed by identity rather than party, and nothing else changes. The notification is
$\op{Update}$'s alone; \cite{BMTZ17} tells the adversary about registrations and departures too,
and extending line~\ref{ln:clocknotify}'s pattern to the other two mutators is mechanical. And
$\opl{Leak}$ returns the whole state, a clock keeping no secrets: the counter is anyone's to read
anyway, and the tables say only who is in the round and who is done with it.

Set against $\Fsig$, most of what the signature example needed is still absent, each absence
legible in the code. No value is the simulator's to choose --- the notification's answer is
discarded --- so there is nothing to sanitize. And where $\Fsig$ reached the slot through
$\Adv(\cdot)$ and re-tested its tables against what could happen while it hung, the clock's one
call is responsive by construction: every mutation of $\op{Update}$ precedes
line~\ref{ln:clocknotify}, and $\opl{Respond}$ keeps the token from leaving, so no instance is
ever re-entered mid-operation and there is nothing to re-test --- enforced against every occupant
of the slot, where Chapter~\ref{sec:fsig} could only advise a simulator class. What the
notification costs is recorded where it bites: Section~\ref{sec:clockprops} plays two properties
over the clock and prices it there.

\section{ویژگی‌ها}\label{sec:clockprops}

Section~\ref{sec:sigprops} instanced Definition~\ref{def:prop} once; the clock is worth a
second, not least because time is the one thing in this paper whose \emph{order} is the content.
We give one game $\Gm$ over $\Gclock$ carrying two properties: $\opl{FinStead}$,
\emph{steadiness}, saying that the counter
moves as the code promises --- forward, by one, and only when pushed --- and $\opl{FinUna}$,
\emph{unanimity}, saying that it moves only when every honest registered party has pushed. Both
take the search reading of Definition~\ref{def:propadv}, winning either being meant to be
impossible rather than hard.

They are played over the game system $\gamma := \{\Gm\} \cup \{\Gclock\}$, with the clock taken
as in Chapter~\ref{sec:gclock} save one field: $\admits := \{\Gm.\PID\}$, pinned. The change is
forced. Tightness at $\{\Gm\}$, which Definition~\ref{def:prop} requires, asks the members under
test to admit whole process ids, and a bare $\Stdpid$ is exactly what Remark~\ref{rem:tight}
refuses; Definition~\ref{def:propenv} says why in operational terms --- an environment shares a
global functionality, touching it as anyone may, and a property of a global $\F$ is a property
against an adversary that shares it. For these two properties sharing is not a nuisance but the
negation of the claim: steadiness and unanimity are assertions about \emph{all} the ticks there
are, and a game can certify them only where its trace is all there is. So the clock is played
privately, and the properties certified are of its \emph{code}, which sharing does not change ---
the situation of Section~\ref{sec:sigprops} exactly, where the $\Fsig$ that ships names its
protocol and the $\Fsig$ under test names the game.

Pinned, the clock is no longer global, so line~\ref{ln:session} of $\opl{Guard}$ asks for one
session; the experiment sits in session $0$ and the clock's process id carries $s = 0$, so the
check passes. And no caller but $\Gm$ now passes line~\ref{ln:caller} at the clock: not the
environment or a machine it hosts, whose claims lie outside a pinned $\admits$, and not the
adversary, whose $A$-identities are refused before line~\ref{ln:clockaz} could even answer them
the time. So the counter changes only inside line~\ref{ln:gtick} below and the membership only
inside the $\op{Register}$ and $\op{Deregister}$ relays beside it, and every change is
recorded --- a game whose verdict is ``the time moved when no one pushed'' has no other way to
know. Inside an interface at $\id$ we write $\id_{\Gclock}$ for $(\Gclock.\PID,\id.P)$, and every
trace entry carries the answer $\tau$ the clock gave.

One line has no analogue in Section~\ref{sec:sigprops}: the relays refuse the root. There the
relays at the root addressed a party $\Fsig$ does not serve and were answered $\rej$ at
line~\ref{ln:interface}, which the verdicts, existential over well-shaped tuples, never read.
Steadiness reads \emph{every} entry, so a $\rej$ of the game's own making cannot be allowed into
the trace; refused at the relay, nothing is appended, and every entry in $\V{tr}$ is an answer the
clock actually gave. Against a replacement this bites, and rightly: a $\pi$ answering an honest
read anything but a number loses steadiness on the spot, a clock that cannot say the time being
broken however it moves.

That line is the first thing a shared shell asks for that a single property would not. Unanimity
does not need it --- its verdict is existential, like $\Fsig$'s, and would read past a $\rej$ as
those did --- so the refusal is steadiness's demand alone, and the shell carries it because both
finalizations hang off the shell. A shell is built to the strictest of the verdicts it serves, and
the cost of that is paid by the others: unanimity is played on a trace cleaner than it needs. It
is a cost worth naming, because it is the one thing sharing a shell can take away. Where a
property would be \emph{weakened} by another's demand rather than merely over-served, the answer
is a second game and not a second finalization.

\begin{interface}[nofloat]{پوستهٔ بازی روی $\Gclock$؛ نهایی‌سازی‌ها در پی می‌آیند \\
\params{$\PID := (\Gm.F,0,0)$, \quad $\Ps := \Gclock.\Ps \cup \{Z\}$, \quad $\admits := \Zpid$, \quad
$\uses := \{(\Gclock,\serves_{\Gm})\}$, \quad $\pars := \none$}}

{\scriptsize
\setlength{\columnsep}{1.4em}
\begin{multicols}{2}
\opsig{$\op{Initialize}()$}:
\begin{algorithmic}[1]
\setcounter{contline}{0}\algcont
\State $\done \gets \false$
\State $\V{tr} \gets (\,)$
  \Comment{رد، یک دنباله}
\algsave
\end{algorithmic}
\vspace{1.4ex}

\opsig{$\id.\op{Register}()$} \quad از $\id'$
\begin{algorithmic}[1]
\algcont
\Req{$\neg\done$}
\Req{$\id.P \neq Z$}
  \Comment{ریشه تنها نهایی می‌کند}
\State $\tau \gets \id_{\Gclock}.\fopl{Register}() \text{ as } \id$ \label{ln:greg}
\State $\V{tr} \gets \V{tr} \cdot (\atom{reg},\id.P,\tau)$
\State \textbf{return} $\tau$
\algsave
\end{algorithmic}
\vspace{1.4ex}

\opsig{$\id.\opdef{Deregister}()$} \quad از $\id'$
\begin{algorithmic}[1]
\algcont
\Req{$\neg\done$}
\Req{$\id.P \neq Z$}
\State $\tau \gets \id_{\Gclock}.\fopl{Deregister}() \text{ as } \id$
\State $\V{tr} \gets \V{tr} \cdot (\atom{dereg},\id.P,\tau)$
\State \textbf{return} $\tau$
\algsave
\end{algorithmic}
\vspace{1.4ex}

\opsig{$\id.\opdef{Tick}()$} \quad از $\id'$
\begin{algorithmic}[1]
\algcont
\Req{$\neg\done$}
\Req{$\id.P \neq Z$}
\State $\tau \gets \id_{\Gclock}.\fopl{Update}() \text{ as } \id$ \label{ln:gtick}
\State $\V{tr} \gets \V{tr} \cdot (\atom{tick},\id.P,\tau)$
\State \textbf{return} $\tau$
\algsave
\end{algorithmic}
\vspace{1.4ex}

\columnbreak

\opsig{$\id.\opdef{Read}()$} \quad از $\id'$
\begin{algorithmic}[1]
\algcont
\Req{$\neg\done$}
\Req{$\id.P \neq Z$}
\State $\tau \gets \id_{\Gclock}.\fopl{Read}() \text{ as } \id$
\State $\V{tr} \gets \V{tr} \cdot (\atom{read},\id.P,\tau)$
\State \textbf{return} $\tau$
\algsave
\end{algorithmic}
\vspace{1.4ex}

\opsig{$\id.\op{Leak}()$} \quad از $\id'$
\begin{algorithmic}[1]
\algcont
\State \textbf{return} $\none$
\algsave
\end{algorithmic}
\vspace{1.4ex}

\end{multicols}}
\end{interface}

\begin{interface}[nofloat]{The finalizations of $\Gm$; the shell over $\Gclock$ is shared}
\opsig{$\id.\opdef{FinStead}()$} \quad from $\id'$ \quad (steadiness)
\begin{algorithmic}[1]
\algcont
\Req{$\neg\done$}
\Req{$\id.P = Z$}
\State $\done \gets \true$
\State parse $\V{tr}$ as $\bigl((\mathit{op}_j,P_j,\tau_j)\bigr)_{j=1}^{n}$; \ $\tau_0 := 0$
\State $w \gets 1$ if $\exists\, j \in \{1,\dots,n\} :\ \
       \tau_j \notin \{\tau_{j-1},\,\tau_{j-1}+1\}$ \label{ln:gstead}
\State \quad ${}\vee\ \bigl(\tau_j = \tau_{j-1}+1 \ \wedge\ \mathit{op}_j \neq \atom{tick}\bigr)$;
\State else $w \gets 0$
\State \textbf{return} $w$
\algsave
\end{algorithmic}
\vspace{1.4ex}

\opsig{$\id.\opdef{FinUna}()$} \quad from $\id'$ \quad (unanimity)
\begin{algorithmic}[1]
\algcont
\Req{$\neg\done$}
\Req{$\id.P = Z$}
\State $\done \gets \true$
\State $\Cs \gets (C,Z).\fopl{Status}() \text{ as } \id$
\State $w \gets 1$ if $\exists\, t \geq 0
       \ \ \exists\, P \in \Gclock.\Ps \setminus \Cs :$ \label{ln:guna}
\State \quad $\bigl(\exists\, \tau_r \leq t : (\atom{reg},P,\tau_r) \in \V{tr}\bigr)
       \ \wedge\ (\atom{dereg},P,\cdot) \notin \V{tr}$
\State \quad ${}\wedge\ \bigl(\exists\, \tau > t : (\cdot,\cdot,\tau) \in \V{tr}\bigr)
       \ \wedge\ \bigl(\forall\, \tau' \in \{t,\,t+1\} :
       (\atom{tick},P,\tau') \notin \V{tr}\bigr)$;
\State else $w \gets 0$
\State \textbf{return} $w$
\algsave
\end{algorithmic}
\end{interface}

The verdicts, in words. The finalization $\op{FinStead}$ returns $1$ on an answer below its predecessor, an answer
more than one above it, or an increase anywhere but at a tick: reads, registrations and
departures must find the time where the last answer left it. A tick answered $\tau_{j-1}+1$ is
allowed --- it is the tick that closed the round --- and $\tau_0 := 0$ is $\op{Initialize}$'s
word taken as the first observation, so a first read answered $1$ is already a counter that moved
before anyone pushed. The finalization $\op{FinUna}$ returns $1$ if the counter passed some $t$ although some party
--- honest at the close, registered by round $t$ and never released --- has no tick answered $t$
or $t+1$: no tick placed during round $t$, that is, the answer $t+1$ covering the one that closed
it. The registration conjuncts read the trace, not the clock's tables, so they too name only what
the game itself did.

Two contrasts with Section~\ref{sec:sigprops}, both instructive. There the trace was a sequence
whose order no verdict read; here $\op{FinStead}$ reads little else, time being an order and
not much besides. And there both verdicts had to ask $P \notin \Cs$ of every party they named;
$\op{FinStead}$ asks nothing of anyone, because at a corrupt party the game's own wrapper
intercepts before the core runs and records (Remark~\ref{rem:gmcorr}), so every entry in $\V{tr}$
is an honest core's answer already. The one verdict that must name honesty is $\op{FinUna}$'s,
whose claim concerns parties that did \emph{not} call, and it names it at the close ---
monotonicity (Section~\ref{sec:fullinterface}) again making that the stronger reading. That the
two differ so --- one reading the register, the other not; one reading order, the other
membership --- is what a shared shell is for: it is the oracles they have in common, not the
verdicts, that made them one game.

\begin{remark}[Agreement is steadiness read across parties]\label{rem:clockagree}
The property a shared clock is usually wanted for is \emph{agreement}: honest parties hold a
common view of the time. Here it is not a third game but a corollary of the verdict above, and by
design. Line~\ref{ln:gstead} quantifies over the trace \emph{interleaved}, not party by party:
between two consecutive $\atom{tick}$ entries every answer must equal its predecessor, whoever
asked, so two honest parties reading between the same two ticks are answered the same value or
the game is won. Stated per party --- each party's own subsequence steady --- the verdict would
be strictly weaker, and agreement a separate property: a replacement could then answer $P$ with
$5$ and $Q$ with $3$ in adjacent activations and fail neither property. What lets one verdict carry
both claims is the execution model: a single token totally orders the trace, so ``simultaneous
reads'' do not exist and a common view reduces to the consistency of adjacent entries. Two limits
are worth naming. Agreement is among honest parties only, a corrupt party's answers being the
adversary's (Remark~\ref{rem:gmcorr}); and it says nothing of a party that does not ask, whose
view may grow stale while the clock runs ahead --- staleness of the party, not disagreement of
the clock, and nothing a functionality could promise a caller that never calls. A \emph{drift}
variant --- views agreeing within some $\Delta$, the loosely synchronized clocks
of~\cite{KMTZ13} --- would be a genuine weakening of line~\ref{ln:gstead}, and under the present
shape it is a third finalization rather than a third game: it reads the same trace of the same
relays and differs in the verdict alone, which is exactly the test for hanging it off this shell.
\end{remark}

\begin{proposition}[$\Gclock$ has both properties]\label{prop:clockprops}
For $\op{Fin}$ either of $\op{FinStead}$ and $\op{FinUna}$, for every
$\Zenv \in \ZenvSet_{\gamma,\gamma}$ and every occupant $\mathcal{X}$ of the adversary slot,
$\Pr[\Win_{\op{Fin}}] = 0$. So $\{\Gclock\}$ has each property within $0$ against
$\bigl(\ZenvSet_{\gamma,\gamma},\Xset\bigr)$ in the sense of Definition~\ref{def:propadv}, for
$\Xset$ the class of all occupants.
\end{proposition}

\begin{proof}
% BEGIN overview prop:clockprops
One invariant carries both properties: the counter changes only inside a tick the trace records and
by exactly one, and membership only inside recorded registrations and departures. Establishing it is
the first paragraph's work, and it turns on three things --- tightness, so no caller but the game
reaches the clock's cores; the game reaching them only at honest parties, so no answer is mediated;
and no core of $\gamma$ handing the token to anything that could run one, the single call the
clock's cores place being responsive. Entries therefore append in execution order, each carrying the
counter as its call returned. Steadiness then follows by comparing consecutive entries, since
between them nothing runs. Unanimity is the longer half: it assumes the four conjuncts of the
verdict and derives a contradiction from the invariant, using monotonicity of corruption to keep the
party honest throughout.
% END overview prop:clockprops
Everything rests on one invariant: the counter changes only inside a tick the trace records, and
by exactly one; the membership only inside recorded registrations and departures. Both tables are
written only in the cores of Chapter~\ref{sec:gclock} --- the counter at one line of
$\op{Update}$, stepping by one when the test of line~\ref{ln:advance} passes, $\V{reg}$ in
$\op{Register}$ and $\op{Deregister}$ --- and those cores run only on calls passing
$\opl{Guard}$: with $\admits$ pinned no caller but $\Gm$ passes line~\ref{ln:caller}, and $\Gm$
calls only through the relays, at honest parties --- at a corrupt one its own wrapper intercepted
a step earlier and no core of $\gamma$ ran --- so no call is mediated and each core's answer is
the answer recorded. Moreover no core of $\gamma$ hands the token to anything that could run one:
the one call the clock's cores place is the notification of line~\ref{ln:clocknotify}, which is
responsive --- $\opl{Respond}$ refuses every call its answerer would place, whatever occupies the
slot, and every mutation of $\op{Update}$ precedes it in any case --- and the relays' one call is
on the clock, that chain returning with no other machine of $\gamma$ having run in between.
Entries therefore append in execution order, each carrying the counter as its call returned.

For steadiness: before the first entry the counter is $\op{Initialize}$'s $0 = \tau_0$. Between
consecutive entries no core of $\gamma$ runs, so the counter at entry $j$ differs from that at
entry $j-1$ by the steps taken inside entry $j$'s own call: none for a read, a registration or a
departure, none or one for a tick. Hence $\tau_j \in \{\tau_{j-1}, \tau_{j-1}+1\}$ throughout,
with the increment only at a tick, and line~\ref{ln:gstead} gives $0$.

For unanimity: suppose the four conjuncts of line~\ref{ln:guna} hold of some $t$ and some
$P \in \Gclock.\Ps \setminus \Cs$ at the close --- honest throughout, by monotonicity. A recorded
registration answered $\tau_r \leq t$ set $\V{reg}[P] \gets 1$; only the cores of $\op{Register}$
and $\op{Deregister}$ write $\V{reg}$, the game placed no $\op{Deregister}$ at $P$ --- there is
no $\atom{dereg}$ entry --- and the observers cannot reach an honest party, so $\V{reg}[P] = 1$
from that moment on. Some entry answers $\tau > t$, and by the invariant the answers are values
the counter took, so it stepped from $t$ to $t+1$ inside some recorded tick; its
line~\ref{ln:advance} read the flags of every party honest at that moment and registered then ---
$P$ among them, the register only ever growing and $\tau_r \leq t$ --- so $\V{tick}[P] = 1$. That
flag is written only by the core of $\op{Update}$ at $P$ and cleared at every step of the counter
and at every re-registration, so it was set while the counter stood at $t$, by a tick of the game
at $P$; that tick answered $t$, or $t+1$ if it was itself the one that closed the round, and it
was recorded. So $(\atom{tick},P,\tau') \in \V{tr}$ for some $\tau' \in \{t,t+1\}$, contradicting
the fourth conjunct, and line~\ref{ln:guna} gives $0$.
\end{proof}

As with $\Fsig$, the bound being $0$ says only that the specification was written to have its
properties; the content is the transfer, and here the notification of line~\ref{ln:clocknotify}
collects its price. Corollary~\ref{cor:proptransfer2} asks that the cores of the game system
place no responsive call, and that line is one, so the transfer to a system $\pi$ put in place of
$\{\Gclock\}$ must be argued directly, for the reason Remark~\ref{rem:respdummy} gives --- and
the same hypothesis guards Theorem~\ref{thm:dummy}, so a system built over the notifying clock
forgoes dummy completeness with it. A clock without the notification --- delete
line~\ref{ln:clocknotify} and the $\uses$ entry --- has both properties by the same proof and
inherits them through the corollary with nothing to discharge.

The choice, note, is only \emph{which} price, not whether. A notification through the ordinary
$\Adv(\cdot)$, the route $\Fsig$ took, would hand the token to an adversary free to drive nested
ticks before answering; answers would then append out of trace order, and steadiness would fail
of the ideal clock itself --- the two verdicts above are properties of $\Gclock$ only because
line~\ref{ln:clocknotify} cannot be acted on before $\op{Update}$ returns. So where $\Fsig$ chose
inheritance and paid with locality, a notifying clock must choose locality and pay with
inheritance; the clock that pays neither is the one that tells the adversary nothing.

%% ==================================================================
%% ==================================================================
\chapter{\texorpdfstring{$\Delta$}{Δ}-Delayed Network}\label{sec:fnet}
\chapterhead{Delta-Delayed Network}

\section{کارکرد}\label{sec:fnetfunc}

The clock exists to be read, and the first reader is a network. Communication with a bounded
delay is~\cite{KMTZ13}'s; the diffusion form, delays measured on a shared clock, is the network
the composable blockchain analyses run over~\cite{BMTZ17}. The network $\Fnet$ below is that functionality
in the language of this paper --- and the first example to \emph{use} the clock rather than be
it. Chapter~\ref{sec:functionalities} named a diffusion channel $\Fdiffuse$ and owed no code;
$\Fnet$ is what such a channel looks like with time attached.

What it does is quickly said. The operation $\opl{Send}$ timestamps the message by the clock --- through
$\id_{\Gclock} := (\Gclock.\PID,\id.P)$, the clock's interface at the party being served, the
notation of Section~\ref{sec:clockprops} --- stores it, tells the adversary --- the network is
public --- and answers a message id. The operation $\opl{Fetch}$ asks the
adversary what to reveal early, then answers everything released to the fetching party together
with everything $\Delta$ old. The division of power is the point: the adversary chooses
\emph{when} within $\Delta$, and to whom first, and may add traffic of its own; what it cannot do
is withhold an honest message past its time, and that bound sits in $\op{Fetch}$'s last line,
where no answer of the slot reaches it. A message sent at time $\tau$ is in every fetch from
time $\tau + \Delta$ on, whoever the adversary favoured first.

\heading{The fields} They follow $\Fsig$'s pattern, not the clock's. $\PID$, $\Ps$ and $\admits$
are parameters: a network is local, one instance per session, its $\admits$ naming the protocol
it serves in the local pattern of Chapter~\ref{sec:functionalities} --- and every instance reads
the one shared clock, which is what makes ``time $\tau$'' comparable across sessions.
$\uses := \{(\Gclock,\serves),(\Adv,\serves)\}$: the first entry requires the clock to serve
every party the network does and is met, whatever session the network sits in, at the shared copy
by the global disjunct of Definition~\ref{def:wfsys} --- a system containing $\Fnet$ is
well-formed only with the clock alongside --- and the second covers the two slot calls below.
$\pars := \Delta$, a parameter a line of code actually reads, as $\Frand$'s $n$ is
(Chapter~\ref{sec:frand}).

\heading{Scheduling is answered, not called} \cite{BMTZ17} lets the adversary deliver a message
early by acting on the network of its own accord. Transcribed literally that is an interface the
adversary \emph{calls}, so a local $\admits$ containing $\Apid$ --- reachable from the slot in
session $0$ and in no other, the one shape Definition~\ref{def:sessunif} excludes and relocation
cannot carry. So the slot is asked instead: $\op{Fetch}$ hands the adversary the time and takes
back a pair $(S,I)$ --- releases among the stored ids, injections of fresh traffic --- sanitized
inline, the good repair being canonical, as in Chapter~\ref{sec:fsig}. Both slot calls are
responsive, the clock's discipline. The operation $\op{Send}$ mutates before it notifies; $\op{Fetch}$ cannot,
its mutations reading the answer, and responsiveness is what makes the difference immaterial:
the token comes straight back, nothing else runs inside either call, and $\Fsig$'s re-entrance
re-tests stay unnecessary here too.

\heading{Corruption, without a line of code} A corrupt party's $\op{Send}$ and $\op{Fetch}$ are
intercepted by line~\ref{ln:medin} of the full interface before the cores run, so the buffer
holds honest sends only, and the cores' clock reads are never mediated --- a timestamp is always
the clock's own word. What a corrupt party ``sends'', the adversary supplies from the receiving
side instead: the injections of line~\ref{ln:netinject} enter the buffer with sender $\none$ and
timestamp $\infty$, released to the fetching party alone --- no age, and so no ripening into
line~\ref{ln:netavail}'s guarantee, which is for the honest and never conscripts the adversary's
own traffic.

\begin{interface}[nofloat]{کارکردِ $\Fnet$ \\
\params{$\PID$, \quad $\Ps$, \quad $\admits$, \quad
$\uses := \{(\Gclock,\serves),\,(\Adv,\serves)\}$, \quad $\pars := \Delta$}}

{\scriptsize
\setlength{\columnsep}{1.4em}
\begin{multicols}{2}
\opsig{$\op{Initialize}()$}:
\begin{algorithmic}[1]
\setcounter{contline}{0}\algcont
\State $\V{ctr} \gets 0$
\State $\V{buf} \gets \emptyset$
  \Comment{درایه‌ها $(m, P, \msg, \tau)$}
\State $\V{rel} : \mathbb{N} \times \Fnet.\Ps \to \{0,1\}$
\State $\V{rel}[*,*] \gets 0$
  \Comment{چه چیزی به چه کسی رها شده}
\algsave
\end{algorithmic}
\vspace{1.4ex}

\opsig{$\id.\op{Send}(\msg)$} \quad از $\id'$
\begin{algorithmic}[1]
\algcont
\State $\tau \gets \id_{\Gclock}.\fopl{Read}() \text{ as } \id$
\State $m \gets \V{ctr}$; \ $\V{ctr} \gets \V{ctr} + 1$
\State $\V{buf} \gets \V{buf} \cup \{(m, \id.P, \msg, \tau)\}$
\State $\Adv^{!}\bigl(\id.\op{Send},\ m,\ \msg\bigr)$ \label{ln:netsend}
  \Comment{شبکه عمومی است}
\State \textbf{return} $m$
\algsave
\end{algorithmic}
\vspace{1.4ex}

\opsig{$\id.\op{Leak}()$} \quad از $\id'$
\begin{algorithmic}[1]
\algcont
\State \textbf{return} $(\V{buf},\ \V{rel})$
\algsave
\end{algorithmic}
\vspace{1.4ex}

\columnbreak

\opsig{$\id.\op{Fetch}()$} \quad from $\id'$
\begin{algorithmic}[1]
\algcont
\State $\tau \gets \id_{\Gclock}.\fopl{Read}() \text{ as } \id$
\State $(S, I) \gets \Adv^{!}\bigl(\id.\op{Fetch},\ \tau\bigr)$ \label{ln:netquery}
  \Comment{a non-pair reads as $(\emptyset,\emptyset)$}
\State $S \gets S \cap \{m : (m,\cdot,\cdot,\cdot) \in \V{buf}\}$;
\State $I \gets I \cap \Msgs$
  \Comment{sanitized inline}
\State $\V{rel}[m, \id.P] \gets 1$ \ for each $m \in S$
\State for each $\msg \in I$: \label{ln:netinject}
\State \quad $\V{buf} \gets \V{buf} \cup \{(\V{ctr}, \none, \msg, \infty)\}$;
\State \quad $\V{rel}[\V{ctr}, \id.P] \gets 1$; \ $\V{ctr} \gets \V{ctr} + 1$
\State \textbf{return} $\bigl\{(m, \msg) : (m, P', \msg, \tau_s) \in \V{buf}$
\State \quad ${}\wedge\ \bigl(\V{rel}[m, \id.P] = 1 \ \vee\ \tau_s + \Delta \leq \tau\bigr)\bigr\}$ \label{ln:netavail}
\algsave
\end{algorithmic}
\end{multicols}}
\end{interface}

Three remarks on the code. Fetched pairs carry no sender: the network is an unauthenticated
diffusion, \cite{BMTZ17}'s, and the injections are what unauthenticated \emph{means} --- a fetch
may contain messages no one sent, the wire being the adversary's to write. An authenticated
\emph{diffusion} is one field and one line away --- return the sender the buffer already keeps,
drop the injections. Authentication with addressing is a further step and a different object;
Chapter~\ref{sec:fac} gives it. The table $\V{rel}$ is sticky: what one fetch reveals no later fetch may hide, so a party's view
of the network only grows --- and line~\ref{ln:netavail} reads $\V{rel}$ \emph{or} the deadline,
so the adversary's silence past $\Delta$ is simply ignored rather than punished. And
$\opl{Leak}$ returns the whole state, the network keeping no secrets: its contents were told to
the adversary at line~\ref{ln:netsend} as they arrived, and the release table is the slot's own
doing.

\section{ویژگی‌ها}\label{sec:netprops}

Section~\ref{sec:fnetfunc} left three properties unformulated: delivery by $\Delta$, monotone
views, and no delivery before sending. The last needs no game at all: an id not yet in $\V{buf}$
is nothing $S$ can name, line~\ref{ln:netavail} intersecting with $\V{buf}$ before anything else
runs, so no occupant of the slot can ever release what has not yet been sent. Monotone views we
still leave unformulated, being a property of $\V{rel}$ alone and orthogonal to timing. What we
formalize here is the first, under the name $\opl{FinLive}$, \emph{liveness}: a message ripened
by $\Delta$ is missing from no later honest fetch. And we formalize a property $\Fnet$ was never
claimed to have, $\opl{FinAuth}$, \emph{authentication}: a message an honest party fetches was
sent by someone. Both take the search reading of Definition~\ref{def:propadv}.

They are two finalizations of one game, and this is where that shape earns the most. The network $\Fnet$ has
the first outright and, being an unauthenticated diffusion by design ---
Section~\ref{sec:fnetfunc}'s own remark on the injections of line~\ref{ln:netinject} ---
provably lacks the second: one shell, one set of oracles, one trace, and two verdicts read off it
of which one is unwinnable and the other is won with probability $1$. Under two separate games
that would be two statements about two objects; here it is one object measured twice. What keeps
the measurements apart is Remark~\ref{rem:gmonce}'s single flag: an environment that calls
$\op{FinAuth}$ has closed the game against $\op{FinLive}$, so no run is judged on both, and
nothing about the pair is a conjunction.

Both are played over the game system $\gamma := \{\Gm\} \cup \{\Fnet,\Gclock\}$, with
$\Fnet.\admits := \{\Gm.\PID\}$ pinned and $\Fnet.\Ps \subseteq \Gm.\Ps$ ---
Section~\ref{sec:sigprops}'s tightness argument transplanted whole, since $\Fnet$'s own fields
already follow $\Fsig$'s pattern rather than the clock's (Section~\ref{sec:fnetfunc}). The
interface box below states $\Gm.\uses$ in full. No caller but $\Gm$ passes $\opl{Guard}$ at $\Fnet$: not the
environment, not the adversary, not a machine hosted at an invented instance of the game's name.
So every message $\Fnet$'s core ever stores, and every answer it ever returns to a $\op{Fetch}$,
passes through the relays below and is recorded --- a game whose verdict names ``sent by someone''
has no other way to know it. $\Fnet$'s own dependence on the clock and the slot is unaffected by
any of this and is met exactly as it was for $\Fnet$ itself, the clock at the shared copy by the
global disjunct of Definition~\ref{def:wfsys}. Inside an interface at $\id$ we write $\id_{\Fnet}$
for $(\Fnet.\PID,\id.P)$ and $\id_{\Gclock}$ for $(\Gclock.\PID,\id.P)$, the second the notation of
Section~\ref{sec:clockprops}.

\begin{interface}[nofloat]{پوستهٔ بازی روی $\Fnet$ و $\Gclock$؛ نهایی‌سازی‌ها در پی می‌آیند \\
\params{$\PID := (\Gm.F,0,0)$, \quad $\Ps := \Fnet.\Ps \cup \{Z\}$, \quad $\admits := \Zpid$, \quad
$\uses := \{(\Fnet,\serves_{\Gm}),\,(\Gclock,\serves)\}$, \quad $\pars := \none$}}

{\scriptsize
\setlength{\columnsep}{1.4em}
\begin{multicols}{2}
\opsig{$\op{Initialize}()$}:
\begin{algorithmic}[1]
\setcounter{contline}{0}\algcont
\State $\done \gets \false$
\State $\V{tr} \gets (\,)$
  \Comment{رد، یک دنباله}
\algsave
\end{algorithmic}
\vspace{1.4ex}

\opsig{$\id.\opdef{Send}(\msg)$} \quad از $\id'$
\begin{algorithmic}[1]
\algcont
\Req{$\neg\done$}
\State $\tau \gets \id_{\Gclock}.\fopl{Read}() \text{ as } \id$
\State $m \gets \id_{\Fnet}.\fopl{Send}(\msg) \text{ as } \id$ \label{ln:gnetsend}
\State $\V{tr} \gets \V{tr} \cdot (\atom{send},\id.P,m,\tau,\msg)$
\State \textbf{return} $m$
\algsave
\end{algorithmic}
\vspace{1.4ex}

\columnbreak

\opsig{$\id.\opdef{Fetch}()$} \quad از $\id'$
\begin{algorithmic}[1]
\algcont
\Req{$\neg\done$}
\State $\tau \gets \id_{\Gclock}.\fopl{Read}() \text{ as } \id$
\State $M \gets \id_{\Fnet}.\fopl{Fetch}() \text{ as } \id$ \label{ln:gnetfetch}
\State $\V{tr} \gets \V{tr} \cdot (\atom{fetch},\id.P,\tau,M)$
\State \textbf{return} $M$
\algsave
\end{algorithmic}
\vspace{1.4ex}

\opsig{$\id.\op{Leak}()$} \quad از $\id'$
\begin{algorithmic}[1]
\algcont
\State \textbf{return} $\none$
\algsave
\end{algorithmic}
\vspace{1.4ex}

\end{multicols}}
\end{interface}

\begin{interface}[nofloat]{The finalizations of $\Gm$; the shell over $\Fnet$ and $\Gclock$ is shared}
\opsig{$\id.\opdef{FinLive}()$} \quad from $\id'$ \quad (liveness)
\begin{algorithmic}[1]
\algcont
\Req{$\neg\done$}
\Req{$\id.P = Z$}
\State $\done \gets \true$
\State $\Cs \gets (C,Z).\fopl{Status}() \text{ as } \id$
\State $w \gets 1$ if $\exists\, P,Q,m,\tau_s,\msg,\tau,M :\ \
       P \notin \Cs \ \wedge\ Q \notin \Cs$ \label{ln:glive}
\State \quad ${}\wedge\ (\atom{send},P,m,\tau_s,\msg) \in \V{tr}
       \ \wedge\ (\atom{fetch},Q,\tau,M) \in \V{tr}$
\State \quad ${}\wedge\ \tau \geq \tau_s + \Delta
       \ \wedge\ (m,\msg) \notin M$;
\State else $w \gets 0$
\State \textbf{return} $w$
\algsave
\end{algorithmic}
\vspace{1.4ex}

\opsig{$\id.\opdef{FinAuth}()$} \quad from $\id'$ \quad (authentication)
\begin{algorithmic}[1]
\algcont
\Req{$\neg\done$}
\Req{$\id.P = Z$}
\State $\done \gets \true$
\State $\Cs \gets (C,Z).\fopl{Status}() \text{ as } \id$
\State $w \gets 1$ if $\exists\, Q,\tau,M,m,\msg :$ \label{ln:gauth}
\State \quad $Q \notin \Cs \ \wedge\ (\atom{fetch},Q,\tau,M) \in \V{tr} \ \wedge\ (m,\msg) \in M$
\State \quad ${}\wedge\ \neg\exists\, P,\tau_s : (\atom{send},P,m,\tau_s,\msg) \in \V{tr}$;
\State else $w \gets 0$
\State \textbf{return} $w$
\algsave
\end{algorithmic}
\end{interface}

Three things about the code. Neither relay reads $\op{Send}$'s or $\op{Fetch}$'s own return value
for a timestamp --- there is none to read --- so both read the clock themselves, one line above
the call to $\Fnet$; the two readings coincide because no core of $\gamma$ runs between them, not
$\Fnet$'s own guard and corruption check, which only reads $\Cs$ and finds $\id.P$ honest, the
relay having run at all, and not $\op{Read}$ itself, which places no call
(Section~\ref{sec:gclockfunc}). Neither relay excludes $\id.P = Z$, unlike
Section~\ref{sec:clockprops}'s: $\Fnet$'s fields follow $\Fsig$'s pattern and not the clock's, and
neither verdict reads any the worse for the root sending or fetching too --- so the shell is spared
the demand steadiness made of the clock's, and both finalizations get the shell they would have
asked for alone. And authentication's
verdict asks nothing of a sender, because $\op{Fetch}$ returns none to ask about --- the fetched
pair is $(m,\msg)$ alone, and whether some $m$ was ever sent, by whom, is exactly the fact no
caller of $\Fnet$ can otherwise recover, which is what the whole game is testing.

\begin{proposition}[$\Fnet$ has liveness]\label{prop:netlive}
For every $\Zenv \in \ZenvSet_{\gamma,\gamma}$ and every occupant $\mathcal{X}$ of the adversary
slot, $\Pr[\Win_{\op{FinLive}}] = 0$. So $\{\Fnet,\Gclock\}$ has liveness within $0$ against
$\bigl(\ZenvSet_{\gamma,\gamma},\Xset\bigr)$ in the sense of Definition~\ref{def:propadv}, for
$\Xset$ the class of all occupants.
\end{proposition}

\begin{proof}
Suppose $(\atom{send},P,m,\tau_s,\msg) \in \V{tr}$, written at some activation of the $\op{Send}$
relay. Its own clock read and $\Fnet.\op{Send}$'s internal one (Section~\ref{sec:fnetfunc})
answer the same value, by the argument just given, so the core stores exactly $(m,P,\msg,\tau_s)$
in $\V{buf}$.

Now suppose also $(\atom{fetch},Q,\tau,M) \in \V{tr}$ with $\tau \geq \tau_s + \Delta$.
$\Fnet.\op{Fetch}$'s core returns every pair $(m',\msg')$ with $(m',P',\msg',\tau_{s'}) \in \V{buf}$
and $\V{rel}[m',Q] = 1 \vee \tau_{s'} + \Delta \leq \tau$ (line~\ref{ln:netavail}), $\tau$ its own
reading of the clock and, again, the relay's. Taking $m' = m$: the buffer entry is
$(m,P,\msg,\tau_s)$ and $\tau_s + \Delta \leq \tau$ holds by hypothesis, so the second disjunct is
met whatever $\V{rel}[m,Q]$ is and whatever the slot answered at line~\ref{ln:netquery}, which
settles only $S$ and $I$ and neither disjunct. Hence $(m,\msg) \in M$, and line~\ref{ln:glive}
gives $0$.
\end{proof}

\begin{proposition}[$\Fnet$ does not have authentication]\label{prop:netauth}
There are $\Zenv^{*} \in \ZenvSet_{\gamma,\gamma}$ and an occupant $\mathcal{X}^{*}$ of the
adversary slot with $\Pr[\Win_{\op{FinAuth}}] = 1$: no $\varepsilon < 1$ bounds
$\{\Fnet,\Gclock\}$'s
authentication against any class containing them.
\end{proposition}

\begin{proof}
% BEGIN overview prop:netauth
A counterexample in two machines. The occupant $\mathcal{X}^{*}$ answers $\Fnet$'s fetch query with
one injection and places no call --- which is exactly why it is not the dummy, whose only way to
answer is to relay to the environment, a call $\opl{Respond}$ refuses. That difference is the whole
point of stating these properties over every occupant of the slot rather than the dummy alone: read
at the dummy, $\op{FinAuth}$ would come out satisfied, and only vacuously. The environment
$\Zenv^{*}$ then fetches once and finalizes on $\op{FinAuth}$ and nothing else. The injection
survives sanitizing, reaches the trace, and no $\op{Send}$ relay ever ran, so the verdict fires with
advantage $1$.
% END overview prop:netauth
Let $\mathcal{X}^{*}$ have $\Adv$'s fields but place no call. On the query
$\id.\op{A}(\id''.\op{Fetch},\tau)$ with $\id''.F = \Fnet.\PID$ --- the shape of
line~\ref{ln:netquery} --- it returns $(\emptyset,\{\msg^{*}\})$ for one fixed $\msg^{*} \in
\Msgs$; on anything else it returns $\none$. Placing no call, it is never refused by
$\opl{Respond}$ at either of $\Fnet$'s slot calls, unlike the dummy, whose only way to answer is
to relay to $\Zenv$
(line~\ref{ln:dumin}), a call $\opl{Respond}$ does refuse (Remark~\ref{rem:respdummy}). This is
exactly where $\mathcal{X}^{*}$ differs from $\Dum$, and exactly why
Propositions~\ref{prop:sigprops}, \ref{prop:clockprops} and~\ref{prop:netlive} are stated over
every occupant of the slot and not the dummy alone --- the dummy's case
of Definition~\ref{def:propadv}, with $\mathcal{X}$ dropped from the subscript, would read
$\op{FinAuth}$ as satisfied, and only vacuously,
$\Dum$ being unable to inject at all.

Let $\Zenv^{*}$ claim $(\Gm.\PID,Q)$ for $Q \in \Fnet.\Ps$, call $\op{Fetch}()$ once, then claim
$(\Gm.\PID,Z)$ and call $\op{FinAuth}()$ --- and $\op{FinAuth}$ and no other, which is the
whole of what it takes to be judged on this property and not the one beside it. It corrupts no
one, hosts nothing, and claims only
identities of $\Gm$'s own name, so $\Zenv^{*} \in \ZenvSet_{\gamma,\gamma}$.

At that fetch, $I = \{\msg^{*}\}$ survives sanitizing ($I \gets I \cap \Msgs$), so the loop of
line~\ref{ln:netinject} adds some fresh $(m,\none,\msg^{*},\infty)$ to $\V{buf}$ and sets
$\V{rel}[m,Q] \gets 1$; $(m,\msg^{*})$ is returned, and $(\atom{fetch},Q,\tau,M) \in \V{tr}$ with
$(m,\msg^{*}) \in M$. No $\op{Send}$ relay ever ran, so $\V{tr}$ carries no entry
$(\atom{send},\cdot,m,\cdot,\msg^{*})$ --- indeed no $\atom{send}$ entry naming $m$ at all, $m$ being
fresh to the injection. So $Q \notin \Cs$ throughout, and line~\ref{ln:gauth} gives $1$.

Nothing about that run touches Proposition~\ref{prop:netlive}. The environment $\Zenv^{*}$ closed the game at
$\op{FinAuth}$, so $\op{FinLive}$ was never answered and $\Win_{\op{FinLive}}$ did not occur:
the two events are disjoint by Remark~\ref{rem:gmonce}, and the same shell can therefore host a
property it always keeps beside one it always loses without the two saying anything about each
other.
\end{proof}

As with $\Gclock$, liveness transfers by Corollary~\ref{cor:proptransfer2} only if the cores of the
game system place no responsive call, and both of $\Fnet$'s do, lines~\ref{ln:netsend}
and~\ref{ln:netquery}; a system $\pi$ put in place of $\{\Fnet,\Gclock\}$ must then have the
transfer argued directly, for the reason Remark~\ref{rem:respdummy} gives, exactly as the clock's
own notification cost Section~\ref{sec:clockprops} the same thing. Tightness, note, did no work
in either proof above beyond licensing $\gamma$ as a game system: unlike correctness and
unforgeability, neither verdict here needs to rule out a message entering $\V{buf}$ by some route
$\Gm$ never sees, only to read off what happens to what $\Gm$ itself placed --- so both
propositions would survive a weaker dependence, were one wanted, though nothing above asks for it.

Authentication transfers nothing, having nothing to transfer: Proposition~\ref{prop:netauth} is
not a bound to inherit but a fact about the specification, exactly the fact
Section~\ref{sec:fnetfunc} already asserted in prose --- proved now rather than merely said, and
by an adversary no stranger than one that ignores its mail. No replacement of $\{\Fnet,\Gclock\}$
is asked to improve on it. What does have the analogous property is Chapter~\ref{sec:fac}'s
$\Fac$: a corrupt sender's traffic there enters through $\op{Send}$ itself, under its own identity,
rather than through an unaccountable slot answer (Chapter~\ref{sec:fac}, ``Authentication is the
adversary's own interface''), so the fetched pair carries a sender for a verdict to name, and there
the transplanted win condition --- fetched, addressed, and never sent --- is not met.

%% ==================================================================
%% ==================================================================
\chapter{\texorpdfstring{$\Delta$}{Δ}-Delayed Authenticated Channel}\label{sec:fac}
\chapterhead{Delta-Delayed Authenticated Channel}

\section{کارکرد}\label{sec:facfunc}

$\Fnet$ diffuses and nobody signs. The complementary object is the one where every message
arrives with its sender attached and none arrives from nobody --- the authenticated channel
of~\cite{Canetti20}, given the same $\Delta$ and the same clock. The channel $\Fac$ below is that
functionality, and the two are worth reading side by side: they differ in a good deal more than
the sender field, and the difference is where authentication actually lives.

What it does is quickly said. The operation $\op{Send}$ takes a recipient as well as a message, stamps the pair
by the clock, files it under a fresh index, and tells the adversary --- authentication is not
secrecy, and this channel hides nothing at all. The operation $\op{Fetch}$ gathers everything addressed to the
fetching party that has come of age, asks the adversary which of the rest to release early,
answers the union, and \emph{empties} the entries it answered. What the receiver gets is the whole
record: sender, recipient, send time, message. The division of power is $\Fnet$'s --- the
adversary chooses when within $\Delta$ and in what order --- and so is the bound: an entry
$\Delta$ old is in the next fetch of the party it names, whatever the slot answers.

\heading{The fields} They are $\Fnet$'s, for $\Fnet$'s reasons. $\PID$, $\Ps$ and $\admits$ are
parameters, one instance per session, with $\admits$ naming the protocol served in the local
pattern of Chapter~\ref{sec:functionalities}, and $\pars := \Delta$. The $\uses$ field is
$\{(\Gclock,\serves),(\Adv,\serves)\}$, its clock entry met at the shared copy by the global
disjunct of Definition~\ref{def:wfsys}. Both conjuncts of $\serves$ are needed and neither is a
formality: the first asks the clock to serve every party the channel does, and without the second
the channel's own name would not be admitted at the clock, so $\opl{Guard}$ would refuse every
read on line~\ref{ln:acsendread}.

\heading{Authentication is the adversary's own interface} $\Fnet$ needs
line~\ref{ln:netinject} because a diffusion channel gives adversarial traffic no name to enter
under: the wire is anonymous, so forgery is modelled by letting the slot write on it directly.
Here there is a name, and the full interface already delivers it. At a corrupt sender the
adversary reaches the core --- lines~\ref{ln:corrgate} and~\ref{ln:medin} admit $\Adv$ at a
corrupt party and neither mediation hook fires on its calls, both testing $\id'.F \neq A$ --- so
corrupt traffic enters through $\op{Send}$ under the sender's own identity, timestamped by the
clock like anyone's. The channel $\Fac$ therefore has no injections and needs none, and the sender field of an
entry is not a report about the wire but a fact about which interface the call came through. That
is what an authenticated channel is: forgery is impossible because there is nowhere to forge from.

The price is that adversarial traffic \emph{ripens}. A message the adversary sends from a corrupt
party is in the list with a real timestamp, so line~\ref{ln:acripe} will deliver it by $\Delta$
and no later answer can retract it. The network $\Fnet$ refuses this deliberately, filing injections at
$\infty$ so the guarantee is for the honest and never conscripts the slot's own traffic. The
difference is a consequence of addressing rather than a choice made against it: a message with a
named sender and a named recipient is a message the sender is on the record for, and holding the
adversary to the deadline it accepted when it spoke as a corrupt party is the reading that keeps
the sender field meaning what it says.

\heading{Fetching empties} Each entry names one recipient, so answering it discharges it, and
line~\ref{ln:acdrain} clears what line~\ref{ln:acripe} and the slot between them selected. The network $\Fnet$
cannot do this --- a diffused message is owed to everyone --- which is why it carries the sticky
$\V{rel}$ table instead and re-answers what it has already answered. The two guarantees are
different in kind and neither implies the other: $\Fnet$'s is that a party's view only grows,
$\Fac$'s that a message is delivered exactly once. A protocol over $\Fnet$ must deduplicate and
one over $\Fac$ must not lose what it fetched, and each obligation is the other's absence.

\heading{Sanitizing the release} The slot's answer is a set of indices to release early, and the
only thing wrong it can be is a set naming entries that are not the fetching party's --- or not
entries at all. Chapter~\ref{sec:sanitization}'s procedure repairs exactly this kind of fault, so
line~\ref{ln:acsan} runs the answer through it under
\[
  \Clean_{\op{addr}}(S;\ \V{list}, P) \;:\iff\;
  S \subseteq \{\, i : \V{list}[i] = (\cdot,\, P,\, \cdot,\, \cdot) \,\} ,
\]
with the list and the fetching party as the read-only context. Two things follow that are worth
saying, because they are what makes $\San$ the right instrument here rather than an assertion. A
good value always exists, so the footnote of Chapter~\ref{sec:sanitization} has nothing to prove
in this case: $\emptyset$ satisfies $\Clean_{\op{addr}}$ whatever the list holds, and being the
empty set it is also the lexicographically first such, so a failing answer is repaired to no
early release at all. And that repair costs the adversary nothing it could not have had, since an
adversary wanting some particular good subset released may submit that subset and be returned it
untouched. An assertion in $\San$'s place would read the same on good answers and hand the slot a
denial of service on bad ones --- one stray index and the honest party's fetch dies --- which is
the failure Chapter~\ref{sec:sanitization} exists to rule out.

\heading{Both slot calls are responsive} The clock's discipline, and here it is load-bearing in a
way worth spelling out. The operation $\op{Fetch}$ reads the time on line~\ref{ln:acfetchread} and asks the slot
on line~\ref{ln:acquery}. Were that call not responsive the adversary would hold the token in
between and could drive the clock forward before answering; $\tau$ would then be stale by the time
line~\ref{ln:acripe} used it, fewer entries would clear $\tau_i + \Delta \leq \tau$, and a message
past its deadline would be lawfully withheld. Responsiveness is what makes the read and the test
one moment. The operation $\op{Send}$ mutates before it notifies and needs no such argument, but is responsive
too, for $\Fnet$'s reason: nothing else may run inside either call.

\begin{interface}[nofloat]{کارکردِ $\Fac$ \\
\params{$\PID$, \quad $\Ps$, \quad $\admits$, \quad
$\uses := \{(\Gclock,\serves),\,(\Adv,\serves)\}$, \quad $\pars := \Delta$}}

{\scriptsize
\setlength{\columnsep}{1.4em}
\begin{multicols}{2}
\opsig{$\op{Initialize}()$}:
\begin{algorithmic}[1]
\setcounter{contline}{0}\algcont
\State $\V{ctr} \gets 0$
\State $\V{list} : \mathbb{N} \to (\Fac.\Ps^{2} \times \mathbb{N} \times \Msgs) \cup \{\none\}$
\State $\V{list}[*] \gets \none$
  \Comment{درایه‌ها $(P, Q, \tau, \msg)$}
\algsave
\end{algorithmic}
\vspace{1.4ex}

\opsig{$\id.\op{Send}(Q, \msg)$} \quad از $\id'$
\begin{algorithmic}[1]
\algcont
\State $\tau \gets \id_{\Gclock}.\fopl{Read}() \text{ as } \id$ \label{ln:acsendread}
\State $\V{ctr} \gets \V{ctr} + 1$
\State $\V{list}[\V{ctr}] \gets (\id.P,\, Q,\, \tau,\, \msg)$
\State $\Adv^{!}\bigl(\id.\op{Send},\ \V{ctr},\ Q,\ \tau,\ \msg\bigr)$ \label{ln:acsend}
\State \textbf{return} $\ok$
\algsave
\end{algorithmic}
\vspace{1.4ex}

\opsig{$\id.\op{Leak}()$} \quad از $\id'$
\begin{algorithmic}[1]
\algcont
\State \textbf{return} $\V{list}$
\algsave
\end{algorithmic}

\columnbreak

\opsig{$\id.\op{Fetch}()$} \quad from $\id'$
\begin{algorithmic}[1]
\algcont
\State $\tau \gets \id_{\Gclock}.\fopl{Read}() \text{ as } \id$ \label{ln:acfetchread}
\State $R \gets \{m : \V{list}[m] = (\cdot,\, \id.P,\, \tau_m,\, \cdot)$
\State \quad ${}\wedge\ \tau_m + \Delta \leq \tau\}$ \label{ln:acripe}
  \Comment{come of age}
\State $S \gets \Adv^{!}\bigl(\id.\op{Fetch},\ \tau\bigr)$ \label{ln:acquery}
  \Comment{released early}
\State $S \gets \San[\Clean_{\op{addr}}]\bigl(S;\ \V{list},\, \id.P\bigr)$ \label{ln:acsan}
\State $M \gets \{\V{list}[m] : m \in R \cup S\}$
\State $\V{list}[m] \gets \none$ \ for each $m \in R \cup S$ \label{ln:acdrain}
\State \textbf{return} $M$
\algsave
\end{algorithmic}
\end{multicols}}
\end{interface}

Three remarks on the code. The recipient is not checked against $\Ps$ on $\op{Send}$: a message
addressed outside the channel's parties is filed and never fetched, which is the same silence a
party who declines to fetch would produce, and the guarantee below is a statement about entries
someone fetches. The send time travels with the message, so a receiver can date what it receives
and check the deadline itself --- $\Fnet$'s fetch returns no such thing, and the difference is
that $\Fac$'s $\Delta$ is locally auditable where $\Fnet$'s is only a fact about the
functionality. And $\opl{Leak}$ returns the list entire, the channel keeping no secrets: its
contents were told to the adversary at line~\ref{ln:acsend} as they arrived, and what has since
been drained is what the slot itself watched leave.

\section{ویژگی‌ها}\label{sec:facprops}

The paragraph above named two claims: delivery by $\Delta$, and the one $\Fnet$ cannot make ---
every message fetched was sent, by the party it names. A third, no delivery before sending, again
needs no game: an index not yet in $\V{list}$, or already drained, fails $\Clean_{\op{addr}}$, so
$\San$ repairs $S$ to $\emptyset$ and nothing not yet real is released early either
(Section~\ref{sec:netprops} made the same point for $\Fnet$'s buffer). The finalizations $\op{FinLive}$ and
$\op{FinAuth}$ below formalize the other two, transplanting Section~\ref{sec:netprops}'s pair
rather than starting over --- the shell of relay, log, finalize does not change between a diffusion
channel and an addressed one, and neither do the names, these being the same two properties asked
of a different channel. Both take the search reading. What changes is which of the two $\Fac$ can
keep: both, this time, and the second is the point of the whole design. That the transplant is
\emph{name for name} is the sharpest way to put the difference between the two channels --- one
game, asked of $\Fnet$ and of $\Fac$, answered differently.

Both are played over the game system $\gamma := \{\Gm\} \cup \{\Fac,\Gclock\}$, with
$\Fac.\admits := \{\Gm.\PID\}$ pinned and $\Fac.\Ps \subseteq \Gm.\Ps$ --- the same tightness
argument as Section~\ref{sec:netprops}, and for the same reason, $\Fac$'s fields already
following $\Fnet$'s pattern (Section~\ref{sec:facfunc}). No caller but $\Gm$ passes $\opl{Guard}$
at $\Fac$, so every entry $\op{Send}$ ever places in $\V{list}$, and every record $\op{Fetch}$
ever returns, passes through the relays below and is recorded. The interface box states
$\Gm.\uses$ in full. Inside an interface at $\id$ we write $\id_{\Fac}$ for $(\Fac.\PID,\id.P)$
and $\id_{\Gclock}$ for $(\Gclock.\PID,\id.P)$, the second the notation of
Section~\ref{sec:clockprops}.

\begin{interface}[nofloat]{پوستهٔ بازی روی $\Fac$ و $\Gclock$؛ نهایی‌سازی‌ها در پی می‌آیند \\
\params{$\PID := (\Gm.F,0,0)$, \quad $\Ps := \Fac.\Ps \cup \{Z\}$, \quad $\admits := \Zpid$, \quad
$\uses := \{(\Fac,\serves_{\Gm}),\,(\Gclock,\serves)\}$, \quad $\pars := \none$}}

{\scriptsize
\setlength{\columnsep}{1.4em}
\begin{multicols}{2}
\opsig{$\op{Initialize}()$}:
\begin{algorithmic}[1]
\setcounter{contline}{0}\algcont
\State $\done \gets \false$
\State $\V{tr} \gets (\,)$
  \Comment{رد، یک دنباله}
\algsave
\end{algorithmic}
\vspace{1.4ex}

\opsig{$\id.\op{Send}(Q,\msg)$} \quad از $\id'$
\begin{algorithmic}[1]
\algcont
\Req{$\neg\done$}
\State $\tau \gets \id_{\Gclock}.\fopl{Read}() \text{ as } \id$
\State $r \gets \id_{\Fac}.\fopl{Send}(Q,\msg) \text{ as } \id$ \label{ln:gacsend}
\State $\V{tr} \gets \V{tr} \cdot (\atom{send},\id.P,Q,\tau,\msg)$
\State \textbf{return} $r$
\algsave
\end{algorithmic}
\vspace{1.4ex}

\columnbreak

\opsig{$\id.\op{Fetch}()$} \quad از $\id'$
\begin{algorithmic}[1]
\algcont
\Req{$\neg\done$}
\State $\tau \gets \id_{\Gclock}.\fopl{Read}() \text{ as } \id$
\State $M \gets \id_{\Fac}.\fopl{Fetch}() \text{ as } \id$ \label{ln:gacfetch}
\State $\V{tr} \gets \V{tr} \cdot (\atom{fetch},\id.P,\tau,M)$
\State \textbf{return} $M$
\algsave
\end{algorithmic}
\vspace{1.4ex}

\opsig{$\id.\op{Leak}()$} \quad از $\id'$
\begin{algorithmic}[1]
\algcont
\State \textbf{return} $\none$
\algsave
\end{algorithmic}
\vspace{1.4ex}

\end{multicols}}
\end{interface}

\begin{interface}[nofloat]{The finalizations of $\Gm$; the shell over $\Fac$ and $\Gclock$ is shared}
\opsig{$\id.\op{FinLive}()$} \quad from $\id'$ \quad (liveness)
\begin{algorithmic}[1]
\algcont
\Req{$\neg\done$}
\Req{$\id.P = Z$}
\State $\done \gets \true$
\State $\Cs \gets (C,Z).\fopl{Status}() \text{ as } \id$
\State $w \gets 1$ if $\exists\, P,Q,\tau_s,\msg,\tau,M :\ \
       P \notin \Cs \ \wedge\ Q \notin \Cs$ \label{ln:gacliv}
\State \quad ${}\wedge\ (\atom{send},P,Q,\tau_s,\msg) \in \V{tr}
       \ \wedge\ (\atom{fetch},Q,\tau,M) \in \V{tr}$
\State \quad ${}\wedge\ \tau \geq \tau_s + \Delta
       \ \wedge\ \neg\exists\, \tau',M' :$
\State \qquad $(\atom{fetch},Q,\tau',M') \in \V{tr} \ \wedge\ (P,Q,\tau_s,\msg) \in M'$;
\State else $w \gets 0$
\State \textbf{return} $w$
\algsave
\end{algorithmic}
\vspace{1.4ex}

\opsig{$\id.\op{FinAuth}()$} \quad from $\id'$ \quad (authentication)
\begin{algorithmic}[1]
\algcont
\Req{$\neg\done$}
\Req{$\id.P = Z$}
\State $\done \gets \true$
\State $\Cs \gets (C,Z).\fopl{Status}() \text{ as } \id$
\State $w \gets 1$ if $\exists\, Q,\tau,M,P,\tau_s,\msg :$ \label{ln:gacauth}
\State \quad $Q \notin \Cs \ \wedge\ (\atom{fetch},Q,\tau,M) \in \V{tr} \ \wedge\ (P,Q,\tau_s,\msg) \in M$
\State \quad ${}\wedge\ (\atom{send},P,Q,\tau_s,\msg) \notin \V{tr}$;
\State else $w \gets 0$
\State \textbf{return} $w$
\algsave
\end{algorithmic}
\end{interface}

Three things about the code, two of them familiar. Both relays read the clock themselves for the
same reason as Section~\ref{sec:netprops}'s: the two readings coincide because no core of $\gamma$
runs between them, $\op{Read}$ placing no call of its own. $\op{Send}$'s relay keeps
$\Fac.\op{Send}$'s own return, $\ok$, having nothing else to report --- unlike $\Fnet$'s, which had
a message id to pass back and $\Fac$ never hands out. And authentication's verdict needs no search
over senders or times, unlike its namesake's: a fetched record already names $P$ and $\tau_s$, so
the win condition is a direct check that the exact tuple was never sent, not an existential one
over what might have produced it --- the address $\Fac$ carries is doing, in the verdict, exactly
what it does at the interface.

\begin{proposition}[$\Fac$ has liveness]\label{prop:acliv}
For every $\Zenv \in \ZenvSet_{\gamma,\gamma}$ and every occupant $\mathcal{X}$ of the adversary
slot, $\Pr[\Win_{\op{FinLive}}] = 0$. So $\{\Fac,\Gclock\}$ has liveness within $0$ against
$\bigl(\ZenvSet_{\gamma,\gamma},\Xset\bigr)$ in the sense of Definition~\ref{def:propadv}, for
$\Xset$ the class of all occupants.
\end{proposition}

\begin{proof}
Suppose $(\atom{send},P,Q,\tau_s,\msg) \in \V{tr}$, written at some activation of the $\op{Send}$
relay (line~\ref{ln:gacsend}). Its own clock read and $\Fac.\op{Send}$'s internal one
(line~\ref{ln:acsendread}) agree, by the argument of Section~\ref{sec:netprops}: no core of
$\gamma$ runs between them, $\Fac$'s guard and corruption check finding $\id.P$ honest since the
relay ran at all, and $\op{Read}$ itself placing no call. So the core wrote some fresh $m$ with
$\V{list}[m] = (P,Q,\tau_s,\msg)$ at that moment.

Suppose too $(\atom{fetch},Q,\tau,M) \in \V{tr}$ with $\tau \geq \tau_s + \Delta$. If
$\V{list}[m]$ still holds $(P,Q,\tau_s,\msg)$ when that fetch's core runs, then $m \in R$ by
line~\ref{ln:acripe} --- the entry names $Q$, and $\tau_s + \Delta \leq \tau$ by hypothesis, for
$\tau$ its own reading of the clock and, by the same argument again, the relay's --- so
$(P,Q,\tau_s,\msg) \in M$ whatever the slot answers at line~\ref{ln:acquery} and whatever $\San$
does with it, neither touching $R$. If instead $\V{list}[m]$ had already been cleared, that
happened only at line~\ref{ln:acdrain} of some earlier fetch by $Q$ --- the one line that ever
clears an entry --- and only on indices already placed in that fetch's own $R \cup S$, hence
already present in that earlier fetch's own $M$. Either way $(P,Q,\tau_s,\msg)$ appears in some
fetch of $Q$'s, and the verdict at line~\ref{ln:gacliv} gives $0$.
\end{proof}

\begin{proposition}[$\Fac$ has authentication]\label{prop:acauth}
For every $\Zenv \in \ZenvSet_{\gamma,\gamma}$ and every occupant $\mathcal{X}$ of the adversary
slot, $\Pr[\Win_{\op{FinAuth}}] = 0$. So $\{\Fac,\Gclock\}$ has authentication within $0$ against
$\bigl(\ZenvSet_{\gamma,\gamma},\Xset\bigr)$ in the sense of Definition~\ref{def:propadv}, for
$\Xset$ the class of all occupants --- the property Section~\ref{sec:netprops} proved $\Fnet$
lacks, held here instead.
\end{proposition}

\begin{proof}
Suppose $(\atom{fetch},Q,\tau,M) \in \V{tr}$ with $Q \notin \Cs$ and $(P,Q,\tau_s,\msg) \in M$. By
$\Fac.\op{Fetch}$'s code, $M = \{\V{list}[m] : m \in R \cup S\}$, and both index sets draw only on
$\V{list}$ as it already stands: $R$ by its own pattern match (line~\ref{ln:acripe}), and $S$ after
sanitizing, since $\Clean_{\op{addr}}$ admits only indices $i$ with
$\V{list}[i] = (\cdot,Q,\cdot,\cdot)$ already and line~\ref{ln:sanpick}
(Chapter~\ref{sec:sanitization}) picks among values meeting that predicate --- neither
manufactures an index $\V{list}$ does not hold. So $(P,Q,\tau_s,\msg) \in M$ only if
$\V{list}[m] = (P,Q,\tau_s,\msg)$ for some $m$ at some point.

The only line that ever writes $\V{list}[m]$ to a value other than $\none$ is $\op{Send}$'s own
core, and by tightness --- $\Fac.\admits = \{\Gm.\PID\}$ pinned --- that core runs only on calls
$\Gm$'s relay places: not the environment, not the adversary, not a machine hosted at an invented
instance of the game's name. So the write producing $(P,Q,\tau_s,\msg)$ happened at some
activation of the $\op{Send}$ relay, and by the coinciding-read argument of the previous proof it
was logged there as exactly $(\atom{send},P,Q,\tau_s,\msg)$. Hence
$(\atom{send},P,Q,\tau_s,\msg) \in \V{tr}$, and the verdict at line~\ref{ln:gacauth} gives $0$.
\end{proof}

As with $\Gclock$ and $\Fnet$, both properties transfer by Corollary~\ref{cor:proptransfer2} only
if the game system's cores place no responsive call --- and here, as there, neither does:
Section~\ref{sec:facfunc}'s ``Both slot calls are responsive'' covers $\op{Send}$'s notification
and $\op{Fetch}$'s query alike. A system $\pi$ put in place of $\{\Fac,\Gclock\}$ must then have
both transfers argued directly, for the reason Remark~\ref{rem:respdummy} gives, exactly as
$\Fnet$'s own notification cost Section~\ref{sec:netprops} the same thing.

Unlike the network, nothing here is left standing only as a negative result.
Proposition~\ref{prop:acauth} is the win condition Section~\ref{sec:netprops} transplanted from
$\Fnet$ and could not close, closed here instead --- and the two proofs, laid side by side, locate
the whole difference in one fact: whether tightness can trace a fetched record back to a
$\op{Send}$ that produced it. For $\Fac$ it always can, $\V{list}$ having no other way to gain an
entry; for $\Fnet$ it cannot, line~\ref{ln:netinject} being exactly the other way in. Two
chapters, one game transplanted twice, and the address is the entire difference between them.

That the transplant keeps the names is what makes the comparison sayable at all. The finalization $\op{FinAuth}$
is not two properties that resemble each other but one win condition --- fetched, addressed, never
sent --- asked of two channels, and the answer is $1$ for $\Fnet$ and $0$ for $\Fac$. Under the
shape Chapter~\ref{sec:properties} replaced, where each property carried its own game, the two
would have been $\Gm_{\op{Auth}}$ and $\Gm_{\op{Auth}}'$ and the sameness would have been a remark;
here it is the notation.

%% ==================================================================
%% ==================================================================
\chapter{تحققِ کانالِ اصالت‌دار}\label{sec:realac}

The channel of Chapter~\ref{sec:fac} was postulated. This chapter earns it: a protocol $\rho$
over three ideal subroutines --- a public-key directory, the unauthenticated network of
Chapter~\ref{sec:fnet}, and the signatures of Chapter~\ref{sec:fsig} --- that UC-realizes an
authenticated channel, and does so \emph{exactly}, within $0$, every subroutine being ideal and
the whole reduction information-theoretic. It is the largest construction in the paper, and the
one that exercises the most of it at once: composition, the responsive discipline, sanitization,
tightness, and the fundamental lemma of Chapter~\ref{app:iub} all appear, each doing a job no
other could.

\heading{What $\rho$ does} A party sending for the first time draws a signing key from $\Fsig$ and
registers it in the directory $\Fpki$. To send $\msg$ to $Q$ it reads the clock for a timestamp
$\tau$, increments a per-party sequence number $s$, signs the tuple $(\P,Q,\tau,s,\msg)$, and
diffuses the tuple with its signature over $\Fnet$. To fetch, it collects what $\Fnet$ delivers,
and keeps a diffused item only if it is addressed to the fetcher, carries a signature that
verifies under the \emph{claimed} sender's directory key, and has a $(\text{sender},s)$ pair not
seen before. What survives it hands up as $(\P,Q,\tau,\msg)$. Diffusion carries no addresses and
no authenticity; the signature supplies the second and the address inside the signed tuple the
first, so that a receiver trusts a sender's name exactly because that name is under the sender's
own signature.

\heading{Why the target is $\Facz$, not $\Fac$} A corrupt sender is the obstruction, and it is
instructive rather than incidental. Over a diffusion network the adversary writes on the wire
whenever it likes, in particular \emph{during} a fetch --- $\Fnet$'s injections
(line~\ref{ln:netinject}) are chosen inside the responsive window of the fetch that carries them.
A message so injected, signed under a key the adversary registered for a corrupt party, is
accepted by the receiver and delivered; but no prior $\op{Send}$ ever placed it, so $\Fac$, which
delivers only what its own $\op{Send}$ stored, cannot produce it, and no simulator can call
$\Fac.\op{Send}$ inside the fetch's responsive window to repair the gap. Worse, the corrupt
sender's timestamp is its own to choose --- the signature validates whatever $\tau$ the payload
carries --- where $\Fac$ stamps the true time of a call it witnessed. Neither gap is a defect in
$\rho$; both are what an authenticated \emph{diffusion} unavoidably gives a corrupt party, and
Chapter~\ref{sec:realac} closes with the argument that no protocol over these subroutines does
better. So the realized object is $\Facz$: $\Fac$ with a corrupt sender's traffic admitted at the
fetch, unstored and freely timestamped, which is exactly the guarantee level
of~\cite{Canetti20}'s authenticated channel and no less. The honest-sender guarantees ---
delivery by $\Delta$, and authenticity, that an honest party's fetched record was genuinely sent
by the honest party it names --- survive whole. And against an adversary that corrupts no one the
two objects coincide: the fundamental lemma prices the difference at $\Pr[\Bad]$ for a $\V{bad}$
that only a corrupt sender can raise, so $\rho$ realizes $\Fac$ itself, exactly, whenever nobody
is corrupt (Corollary~\ref{cor:realac}).

\heading{Two subroutines are not quite off the shelf} The signatures are the responsive variant
$\Fsig^{!}$ of Chapter~\ref{app:responsive} --- $\Adv^{!}$ in $\opl{Gen}$, $\opl{Sign}$ and
$\opl{Verify}$ --- and the reason is the token. $\Fac$'s fetch is atomic: its one slot call is
responsive, so no occupant drives the execution between the call and its answer. Were $\rho$'s
fetch to reach $\Fsig$ through the ordinary $\Adv(\cdot)$, a verification would hand the adversary
the token mid-fetch and let it act elsewhere before returning, an interleaving $\Facz$'s atomic
fetch cannot show and no simulator could reproduce. The responsive variant $\Fsig^{!}$ removes the escape; $\Fnet$'s and
$\Facz$'s calls are already responsive, so every send and every fetch is atomic on both sides.
And the directory is \emph{local}, not the global $\Gpki$ of Chapter~\ref{sec:pki}: a shared
directory is written by an honest $\rho.\op{Send}$ in the real world and read straight back by the
environment through $\opl{Retrieve}$, where in the ideal world nothing writes it and the simulator
holds only the peek $\Gpki$ grants at honest parties --- an unsimulatable difference, the
commitment problem of a global setup in this paper's coordinates. A local $\Fpki$, its identities
private to the session like any hybrid, closes it; $Z$ is kept out of its $\admits$ for the same
reason, a direct retrieval by the environment being a call the ideal world has no machine to
serve.

\heading{The directory forgets bindings on purpose} $\Fpki$ sanitizes a registered key to be a
key and nothing more --- \emph{per slot}, dropping $\Gpki$'s cross-party freshness
(Chapter~\ref{sec:pki}, $\Clean_{\op{reg}}$). The freshness would be a weapon: the adversary
chooses an honest party's signing key, it being the value $\Fsig.\op{Gen}$ returns, so it could
pre-register that key at a corrupt slot and have the honest party's own registration
\emph{repaired} to a different one, leaving directory and signing key disagreeing and the party's
every message rejected for ever --- a denial of the liveness $\Facz$ promises unconditionally.
Per-slot fixity is enough because attribution never rested on binding: the sender's name travels
inside the signed tuple, the receiver looks that name up, and a signature valid under an honest
party's key on a tuple naming a \emph{different} sender is a forgery $\Fsig$ suppresses. Binding
buys nothing here and costs liveness, so it goes.

\begin{interface}[nofloat]{دفترچهٔ محلیِ $\Fpki$ \\
\params{$\PID$, \quad $\Ps$, \quad $\admits$, \quad $\uses := \emptyset$, \quad $\pars := \none$}}

{\scriptsize
\setlength{\columnsep}{1.4em}
\begin{multicols}{2}
\opsig{$\op{Initialize}()$}:
\begin{algorithmic}[1]
\setcounter{contline}{0}\algcont
\State $\V{VK} : \Fpki.\Ps \to \Keys \cup \{\unset\}$
\State $\V{VK}[*] \gets \unset$
\algsave
\end{algorithmic}
\vspace{1.4ex}

\opsig{$\id.\op{Register}(\vk)$} \quad از $\id'$
\begin{algorithmic}[1]
\algcont
\If{$\id'.F \in \{A,Z\} \ \wedge\ \id.P \notin \Cs$}
  \State \textbf{return} $\V{VK}[\id.P]$ \label{ln:rpkipeek}
    \Comment{یک نگاه، نه یک نوشتن}
\EndIf
\If{$\V{VK}[\id.P] \neq \unset$}
  \State \textbf{return} $\V{VK}[\id.P]$ \label{ln:rpkifixed}
\EndIf
\State $\vk \gets \San[\Clean_{\op{key}}](\vk)$
\State $\V{VK}[\id.P] \gets \vk$ \label{ln:rpkiwrite}
\State \textbf{return} $\vk$
\algsave
\end{algorithmic}
\vspace{1.4ex}

\columnbreak

\opsig{$\id.\op{Retrieve}(P)$} \quad از $\id'$
\begin{algorithmic}[1]
\algcont
\State \textbf{return} $\V{VK}[P]$
\algsave
\end{algorithmic}
\vspace{1.4ex}

\opsig{$\id.\op{Leak}()$} \quad از $\id'$
\begin{algorithmic}[1]
\algcont
\State \textbf{return} $\V{VK}$
\algsave
\end{algorithmic}
\vspace{1.4ex}

\opsig{$\Clean_{\op{key}}(\vk)$}:
\begin{algorithmic}[1]
\algcont
\State \textbf{return} $\vk \in \Keys$
\algsave
\end{algorithmic}
\end{multicols}}
\end{interface}

$\Facz$ is $\Fac$ with one operation changed: $\op{Fetch}$ takes back a pair $(S,I)$ from the
slot, $S$ the early releases among stored ids as before, $I$ a set of \emph{injected} records the
slot delivers unstored. Sanitizing pins each injection to a corrupt sender and to the fetching
party, so an injection can only ever put a corrupt party's word into the fetch that asked for it;
$\Fac$'s own $\op{Send}$, $\op{Leak}$ and $\op{Initialize}$ are unchanged.

\begin{interface}[nofloat]{Functionality $\Facz$: the fetch of $\Fac$, replaced \\
\params{fields as $\Fac$ (Section~\ref{sec:facfunc})}}

{\scriptsize
\opsig{$\id.\op{Fetch}()$} \quad from $\id'$
\begin{algorithmic}[1]
\State $\tau \gets \id_{\Gclock}.\fopl{Read}() \text{ as } \id$
\State $\Cs \gets (C,\id.P).\fopl{Status}() \text{ as } \id$
\State $R \gets \{m : \V{list}[m] = (\cdot,\, \id.P,\, \tau_m,\, \cdot) \ \wedge\ \tau_m + \Delta \leq \tau\}$
\State $(S, I) \gets \Adv^{!}\bigl(\id.\op{Fetch},\ \tau\bigr)$ \label{ln:aczquery}
  \Comment{releases and injections}
\State $S \gets \San[\Clean_{\op{addr}}]\bigl(S;\ \V{list},\, \id.P\bigr)$
\State $I \gets \San[\Clean_{\op{inj}}]\bigl(I;\ \Cs,\, \id.P\bigr)$ \label{ln:aczsan}
\State $M \gets \{\V{list}[m] : m \in R \cup S\} \ \cup\ I$ \label{ln:aczdeliver}
\State $\V{list}[m] \gets \none$ \ for each $m \in R \cup S$
\State \textbf{return} $M$
\end{algorithmic}
\vspace{1.0ex}
\opsig{$\Clean_{\op{inj}}(I;\ \Cs, P_f)$}:
\begin{algorithmic}[1]
\State \textbf{return}\ $I \subseteq \{(P,Q,\tau,\msg) : P \in \Cs \ \wedge\ Q = P_f\}$
\end{algorithmic}}
\end{interface}

The protocol reads its four subroutines through the interfaces at the party being served, the
notation $\id_{\F} := (\F.\PID,\id.P)$ of Chapter~\ref{sec:fnet}. Its $\uses$ names all four,
the clock among them, since $\rho$ stamps its own payloads and does not recover $\Fnet$'s stamp,
which $\Fnet$ never returns.

\begin{interface}[nofloat]{پروتکلِ $\rho$ روی $\{\Fpki,\Fnet,\Fsig^{!}\}$ \\
\params{$\PID$, \quad $\Ps$, \quad $\admits$, \quad
$\uses := \{(\Fpki,\serves),\allowbreak(\Fnet,\serves),\allowbreak
(\Fsig^{!},\serves),\allowbreak(\Gclock,\serves)\}$, \quad
$\pars := \none$}}

{\scriptsize
\setlength{\columnsep}{1.4em}
\begin{multicols}{2}
\opsig{$\op{Initialize}()$}:
\begin{algorithmic}[1]
\setcounter{contline}{0}\algcont
\State $\V{gen} : \Ps \to \{0,1\}$; \ $\V{gen}[*] \gets 0$
\State $\V{seq} : \Ps \to \mathbb{N}$; \ $\V{seq}[*] \gets 0$
\State $\V{seen} : \Ps \to \mathcal{P}(\Ps \times \mathbb{N})$
\State $\V{seen}[*] \gets \emptyset$
\algsave
\end{algorithmic}
\vspace{1.4ex}

\opsig{$\id.\op{Send}(Q,\msg)$} \quad از $\id'$
\begin{algorithmic}[1]
\algcont
\If{$\V{gen}[\id.P] = 0$} \label{ln:rgen}
  \State $\vk \gets \id_{\Fsig}.\fopl{Gen}() \text{ as } \id$
  \State $\id_{\Fpki}.\fopl{Register}(\vk) \text{ as } \id$
  \State $\V{gen}[\id.P] \gets 1$
\EndIf
\State $\tau \gets \id_{\Gclock}.\fopl{Read}() \text{ as } \id$
\State $\V{seq}[\id.P] \gets \V{seq}[\id.P] + 1$
\State $x \gets (\id.P,\, Q,\, \tau,\, \V{seq}[\id.P],\, \msg)$
\State $\sigma \gets \id_{\Fsig}.\fopl{Sign}(x) \text{ as } \id$
\State $\id_{\Fnet}.\fopl{Send}\bigl((x,\sigma)\bigr) \text{ as } \id$ \label{ln:rsend}
\State \textbf{return} $\ok$
\algsave
\end{algorithmic}
\vspace{1.4ex}

\columnbreak

\opsig{$\id.\op{Fetch}()$} \quad from $\id'$
\begin{algorithmic}[1]
\algcont
\State $N \gets \id_{\Fnet}.\fopl{Fetch}() \text{ as } \id$
\State $M \gets \emptyset$
\For{$(m, y) \in N$} \label{ln:rloop}
  \State parse $y$ as $(x,\sigma)$, \ $x = (P,Q,\tau,s,\msg)$
    \Comment{else skip}
  \State \textbf{if} $Q \neq \id.P$ \textbf{then} skip
  \State $\vk \gets \id_{\Fpki}.\fopl{Retrieve}(P) \text{ as } \id$
  \State \textbf{if} $\vk = \unset$ \textbf{then} skip
  \State $b \gets \id_{\Fsig}.\fopl{Verify}(\vk, x, \sigma)$ \label{ln:rverify}
  \Statex \quad $\text{as } \id$
  \State \textbf{if} $b \neq 1$ \textbf{then} skip
  \State \textbf{if} $(P,s) \in \V{seen}[\id.P]$ \textbf{then} skip \label{ln:rdedup}
  \State $\V{seen}[\id.P] \gets \V{seen}[\id.P] \cup \{(P,s)\}$
  \State $M \gets M \cup \{(P,Q,\tau,\msg)\}$
\EndFor
\State \textbf{return} $M$
\algsave
\end{algorithmic}
\end{multicols}}
\end{interface}

\heading{The simulator} Fix an occupant $\Adv$ of the adversary slot. In the real world $\Adv$
faces the three hybrids; in the ideal world they are gone, and $\Sim_{\Adv}$ must present their
faces from inside itself. It is an absorption (Definition~\ref{def:zadv}, read on the adversary
side): $\Sim_{\Adv}$ bundles $\Adv$ with one internal copy each of $\Fpki$, $\Fnet$ and
$\Fsig^{!}$, runs $\Adv$ against them exactly as the real execution would, and drives them from
what $\Facz$ tells it. On an honest $\op{Send}$ notification $\Facz$ hands it $(\P,Q,\tau,\msg)$;
it advances the internal $\V{seq}$ for $\P$, forms $x = (\P,Q,\tau,s,\msg)$, and replays $\rho$'s
own steps into the internal hybrids --- an internal $\Fsig^{!}.\op{Sign}$, whose $\Adv^{!}$ query
the hosted $\Adv$ answers, and an internal $\Fnet.\op{Send}$, whose notification it delivers ---
so the hosted $\Adv$ sees precisely the queries the real world would place. When $\Facz$ makes its
responsive fetch query, $\Sim_{\Adv}$ runs $\rho.\op{Fetch}$'s filter over the internal $\Fnet$'s
delivery against the hosted $\Adv$, learns exactly which records $\rho$ would hand up, and answers
$(S,I)$: the ripe or early-released honest entries by their $\Facz$ ids in $S$, and the
corrupt-sender records --- those verifying under a key registered at a corrupt slot --- in $I$.
Corruption it passes straight through, $\op{Corrupt}$ surviving the bundle as
Remark~\ref{rem:midcall}'s does.

Made explicit: two pieces of state the paragraph above passes over silently, both forced by the
code just above it rather than chosen. $\Facz$'s own $\op{Send}$ (inherited from $\Fac$,
Section~\ref{sec:facfunc}) notifies the slot with $(\op{Send},m,Q,\tau,\msg)$, $m$ its own list
index --- not $(\P,Q,\tau,\msg)$ alone, the sender $\P$ read off the identity the query arrives at
(Definition~\ref{def:respcall}). The simulator $\Sim_{\Adv}$ must keep that index, $\V{fid}$ below, to name the
entry later; it cannot recover it at fetch time, having no access to $\Facz.\V{list}$ itself.
And $\rho$'s own tables --- $\V{gen}$ at line~\ref{ln:rgen}, $\V{seq}$, and the $\V{seen}$ its
freshness test reads at line~\ref{ln:rdedup} --- are tables $\Sim_{\Adv}$ does not host, $\rho$
being no part of the ideal world and only its subroutines internal here, so reproducing its
filter needs shadows of its own, kept equal to $\rho$'s by construction rather than by reading
them.

Two routing questions the prose leaves open, and both decide whether the simulator works at all.
The first is what becomes of a call the hosted $\Adv$ places on $\rho$'s \emph{own} identity.
Line~\ref{ln:corrgate} admits $\Adv$ at a corrupt party and neither mediation hook fires on its
calls, so in the real world such a call reaches $\rho$'s core --- and every subroutine call that
core then places is mediated straight back to $\Adv$ (line~\ref{ln:medin}, the caller being $\rho$
and the party corrupt), so nothing of it ever reaches $\Fnet$'s buffer. Relayed outward it would
instead reach $\Facz$, whose $\op{Send}$ files a record that \emph{ripens} and is delivered by
$\Delta$ whatever the slot answers (Section~\ref{sec:facfunc}) --- a delivery with no real-world
counterpart. So these calls are not relayed: Definition~\ref{def:bundle}'s explicit routing keeps
them inside, where the last handler below runs $\rho$'s core against the hosted $\Adv$ exactly as
$\opl{Mediate}$ does really. The outward set is declared accordingly.

The second is how far a responsive answer reaches. Both handlers answer a responsive call, and
$\opl{Respond}$ (Definition~\ref{def:respond}) refuses every call the responder \emph{places};
calls between two hosted machines are not placed at all (Definition~\ref{def:bundle}), so the
internal copies are reachable, but nothing outside the bundle is. Two things outside it would
otherwise be reached, and $\Sim_{\Adv}$'s code avoids each.

A hosted copy's \emph{full} interface opens by reading the corruption register, which
Definition~\ref{def:respcall} refuses outright --- a responsive answer is computed from what the
responder already knows. So the handlers drive the copies through their \emph{cores}, written
$\opl{Gen}$ rather than $\fopl{Gen}$ above, and nothing is lost by it: a party corrupt at the send
or the fetch is mediated before $\Facz$'s own core runs, so no slot query is ever issued on its
behalf and no handler ever serves one, and at an honest party the full interface adds only its
silencer --- the gate passes, and neither hook fires. And $\Fnet$ opens both $\op{Send}$ and
$\op{Fetch}$ by reading $\Gclock$, which is shared rather than hosted. That read is answered with
the $\tau$ the notification or query already carries, which is the right value and not merely a
convenient one: on the send it is the stamp the proof shows $\rho$'s and $\Fnet$'s own reads
share, and on the fetch it is $\Facz$'s read, which no responsive call can have advanced past
(Proposition~\ref{prop:respcorr}).

\begin{interface}[nofloat]{Simulator $\Sim_{\Adv}$ \\
\params{occupies $\IDs(\Adv)$; hosts $\Adv$ and one internal copy each of $\Fpki$, $\Fnet$,
$\Fsig^{!}$ (Definition~\ref{def:zadv}, read on the adversary side); outward set the identities
the ideal execution still serves, less $\rho$'s own, which the last handler routes by hand
(Definition~\ref{def:bundle})}}

{\scriptsize
\setlength{\columnsep}{1.4em}
\begin{multicols}{2}
\opsig{$\op{Initialize}()$}:
\begin{algorithmic}[1]
\setcounter{contline}{0}\algcont
\State $\V{gen} : \Ps \to \{0,1\}$; \ $\V{gen}[*] \gets 0$
\State $\V{seq} : \Ps \to \mathbb{N}$; \ $\V{seq}[*] \gets 0$
\State $\V{seen} : \Ps \to \mathcal{P}(\Ps \times \mathbb{N})$; \ $\V{seen}[*] \gets \emptyset$
\State $\V{fid} : \Ps \times \mathbb{N} \to \mathbb{N} \cup \{\unset\}$; \ $\V{fid}[*] \gets \unset$
  \Comment{ثبتِ خودی $\mapsto$ اندیسِ $\Facz$}
\algsave
\end{algorithmic}
\vspace{1.4ex}

\opsig{on $\Facz$'s $\op{Send}$ note $(m,Q,\tau,\msg)$ at $\P$}
\begin{algorithmic}[1]
\algcont
\If{$\V{gen}[\P] = 0$} \label{ln:simgen}
  \State $\vk \gets$ internal $\id_{\Fsig}.\opl{Gen}()$
  \State internal $\id_{\Fpki}.\opl{Register}(\vk)$
  \State $\V{gen}[\P] \gets 1$
\EndIf
\State $\V{seq}[\P] \gets \V{seq}[\P] + 1$; \quad $s \gets \V{seq}[\P]$
\State $x \gets (\P,\, Q,\, \tau,\, s,\, \msg)$
\State $\sigma \gets$ internal $\id_{\Fsig}.\opl{Sign}(x)$
  \Comment{hosted $\Adv$ answers the query}
\State internal $\id_{\Fnet}.\opl{Send}\bigl((x,\sigma)\bigr)$
  \Comment{its clock read answered $\tau$}
\State $\V{fid}[\P,s] \gets m$
\algsave
\end{algorithmic}
\vspace{1.4ex}

\columnbreak

\opsig{on $\Facz$'s responsive $\op{Fetch}$ query $(\tau)$ at $Q$}
\begin{algorithmic}[1]
\algcont
\State $N \gets$ internal $\id_{\Fnet}.\opl{Fetch}()$
  \Comment{clock answered $\tau$; $\Adv$ answers the release query}
\State $S, I \gets \emptyset, \emptyset$
\For{$(m,y) \in N$} \label{ln:simloop}
  \State parse $y$ as $(x,\sigma)$, \ $x = (\P,Q',\tau_x,s,\msg)$
    \Comment{else skip}
  \State \textbf{if} $Q' \neq Q$ \textbf{then} skip
  \State $\vk \gets$ internal $\id_{\Fpki}.\opl{Retrieve}(\P)$
  \State \textbf{if} $\vk = \unset$ \textbf{then} skip
  \State $b \gets$ internal $\id_{\Fsig}.\opl{Verify}(\vk,x,\sigma)$
  \State \textbf{if} $b \neq 1$ \textbf{then} skip
  \State \textbf{if} $(\P,s) \in \V{seen}[Q]$ \textbf{then} skip \label{ln:simdedup}
  \State $\V{seen}[Q] \gets \V{seen}[Q] \cup \{(\P,s)\}$
  \State \textbf{if} $\V{fid}[\P,s] \neq \unset$
  \State \quad \textbf{then} $S \gets S \cup \{\V{fid}[\P,s]\}$ \label{ln:simhonest}
  \State \quad \textbf{else} $I \gets I \cup \{(\P,Q,\tau_x,\msg)\}$ \label{ln:siminject}
\EndFor
\State \textbf{return} $(S,I)$
\algsave
\end{algorithmic}
\vspace{1.4ex}

\opsig{on $\Adv$'s $\op{Send}$ or $\op{Fetch}$ at a corrupt $\P$}
\begin{algorithmic}[1]
\algcont
\State run $\rho$'s own core for that operation (the protocol box above) on $\V{gen}$, $\V{seq}$,
  $\V{seen}$, with every subroutine call it places --- to $\Fsig^{!}$, $\Fpki$, $\Fnet$ or
  $\Gclock$ --- handed to the hosted $\Adv$ rather than placed, as the full interface's
  line~\ref{ln:medin} hands it really; return the core's own value \label{ln:simcorr}
\State never relayed outward: no $\Facz.\V{list}$ entry, and none drained
\algsave
\end{algorithmic}
\vspace{1.4ex}

\opsig{Routing, and $\op{Corrupt}$}
\begin{algorithmic}[1]
\algcont
\State any other call between the hosted $\Adv$ and an internal hybrid: served inside, the
  bundle's ordinary routing (Definition~\ref{def:bundle}). A $\op{Corrupt}$ leaves under the
  hosted $\Adv$'s own claim, the bundle occupying $\IDs(\Adv)$, so $\op{Corrupt}$'s own gate
  $\id'.F = A$ (its line~\ref{ln:corronlya}) is met without rewriting
\algsave
\end{algorithmic}
\end{multicols}}
\end{interface}

Three things about $\Sim_{\Adv}$'s code, all used by the proof but easy to read past. First,
$\V{fid}[\P,s] \neq \unset$ at line~\ref{ln:simhonest} is exactly ``an undrained $\V{list}$ entry
matches the record'' in the proof's own words: $\Sim_{\Adv}$ recognizes a record precisely when it
once relayed the honest $\op{Send}$ that produced it, which is the only way an entry gets into
$\V{fid}$ at all.

Second, $\V{gen}$ is a shadow of $\rho$'s own flag and not a reading of the internal $\Fpki$'s
$\V{VK}$, though at an honest party the two agree --- $\Fpki$'s $\op{Register}$ turns an $A$- or
$Z$-caller's call there into a peek at its line~\ref{ln:rpkipeek}, so only $\rho$ ever writes that
entry. They part company at a corrupt one, which is why the shortcut is not taken: there $\rho$'s
$\op{Register}$ is mediated away before $\Fpki$'s core runs, leaving $\V{VK}[\P]$ unset while
$\rho$ sets $\V{gen}[\P] \gets 1$ regardless, $\rho.\op{Send}$'s line~\ref{ln:rgen} not reading
the answer. A simulator
testing $\V{VK}$ would re-run $\op{Gen}$ and $\op{Register}$ on the party's every later send and
show the hosted $\Adv$ a pair of queries the real $\Adv$ never sees.

Third, line~\ref{ln:simhonest} puts \emph{every} matching honest record into $S$, ripe ones
included, where $\Facz$ would have taken the ripe ones through $R$ on its own. This is deliberate
and costs nothing: $\Clean_{\op{addr}}$ admits any undrained entry addressed to $Q$, and
$\Facz.\op{Fetch}$ unions $R$ with $S$ at its line~\ref{ln:aczdeliver} and drains their union, so
naming a ripe entry twice delivers and drains it once. Naming it is also all $\Sim_{\Adv}$ can do, having no clock read of
its own inside the query with which to tell ripe from early.

\begin{theorem}[$\rho$ realizes the authenticated channel]\label{thm:realac}
Let $\pi := \{\rho,\Fpki,\Fnet,\Fsig^{!}\}$ with parameters matched so that $\pi \cup \{\Gclock\}$
is well-formed, $\rho.\admits = \Facz.\admits$, and $\Fnet.\pars = \Facz.\pars = \Delta$. Then
$\pi$ UC-emulates $\{\Facz\}$ within $0$: for every occupant $\Adv$ of the adversary slot the
simulator $\Sim_{\Adv}$ above gives, for every $\Zenv \in \ZenvSet_{\pi,\{\Facz\}}$,
\[
  \opl{Exec}\bigl(\pi,\Zenv,\Adv,\Corr\bigr)
  \;\equiv\;
  \opl{Exec}\bigl(\{\Facz\},\Zenv,\Sim_{\Adv},\Corr\bigr) ,
\]
equal as distributions. The quantifier is over \emph{every} adversary, not the dummy alone: every
slot call in either world is responsive, so Theorem~\ref{thm:dummy} does not reach this and the
simulator is exhibited against each $\Adv$ directly.
\end{theorem}

\begin{proof}
% BEGIN overview thm:realac
One induction on activations, carrying an invariant in three parts. The coupling pairs machine to
machine as in Lemma~\ref{lem:absorb}, so the routing, claims and refusal arguments are that
lemma's; what is new is the state carried across the $\rho$/$\Facz$ cut, and the invariant on it is
the whole of the argument. \emph{The invariant} states the three parts: the internal hybrids hold
what the real ones hold, $\Sim_{\Adv}$'s shadows hold what $\rho$ holds, and $\Facz.\V{list}$ holds
exactly the undrained honest sends. Three cases preserve it. \emph{Send} shows the two worlds stamp
a send at the same time, and that replaying $\rho$'s own steps into the internal copies keeps them
equal. \emph{Fetch} carries the weight: it splits not on the sender's current status but on whether
an undrained $\V{list}$ entry matches the record, and shows that the unmatched case forces a
corrupt sender --- which is what unforgeability buys here, and why $\Clean_{\op{inj}}$ never trips.
\emph{Corrupt-party traffic} covers the one activation the declared outward set withholds, and
shows that withholding it is what the invariant needs.
% END overview thm:realac
Couple the two executions on equal tapes, paired machine to paired machine as in
Lemma~\ref{lem:absorb}: $\Zenv$, $\Corr$ and $\Gclock$ each to itself, the hosted $\Adv$ and the
internal hybrids of $\Sim_{\Adv}$ to the real $\Adv$ and the real hybrids, and $\rho$ together
with what the hybrids hold to $\Facz$. Everything below is one induction on activations, in the
shape of Definition~\ref{def:regcorr}; the claims and refusal paragraphs are
Lemma~\ref{lem:absorb}'s verbatim, and its routing paragraph holds of every call but one ---
$\Sim_{\Adv}$ is a bundle, and the calls that cross its boundary are the ordinary outward and
inward relays, save those the hosted $\Adv$ places on $\rho$'s own identities, which the declared
outward set withholds and the last handler serves inside --- so only the state carried across the
$\rho$/$\Facz$ cut is new, and the invariant on it is the whole of the argument.

\emph{The invariant.} At every corresponding pair of configurations: (i) the internal $\Fpki$,
$\Fnet$ and $\Fsig^{!}$ of $\Sim_{\Adv}$ hold the same tables as the real hybrids, having been
initialised alike and driven by the same calls in the same order; (ii) $\Sim_{\Adv}$'s shadow
$\V{gen}$, $\V{seq}$ and $\V{seen}$ equal $\rho$'s; and (iii) $\Facz.\V{list}$ holds exactly the
records $(P,Q,\tau,\msg)$ of the sends $\rho$ made through a party's own $\op{Send}$ while that
party was honest --- one entry per such $\rho.\op{Send}$, at the timestamp it read --- that no
fetch has since drained, and no others. Later corruption of the sender does not disturb the entry,
corruption being monotone and the record already placed. (i) and (ii) hold because paired copies
run the same code on equal tapes under the same driving: at an honest party $\rho$'s core drives
the real hybrids and $\Sim_{\Adv}$'s handlers drive their copies through the same steps, and at a
corrupt one both sides mediate the same subroutine calls away before any callee's core runs, as
the third case below checks. (iii) is what the three cases preserve.

\emph{Send.} $\Zenv$ calls $\op{Send}(Q,\msg)$ at an honest $\P$. On the ideal side
$\Facz.\op{Send}$ reads the clock, appends $(\P,Q,\tau,\msg)$ to $\V{list}$, and notifies
$\Sim_{\Adv}$ responsively. On the real side $\rho.\op{Send}$ reads the clock, and its internal
$\Fnet.\op{Send}$ reads it again; the two clock reads coincide, by the argument of
Section~\ref{sec:netprops} --- no core of either execution runs between them, $\Gclock.\op{Read}$
placing no call --- so the send is stamped $\tau$ on both sides. $\Sim_{\Adv}$, notified, replays
exactly the $\op{Gen}$/$\op{Register}$/$\op{Sign}$/$\op{Send}$ steps $\rho$ runs, into copies that
(i) keeps equal to $\rho$'s --- the internal $\Fnet$'s own clock read answered with the $\tau$ the
notification carried, which is that same stamp, since the responder may place no call of its own
to make it --- and the hosted $\Adv$ answers each $\Adv^{!}$ query as the real $\Adv$ does over
equal tapes. Both return $\ok$; one new $\V{list}$ entry matches the one send; (iii) holds.

\emph{Fetch.} $\Zenv$ calls $\op{Fetch}()$ at $Q$. Everything turns on the two worlds delivering
the same set $M$, and on $\Sim_{\Adv}$ being \emph{able} to name it within $\Clean_{\op{inj}}$.
Run $\rho.\op{Fetch}$'s filter, which $\Sim_{\Adv}$ mirrors internally: a delivered record
$(P,Q,\tau,\msg)$ is one whose diffused tuple $(x,\sigma)$, $x = (P,Q,\tau,s,\msg)$, is addressed
to $Q$, carries a signature verifying under $\V{VK}_{\Fpki}[P]$, and is fresh in $\V{seen}[Q]$.
Split not on $P$'s current status but on whether an undrained $\V{list}$ entry matches the
record --- the honest case and the injected case being exactly these two.

If one matches: by (iii) it is a send $\rho$ made for $P$ while $P$ was honest, and $\Facz$
delivers it from store --- through $R$ if ripe, and through $S$, its id admitted by
$\Clean_{\op{addr}}$ as addressed to $Q$, if $\Sim_{\Adv}$ releases it early to match $\Fnet$'s
early release. Whether or not the sender is corrupt by now, the record is delivered from
$\V{list}$ and never injected, so none is delivered twice.

If none matches: then $P$ is corrupt at the fetch, and this is where unforgeability earns its
place. Were $P$ honest at the fetch it would be honest throughout, by monotonicity, so
$\V{VK}_{\Fpki}[P]$ would be $P$'s own $\Fsig$ key --- $\rho$ registers exactly what $\op{Gen}$
returned --- and a signature verifying under an honest party's key is never a forgery:
$\op{Verify}$'s honest-key test (Section~\ref{sec:fsigfunc}) sets $b \gets 0$ on any fresh triple
under such a key, so $b = 1$ forces the triple recorded, and only an honest $\op{Sign}$ records
one, which by (iii) would leave a matching $\V{list}$ entry --- against the case. So $P \in \Cs$;
the record is the adversary's to have produced, an $\Fnet$ injection or a signature made as a
corrupt party, and $\Sim_{\Adv}$ places $(P,Q,\tau,\msg)$ in $I$, which $\Clean_{\op{inj}}$ admits,
$P$ being corrupt and $Q$ the fetcher. The channel $\Facz$ delivers it unstored, as $\rho$ hands it up. This
is what unforgeability buys: no record lacking a genuine send behind it can ever wear an honest
name, so the injection channel is only ever asked for corrupt senders, and $\Clean_{\op{inj}}$
never trips.

Freshness matches on both sides by (ii), and draining matches: a stored record delivered is
drained from $\V{list}$ ideally and marked seen really, so $\Fnet$'s re-delivery of the same tuple
is filtered by $\V{seen}$ really and finds no $\V{list}$ entry ideally, falling into neither case a
second time. So the two $\op{Fetch}$es return the same $M$ and step to corresponding
configurations; (iii) is preserved, its stored entries drained in lockstep.

\emph{Corrupt-party traffic.} $\Adv$ calls $\op{Send}$ or $\op{Fetch}$ at a corrupt $\P$, on
$\rho$'s own identity --- the one activation the outward set withholds, and the case the two
above do not reach. The full interface admits it at line~\ref{ln:corrgate}, $\P$ being corrupt,
and neither hook fires on an $A$-caller, so on the real side $\rho$'s core runs. Every subroutine
call that core then places carries $\rho$'s name at a corrupt party, so line~\ref{ln:medin} of
that same interface hands each to $\Adv$ before the callee's core runs: nothing of the activation
reaches $\Fnet$'s buffer, $\Fpki$'s directory or $\Fsig$'s tables, and its very timestamp is
$\Adv$'s to supply. $\Sim_{\Adv}$'s line~\ref{ln:simcorr}
runs that same core against the hosted $\Adv$ on the same shadows, so both sides move $\V{gen}$,
$\V{seq}$ and $\V{seen}$ alike and leave the hybrids untouched, preserving (i) and (ii); and
neither side touches $\V{list}$ --- ideally because the call never leaves $\Sim_{\Adv}$, really
because there is no $\Facz$ to touch --- so (iii) is preserved with nothing to check.

Withholding these calls is what (iii) needs, and it is worth saying why rather than leaving it to
the routing. Relayed outward the call would reach $\Facz.\op{Send}$, which files a record at the
true time; the ripeness test $\Facz$ inherits at $\Fac.\op{Fetch}$'s line~\ref{ln:acripe} would
then deliver it to $Q$ within $\Delta$ whatever the slot answered, where the real world delivers
nothing at all unless $\Adv$ injects on its own account.
The entry would break (iii) the moment it was filed, and no later answer of $\Sim_{\Adv}$'s could
retract it --- $R$ being $\Facz$'s to compute. This is the ripening that
Section~\ref{sec:facfunc} charges to a corrupt sender speaking through $\Fac.\op{Send}$, and a
protocol whose corrupt parties speak through $\Fnet$ instead does not pay it.

No other activation touches the cut: a call among $\Zenv$, $\Gclock$, $\Corr$ and $\Adv$ that
targets none of $\rho$'s identities is routed and claimed identically on both sides by
Lemma~\ref{lem:absorb}, and reads no state the invariant governs. Corresponding halting
configurations return $\Zenv$'s one output. The
executions are equal as distributions, and the advantage is $0$ --- not a bound but an identity,
every step deterministic once the tapes are fixed and no forgery ever slipping through to force
an injection the sanitizer would refuse.
\end{proof}

The bound is $0$, and it is worth being clear on why nothing weaker was available and nothing
stronger was needed. Every subroutine is ideal, so there is no computational gap to charge for;
the reduction is an identity of distributions, as the composition theorem's is. Transfer, note,
does not run here: Corollary~\ref{cor:proptransfer2} wants cores placing no responsive call, and
$\rho$'s reach $\Fsig^{!}$, $\Fnet$ and the clock all responsively, so a property of $\Facz$ does
not descend to $\pi$ for free --- it must be re-proved over $\pi$, exactly as
Section~\ref{sec:clockprops} warned for anything built over the notifying clock. What the theorem
gives instead is the stronger thing for a realization: not a bound against the dummy but an
equality against every adversary.

\begin{corollary}[Against non-corrupting classes, $\rho$ realizes $\Fac$ itself]\label{cor:realac}
Instrument $\Facz$ to set $\V{bad} \gets \true$ at line~\ref{ln:aczdeliver} when $I \neq
\emptyset$, and let $\Fac^{\flat}$ be the pair-shaped $\Fac$ that ignores $I$ --- its fetch
$\Fac$'s, taking $(S,I)$ and using $S$ alone. Now $\Facz$ and $\Fac^{\flat}$ are identical until bad
(Definition~\ref{def:iub}), the injection the one bad-guarded difference. So for occupants and
environments that corrupt nobody, $\Clean_{\op{inj}}$ forces $I = \emptyset$ and $\Pr[\Bad] = 0$,
whence by Lemma~\ref{lem:iub} the two are indistinguishable and $\pi$ UC-emulates
$\{\Fac^{\flat}\}$, hence $\{\Fac\}$, within $0$ against the non-corrupting classes. Corrupt a
party and the guarantee is $\Facz$'s, no less: the honest-sender half of $\op{FinAuth}$
(Section~\ref{sec:facprops}) transfers, its $\op{Send}$-tracing argument reading only honest
records, while the full $\op{FinAuth}$ fails of $\Facz$ exactly as it does of $\Fnet$, and for the
same reason.
\end{corollary}

\begin{remark}[The weakening is forced]\label{rem:acimposs}
No protocol over $\{\Fpki,\Fnet,\Fsig^{!}\}$ realizes $\Fac$ as shipped, so $\Facz$ is not a
convenience but the true target. Suppose one did. A corrupt party diffuses, inside a fetch's
responsive window, a well-formed item its receiver accepts --- $\Fnet$ grants exactly this, the
wire being the adversary's to write (Section~\ref{sec:fnetfunc}) --- carrying a timestamp
$\tau^{*}$ of the adversary's choosing. On the ideal side the accepting fetch is atomic, so the
simulator must have called $\Fac.\op{Send}$ strictly earlier to place a matching entry; but $\Fac$
stamps that call the true time it was made, not $\tau^{*}$, and the receiver reads $\tau^{*}$ off
the delivered record. An environment that sends at one time and has its corrupt confederate claim
another distinguishes with certainty. The freedom is $\Fnet$'s to give and no receiver-side check
removes it, since the signature validates whatever time the payload carries; only weakening the
ideal object to admit it --- which is $\Facz$ --- closes the gap. This is the diffusion analogue
of the impossibility that forces $\Fnet$'s own injections to sit outside its delivery guarantee.
\end{remark}

This is the paper's longest single construction, and the point of running it in full is that
every layer of the framework shows up load-bearing. Composition supplies the cut between $\rho$
and its hybrids; the responsive discipline of Chapter~\ref{app:responsive} makes both fetches
atomic so the schedules can be coupled at all; sanitization confines a corrupt slot to corrupt
senders; tightness would be what a property of $\Facz$ needs to descend, and its absence under
responsive cores is why the realization is proved directly rather than inherited; and the
fundamental lemma of Chapter~\ref{app:iub} is what turns the exact realization of the weaker
object into the exact realization of the stronger one wherever no corruption is in play. The
authenticated channel, postulated in Chapter~\ref{sec:fac}, is now a theorem.

%% ==================================================================
%% ==================================================================
\chapter{\texorpdfstring{$\Fsig$}{F\_Sig}: A Deeper Dive}\label{sec:sigdeep}
\chapterhead{F-Sig: A Deeper Dive}

Chapter~\ref{sec:fsig} gave a signature functionality and proved two properties of it. This
chapter asks the three questions that leaves open. What does a \emph{scheme} have to satisfy for
a protocol built on it to inherit those properties --- and what does the protocol look like, drawn
in this framework's own terms, with its randomness and its state held where
Section~\ref{sec:leakage} says they must be? Does the converse hold: are the two properties
\emph{enough} to characterize $\Fsig$, so that anything with its interfaces satisfying them
realizes it? And is $\Fsig$ the strongest object of its shape, or does a realization of it satisfy
something it does not?

The answers are: yes with a scheme's own security as the bound; no, and instructively so; and no
--- a realization can be strictly stronger, which is worth knowing about ideal functionalities in
general.

\section{طرح‌های امضا و بازی‌هایشان}\label{sec:dssyntax}

A \emph{signature scheme} is a triple of algorithms $\DS = (\op{Gen},\op{Sign},\op{Ver})$ over a
key space $\Keys$, a message space $\Msgs$, a signature space $\Sigs$ and a coin space
$\{0,1\}^{n}$, with $\op{Gen}$ and $\op{Sign}$ taking their coins explicitly:
\[
  \begin{gathered}
  (\vk, \V{sk}) \gets \op{Gen}(r) , \qquad
  \sigma \gets \op{Sign}(\V{sk}, \msg; r) , \\[3pt]
  b \gets \op{Ver}(\vk, \msg, \sigma) ,
  \end{gathered}
\]
$\op{Ver}$ deterministic. The coins are arguments rather than an implicit tape because $\psig$
below draws them through $\Frand$, and a scheme whose randomness is hidden inside it could not be
wired to a randomness functionality at all.

Its three security notions are games in the ordinary property-based sense, each played by an
adversary $\mathcal{B}$ against $\DS$ alone --- no framework, no execution, the notions a
cryptographer would write down anyway, unforgeability under chosen-message attack chief among
them~\cite{GMR88}. One shell, one signing oracle, shared by all three in the
same idiom the paper's own game systems share a shell across several finalizations
(Definition~\ref{def:prop}): $\mathcal{B}$ receives $\vk$, queries $\op{Sign}$ as it likes, then
halts by calling exactly one of the three procedures below, which decides whether it wins.

\begin{procedure}{بازی برای طرحِ امضای $\DS$: پوسته}
\opsig{$\op{Initialize}()$}:
\begin{algorithmic}[1]
\setcounter{contline}{0}\algcont
\State $(\vk,\V{sk}) \gets \DS.\op{Gen}(r)$ \Comment{$r$ یکنواخت} \label{ln:dsgen}
\State \textbf{return} $\vk$
\algsave
\end{algorithmic}
\vspace{1.4ex}

\opsig{$\op{Sign}(\msg)$}:
\begin{algorithmic}[1]
\algcont
\State $\sigma \gets \DS.\op{Sign}(\V{sk},\msg;r)$ \Comment{$r$ یکنواخت} \label{ln:dssign}
\State $\V{tr} \gets \V{tr} \cdot (\msg,\sigma)$
\State \textbf{return} $\sigma$
\algsave
\end{algorithmic}
\end{procedure}

\begin{procedure}{Finalizations of the game, one per notion; the shell above is shared}
\opsig{$\op{Cor}(\msg)$} \quad (correctness)
\begin{algorithmic}[1]
\algcont
\State $\sigma \gets \op{Sign}(\msg)$
  \Comment{a fresh call, on the game's own oracle}
\State $w \gets 1$ if $\DS.\op{Ver}(\vk,\msg,\sigma) \neq 1$; \quad else $w \gets 0$ \label{ln:dscor}
\State \textbf{return} $w$
\algsave
\end{algorithmic}
\vspace{1.4ex}

\opsig{$\op{Uf}(\msg,\sigma)$} \quad (unforgeability)
\begin{algorithmic}[1]
\algcont
\State $w \gets 1$ if $\DS.\op{Ver}(\vk,\msg,\sigma) = 1 \ \wedge\ (\msg,\sigma) \notin \V{tr}$;
       \quad else $w \gets 0$ \label{ln:dsuf}
\State \textbf{return} $w$
\algsave
\end{algorithmic}
\vspace{1.4ex}

\opsig{$\op{Col}(\msg,\msg',\sigma)$} \quad (collision resistance)
\begin{algorithmic}[1]
\algcont
\State $w \gets 1$ if $\msg \neq \msg' \ \wedge\ \DS.\op{Ver}(\vk,\msg,\sigma) =
       \DS.\op{Ver}(\vk,\msg',\sigma) = 1$; \quad else $w \gets 0$ \label{ln:dscol}
\State \textbf{return} $w$
\algsave
\end{algorithmic}
\end{procedure}

Write $\varepsilon_{\op{cor}}(\mathcal{B})$, $\varepsilon_{\op{uf}}(\mathcal{B})$ and
$\varepsilon_{\op{col}}(\mathcal{B})$ for the probabilities that $\op{Cor}$, $\op{Uf}$ and
$\op{Col}$ return $1$. Perfect correctness is $\varepsilon_{\op{cor}} = 0$, which most schemes
have and none of what follows needs. Unforgeability is \emph{strong}~\cite{ADR02}: line~\ref{ln:dsuf} excludes
the pair the oracle returned rather than the message alone, matching $\opl{FinUF}$'s reading in
Section~\ref{sec:sigprops}. Collision resistance is not standard and is not implied by the other
two; Section~\ref{sec:sigcoll} is where it earns its place.

\section{The protocol \texorpdfstring{$\psig$}{π\_Sig}}\label{sec:psig}

$\psig$ carries $\Fsig$'s three interfaces and holds nothing of its own: every coin is drawn
through $\Frand$ and every value that outlives an invocation sits in $\Fstore$, which is
Section~\ref{sec:leakage}'s \emph{leakage-well-formed} shape and the first protocol in this paper
written to it. That is not decoration. Leakage is gated on the party being corrupt
(line~\ref{ln:leakcorr}) but a functionality's $\opl{Leak}$ hands over its state \emph{entire}, so
a single $\Frand$ shared across parties would surrender every party's coins the moment any one
party was corrupted, and with them every signing key --- no unforgeability could survive it.
The protocol $\psig$ therefore addresses \emph{one instance of each per party}, $\Frand^{P}$ and $\Fstore^{P}$
with $\Ps = \{P\}$, which Remark~\ref{rem:gminstances} already allows and the guard's blindness to
the instance number already reaches. And it erases each coin the moment it has been used, which is
what $\Frand$'s $\op{Erase}$ was put there for (Section~\ref{sec:frandfunc}): what has been
erased cannot leak, so a corrupt party surrenders its own signing key and nothing else --- not
even its own past coins.

\begin{interface}[nofloat]{پروتکلِ $\psig$ روی $\{\Frand^{P},\Fstore^{P}\}_{P \in \Ps}$ \\
\params{$\PID$, \quad $\Ps$, \quad $\admits$, \quad
$\uses := \{(\Frand,\serves),(\Fstore,\serves)\}$, \quad $\pars := \DS$}}

{\scriptsize
\setlength{\columnsep}{1.4em}
\begin{multicols}{2}
\opsig{$\id.\op{Gen}()$} \quad از $\id'$
\begin{algorithmic}[1]
\setcounter{contline}{0}\algcont
\State $\vk \gets \id_{\Fstore}.\fopl{Get}(\atom{vk}) \text{ as } \id$ \label{ln:psgetvk}
\If{$\vk \neq \rej$}
  \State \textbf{return} $\vk$
    \Comment{یک کلید به‌ازای هر طرف}
\EndIf
\State $(i,r) \gets \id_{\Frand}.\fopl{Rnd}() \text{ as } \id$
\State $(\vk,\V{sk}) \gets \DS.\op{Gen}(r)$ \label{ln:psgen}
\State $\id_{\Frand}.\fopl{Erase}(i) \text{ as } \id$ \label{ln:pserase}
  \Comment{سکه‌ها نمی‌توانند نشت کنند}
\State $\id_{\Fstore}.\fopl{Put}(\atom{vk},\vk) \text{ as } \id$
\State $\id_{\Fstore}.\fopl{Put}(\atom{sk},\V{sk}) \text{ as } \id$
\State \textbf{return} $\vk$
\algsave
\end{algorithmic}
\vspace{1.4ex}

\opsig{$\id.\op{Verify}(\vk,\msg,\sigma)$} \quad از $\id'$
\begin{algorithmic}[1]
\algcont
\State \textbf{return} $\DS.\op{Ver}(\vk,\msg,\sigma)$ \label{ln:psver}
\algsave
\end{algorithmic}

\columnbreak

\opsig{$\id.\op{Sign}(\msg)$} \quad از $\id'$
\begin{algorithmic}[1]
\algcont
\State $\V{sk} \gets \id_{\Fstore}.\fopl{Get}(\atom{sk}) \text{ as } \id$
\If{$\V{sk} = \rej$}
  \State \textbf{return} $\none$
    \Comment{هنوز کلیدی نیست}
\EndIf
\State $(i,r) \gets \id_{\Frand}.\fopl{Rnd}() \text{ as } \id$
\State $\sigma \gets \DS.\op{Sign}(\V{sk},\msg;r)$ \label{ln:pssign}
\State $\id_{\Frand}.\fopl{Erase}(i) \text{ as } \id$
  \Comment{به‌محضِ استفاده پاک می‌شود}
\State \textbf{return} $\sigma$
\algsave
\end{algorithmic}
\vspace{1.4ex}

\opsig{$\id.\op{Leak}()$} \quad از $\id'$
\begin{algorithmic}[1]
\algcont
\State \textbf{return} $\none$
  \Comment{چیزی ندارد که نشت کند}
\algsave
\end{algorithmic}
\end{multicols}}
\end{interface}

The draws come back as pairs because $\Frand$'s $\op{Rnd}$ answers with the index it wrote
beside the string (line~\ref{ln:rndreturn}), which is what lets $\psig$ name the entry it then
erases; $\op{FinUnp}$ reads the strings and is indifferent to the index travelling beside them.

Three remarks. $\psig$'s own $\opl{Leak}$ returns $\none$ because it has no state to give: the key
is $\Fstore^{P}$'s, and a corrupt party's adversary reads it there, which is precisely what
leakage-well-formedness asks --- the leak is where the state is, and nowhere else. The operation $\op{Verify}$
places no call at all, $\op{Ver}$ being deterministic, so a verification is local and neither
draws nor stores; that is why $\psig$ is the one protocol here whose $\op{Verify}$ cannot be
distinguished from $\Fsig$'s by the token's movement. And the key is fetched from the store rather
than cached in $\psig$, which is what makes ``one key per party'' a fact about $\Fstore$'s
per-slot fixity (Section~\ref{sec:storeprops}, $\opl{FinKeep}$) rather than a promise $\psig$
makes about itself.

\begin{proposition}[$\psig$ inherits the two properties]\label{prop:psigprops}
Let $\gamma := \{\Gm\} \cup \psig \cup \{\Frand^{P},\Fstore^{P}\}_{P}$ be a game system in the
sense of Definition~\ref{def:prop}, with $\Gm$ the shell of Section~\ref{sec:sigprops}. For every
$\Zenv \in \ZenvSet_{\gamma,\gamma}$ placing at most $q$ calls on the relays and every occupant
$\mathcal{X}$ of the adversary slot there are adversaries $\mathcal{B}_{\op{cor}}$ and
$\mathcal{B}_{\op{uf}}$ against $\DS$, each running $\Zenv$ and $\mathcal{X}$ once with $O(q)$
overhead, such that
\[
  \begin{gathered}
  \Pr[\Win_{\op{FinCor}}] \;\leq\; q \cdot \varepsilon_{\op{cor}}(\mathcal{B}_{\op{cor}}) , \\[3pt]
  \Pr[\Win_{\op{FinUF}}] \;\leq\; q \cdot \varepsilon_{\op{uf}}(\mathcal{B}_{\op{uf}})
  \;+\; q^{2} 2^{-n} .
  \end{gathered}
\]
\end{proposition}

\begin{proof}
% BEGIN overview prop:psigprops
One preliminary paragraph and two reductions. The preliminary collects what the subroutines already
give: tightness makes every trace entry an answer $\psig$'s own core gave at an honest party,
Proposition~\ref{prop:storeprops} makes the store fixed and faithful so one key is issued per party
for good, and Proposition~\ref{prop:randunp} makes the coins uniform and, being erased, unreadable.
Each property then reduces to the scheme's own game by the same move: assume the verdict fires,
identify the triple it names as a failure of $\DS$ outright, and have the reduction guess in advance
which of the at most $q$ calls is the failing one --- which is where the factor $q$ comes from, in
both bounds.
% END overview prop:psigprops
By tightness no caller but $\Gm$ reaches $\psig$, so every $\atom{gen}$, $\atom{sign}$ and $\atom{ver}$
entry of $\V{tr}$ is an answer $\psig$'s own core gave at an honest party --- at a corrupt one the
game's wrapper intercepts and records nothing (Remark~\ref{rem:gmcorr}). By
Proposition~\ref{prop:storeprops} the store is fixed and faithful: the $\vk$ and $\V{sk}$ a party
puts are what it later gets, and $\op{Gen}$'s guard on line~\ref{ln:psgetvk} therefore issues one
key per party for good. By Proposition~\ref{prop:randunp} the coins are uniform and, being erased
at line~\ref{ln:pserase} and its counterpart in $\op{Sign}$, are read by nobody: $\Frand^{P}$'s
leak is empty of them and $\Frand^{P}$ serves $P$ alone, so no corruption of another party reaches
them.

For correctness, suppose line~\ref{ln:gcor} fires: honest $P$, $Q$, a $(\atom{gen},P,\vk)$, a
$(\atom{sign},P,\msg,\sigma)$ with $\sigma \in \Sigs$, and a $(\atom{ver},Q,\vk,\msg,\sigma,0)$. The
verification is $\DS.\op{Ver}(\vk,\msg,\sigma)$ by line~\ref{ln:psver}, and the signature is
$\DS.\op{Sign}(\V{sk},\msg;r)$ under the $\V{sk}$ paired with that very $\vk$ by
line~\ref{ln:psgen}, on coins uniform by the paragraph above. So the triple is a correctness
failure of $\DS$ outright. The adversary $\mathcal{B}_{\op{cor}}$ runs the execution, guesses in advance which
of the at most $q$ signing calls is the failing one, and outputs its message; it wins whenever the
game does and the guess is right, which is the factor $q$.

For unforgeability, suppose line~\ref{ln:guf} fires: honest $P$, $Q$, $(\atom{gen},P,\vk)$,
$(\atom{ver},Q,\vk,\msg,\sigma,1)$ and no $(\atom{sign},P,\msg,\sigma)$. Two cases. Either $\vk$ is
the key $\op{Gen}$ issued to $P$, and then $\op{Ver}(\vk,\msg,\sigma) = 1$ with $(\msg,\sigma)$
returned by no $\op{Sign}$ at $P$ --- a strong forgery, which $\mathcal{B}_{\op{uf}}$ harvests by
planting its challenge key at a guessed one of the at most $q$ generating parties and answering
that party's signing calls from its oracle, every other party it runs itself; or $\vk$ is
\emph{another} honest party's key that happens to equal $P$'s, and then two independent
$\op{Gen}(r)$ draws on uniform $n$-bit coins collided, which over at most $q$ generations is at
most $q^{2}2^{-n}$. Note where the erasure is spent: without it, corrupting any party would hand
$\mathcal{X}$ the coins of $P$'s own $\op{Gen}$ and $\op{Sign}$, and $\mathcal{B}_{\op{uf}}$ could
not answer with an oracle it does not have the coins for.
\end{proof}

This is the first bound in the paper that is neither $0$ nor information-theoretic, and it is
worth saying what has changed. The functionality $\Fsig$ has its properties because it was written to
(Proposition~\ref{prop:sigprops}); $\psig$ has them because $\DS$ does, and the two $q$-factors
are the ordinary price of a guessing reduction. Corollary~\ref{cor:proptransfer2} is what makes
the exchange legitimate in the other direction --- a property of $\Fsig$ descends to any
realization of it --- and this proposition is the same statement approached from the realization's
side, proved directly because $\psig$ is not obtained from $\Fsig$ by emulation but built from a
scheme.

\section{عکسِ قضیه، و آنچه ویژگی‌ها به آن نمی‌رسند}\label{sec:sigconv}

Proposition~\ref{prop:psigprops} runs from a scheme to the properties. The converse would run the
other way and much further: that the two properties \emph{characterize} $\Fsig$, so that any
functionality carrying its interfaces and satisfying them realizes it. It is a natural conjecture
--- properties are what we ask of an ideal object, so an object satisfying them ought to be as good
as the one we asked them of --- and it is false.

\begin{proposition}[The two properties do not characterize $\Fsig$]\label{prop:convfails}
There is a standard functionality $\F'$ with $\Fsig$'s fields and interfaces, having both
$\opl{FinCor}$ and $\op{FinUF}$ within $0$ against every environment and every occupant of the
slot, such that $\{\F'\}$ does not UC-emulate $\{\Fsig\}$ within any $\varepsilon < \tfrac12$.
\end{proposition}

\begin{proof}
Let $\F'$ be $\Fsig$ with one change: in $\op{Verify}$, when no honest party's key equals the
queried $\vk$ --- that is, when $\vk \neq \V{VK}[P']$ for every $P' \notin \Cs$ --- the answer is
a fresh uniform bit, recorded nowhere, so that a later query on the same triple is answered
independently. On every other triple $\F'$ is $\Fsig$ exactly, memoization included.

$\F'$ has both properties within $0$. Each of lines~\ref{ln:gcor} and~\ref{ln:guf} is existential
over a $(\atom{gen},P,\vk)$ entry with $P \notin \Cs$, so a triple whose $\vk$ is no honest party's
key witnesses neither, and on all other triples $\F'$ answers as $\Fsig$ does, whose verdicts are
$0$ by Proposition~\ref{prop:sigprops}.

It does not emulate. Let $\Zenv$ pick any $\vk^{*}$, $\msg$, $\sigma$ with $\vk^{*}$ generated by
nobody, corrupt no one, and call $\op{Verify}(\vk^{*},\msg,\sigma)$ twice at an honest party,
returning $1$ if the two answers differ. Against $\F'$ the answers are independent uniform bits,
so it returns $1$ with probability $\tfrac12$. Against $\Fsig$ the first call fixes
$\V{Ver}[\vk^{*},\msg,\sigma]$ and the second returns it, so the answers agree and it returns $1$
with probability $0$ --- whatever the simulator does, the memoization being $\Fsig$'s own code and
no value of the slot's able to unfix it. The advantage is $\tfrac12$.
\end{proof}

The diagnosis is more useful than the counterexample. Emulation pins an object's whole observable
behaviour; the two properties pin two things about it --- that honestly produced signatures verify,
and that forgeries under honest keys do not. Everything else is out of their reach, and the gap
has a shape: both verdicts are existential over $(\atom{gen},P,\vk)$ with $P$ honest, so
\emph{nothing whatever} is said about keys no honest party generated, and nothing is said about
consistency, one answer being all either verdict inspects. What is missing is not a stronger
notion of unforgeability but two further properties, and they are the ones $\Fsig$'s code makes
obvious once looked for: its $\op{Verify}$ memoizes, and its $\op{Gen}$ issues one key per party.

\begin{interface}[nofloat]{Two further finalizations of the shell of Section~\ref{sec:sigprops}}
\opsig{$\id.\opdef{FinCons}()$} \quad from $\id'$ \quad (consistency)
\begin{algorithmic}[1]
\setcounter{contline}{0}\algcont
\Req{$\neg\done$}; \ \Req{$\id.P = Z$}; \ $\done \gets \true$
\State $w \gets 1$ if $\exists\, Q,Q',\vk,\msg,\sigma :\ \
       (\atom{ver},Q,\vk,\msg,\sigma,0) \in \V{tr}$ \label{ln:gcons}
\State \quad ${}\wedge\ (\atom{ver},Q',\vk,\msg,\sigma,1) \in \V{tr}$; \quad else $w \gets 0$
\State \textbf{return} $w$
\algsave
\end{algorithmic}
\vspace{1.4ex}

\opsig{$\id.\opdef{FinOne}()$} \quad from $\id'$ \quad (one key per party)
\begin{algorithmic}[1]
\algcont
\Req{$\neg\done$}; \ \Req{$\id.P = Z$}; \ $\done \gets \true$
\State $\Cs \gets (C,Z).\fopl{Status}() \text{ as } \id$
\State $w \gets 1$ if $\exists\, P \notin \Cs,\ \vk \neq \vk' :\ \
       (\atom{gen},P,\vk) \in \V{tr} \ \wedge\ (\atom{gen},P,\vk') \in \V{tr}$; \label{ln:gone}
\State \quad else $w \gets 0$
\State \textbf{return} $w$
\algsave
\end{algorithmic}
\end{interface}

Both take the search reading, both are attained at $0$ by $\Fsig$ --- $\opl{FinCons}$ because
$\V{Ver}$ is written once per triple and read thereafter, $\opl{FinOne}$ because $\V{VK}[P]$ is
written once and read thereafter, each by the same first-write-wins argument
Proposition~\ref{prop:pkifix} makes for the directory --- and both are attained at $0$ by $\psig$,
by $\op{Ver}$'s determinism and $\Fstore$'s $\op{FinKeep}$ respectively. Neither asks anything of
a corrupt party: $\op{FinCons}$ names no party at all, the game's own wrapper having already kept
corrupt answers out of $\V{tr}$.

\begin{theorem}[The converse, from the four properties]\label{thm:sigconv}
Let $\F'$ be a standard functionality with $\Fsig$'s fields and interfaces, whose $\op{Sign}$
returns $\none$ exactly when no key has been generated at the calling party, and let $\F'$ have
$\op{FinCor}$, $\op{FinUF}$, $\op{FinCons}$ and $\op{FinOne}$ within $0$ against
$\ZenvSet_{\gamma,\gamma}$ and every occupant of the slot, for $\gamma$ the game system of
Section~\ref{sec:sigprops} with $\F'$ in place of $\Fsig$. Then $\{\F'\}$ UC-emulates
$\{\Fsig\}$ within $0$: for every occupant $\Adv$ there is a simulator $\Sim_{\Adv}$ with
\[
  \opl{Exec}\bigl(\{\F'\},\Zenv,\Adv,\Corr\bigr)
  \;\equiv\;
  \opl{Exec}\bigl(\{\Fsig\},\Zenv,\Sim_{\Adv},\Corr\bigr)
\]
for every $\Zenv \in \ZenvSet_{\{\F'\},\{\Fsig\}}$.
\end{theorem}

\begin{proof}
% BEGIN overview thm:sigconv
The simulator is the obvious one --- $\Adv$ bundled with an internal copy of $\F'$, answering each
of $\Fsig$'s three slot queries by putting the same call to that copy --- so the whole of the
argument is that the two sides never disagree on a returned value. Since $\Fsig$'s answer \emph{is}
the slot's answer unless its own code overrides, that reduces to showing no override ever fires,
and there are exactly four: $\op{Gen}$'s memoization, $\op{Gen}$'s sanitizer, $\op{Sign}$'s
sanitizer, and $\op{Verify}$'s honest-key test. Each is taken in turn and excluded by one of the
four hypothesised properties at $0$, the sanitizer cases splitting further by which repair the
sanitizer would make. That is where all four properties are spent, and it is why the theorem needs
exactly them.
% END overview thm:sigconv
The simulator $\Sim_{\Adv}$ bundles $\Adv$ with an internal copy of $\F'$ and runs it on the calls $\Fsig$
reports. The functionality $\Fsig$ asks the slot for exactly the three values it does not pin --- the key at
$\op{Gen}$, the signature at $\op{Sign}$, the bit at $\op{Verify}$ --- and $\Sim_{\Adv}$ answers
each by putting the same call to its internal $\F'$ at the same identity and forwarding what comes
back, letting the hosted $\Adv$ see, and answer, whatever $\F'$'s own code would put to a slot.
Leakage and corruption instructions it serves from the internal copy likewise; note that this is
why $\F'$'s $\opl{Leak}$ need not resemble $\Fsig$'s at all, the environment reaching it only
through the slot, which $\Sim_{\Adv}$ occupies.

Couple the two executions on equal tapes as in Lemma~\ref{lem:iub}. Everything reduces to one
claim: at every call, $\Fsig$ returns what the internal $\F'$ returned, so that the two sides
agree on the value the environment sees. Since $\Fsig$'s answer is the slot's answer \emph{unless}
its own code overrides, it is enough that no override ever fires, and there are exactly four.

$\op{Gen}$'s memoization and $\op{Gen}$'s sanitizer. If $\V{VK}[\id.P]$ is already set $\Fsig$
returns it; $\op{FinOne}$ at $0$ says the internal $\F'$ returns the same key at that party too,
so the two agree. Otherwise $\San[\Clean_{\vk}]$ replaces the offered $\vk$ if it is not a key, or
duplicates one already issued, or has a signature already recorded valid under it. Not a key: then
$\F'$ answered $\op{Gen}$ with a non-key, and a later $\op{Sign}$ at that party would carry
$\sigma \in \Sigs$ with $(\atom{ver},\cdot,\vk,\cdot,\cdot,0)$ available to the environment ---
excluded by $\op{FinCor}$ at $0$ once the environment signs and verifies, which it may. Duplicate:
two honest parties then share a key, and signing at one and verifying at the other wins
$\op{FinUF}$, excluded. Pre-signed: some $\V{Ver}[\vk,\msg,\sigma] = 1$ stands before $\vk$ is
issued, and $\Fsig$ records a $1$ only at $\op{Sign}$, so an honest party signed under a key not
yet generated --- which $\Fsig$'s $\op{Sign}$ refuses, returning $\none$, and which the
hypothesis on $\F'$'s $\op{Sign}$ refuses too. So the sanitizer is idle.

$\op{Sign}$'s sanitizer. The sanitizer $\San[\Clean_{\sigma}]$ replaces $\sigma$ if it is not a signature or if
$\V{Ver}[\vk,\msg,\sigma] = 0$ already. Not a signature: $\op{FinCor}$'s $\sigma \in \Sigs$
conjunct is exactly the reading that lets this case be excluded, the environment verifying the
returned value and $\op{FinCor}$ forbidding a $0$. Already $0$: then the same triple verified $0$
earlier and verifies $1$ now --- the internal $\F'$ having produced a signature it must accept ---
which is $\op{FinCons}$, excluded at $0$. So it too is idle.

$\op{Verify}$'s memoization and its honest-key suppression. The first returns the recorded verdict
where one exists, and $\op{FinCons}$ at $0$ says the internal $\F'$ answers the same triple the
same way, so they agree; this is the override the counterexample of
Proposition~\ref{prop:convfails} exploited, and it is exactly what $\op{FinCons}$ closes. The
second sets $b \gets 0$ on a fresh triple under an honest party's key, and $\op{FinUF}$ at $0$ says
the internal $\F'$ answers $0$ there too --- a $1$ being its win condition. So both agree.

No override fires, every value the environment sees is the internal $\F'$'s own, and the coupled
runs proceed in step to the same halt, exactly as in Lemma~\ref{lem:iub}'s no-flag case. The
distributions are equal and the advantage is $0$.
\end{proof}

So the converse holds, and what it took to hold is the point. Two properties were not a
characterization; four are, and the two added are not deep facts about signing but the two places
$\Fsig$'s code \emph{commits} --- a table written once and read thereafter, twice over. A
functionality's properties characterize it exactly when they cover every such commitment, and
reading a specification for its commitments is a better way to find the missing property than
asking what else one would like to be true.

\section{A property of \texorpdfstring{$\psig$}{π\_Sig} that \texorpdfstring{$\Fsig$}{F\_Sig} lacks}\label{sec:sigcoll}

The last question is whether $\Fsig$ is the strongest object of its shape. It is not, and the
gap is not an oversight: it is what sanitization costs.

\begin{interface}[nofloat]{A further finalization: collision resistance}
\opsig{$\id.\opdef{FinColl}()$} \quad from $\id'$ \quad (collision resistance)
\begin{algorithmic}[1]
\setcounter{contline}{0}\algcont
\Req{$\neg\done$}; \ \Req{$\id.P = Z$}; \ $\done \gets \true$
\State $\Cs \gets (C,Z).\fopl{Status}() \text{ as } \id$
\State $w \gets 1$ if $\exists\, P \notin \Cs,\ Q,Q' \notin \Cs,\ \vk,\ \msg \neq \msg',\ \sigma :$ \label{ln:gcoll}
\State \quad $(\atom{gen},P,\vk) \in \V{tr}
       \ \wedge\ (\atom{ver},Q,\vk,\msg,\sigma,1) \in \V{tr}$
\State \quad ${}\wedge\ (\atom{ver},Q',\vk,\msg',\sigma,1) \in \V{tr}$; \quad else $w \gets 0$
\State \textbf{return} $w$
\algsave
\end{algorithmic}
\end{interface}

It takes the search reading, winning being meant to be hard rather than no better than a guess.

\begin{proposition}[$\psig$ has it, $\Fsig$ does not]\label{prop:sigcoll}
$\psig$ has $\opl{FinColl}$ within $q \cdot \varepsilon_{\op{col}}(\mathcal{B}_{\op{col}})$ for an
adversary $\mathcal{B}_{\op{col}}$ against $\DS$ built as in
Proposition~\ref{prop:psigprops}. The functionality $\Fsig$ does not have it within any $\varepsilon < 1$: there
are $\Zenv^{*}$ and an occupant $\mathcal{X}^{*}$ with $\Pr[\Win_{\op{FinColl}}] = 1$.
\end{proposition}

\begin{proof}
For $\psig$: a winning trace exhibits one $\sigma$ verifying under an honest party's $\vk$ on two
distinct messages, and verification is $\DS.\op{Ver}$ by line~\ref{ln:psver}, so the pair is a
collision for $\DS$ under a key drawn on uniform coins. The adversary $\mathcal{B}_{\op{col}}$ plants its
challenge key at a guessed one of the at most $q$ generating parties, answers that party's signing
from its oracle and runs the rest itself, and outputs the two messages and the signature.

For $\Fsig$: let $\mathcal{X}^{*}$ answer every $\op{Gen}$ query with a fresh key, every
$\op{Verify}$ query with $0$, and both $\op{Sign}$ queries with one fixed $\sigma^{*} \in \Sigs$.
Let $\Zenv^{*}$ corrupt nobody, call $\op{Gen}()$ then $\op{Sign}(\msg)$ and $\op{Sign}(\msg')$
for $\msg \neq \msg'$ at an honest party, then $\op{Verify}$ on both triples, then finalize. The
first signature is accepted, $\Clean_{\sigma}$ asking only that $\sigma^{*}$ be a signature and
that $\V{Ver}[\vk,\msg,\sigma^{*}] \neq 0$, which holds, it being unset; $\V{Ver}[\vk,\msg,
\sigma^{*}] \gets 1$. The second is accepted for the same reason at the \emph{other} message,
$\V{Ver}[\vk,\msg',\sigma^{*}]$ being unset as well, and $\V{Ver}[\vk,\msg',\sigma^{*}] \gets 1$.
Both verifications then return the recorded $1$, and the verdict at line~\ref{ln:gcoll} gives $1$
with certainty.
\end{proof}

The attack is not a defect to be patched --- or rather, it can be patched, by adding to
$\Clean_{\sigma}$ the conjunct that $\sigma$ be recorded valid for no other message under $\vk$,
and the interest is in what that would cost. Sanitization exists so that a functionality's
guarantees do not depend on the slot answering helpfully (Chapter~\ref{sec:sanitization}); each
conjunct of a $\Clean$ predicate is a guarantee bought, and each is also a value the simulator may
no longer choose. The functionality $\Fsig$ as written buys correctness and unforgeability and declines to buy
collision resistance, which is a defensible reading of what a signature functionality is for --- and
the consequence, made precise by Proposition~\ref{prop:sigcoll}, is that a realization may be
strictly stronger than the ideal object it realizes.

That is worth sitting with, because it inverts the usual intuition. An ideal functionality is
often read as the best possible behaviour, with realizations approximating it from below; here
$\psig$ satisfies a property $\Fsig$ provably fails, and no amount of emulation transfers it
downward --- Corollary~\ref{cor:proptransfer2} carries properties from the ideal object to the
realization, never the other way. What an ideal functionality fixes is what a protocol may be
\emph{relied on} for, not what it happens to achieve; anything a realization achieves beyond that
is real, and is invisible to every argument that goes through the ideal object.


%% ==================================================================
%% Pending issues. Unnumbered, like the introduction, so it does not take
%% an chapter letter: it is a closing note about the framework rather
%% than part of it. Technical items only -- open questions the body
%% already raises, gathered in one place, plus what the review of this
%% paper has and has not covered.
%% ==================================================================
\phantomsection
\addcontentsline{toc}{chapter}{Pending Issues}
\chapter*{Pending Issues}
\chapterheadstar{Pending Issues}

What this paper leaves open, gathered from the places the body raises it, together with what its
review does and does not cover. Presentation and typesetting are not listed.

\heading{Dummy completeness stops at responsive calls} Theorem~\ref{thm:dummy} hypothesises systems
whose cores place no responsive call. The dummy's relay is refused at the first such call and
answers $\none$ where a real adversary would answer with substance, so the regrouping fails.
Recovering completeness needs a matching notion on the environment's side, and we do not give one.

\heading{Games have no responsive theory} All three property-transfer results inherit that
condition from Theorem~\ref{thm:dummy}. A game system whose $\F$ reaches the slot through
$\Adv^{!}$ sits outside all three and must have its dummy-only hypothesis established directly.
Four of this paper's own functionalities decline the hypothesis for exactly this reason and argue
their transfers by hand.

\heading{The quantifier order has no converse} Definition~\ref{def:uc} fixes the simulator before
the environment, and implies the weaker reading in which the simulator may depend on both. The
converse we do not prove. Whether the two notions separate is open.

\heading{Blocked pairs are not characterized} Tightness keeps a blocked comparison faithful for the
private members, but it constrains no member of $I$, and an interface member admitting a bare used
name is reachable by a hosted fake instance. Pinning those, or admitting only
$\Zpid \cup \Apid$ where the environment reaches them directly, would close it. We give no general
sufficient condition and no characterization.

\heading{What spawning would cost} Admissibility withholds the whole name closure, which is more
than a comparison strictly needs. Letting a simulator spawn functionalities at identities $\varphi$
does not occupy would leave only $\IDs(\varphi) \setminus \IDs(\pi)$ to withhold --- empty for the
usual pairing. What that relaxation costs in composition is not worked out.

\heading{Trace atomicity is not temporal consistency} Every trace entry is assembled atomically,
and every append in the paper was checked. Whether an entry's timestamp and the set fetched
alongside it agree \emph{in time} is a different property, secured by responsiveness rather than by
atomicity. The two are easy to conflate and the paper nowhere separates them in one place.

\heading{What the review covers} Six of twenty-one audit units were examined in full, chosen
leaves-first by criticality, and in a single pass rather than the two independent ones the protocol
asks for. Every defect found clustered on one mechanism --- what a bundle places outward, and
whether it may use its hosted members while answering a responsive call --- and all were fixed.
The exposure is therefore in what one pass did not look for; the realization of the authenticated
channel is where a further pass should go, being the unit that actually contained defects.

%% ==================================================================
%% book's thebibliography opens a \chapter* and, like every starred
%% heading, adds nothing to the contents. The other two in this document
%% carry their own \addcontentsline; this one needs the same.
\cleardoublepage
\phantomsection
\addcontentsline{toc}{chapter}{Bibliography}
\begin{thebibliography}{99}
\chapterheadstar{Bibliography}

\bibitem{ADR02}
J.~H.~An, Y.~Dodis, and T.~Rabin.
\newblock On the security of joint signature and encryption.
\newblock In \emph{Advances in Cryptology --- EUROCRYPT 2002}, volume 2332 of \emph{LNCS}, pages
83--107. Springer, 2002.
\newblock \url{https://eprint.iacr.org/2002/046}

\bibitem{Bellare97}
M.~Bellare.
\newblock Practice-oriented provable-security.
\newblock In \emph{Proceedings of the First International Workshop on Information Security (ISW
'97)}, volume 1396 of \emph{LNCS}, pages 221--231. Springer, 1998.
\newblock \url{https://cseweb.ucsd.edu/~mihir/papers/isw-pops.pdf}

\bibitem{BMTZ17}
C.~Badertscher, U.~Maurer, D.~Tschudi, and V.~Zikas.
\newblock Bitcoin as a transaction ledger: A composable treatment.
\newblock In \emph{Advances in Cryptology --- CRYPTO 2017}, volume 10401 of \emph{LNCS}, pages
324--356. Springer, 2017.
\newblock \url{https://eprint.iacr.org/2017/149}

\bibitem{BP15}
A.~Bauer and M.~Pretnar.
\newblock Programming with algebraic effects and handlers.
\newblock \emph{Journal of Logical and Algebraic Methods in Programming}, 84(1):108--123, 2015.
\newblock \url{https://arxiv.org/abs/1203.1539}

\bibitem{BR06}
M.~Bellare and P.~Rogaway.
\newblock The security of triple encryption and a framework for code-based game-playing proofs.
\newblock In \emph{Advances in Cryptology --- EUROCRYPT 2006}, volume 4004 of \emph{LNCS}, pages
409--426. Springer, 2006.
\newblock \url{https://eprint.iacr.org/2004/331}

\bibitem{Canetti04}
R.~Canetti.
\newblock Universally composable signature, certification, and authentication.
\newblock In \emph{17th IEEE Computer Security Foundations Workshop (CSFW 2004)}, pages 219--233.
IEEE, 2004.
\newblock \url{https://eprint.iacr.org/2003/239}

\bibitem{Canetti20}
R.~Canetti.
\newblock Universally composable security.
\newblock \emph{Journal of the ACM}, 67(5):28:1--28:94, 2020.
\newblock \url{https://eprint.iacr.org/2000/067}

\bibitem{CDPW07}
R.~Canetti, Y.~Dodis, R.~Pass, and S.~Walfish.
\newblock Universally composable security with global setup.
\newblock In \emph{Theory of Cryptography --- TCC 2007}, volume 4392 of \emph{LNCS}, pages 61--85.
Springer, 2007.
\newblock \url{https://eprint.iacr.org/2006/432}

\bibitem{FKKKW26}
P.~Farshim, M.~Karvonen, A.~Knispel, M.~Kohlweiss, and P.~Wadler.
\newblock UC, categorically: Rigorous diagrammatic proofs.
\newblock \emph{arXiv preprint arXiv:2608.04521}, 2026.
\newblock \url{https://arxiv.org/abs/2608.04521}

\bibitem{GMR88}
S.~Goldwasser, S.~Micali, and R.~Rivest.
\newblock A digital signature scheme secure against adaptive chosen-message attacks.
\newblock \emph{SIAM Journal on Computing}, 17(2):281--308, 1988.
\newblock \url{https://people.csail.mit.edu/silvio/Selected\%20Scientific\%20Papers/Digital\%20Signatures/A_Digital_Signature_Scheme_Secure_Against_Adaptive_Chosen-Message_Attack.pdf}

\bibitem{KMTZ13}
J.~Katz, U.~Maurer, B.~Tackmann, and V.~Zikas.
\newblock Universally composable synchronous computation.
\newblock In \emph{Theory of Cryptography --- TCC 2013}, volume 7785 of \emph{LNCS}, pages
477--498. Springer, 2013.
\newblock \url{https://eprint.iacr.org/2011/310}

\bibitem{PP09}
G.~Plotkin and M.~Pretnar.
\newblock Handlers of algebraic effects.
\newblock In \emph{18th European Symposium on Programming (ESOP 2009)}, volume 5502 of
\emph{LNCS}, pages 80--94. Springer, 2009.
\newblock \url{https://homepages.inf.ed.ac.uk/gdp/publications/Effect_Handlers.pdf}

\bibitem{Shoup04}
V.~Shoup.
\newblock Sequences of games: A tool for taming complexity in security proofs.
\newblock \emph{IACR Cryptology ePrint Archive}, Report 2004/332, 2004.
\newblock \url{https://eprint.iacr.org/2004/332}

\end{thebibliography}


\end{document}
